% =====================================================================
%  GENERATED; DO NOT EDIT BY HAND.  Rebuild: sh ebook-method/mindre-arbejder-2/rebuild.sh
%  Harald Høffding, Mindre Arbejder. Anden Række. København: Gyldendalske Boghandel (Nordisk Forlag), 1905.
%  Aarhus Stiftsbogtrykkeri.  237 pp.; seventeen lectures, essays, speeches and reviews of 1891-1905.
%  Diplomatic transcription (HOEFFDING-QUEUE #11, with #5 »Begrebet Analogi«, pp. 33-46), 2026-10-15.
%  ACCURACY (2026-10-16).  Scan: Google/NYPL bitonal 600 ppi (Internet Archive mindrearbejder01hfgoog), 237 pp.,
%  complete; no second copy online.  Readings: our tesseract (tessdata_best dan, 300 dpi; own non-words 0.77 %) as the
%  layer, the Archive's ABBYY as the witness; collated (ebook-method/collate/check.py; decisions.tsv: rules, blind crop
%  reads with two planted controls a sheet and an audit sample of every rule and voter class, adjudication by sense).
%  Later rounds: an edge sweep for the cut margins; the 15 footnotes read whole and placed from a marks table; a
%  structure reading (heads, verse, the asterisms); italics from slant + ABBYY agreement and the readers; the 24 most
%  damaged pages proofread whole.  Blind page check (6 pages sampled, seed 2026101511): 18 word errors in 2,029 words,
%  all on the 3 damaged pages among them (8.9 per 1,000 there), mended; none on the clean pages.
%  The English translation was a second reading: its translators listed 373 doubtful spots; 230 went back to the
%  image (crop sheets, 5 readers: 105 faults of the transcription, 75 the print's own, 31 letters supplied), and 8
%  pages whose margin is cut were read whole again (32, 46, 54, 112, 122, 164, 170, 179): 306 rows, odd-fixes.tsv.
%  That round found faults the earlier ones had missed on damaged pages: expect a few residual errors per damaged page
%  not read whole (lost ø-strokes, smeared line starts read as capitals, cut letters), near none on clean pages.
%  SCAN DAMAGE: the scan cuts the gutter margin on many pages (letters, or short words, lost at line ends of even and
%  line starts of odd pages).  These are SUPPLIED BY SENSE, silently in the text; each is recorded in proof-fixes.tsv or
%  odd-fixes.tsv ("supplied by sense"; "uncertain" where the word cannot be known from its traces).  Pages read whole:
%  32, 46, 49, 54, 72, 90, 97, 100, 110, 112, 115, 116, 118, 120, 122, 126, 131, 147, 150, 152, 156, 162, 164, 170,
%  173, 176, 177, 178, 179, 183, 218, 221, 225, 234.
%  Rules: ø where the bitonal scan lost the stroke (religiose -> religiøse); German ü/ä/ö where it lost the dots
%  (fiir -> für).  The Indhold and the heads differ in places: both kept as printed (see the comment at the Indhold).
% =====================================================================
\documentclass[12pt,a4paper,openany]{book}
\usepackage[utf8]{inputenc}
\usepackage[T1]{fontenc}
\usepackage{libertinus}
\usepackage{libertinust1math}
\usepackage[danish]{babel}
\usepackage[protrusion=true,expansion=true]{microtype}
\usepackage{setspace}
\usepackage[margin=1.2in]{geometry}
\usepackage{marginnote}
\usepackage{hyperref}
\hypersetup{pdftitle={Mindre Arbejder. Anden Række}, pdfauthor={Harald Høffding}, hidelinks}
\setstretch{1.3}
\renewcommand{\thefootnote}{\arabic{footnote})}
\newcommand{\opage}[1]{\marginnote{\footnotesize\textit{[#1]}}}
\newcommand{\essay}[2][]{\chapter*{#2}\addcontentsline{toc}{chapter}{#1}}   % [toc title]{title as printed, with its note}
\newcommand{\streg}{\begin{center}\rule{3em}{0.4pt}\end{center}}
\newcommand{\stjerner}{\begin{center}*\hspace{0.4em}*\\[-0.9ex]*\end{center}}   % the print's asterism (⁂), pp. 50, 51, 53, 55, 221
\begin{document}

\frontmatter
% PDF 11 (half-title), PDF 13 (title), PDF 14 (printer), PDF 15 (Indhold): read by the structure readers (struct-K2.txt S-109..112)
% and the Archive's ABBYY (PDF 11, 14).  The cover (PDF 9) repeats the title with "KØBENHAVN OG KRISTIANIA".
\begin{titlepage}\centering
{\Large MINDRE ARBEJDER}\\[1ex]{\small ANDEN RÆKKE}
\end{titlepage}
\begin{titlepage}\centering\vspace*{2cm}
{\large\textsc{Harald Høffding}}\\[1ex]\rule{3em}{0.4pt}\\[3ex]{\Huge MINDRE ARBEJDER}\\[2ex]{\large ANDEN RÆKKE}\par\vfill
{\small KØBENHAVN}\\[0.5ex]{\small GYLDENDALSKE BOGHANDEL}\\[0.5ex]{\small --- NORDISK FORLAG ---}\\[1ex]1905
\end{titlepage}
\thispagestyle{empty}\vspace*{\fill}\begin{center}{\small AARHUS STIFTSBOGTRYKKERI}\end{center}\clearpage
\chapter*{INDHOLD.}
\streg
\begin{flushright}\small Side\end{flushright}
\noindent\begin{tabular}{@{}p{0.86\textwidth}@{\hspace{0.5em}}r@{}}
Filosofien og Livet. (1901) & 1\\
Psykologiens nuværende Stilling og dens Forhold til dens Nabovidenskaber. (1904) & 21\\
Begrebet Analogi i dets filosofiske Betydning. (1904) & 33\\
Forord og Efterskrift til min Religionsfilosofi. (1905) & 47\\
Om nogle religionsfilosofiske Arbejder fra den nyeste Tid. (1903) & 58\\
Parmenides fra Elea. (1901) & 141\\
Holberg som Filosof. (1903) & 150\\
Henrik Steffens (1902) & 162\\
Henrik Steffens i det danske Aandsliv. (1902) & 172\\
Søren Kierkegaard. En Karakteristik. (1903) & 181\\
Fichtes »Taler til den tyske Nation« og de danske Nordslesvigere (1899) & 192\\
Harnack om Kristendommens Væsen. (1901) & 203\\
Tale for Georg Brandes. (1891) & 209\\
Tale for Henrik Ibsen. (1898) & 213\\
Anmeldelse af Bjørnson's »Poul Lange og Tora Parsberg«. (1898) & 217\\
Julius Langes Breve. (1902) & 224\\
Goethe. (1899) & 226\\
\end{tabular}
% The Indhold and the heads differ (both recorded, neither normalised): »Begrebet Analogi i dets« / »OG DETS«; »Holberg som
% Filosof« / »LUDVIG HOLBERG SOM FILOSOF«; Kierkegaard (1903) / (1902); Fichte »... og de danske Nordslesvigere (1899)« /
% »... OG NORDSLESVIGERNES NATIONALE KAMP. (1898)«; »Tale for Henrik Ibsen« / »TALE TIL HENRIK IBSEN«; the Bjørnson review's
% Indhold title / the head »BJØRNSTJERNE BJØRNSON: POUL LANGE OG TORA PARSBERG.«  Dates in the heads only where printed there.
\mainmatter
% --- p. 1 ---
\setcounter{footnote}{0}%
% CHECK p. 1: centred line: heading or verse
\opage{1}\essay[Filosofien og Livet. (1901)]{FILOSOFIEN OG LIVET.}
\streg

Engang i sine unge Dage mødte Søren Kierkegaard F. C.
Sibbern paa Gammeltorv og spurgte ham, hvorledes Filosofien
forholder sig til Livet i Virkeligheden. Sibbern blev, som han
selv har fortalt, forundret over Spørgsmaalet, fordi al hans Filosofi
gik ud paa at studere Liv og Virkelighed; men han forstod,
at En, der var stærkt paavirket af Hegel, kunde spørge
saaledes. Dog falder Spørgsmaalet ikke bort, fordi Livet og Virkeligheden
blive Genstande for Filosofien. Tvertimod stilles det
netop derved endnu skarpere. Thi Livet skal jo først og fremmest
leves, ikke studeres, og Virkeligheden stiller os først og
fremmest Opgaver for vor Villen og Handlen. Forat studere og
tænke maa man stille sig udenfor Livet, se paa dette som paa
en Genstand. Da Wieland raadede den unge Arthur Schopenhauer
fra at studere Filosofi, svarede denne: »Livet er en mislig
Sag; jeg har besluttet at tilbringe mit med at tænke derover.«
Schopenhauer var altsaa enig med Sibbern i, at Filosofen
skal studere Livet; men han glæder sig over, at dette
Studium kan lægge Beslag paa hele Livet, saa at man slipper
for det »virkelige Livs« Pinagtigheder. Schopenhauer kunde til
sin Filosofis Motto have brugt Goethes Ord: »Was im Leben
uns verdriesst, man im Bilde gern geniesst«. Det kommer da
an paa Forholdet mellem Billedet af Livet og selve Livet. Uenigheden
mellem Sibbern og Schopenhauer bestaar i, at Sibbern
% --- p. 2 ---
\setcounter{footnote}{0}%
\opage{2}er begejstret over, at han gennem Filosofien staar i Forhold til
Livet, medens Schopenhauer vil benytte Beskæftigelsen med Billedet
af Livet til at slippe fra selve Livet.

Klarest vil maaske det hele Spørgsmaal kunne belyses, naar
Vægten lægges paa, at det at filosofere selv er en Form for
aandeligt Liv. Den Strid, der synes at bestaa mellem Videnskaben
og Livet, kommer for en stor Del af, at man fra deres
Side, der mest føre Ordet »Liv« i Munden, overser, at ogsaa
Tænken og Forsken er en Form for aandeligt Liv, medens Tænkerne
og Forskerne ere tilbøjelige til at stille Tankelivet over
eller dog udenfor de andre Former for aandeligt Liv. Ved
Spørgsmaalet om Filosofiens, eller idethele Videnskabens Forhold
til Livet, spørges der altsaa egentligt om Forholdet mellem
to Former af aandeligt Liv.

\section*{I.}

Grunden til, at man ved Filosofien snarere end ved andre
Videnskaber kommer til at spørge om dens Forhold til Livet,
ligger i, at Filosofien særligt beskæftiger sig med Spørgsmaal,
der vedkomme Livet i dets Helhed og ligge paa Grænsen
mellem Tankelivet og de andre Livsformer. Filosofiske Problemer
angaa Erkendelsens Natur, Betingelser og Grænser, Tilværelsens
Væsen og Værd, samt Grundlaget for Vurderingen
af menneskelige Handlinger og Samfundsforhold. Det er Filosofiens
Opgave at samle og bruge Alt, hvad der kan tjene til
at belyse disse Spørgsmaal, og derigennem at orientere os overfor
Livets og Tilværelsens sidste Gaader. I sine Forsøg herpaa
træder Filosofien op som en selvstændig Form for Aandsliv,
der virker bestemmende tilbage paa Aandslivets andre Former.
Saaledes har Vedantafilosofien været af stor Betydning for den
indiske Aandsudvikling; Platonismen har vældigt paavirket den
kristne Kirkes Udvikling i Oldtid og Middelalder, og i det forløbne
Aarhundrede have Spinozas og Kants Ideer baade direkte
og indirekte paavirket den aandelige Udvikling, selv hos dem,
der mente slet ikke at have Noget med Filosofi at gøre. Der
er her stadigt et Tankearbejde at gøre, uden hvilket Aandslivet
vil miste Klarhed og Dybde. Og dette Arbejde kan kun gøres
% --- p. 3 ---
\setcounter{footnote}{0}%
\opage{3}ret, naar det gøres med fuld Frihed, --- naar Tankearbejdet
staar som en selvstændig Form for Livet, ikke blot som et Middel,
men som Formaal, altsaa som det, der har umiddelbar Værdi.
Filosofien skal søge at gøre Regnskabet op med Hensyn til Menneskets
Stilling i Tilværelsen som erkendende, følende og villende,
og en saadan Opgørelse kan kun have Betydning, naar
den gøres uhildet, hensynsløst, med fuld Frihed, --- hvilket Resultat
den saa end fører til. Tanken er den Del af Livet, der
skal kaste Lys over de andre Dele af Livet, og det kan da ikke
nytte, at disse andre Dele af Livet kalde sig selv for »Livet«
og ville hovmesterere Tanken. De Erfaringer, Tanken behøver
forat kunne behandle sine Opgaver, maa den selv søge og selv
anerkende: kun ved Tankens Hjælp kan der skelnes mellem
virkelig Erfaring og blotte Forestillinger eller Illusioner. I sit
Forsøg paa at komme til Klarhed, det vil sige til fuld Bevidsthed
om Menneskets Stilling i Tilværelsen, arbejder Filosofien
Haand i Haand med al anden videnskabelig Forsken og har, i
sine betydningsfuldeste Perioder, ikke skilt sin Sag fra dennes.

Men hvor vanskeligt er det ikke at komme til fuld Bevidsthed!
Bag vor klareste Bevidsthed ligger der stedse ubevidste
Tendenser og Dispositioner, der først senere kunne drages frem.
Vor Bevidsthed arbejder sig stedse frem af et dunkelt Kaos,
hvis spredte Elementer formes ved en uvilkaarlig Sammenfatten.
Selv vor klareste Tænken har derfor sin Akilleshæl, nemlig det
Punkt, hvor de ubevidste Kræfter virke. Denne Akilleshæl kan
siden opdages, og Tænkningen kan maaske undgaa at blive oftere
saaret paa dette Sted. Da Tetis dyppede sin Søn i Ild forat
gøre ham usaarlig, holdt hun ham i Hælen; hun var ikke saa
klog at dyppe ham to Gange og anden Gang holde ham i den
anden Hæl. Denne Klogskab har Filosofien vundet paa sin Vej,
jo mere den har frigjort sig fra Dogmatismen. Ethvert filosofisk
Forsøg hviler paa visse Forudsætninger, som den Filosoferende
ikke kan være sig fuldt bevidst, og som udtrykker den Afhængighed,
Tanken trods al sin Energi stedse staar i til andre Sider
af Livet. Men saadanne Forudsætninger kan saa senere
Tænkning drage frem, og derpaa beror --- afset fra det nye
Stof, den fortsatte ydre og indre Erfaring kan bringe  --- den
% --- p. 4 ---
\setcounter{footnote}{0}%
\opage{4}filosofiske Forskens Fremskridt. Akilles maa stedse paany gøres
usaarlig, da der stedse forestaar nye Kampe, og derfor bliver
der altid en Akilleshæl, saalænge Livet og Udviklingen varer.
Filosofen kan ikke undgaa at være Barn af sin Tid og sit Folk.
I sit mest aktive Tankearbejde ledes han ofte af en uvilkaarlig
Trang til at finde Udtryk for Ejendommeligheder og Tendenser,
han ikke agter paa --- netop fordi de ligge ham saa nær. Fra
denne Side set staar Filosofien i et passivt, modtagende Forhold
til Aandslivets andre Sider. De filosofiske Tanker blive her
Symptomer eller Symboler, saa meget Arbejde deres Udformning
end har kostet. Platons Idelære staar overfor Erfaringsverdenens
Dunkelhed og Uro, ligesom Olympen i sin evige
Glans staar over Menneskelivets Kampe og Omskiftelser, og i
Kants »Kritik der reinen Vernunft« staar den evige Stræben
som det forløsende Ord, ikke mindre end i Goethes »Faust«.
Hvad Myten og Digtningen udtrykke i Billedets Form, udtrykker
Filosofen i sine Ideer. Og hvor det væsentligt er Livets
dunkle Trang, der driver Tanken frem, alt medens denne mener
at følge sine egne Love, --- der kan det blive vanskeligt at afgøre,
hvor Tænkningen hører op, og hvor Poesien begynder.
Maaske kunne de sidste Gaader kun besvares i Poesiens symbolske
Form. Det bliver da Kritikens Sag at agte paa, at Grænsen
mellem Tænkning og Poesi overholdes.

Det er altsaa et dobbeltsidigt Forhold, Filosofien staar i til
Livet, --- et aktivt i det bevidste Tankearbejde paa Grænseproblemernes
Omraade, og et passivt i den uvilkaarlige Symboliseren.
Begge Sider af Forholdet have deres store Betydning, og
begge ville stedse kunne paavises, hvor Filosofiens Gerning øves
i større Stil. Efter de forskellige filosofiske Retningers Ejendommelighed
vil Vægten komme til at ligge snart mere paa den
ene, snart mere paa den anden af disse to Sider. I Romantikens
Filosofi fik --- trods det betydelige Tankearbejde, der øvedes,
og som man nutildags ikke altid tilstrækkeligt anerkender --- Symboliseringen
Overvægten over den egentlige Tænken, medens
Positivismen stræbte at holde sig saa nær som muligt til de bestemte,
givne Erfaringer. Mellem disse to Retninger har Kampen
staaet i det 19. Aarhundredes Filosofi. Vi kunne nu ved
% --- p. 5 ---
\setcounter{footnote}{0}%
\opage{5}Aarhundredets Slutning se tilbage paa denne Kamp uden Partiskhed
og forstaa den udfra de historiske Forhold, under hvilke
det filosofiske Arbejde skulde øves. Hvad der ganske særligt
karakteriserer det nu forløbne Aarhundrede, i Sammenligning
med dets Forgænger, var nu ogsaa netop dets Blik for denne
historiske Betingethed af alt Aandsliv, altsaa ogsaa af Tankelivet.
Det har indskærpet det Afhængighedsforhold, i hvilket den
Enkelte stedse staar til hele Slægten. Romantiken og Positivismen
fremdrog hver paa sin Maade denne Afhængighed. Johann
Gottlieb Fichte, der kan betragtes som Grundlægger af
Romantikens Filosofi, udtalte, at ved Valget af de Principer, der
tilsidst bestemme vor Verdensanskuelse, er det Tænkerens Personlighed,
der gør Udslaget. Et filosofisk System var for ham
ikke et dødt Bohave, man paa udvortes Maade kunde anskaffe
sig en Gang for alle, men det var en Tankeverden, der fødtes
i Menneskets inderste Sjæl. Hvilken Filosofi man har, beror
paa, hvilket Menneske man er. Hegel førte denne Tankegang
videre i sit geniale, men dunkle Skrift »Phänomenologie des
Geistes«, der skulde vise, hvorledes Tænkningens Art, Grad og
Retning bestemmes ved det Udviklingstrin, til hvilket Aandslivet
var naat. I Comtes Lære om de tre Stadier indskærpes
det, at Erkendelsens Metode og Resultater blive forskellige efter
Menneskets Udviklingstrin: i det »teologiske« Stadium anvendes
andre Forklaringsmaader end i det »metafysiske«, og det »positive«
Stadium, den nyere Erfaringsvidenskabs Standpunkt, stiller
sine Spørgsmaal og fordrer sine Svar efter andre Forudsætninger
end de to tidligere Stadier. Hvad der i det ene Stadium
er god Forklaring, er det ikke i det andet. Erkendelsens Fremgangsmaade,
Opgaver og Resultater afhænge saaledes stedse af
Aandslivets og Kulturens almindelige Udvikling. Det mest omfattende
Synspunkt herfor er givet i Udviklingsfilosofien, der
betragter al Erkendelse i Lyset af Kampen for Tilværelsen og
af Tillempelsens Nødvendighed. Livets og Livskampens Vilkaar
bestemme gennem mange Mellemled de Former, i hvilke Problemerne
stilles, og de Løsninger, man mener at finde.

At Filosofien saa klart er bleven opmærksom herpaa, betegner
et stort Fremskridt i Forstaaelsen af Forholdet mellem
% --- p. 6 ---
\setcounter{footnote}{0}%
\opage{6}Tankelivet og Livets andre Former. --- Jeg gaar nu over til at
undersøge dette Forhold mere i det Enkelte.

\section*{II}

Filosofen er først og fremmest Logiker. Man har endog
forsøgt den Definition af Filosofien, at den er Videnskaben om
den menneskelige Erkendelse. Denne Definition er ikke udtømmende,
men den fremhæver et Hovedpunkt. En Undersøgelse
af vore Begrebers, Dommes og Slutningers Natur og Oprindelse
og af hele den Maade, paa hvilken de hænge sammen med vor
Bevidstheds Væsen og Virkemaade, vil altid være af afgørende
Betydning i Filosofien. Det er tilsidst Logikerens Opgave at
sige, hvad det er at \textit{forstaa}. Hvad kræves der til fuld Forstaaelse
af en Sag, og hvorvidt er fuld Forstaaelse mulig? Det
er Bestræbelsen for at forstaa, der fører til Udviklingen af Begreber,
Domme og Slutninger, og naar vi søge at forstaa Noget,
omsætte vi det i vore Begrebers, Dommes og Slutningers Former,
idet vi danne saa sammenhængende Tankerækker, som vi
formaa, paa Grundlag af det i Erfaringen givne Indhold. Men
med hvilken Ret omsætte vi saaledes Tilværelsens Indhold, som
vi jo dog ikke selv have frembragt, i vor Tænknings Former?
Dersom Tilværelsen skulde være fuldt forstaaelig (rationel), maatte
alt dens Indhold kunne fremstilles som et stort System af Begreber,
Domme og Slutninger, og hver enkelt Begivenhed i Verden
maatte kunne stilles i samme Forhold til andre Begivenheder,
som Konklusionen i en Slutning staar i til Præmisserne.
Paa alle Videnskabens Omraader arbejder Tanken utrætteligt
paa gennem Iagttagelse og Analyse at føre Tilværelsens Indhold
tilbage til saa simple Synspunkter, at en rent logisk Formulering
bliver mulig. Erkendelsesproblemet opstilles, naar det
Spørgsmaal opkastes, hvorvidt Tilværelsesindholdet paa udtømmende
Maade kan omsættes i Tankens Former.

For Logiken har kun det Gyldighed, der paa modsigelsesfri
Maade kan indordnes som Led i en Tankerække. Men selvom
en saadan rent logisk Formulering var mulig af Alt i Tilværelsen,
vilde Tilværelsens Indhold og Beskaffenhed ikke derved
være udtømmende karakteriseret. Vi vilde i det Højeste kunne
% --- p. 7 ---
\setcounter{footnote}{0}%
\opage{7}sige, at Tilværelsen kan opfattes eller forstaas, \textit{som om} dens
Tilstande og Begivenheder forholdt sig til hverandre, ligesom
Leddene i en Slutningsrække forholde sig til hverandre. Men
Tilstandene og Begivenhederne indeholde stedse Andet og Mere
end de rent logiske Forhold, efter hvilke de kunne ordnes i vor
Forstaaelse. De specielle Egenskaber (Kvaliteter), hvormed de
optræde, ere ikke forklarede, fordi vi med strengeste Nøjagtighed
kunne paavise deres Forhold til andre Egenskaber. De specielle
Farvekvaliteter f. Ex. ere ikke forklarede i deres Ejendommelighed,
fordi det lovmæssige Forhold mellem Farverækken
og Lysets Brydningsgrader er paavist. Og der bliver i hvert
Tilfælde den store Forskel tilbage mellem en Tankerække og
en Række af Begivenheder eller Tilstande, at i Tankerækken
forholder det ene Led sig til det andet blot som Grund til Følge,
men i Begivenhedernes Række forholder det ene Led sig tillige
til det andet som Aarsag til Virkning. Den Begivenhed, vi kalde
Aarsagen, gaar i Tiden forud for den Begivenhed, vi kalde Virkningen.
Tidsforholdet har her afgørende Betydning, medens
i Tankerækken Konklusionen egentligt er given, saasnart vi have
Præmisserne, hvor lang Tid vi end behøve forat indse det. I
den rene Logik har Tidsforholdet ingen Betydning. Men i Tilværelsen
\textit{sker} der noget; Forskellen mellem Fortid og Fremtid
har her real Betydning. Der gives her et »endnu ikke« og et
»ikke mere«; der opstaar derved en Modsætning mellem Erindring
og Forventning, og det bliver Spørgsmaalet, om vi med
Rette grunde vore Forventninger paa vore Erindringer, og om
de Forventninger, vi saaledes begrunde, blive bekræftede ved
de følgende Begivenheder. Saasnart Filosofien --- med Hume
og Kant --- blev klar over den væsentlige Forskel mellem Forholdet
mellem Grund og Følge paa den ene Side og Forholdet
mellem Aarsag og Virkning paa den anden Side, var Erkendelsesproblemet
stillet for Alvor, og det blev et stort Spørgsmaal,
om Tilværelsen kunde antages at være forstaaelig, rationel
i sig selv.

Som Logiker tager Filosofen kun Hensyn til en eneste aandelig
Fornødenhed, Trangen til Forstaaelse. Det er hans Opgave
% --- p. 8 ---
\setcounter{footnote}{0}%
\opage{8}at afstikke Grænsen mellem \textit{det} Forstaaelige og det Uforstaaelige.
Men der gives andre aandelige Fornødenheder end Trangen
til at forstaa, og det hører til Filosofiens Opgaver ogsaa at
undersøge disse andre Fornødenheder og deres Forhold til hin
rent intellektuelle Trang. Filosofen skal ikke blot være Logiker,
men ogsaa Psykolog. Selve Begrebet Forstaaelse er et psykologisk
Begreb, da al Forstaaelse bliver til ved en Bevidsthedsproces.
Al Forstaaelse er en Ytring af \textit{det personlige} Liv, hvis
mangfoldige Former og Trin Psykologen skal følge Skridt for
Skridt fra dets mest elementære Ytringer til dets højeste. Her
optræder Filosofen som Iagttager og skal øve en Beskrivelsens
Kunst, der snart kan kræve Evner af lignende Art som Digterens,
snart fører til Anvendelse af den experimentale Metode.
De ofte zarte og halvt skjulte Nuancer, i hvilke sjælelige Fænomener
fremtræde, maa opsøges ad højst forskellige Veje. Og
naar Iagttagelsen og Beskrivelsen lykkes, staar endnu Forklaringen
tilbage, Paavisningen af Aarsagsforholdene paa det sjælelige
Omraade. Her har Filosofen det for ham mest nærliggende
Exempel paa et Tilværelsesindhold, som ikke udtømmende kan
fremstilles i rent logisk Form. Og dog arbejder han stadigt paa
en saadan Fremstilling. Selv de højeste sjælelige Tilstande søger
han at finde Love og Betingelser for; selv deres skjulte og stille
Fremvæxt søger han at efterspore. Han søger en Kontinuitet i
det \textit{personlige Livs} Udvikling, der kan gøre det muligt at opfatte
de forudgaaende Tilstande, \textit{som om} de forholdt sig til de
følgende som Præmisser til Konklusion. Men hans egen direkte
Erfaring viser ham her tydeligt, at de virkelige Tilstande ikke
ere udtømmende karakteriserede, fordi deres indbyrdes Forhold
kan fremstilles paa denne Maade. --- ---

Baade som Psykolog og som Logiker kommer Filosofen i
Berøring med andre Videnskaber, idet han benytter Naturvidenskaben
og de historiske Videnskaber til at belyse Sjælelivets
Udvikling, og idet han i de specielle videnskabelige Metoder
finder Anvendelser af de almindelige logiske Principer. Og i
begge Egenskaber vil han ogsaa tilsidst føle sig stillet overfor
\textit{det metafysiske} Problem, Problemet om Tilværelsens inderste
Natur og Væsen. Den Maade, paa hvilken dette Problem be%
% --- p. 9 ---
\setcounter{footnote}{0}%
\opage{9}handles, vil væsentligt afhænge af det Resultat, man er kommen
til ved Behandlingen af Erkendelsesproblemet, og af den Opfattelse,
man har dannet sig om Sjælelivets Natur og om dets Forhold
til andre Sider af Tilværelsen. Og hvis man helt afviser
dette Problem, saa vil det ogsaa væsentligt være af Grunde, der
ere hentede fra Logik eller Psykologi. I denne Sammenhæng vil
det især være af Interesse at se, hvilken Betydning et Forsøg
paa at behandle Tilværelsesproblemet, altsaa paa at danne en
metafysisk Hypotese, kan have.

En saadan Hypotese kan tjene som en Ide eller et Princip,
en ledende Tanke, der stiller os videregaaende Spørgsmaal
og Problemer, hvor en ængstelig Begrænsning til de enkelte
foreliggende Erfaringer vilde stanse Forskningen. Der vil kunne
gives os en større Horizont for vort Arbejde paa de specielle
Omraader, --- en Advarsel mod at stanse for tidligt. Saaledes
vil den Antagelse, at Tilværelsen baade fra sin materielle og
sin aandelige Side danner en stor Kontinuitet, et stort, ubrudt
System af Aarsagsrækker, naturligt opfordre til at vedblive at
søge, indtil man i det Enkelte har fundet de bestemte Mellemled
og derved udfyldt de Huller, der syntes at være i de ydre
eller indre Begivenheders Rækker. Derimod ville spiritualistiske
og indeterministiske Antagelser indeholde en Fristelse til at
stanse Forskningen for tidligt og slaa sig til Ro med »Sjælens«
eller »den frie Villies« Indgriben i Begivenhedernes Gang. Dette
er den ene Betydning, en metafysisk Hypotese kan have.

En anden Betydning har den ved at være et Forsøg paa
at gøre vort intellektuelle Regnskab op. Vi drage i den metafysiske
Hypotese de sidste Konsekvenser af vor Erkendelse og
se, hvorhen de pege. En saadan Opgørelse vil være af stor
Betydning for det hele Aandsliv. Det er en fundamental Orientering,
der foretages. Tænkningens Historie viser os, at der her
er visse Hovedsystemer, mellem hvilke Kampen staar. Deres
Tal og deres specielle Karaktermærker ere ikke de samme til
enhver Tid. Nogle metafysiske Muligheder kunne falde bort som
ugyldige eller ufornødne, og andre komme i deres Sted, eller
de ældre Hypoteser maa omarbejdes for at kunne harmonere med
de Krav, nye Erfaringer stille. Maaske er det en evig Strid, vi
% --- p. 10 ---
\setcounter{footnote}{0}%
\opage{10}her staa overfor, men saa er det dog af Betydning, at vi blive
klare over, hvad Striden angaar, hvorfor den er nødvendig, og
hvorledes den maa føres.

Paa det Grænseomraade, hvor de metafysiske Hypoteser bevæge
sig, og hvor Tænkning og Poesi ofte saa umærkeligt kunne
gaa over i hinanden, vil Filosofens Individualitet, hans »personlige
Ligning« uundgaaeligt komme til at spille en Rolle. Overfor
det menneskelige Hang til Dogmatisme, --- det vil sige, til
at gøre sine Forudsætninger til de eneste rette eller eneste
mulige, --- er det af Vigtighed at være opmærksom herpaa.
Men deraf følger ikke, at ethvert Forsøg paa at opbygge en
Tankeverden paa Grænsen af vor sikre Erkendelse er uberettiget
og unyttigt, naar Forsøget foretages baade med kritisk
Sans og med dyb Følelse af Problemernes Brod. En borneret
Skepsis kan være ligesaa betænkelig for Aandslivets Klarhed
og Sundhed, som en borneret Dogmatisme. Hverken nye Tvivl
eller gamle Dogmer bør hindre os i at se os frit omkring, selv
naar vi ere komne ud i Grænseregionerne. Kun den uafviselige
Fordring gælder, at der maa tilstræbes en Harmoni mellem
den videnskabelige Forstaaelse af Tilværelsen og den Tankeverden,
den Enkelte opbygger paa Videnskabens Grænser. I Længden
kan en Strid mellem Videnskabens og det personlige Livs
Krav ikke bestaa. --- Der er en Ting, man i Regelen overser, naar
man taler om Modsætningen mellem »Tro og Viden«. Videnskabelig
Forsken er ikke et rent upersonligt Anliggende. Det
har den aldrig været for de store Forskere, hvem Videnskabernes
grundlæggende Tanker skyldes. I Bestræbelsen for at finde
Sammenhæng i Naturens Gang og i Aandens Liv, og mellem
begge, ytrer sig en personlig Trang, en aandelig Livsfornødenhed.
Vor Sjæls uvilkaarlige Stræben efter at bevare Enheden
med sig selv under alle Forandringer, der opleves og erfares,
og efter at sammenfatte alle de forskellige Elementer, der røre
sig i den, er den sidste Forklaring, der kan gives af den Utrættelighed,
med hvilken der arbejdes paa at finde Aarsagerne til
Alt, hvad der sker. Jo mere det Nye kan opfattes som Fortsættelse
eller Omsætning af det, der forud kendtes, des bedre
fyldestgøres hin Trang. Ønsket om at finde naturlige, det vil
% --- p. 11 ---
\setcounter{footnote}{0}%
\opage{11}sige: i Erfaringen kendte og paaviselige, Aarsager til Begivenhederne,
er derfor ikke blot paatvunget Erkendelsen udefra, men
udtrykker en væsentlig Side i vort Aandlivs Natur. Videnskabelig Forsken er et Personlighedsarbejde, ligesaa vel som Kunsten
er det, og ligesaa vel som den praktiske Karakterudvikling er
det. Der er en Analogi, og mere end en Analogi, mellem Tænkning
og Karakter: begge fordre Sammenhæng og Konsekvens
og bekæmpe det Sporadiske, det Isolerede og det Fragmentariske.
Kun ved denne Betragtning forstaas den dybe Glæde, der opstaar,
naar Tankearbejdet gaar frem, og naar Forstaaelse vindes.

\section*{III}

Filosofen er ikke blot Logiker og Psykolog; han er ogsaa
Etiker. Det er ikke blot hans Opgave at konstruere Tankerækker
og finde Aarsagsrækker, men ogsaa at prøve Værdier, at
gennemføre en sammenhængende Vurdering af menneskelige
Egenskaber, Handlinger og Samfundsformer. Al Vurdering bestemmes
tilsidst ved en Grundværdi, en umiddelbar Værdi, i
Forhold til hvilken al anden Værdi gælder. Hvad der for Mennesket
har umiddelbar Værdi, det vil sige: medfører umiddelbar
Tilfredsstillelse, søger han at fastholde, og hvis det ikke
umiddelbart kan fastholdes eller opnaas, søger han efter Midler
til at virkeliggøre det og fæstne det. De Formaal, Mennesket
stiller sig, bestemmes ved, hvad der har umiddelbar Værdi for
ham, og af de Formaal, han stiller sig, bestemmes igen de Normer,
den Maalestok, som han anvender ved Vurderingen af Egenskaber,
Handlinger og Livsforhold. Begreberne: Værdi, Formaal
og Norm ligge saaledes i samme Række. Udfra en vis Grundværdi
kan Tænkningen skride frem med logisk Nødvendighed,
naar der raades over den fornødne Erfaring til at faststille den
Maade, paa hvilken en Egenskab, en Handling eller en Samfundsordning
griber ind i det menneskelige Liv. Spørgsmaalet
bliver, om Egenskaben, Handlingen eller Samfundsordningen
ligger i en Række, der i sin Fortsættelse fører til Grundværdiens
Virkeliggørelse.

Den etiske Tænkning er altsaa kun en speciel Anvendelse
af logisk og psykologisk Tænkning. Forholdet mellem Formaal
% --- p. 12 ---
\setcounter{footnote}{0}%
\opage{12}og Middel er egentlig ikke forskelligt fra Forholdet mellem Grund
og Følge og mellem Aarsag og Virkning. Formaalets Værdi
indeholder Grunden til Midlets Værdi, og Midlet er Aarsag til
Formaalets Virkeliggørelse. Den, der har stillet sig et Formaal,
fordi Noget staar for ham som værdifuldt, kan ikke uden at komme
i Modsigelse med sig selv undlade at vælge de Midler, der vise
sig nødvendige til Formaalets Opnaaelse; i hvert Tilfælde maa
han tillægge disse Midler en afledet Værdi. Alt hvad der har
Værdi for os, søge vi uvilkaarligt, dog ikke altid med Bevidsthed,
at hævde, og denne Hævden ytrer sig først og fremmest i
en Anerkendelse af Værdien af de Egenskaber, Handlinger og
Samfundsordninger, der kunne være Midler og Veje til det
Værdifuldes Bestaaen. Fuldkomnest vil denne Hævden være,
naar selve Bestræbelserne for at finde saadanne Midler og Veje,
ja, naar selve disse Midler og Veje kunne faa umiddelbar Værdi
for os. Afstanden mellem Værdi og Norm og mellem Formaal
og Middel bliver da saa ringe som muligt.

Det er af Betydning at fremhæve dette indre Slægtskab
mellem Etik, Logik og Psykologi. Den Fordom er endnu ikke
forsvunden, at man ved Overgangen til Etiken maa afføre sig
sin ellers sædvanlige Maade at tænke og forske paa. Den etiske
Tænkning indfører kun den specielle Forudsætning af en Grundværdi,
men behøver ellers ingen særegne Love eller Principer.
Vanskelighederne for Etiken bero dels paa, at vor Erfarings
Omfang og Sikkerhed ikke altid tillader os at drage ganske nøjagtige
Konsekvenser, dels paa, at enhver etisk Tankegang kun
gælder under Forudsætning af bestemte Grundværdiers Gyldighed.
Den etiske Undersøgelse stanser, naar Spørgsmaalet om
selve disse Grundværdier rejses. Hvorledes vi filosofisk maa
stille os, naar Diskussionen har naat dette Punkt, ville vi i det
Følgende vise.

\section*{IV}

Baade som Logiker, som Psykolog og som Etiker støder
Filosofen paa Vanskeligheder, der hænge sammen med vor Erkendelses
Begrænsning. Det vil være af Interesse at se, hvorledes
han maa stille sig til disse Vanskeligheder.

% --- p. 13 ---
\setcounter{footnote}{0}%
\opage{13}Det er allerede ovenfor bemærket, at Tilværelsen ikke kan
omsættes til et logisk System, den fuldkomneste Form for Forstaaelse,
vi besidde. Den exakte Naturvidenskab nærmer sig til
dette intellektuelle Ideal paa de Omraader, hvor nøjagtig Maaling
og Beregning kan afløse den blotte Iagttagelse og Beskrivelse,
og hvor det derfor er muligt at føre Lovene for Begivenhedernes
Gang tilbage til en Række af Ligninger. Men selvom
en saadan Reduktion af det Kvalitative til det Kvantitative kunde
gennemføres paa alle Omraader, vilde Forskellen mellem den
logiske Form og det reale Indhold dog stedse blive tilbage. Det
reale, kvalitative Indhold maa stedse forudsættes og tages som
givet, men kan ikke udledes af selve den logiske Form: det er
altsaa rigere end denne Form, saa betydningsfuld denne end
maatte være. Der bestaar tilsidst et irrationelt (i matematisk
Forstand irrationelt) Forhold mellem vor Erkendelses Form og
dens Indhold. Tanken er selv en Del af Tilværelsen, og vi benytte
denne Del som en Form, i hvilken hele Tilværelsen skal
udtrykkes. Men denne Opgave kan ikke udtømmende løses, da
Tanken stedse søger og finder nyt Indhold, Indhold, den maa
tage, som det kommer, og ikke kan hexe frem af sine rene
Former. Selvom det Indhold, den hidtil har haft at bearbejde,
nogenlunde villigt har ladet sig udtrykke i Tankens Former,
saa ligger deri ingen Garanti for, at det vil gaa ligesaa let for
det fremtidigt givne Indhold. De rette Linjer, ved hvilke Tanken
beskriver Tilværelsens Bevægelse, kunde maaske være forsvindende
smaa Dele af en uendelig Kurve, hvis Formel det
ikke er den muligt at finde.

Det var dette irrationelle Forhold mellem Erkendelsens Form
og Indhold, Kant paa en noget uheldig Maade udtrykte ved at
opstille det meget omdisputerede »Ding an sich«. I den bestemte
Skikkelse, i hvilken han opstillede dette Begreb, kan det
ikke fastholdes. Men som Udtryk for det irrationelle Forhold
mellem Erkendelsens Form og Indhold kan det ikke forsvinde
fra Filosofien, men vil, saa ivrigt man end vil søge at dække
over det, stikke Hovedet op paa et eller andet Punkt i enhver
filosofisk Opfattelse. Det udtrykker den Kendsgerning, at der
intet Bevis kan føres for Tilværelsens Rationalitet. Det kunde
% --- p. 14 ---
\setcounter{footnote}{0}%
\opage{14}jo være, at hele vor stolte Videnskab kun blev mulig ved, at
Tilværelsen i lange Tider saa at sige kun havde vendt sin
rationelle Side ud imod os, men at vi, hvis den engang vendte
en af sine andre Sider ud imod os, vilde staa overfor en fuldstændig
Umulighed af at opstille sammenhængende Tankerækker,
der kunde anvendes paa den. Tilværelsens logiske Istid vilde
da begynde.

Ogsaa ad en anden Vej kunne de irrationelle Muligheder, en
indgaaende Drøftelse af Erkendelsesproblemet opdager, belyses.
Al Videnskab fremgaar af Livet, af Menneskets Trang og Villie
til at orientere sig i Tilværelsen. Sandheden er forsaavidt et
Villiesprodukt. Men er Sandheden nødvendig? Er den det
Højeste? Kan vor Villie ikke frigøre sig fra de Regler og Normer,
den dog selv har frembragt? --- Det Berettigede i en saadan
Spørgen (der ikke er saa ny, som man i den nyeste Tid
har ment) beror paa, at det, der for os er Sandhed, stedse svarer
til vort bestemte Livs- eller Udviklingstrin, og at der behøves
en ny Sandhed, naar et nyt Livstrin er naat. Selv den mest
exakte Videnskab, vi nu besidde, svarer til et bestemt Livstrin,
hvis Tid engang vil være forbi, saa at den maa afløses af en
ny Erkendelse. Dette følger ligefrem af, at al Erkendelse grunder
sig paa Erfaring. Al vor Videnskab opstaar ved Tydning
af de Erfaringer, der hidtil ere gjorte. Deraf følger nu dog ikke,
at man med tilbørlig god eller ond Villie uden videre kunde
bryde ud og stille sig udenfor den Sandhed, der er naat. Den,
der virkeligt har naat et nyt Livstrin, vil have nok at gøre med
at finde den nye Sandhed, der svarer til dette Trin. Lad ham
saa fremstille den, forat de Andre, der ikke have naat det »højere«
Trin, kunne faa en Anelse om de Begivenheder, der forestaa i Sandhedens
Verden! Kan han ikke det, er det kun Afmagt, naar han pukker
paa sin Villie, og det er et Spørgsmaal, om det virkeligt er en
Villie, der rører sig hos ham. Og de, der endnu staa paa det
»lavere« Trin, ville i hvert Tilfælde prøve paa, om det dog ikke
alligevel skulde være den gamle Logik og de gamle Sandheder,
han benytter i sine Protester mod »Sandheden«, og maaske
endog benytter paa en »altfor menneskelig« Maade.

Hvor der virkeligt naas et nyt Livstrin, vil Erkendelsen
% --- p. 15 ---
\setcounter{footnote}{0}%
\opage{15}strax begynde at virke forat finde en kontinuerlig Sammenhæng
mellem det nye Trin og de tidligere. Enhver ny Religion søger
saaledes at betragte de tidligere Religioner som forberedende
Trin, som Veje, der maatte vandres, før Tidens Fylde kunde
komme. Og den videnskabelige Erkendelse søger ad biologisk,
psykologisk og historisk (sociologisk) Vej at finde Lovene og
Vejene for Overgangen fra en Livsform til en anden. Saaledes
bære vi os ad, naar vi sammenligne Menneskets Erkendelse
med Dyrets, den græske Videnskab med de vilde Folkeslags
Forestillinger, den moderne Videnskab med den græske. Det er
denne Metode, Hegel, Comte og Spencer hver paa sin Maade
have anvendt forat vise Erkendelsens Afhængighed af Livets Trin.
Paa en saadan Maade vil Erkendelsen, naar det ikke længere
er muligt at bygge en sammenhængende Tankerække, stedse
gaa længere tilbage og undersøge, om der ikke ved selve Tankerækkens første Udspring viser sig en Sammenhæng, som baade
kan forklare, hvorledes Tankerækken blev mulig, og hvorfor den
ikke kunde fortsættes længere. ---

Hvad der saaledes gælder for vor Erkendelse idethele, gælder
særligt for den etiske Erkendelse. Etiken opererer, som vi
have set, med Begreberne Værdi og Formaal som grundlæggende.
Men her opstaar da det Spørgsmaal: \textit{hvilke} umiddelbare Værdier,
hvilke højeste Formaal skulle lægges til Grund ved Vurderingen
og bestemme dennes Norm og Maalestok? Der er
Mulighed for ligesaa mange etiske Systemer, som der kan gives
Grundværdier eller Grundformaal, udfra hvilke en konsekvent
gennemført Vurdering af menneskelige Egenskaber, Handlinger
og Samfundsforhold er mulig. Men fra det ene Grundformaal
eller den ene Grundværdi til den anden fører ingen rent intellektuel
Vej; Overgangen sker ved en psykologisk og historisk
Proces, der kan være lang og indviklet og foregaa gennem
mange Forskydninger og Omsætninger. Her gaar Etiken, ligesom
før Erkendelseslæren, tilbage til Biologi, Psykologi og Sociologi.
Forgæves vil man søge at omgaa denne Vanskelighed
ved at definere Etiken som Læren om Samfundets nødvendige
Krav til den Enkelte. Hvis den Enkelte har Magt, Snildhed og
Konsekvens, vil han ikke lade sig imponere heraf. Selvom man
% --- p. 16 ---
\setcounter{footnote}{0}%
\opage{16}maaske kan tvinge ham til at indrømme, at han gør klogest i
at rette sig efter de Krav, Samfundet stiller til ham, vil han
ikke deraf lade sig tvinge til at anerkende, at de ere gode, berettigede
og retfærdige. Denne Anerkendelse opstaar kun hos
ham, naar Samfundets Bestaaen og Udvikling faar umiddelbar
Værdi for ham. Han maa opdage Samfundslivets Værdi. Og
denne Opdagelse sker ikke ved blot Ræsonnement. Atter her
komme vi tilbage til Livets forskellige Trin: hvad der er etisk
konsekvent paa det ene Trin, behøver ikke at være det paa det
andet. Den etiske Tænknings Betydning bestaar i, at den drager
alle Konsekvenser af de Grundværdier, der kundgøre sig
paa et givet Livstrin; men den er ikke tilstrækkelig til direkte
og definitivt at føre fra et Livstrin til et andet.

Men den Sammenhæng, Etiken ikke kan skaffe os, fordi
Grundværdierne, og med dem Udgangspunkterne for de etiske
Tankerækker, kunne staa i Strid med hverandre, kan findes ved
at gaa tilbage til Biologi, Psykologi og Sociologi. Her kunne
vi studere Værdiernes Opstaaen, Udvikling og Forsvinden. Værdiernes
Historie kan have Kontinuitet, selvom de etiske Systemer,
der ere Projektioner udfra de enkelte Livstrin, i deres Principer
staa skarpt imod hverandre. ---

Men --- Nissen flytter med. Thi Biologi, Psykologi og Sociologi
ere jo selv Erkendelsesarter og bygge paa al Erkendelses
sidste Principer. Ogsaa den Maade, paa hvilken en biologisk,
psykologisk og sociologisk Erkendelse til given Tid udvikler sig,
vil bero paa det Livstrin, der er naat. Vi synes da at bevæge
os i en Kreds: selve Videnskaben om Livstrinnenes Udvikling
beror paa det Livstrins Beskaffenhed, til hvilket Udviklingen
har ført! Denne Kreds betegner dog ingen Fejl i vor Tankegang.
Den betegner kun, at vi her staa ved et Vexelvirkningsforhold,
der uafbrudt gør sig gældende: Erkendelsen afhænger
af Livet, og selve den Opfattelse af Livet, man danner sig, afhænger
af den Maade, paa hvilken Erkendelsen bestemmes i
sin Afhængighed af Livet. Denne Kredsgang udelukker ikke en
Fremskriden, --- naar blot den Kreds, vi her stadigt paany maa
beskrive, igen (som Epicykel) bevæger sig henad en større
Kreds, og denne maaske igen henad en endnu større, og saa
% --- p. 17 ---
\setcounter{footnote}{0}%
\opage{17}fremdeles. Vi staa her ved en Proces, der kan fortsættes i det
Uendelige, men hvis Lov vi kunne opdage allerede indenfor vor
Erfaring.

Selvom vi se bort fra dette Hovedproblem, kan der være
Vanskeligheder nok ved at forstaa Udviklingen fra et Livstrin
til et andet. Kontinuiteten kan ikke overalt paavises i det Enkelte.
Der er Punkter, hvor Forskningen atter og atter stanser.
Der er, saa langt Bevægelserne af Kloderne og af Atomerne
end forfølges tilbage, stedse en vis Retning af Bevægelsen og
en vis Fordeling af de virkende Kræfter, der ikke kunne udledes
af Bevægelsens almindelige Love. Ved Stoffernes Sammensætning
opstaa nye Stoffer med nye Egenskaber, der ikke
finde deres Forklaring ved den almindelige Lov om Materiens
Bestaaen under alle Forandringer; tvertimod forudsætter jo
netop denne Lov, at der sker Forandringer, det vil sige, at der
opstaar nye Egenskaber. Livets og Bevidsthedens Opstaaen staar
som en uløst Gaade, og lignende Gaader staa vi ved, naar de
store historiske Personligheders Optræden fører Menneskelivet
ind paa ny Veje. Her er Knuder, der idetmindste foreløbigt
ikke kunne løses. Betyder det, at Kontinuiteten virkeligt er brudt
paa disse Punkter, eller betyder det kun, at Sammenhængen er
saa indviklet, at vi ikke kunne udrede den, eller at der i Knuderne
er indflettet flere Traade end de, vi kende fra vor Erfaring?
--- Saa stor Vægt man end vil lægge paa saadanne
uløste eller uløselige Spørgsmaal, saa lader Forskningen sig dog
ikke stanse af dem. Hvor den ikke kan forklare, iagttager og
beskriver den. Og den glæder sig over Tilværelsens Fylde, der
lover os, at vort intellektuelle Arbejde ikke skal mangle Stof
og Opgaver, hverken paa det aandelige eller paa det materielle
Omraade. Vor Sans kan sløves, og vor Tanke blive træt, men
Tilværelsen er evigt ung, og selv naar det bliver Vinter for os,
kan der for den begynde et nyt Foraar.

\section*{V}

Der venter endnu Filosofen en stor Opgave. Som Logiker
opererer han med Begreberne Grund og Følge, som Psykolog
med Begreberne Aarsag og Virkning og som Etiker med Be%
% --- p. 18 ---
\setcounter{footnote}{0}%
\opage{18}greberne Formaal og Middel. Men tilsidst kunne disse forskellige
Synspunkter ikke falde udenfor hverandre. De maa alle
komme til deres Ret i vor afsluttende Opfattelse af Tilværelsen.
I de metafysiske Systemer har i Reglen et af disse Synspunkter
haft Overvægten. Man har opfattet Tilværelsen som en stor
Logik, eller som en stor Mekanisme, eller som et Værdiernes
Rige. Kunne nu alle disse Synspunkter forenes i én Ide?

Herom føres den store, maaske den evige Strid mellem
Verdensanskuelserne. Kun da har denne Strid filosofisk Interesse,
naar den føres med klare Tanker og paa en saadan Maade,
at der stadigt arbejdes paa et bestemt Grundlag af Erfaring og
videnskabelig Forstaaelse. Det er her, Videnskab, Filosofi og
Religion støde sammen. Det er Religionsfilosofiens Opgave at
undersøge denne Grænsestrids Betingelser og Betydning.

Et opmærksomt Blik paa Problemernes Udvikling gennem
Tiderne fører, efter min Opfattelse, til det Resultat, at det stunder
til Fald baade for den religiøse og den antireligiøse Dogmatisme.
Intet Dogme, det vil sige: ingen Forestilling, der direkte
eller indirekte er udsprungen af den religiøse Følelse,
forklarer Tilværelsen for os, hverken i det Hele eller i det
Enkelte. Der er Huller i vore Tankerækker og i vore Aarsagsrækker;
men de kunne ikke fyldes ved, at man appellerer til
Guders eller Sjæles Indgriben. Et saadant Hul er et uløst Problem,
hverken Mere eller Mindre. Men selvom Hullerne udfyldes,
saa kan, som vi have set ovenfor, Tilværelsens Indhold
ikke udtømmes i nogen Tankerække eller Aarsagsrække. Det
er sikkert et vigtigt Bidrag til Tilværelsens Karakteristik, at
dens Tilstande og dens Begivenheder i stort Omfang kunne beregnes
og forudsiges gennem de Love, vor Erkendelse kan finde.
Dette Moment betone de monistiske Systemer ganske særligt:
da al Forstaaelse beror paa, at der er Enhed og Sammenhæng
i Alt, hvad der sker, en Identitet bag alle Forskellighederne, saa
slutter man, at da Tilværelsen er forstaaelig i saa høj Grad,
som den er det, maa der ogsaa være Enhed, Sammenhæng,
Identitet i dens Væsen. Men da den Mulighed ikke er udelukket,
at Tilværelsen staar i et irrationelt Forhold til vor Er%
% --- p. 19 ---
\setcounter{footnote}{0}%
\opage{19}kendelse, kan ingen Monisme blive Mere end en spekulativ
Hypotese.

Den Enhed, som logisk Tænkning og Aarsagsforklaring
Spore i Tilværelsen, vil være den samme, hvorledes det saa
gaar med det, vi kalde Værdier. Den logiske Konsekvens og
den kausale Sammenhæng er den samme, hvad enten Udgangs- og Slutningspunkterne for de Tanke- og Aarsagsrækker, vi kunne
danne, ere værdifulde eller ikke. Forstaaelsen er forsaavidt uafhængig
af Vurderingen.

Men derved bortfalder ikke den Kendsgerning, at Menneskene
i deres personlige og praktiske Liv arbejde med Begreber som
Formaal og Værdi. Den Kendsgerning, at saadanne Begreber
dannes og maa dannes, maa tages med i Betragtning ved vor
afsluttende Betragtning af Tilværelsen. Det er her, ligesom ved
Erkendelsen, en Del af Tilværelsen, der træder i Forhold til
den øvrige Tilværelse. I den religiøse Følelse udtales en Dom
om Tilværelsen som Helhed betragtet. De Værdier, Mennesket
kender, have deres Udspring i virkelige Forhold. Deres Skæbne
er indvævet i Tilværelsens store Sammenhæng. Den religiøse
Følelse bestemmes ved det Forhold, Mennesket erfarer mellem
Værdierne og Virkeligheden, og den religiøse Tro er en fast
Forvisning om, at hvad der virkeligt har Værdi, ogsaa vil bestaa,
hvilken Bølgegang det saa maa gennemfare, og hvilke Forandringer
og Forskydninger det saa maa undergaa. Forskellen
mellem religiøse Standpunkter vil ifølge Sagens Natur dels kunne
bero paa, hvad det er, man tillægger Værdi, --- dels paa, hvorvidt
man kender Virkeligheden, --- dels paa, hvilket Forhold
man har erfaret mellem det Værdifulde, man har fundet, og
den Virkelighed, man kender. Her er Aarsager nok til Strid.
Og da baade Værdierne, Virkeligheden og Forholdet mellem dem
er i stadig Udvikling, vil Striden atter og atter kunne blusse
op i nye Former.

Religionsfilosofiens Opgave er det at udfinde, hvad Striden
egentligt gælder, og at undersøge den psykologiske Mulighed af
at forene Tro paa Værdiens Bestaaen med den videnskabelige
Forstaaelse af Tilværelsen og med det etiske Arbejde paa Værdiers
Opdagelse og Frembringelse. Filosofi kan ikke frembringe
% --- p. 20 ---
\setcounter{footnote}{0}%
\opage{20}Tro. Tro udtænkes ikke og konstrueres ikke, men voxer frem
af Livet. Tro er Genstand for Filosofien, ikke Frugt af Filosofien.
Filosofien undersøger dens psykologiske Mulighed og dens
etiske Betydning, saavel som dens erkendelsesteoretiske Berettigelse.
I denne sidste Henseende bliver det især Opgaven at
fastsætte Grænsen mellem Tænkning og Poesi.

Maaske er det Højeste kun tilgængeligt for os i Poesiens
Form, og en Livspoesi, i hvilken det personlige Liv fik Udtryk
for sine dybeste Erfaringer, vilde mere og mere vise sig at fyldestgøre
den Trang, der forhen kun kunde tilfredsstilles ved Myter
og Dogmer. Da vil Goethe faa Ret:

\begin{verse}
Und deines Geistes höchster Feuerflug \\
% CHECK p. 20: centred line: heading or verse \\
Hat schon am Gleichniss, hat am Bild genug!
\end{verse}

Skal Livet bevare sin Sundhed og Energi under den fremskridende
Deling af Arbejdet, saa maa de delte Kræfter kunne
forenes til Samvirken. Der maa ske nye Koncentrationer, hvis
Fortidens store Værdier skulle bevares i nye Former og under
de ændrede Vilkaar. Dertil kræves store, nydannende Personligheder,
i hvilke Digteren, Profeten og Tænkeren forenede arbejde
paa samme Værk. I Forhold til et saadant Værk er det
at filosofere kun et ensidigt Arbejde; men det er et Arbejde,
der skal og maa gøres, og som indirekte og ad mange Omveje
kan yde Bidrag til den store Koncentrationsproces, af hvilken
et nyt Livstrin fremgaar.

% --- p. 21 ---
\setcounter{footnote}{0}%
\opage{21}\essay[Psykologiens nuværende Stilling og dens Forhold til dens Nabovidenskaber. (1904)]{PSYKOLOGIENS NUVÆRENDE STILLING\\ OG DENS FORHOLD TIL DENS NABOVIDEN-\\ SKABER\footnote[1]{Et Foredrag paa den internationale videnskabelige Kongres i St. Louis. Den 23. September 1904.}.}
\streg

\section*{I}

I vore Dages Opfattelse af Videnskabens Væsen og Betydning
gør sig særligt den Overbevisning gældende, at vel kan
Virkelighedens store og komplexe Sammenhæng kun forstaas af
os, naar vi paa mere eller mindre kunstig Vis isolere enkelte
Elementer og undersøge deres indbyrdes Forhold, men at de
Elementer, vor Videnskab saaledes skelner imellem, ikke uden
videre tør betragtes som Virkelighedens egne, konstituerende Elementer.
Med andre Ord: Forstaaelsens og Virkelighedens Betingelser
falde ikke sammen. Vor Vej til Forstaaelse er ikke
nødvendigvis Naturens Vej til Frembringelse. Det er den kritiske
Filosofis gamle Grundtanke. Langsomt er den trængt igennem,
især i den nyeste Tid, ikke blot hos Filosofer, men ogsaa
hos Naturforskere, der have drøftet deres Videnskabs første Principer.
Der viser sig da Muligheden af et irrationelt Forhold
mellem Tanken og Realiteten, --- Muligheden af, at Tankens
Analyse ikke kan lade Realitetens store Syntese komme til sin
fulde Ret. Dog rokkes derved ikke ved Videnskabens Gyldighed,
ti de Analyser og Distinktioner, vi foretage for at naa viden%
% --- p. 22 ---
\setcounter{footnote}{0}%
\opage{22}skabelig Forstaaelse, maa Punkt for Punkt være begrundede ved
Iagttagelse af den levende og konkrete Realitet.

Paa intet Erfaringsomraade gør dette Forhold sig saa slaaende
gældende, som i Psykologien.  De Elementer, der paa det materielle
Omraade udsondres af Helheden, ligge dog i Rummet
udenfor hverandre, saa inderlig end Vexelvirkningen mellem
dem antages at være. Men et saadant ydre Sondringsforhold
kan ikke antages paa Sjælelivets Omraade. Her er det Enkelte
saaledes indvævet i det Hele, at det faar hele sin bestemte Karakter
derved, og Helheden kan her ikke betragtes som blot
Produkt af Elementerne. Analyse og Distinktion faa derfor i
Psykologien en mere kunstig Karakter end i Naturvidenskaben.
De psykiske Elementer tør vi endnu mindre tillægge absolut
Realitet end de materielle Atomer. Og dog gives der nu engang
ingen anden Vej til videnskabelig Forstaaelse af Sjælelivet
end Analyse paa Grundlag af Iagttagelse og Experiment.
Desuden viser Iagttagelse os, at Sjælelivet ved sin første Begyndelse
og paa sin Grænse mere og mere antager en sporadisk
Karakter, fremtræder glimtvis og som tilsyneladende isolerede
Fornemmelser. Derved træder det i stærk Modsætning til den
Maade, hvorpaa det ytrer sig, hvor det har naaet sin fulde Udvikling
og Modenhed; her har det afgjort Helhedens, Syntesens
Præg og fremviser en inderlig Sammenhæng med alle enkelte
Elementer. Der opstaar herved en for Psykologiens hele Karakter
betydningsfuld Antinomi. Ti vi formaa ikke at opfatte Sjælelivet
som opstaaet ved blot Kombination af de Elementer, mellem
hvilken Analysen sondrer, ejheller af de sporadiske Glimt,
til hvilke det synes at være reduceret ved sine Grænser; og
dog maa der være en indre Sammenhæng mellem den konkrete
Bevidsthed med dens syntetiske Helhedspræg og de Elementer,
Analyse og Grænseiagttagelse føre tilbage til.

Denne Antinomi har grebet dybt ind i den nyere Psykologis
Udviklingsgang. Den fremtræder særligt i Brydningen mellem
de to store Retninger, den ene grundlagt af Hartley og
Hume og fortsat gennem Associationspsykologien og den Herbartske
Skole, den anden grundlagt af Leibniz og Kant og fortsat
gennem den tyske romantiske Skole og dens Eftervirknin%
% --- p. 23 ---
\setcounter{footnote}{0}%
\opage{23}ger. Hin Retning fører til en psykisk Atomistik, medens denne
fastholder Sjælelivets syntetiske Karakter. Jeg skal ikke her forfølge
denne Brydnings Historie i det Enkelte. Der gives ingen
Psykolog, hvis almindelige Standpunkt og specielle Synsmaader
ikke ere bestemte ved den Stilling, han indtager til Forholdet
mellem Atomistik og konkret Realitet paa Sjælelivets Omraade.
Paa begge Sider er der Fristelse til Dogmatisme, idet man enten
ensidigt slaar det enkelte Synspunkt fast som udtømmende, eller
statuerer en uforsonlig Modsætning mellem Synspunkterne.

Der gives en psykologisk Atomistik, der betragter den
psykologiske Analyses Elementer som absolute Bestanddele, hvis
mekaniske Sammensætning frembringer Sjælelivet som sit Produkt.
Den overser, at hele det psykologiske Problem træder i
Kraft igen ved ethvert af de psykiske Atomer. Ti ligesom det
i Naturvidenskaben har vist sig, at de Atomer, ved hvilke man
troede at kunne stanse, vare hele Verdener, i hvis Indre der
foregik Bevægelser og Strømninger, som gjorde nye Undersøgelser
nødvendige, saaledes viser nærmere Betragtning af de Forhold,
under hvilke vore simpleste Sansefornemmelser optræde,
at de selv maa opfattes som syntetiske Fænomener, som konkrete
Helheder, idet de svare til en mere eller mindre komplex
fysiologisk Proces.

Der gives paa den anden Side en Dogmatisme, for hvilken
Sjælelivet udgør en udelelig Enhed, maaske endog en substantiel
Enhed, der skal trodse al Analyse. Den overser, at Sjælelivet
som alt Liv bestaar under en stadig Kamp med Modsætninger,
indre og ydre, og at der allerede for den praktiske Selviagttagelse
frembyder sig afgørende Forskelle, f. Ex. ved ethvert
Valg mellem Muligheder. De forskellige kæmpende Dispositioner
eller Tendenser kunne kundgøre sig i Bevidstheden paa
eftertrykkelig Maade. Sjælelivet bestaar under en stadig Stræben
efter at harmonisere de stridende Elementer, og Syntesens
Energi fremtræder i mange forskellige Grader af Intensitet. Det
er altsaa ikke først den videnskabelige Iagttagelse og Analyse,
der frembringer Forskellighederne, men disse høre selve Livet
til. Som den videnskabelige Analyse atter og atter ved sin Ensidighed
nøder os til at gaa tilbage til Livets store Syntese,
% --- p. 24 ---
\setcounter{footnote}{0}%
\opage{24}saaledes nøder atter Livets indre Spændingsforhold os til at anerkende
Forskellenes og Modsætningernes Realitet og godkender
forsaavidt Betydningen af den videnskabelige Analyse. Vi kunne
kun analysere, hvad Livet forud har forbundet; men denne Forbindelse
udelukker ikke Forskelsforhold mellem de Elementer,
der ere forbundne.

I den amerikanske Literatur er Forholdet mellem Analyse
og Syntese i Psykologien energisk blevet drøftet af to saa fremragende
Forskere som William James og Hugo Münsterberg.
Maaske er dog især den sidstnævnte Forsker bleven ført til at
hævde en skarpere Modsætning mellem Psykologien og Livet,
end det konsekvent er muligt at fastholde, naar det dog stedse
er Livet, der skal yde Psykologien dens Materiale, og naar det
staar fast, at Selviagttagelse ikke først begynder med videnskabelig
Psykologi, men er en praktisk Fornødenhed, der forudsætter
Muligheden af, at virkelige Forskelligheder existere
indenfor Livets Helhed. ---

Modsætningen melllem Analyse og Syntese paa det psykologiske
Omraade staar i en vis Forbindelse med Modsætningen
mellem Intellektualisme og Voluntarisme. Ti det er indenfor
Fornemmelsernes og Forestillingernes Omraade, at man lettest
kan skelne forskellige Elementer. Det er den mest artikulerede
Side af Sjælelivet, og det er tillige fra denne Side, at Sjælelivet
ligger mest aabent for Iagttagelse og Experiment. I Modsætning
til det frembyder Følelses- og Villieslivet en stærkere
Koncentration; den syntetiske Karakter træder her mere udpræget
frem. Overvurderingen af den analytiske Metodes Resultater
vil let føre til en Undervurdering af Følelses- og Villieslivets
Betydning i Virkeligheden, maaske endog til et Forsøg
paa at betragte Følelses- og Villieslivet som blot Resultat af
Fornemmelsens og Forestillingens Vexelspil. Herimod strider, at
det ikke blot er umuligt at udlede Koncentrationen af absolute 
Elementers Vexelvirkning, men at ved nærmere Undersøgelse
hele Retningen af Fornemmelsernes og Forestillingernes Udvikling
bestemmes ved de Interesser, Værdier og Formaal, Følelsesog
Villieslivet vidne om. Den fra Fichte og Schopenhauer stammende
Voluntarisme, der har fundet en vigtig Støtte i den bio%
% --- p. 25 ---
\setcounter{footnote}{0}%
\opage{25}logiske Teori om Kampen for Livet, vinder ogsaa mere og mere
Tilslutning som det væsentligste psykologiske Synspunkt.

\section*{II}

Baade ved at fremhæve det stedse inkommensurable Forhold
mellem Analyse og Syntese og ved at betone Voluntarismen
i Modsætning til Intellektualismen, forminske vi Tilbøjeligheden
til at antage, at Personlighedens Begreb en Gang for
alle skulde kunne afsluttes, saaledes som den spekulative Filosofi
og Teologien ofte have trot. Positivismen og Erfaringsvidenskaben
faa ofte Skyld for at fornegte Personlighedsbegrebet, og
man stiller i den nyeste Tid ofte historisk Opfattelse og Frihedslære
i Modsætning til Erfaringspsykologien. Jeg tror, at
netop den empiriske, experimentale og analytiske Retning i Psykologien
hviler paa en alvorlig og oprigtig Anerkendelse af
det personlige Livs Realitet. Den hviler paa den Overbevisning,
at Aandslivets Værdi ikke forringes ved at være knyttet til bestemte
Betingelser og underkastede bestemte Love. Derfor undersøger den
frejdigt Aandslivet ad alle de Veje, der staa aabne for Videnskaben.

Forskellen mellem de forskellige psykologiske Skoler beror
egentligt paa, hvor man finder Problemet, og hvorledes Bevisbyrden
fordeles. Er den psykologiske Gaade den, hvorledes der
kan være Enhed og Sammenhæng i Bevidstheden, eller opstaar
det Gaadefulde dér, hvor Bevidsthedslivet antager en sporadisk
Karakter, fremtræder i isolerede Glimt? Dette er Hovedspørgsmaalet.
Det deler sig naturligvis i mange enkelte Spørgsmaal.
Fordi jeg er overbevist om, at Bevidsthedens Natur ytrer sig i
en indre Sammenhæng og Sammenfatten, stiller der sig alligevel
Opgaver nok: nemlig at finde de specielle Former for Sammenfatten,
der ingenlunde kunne udledes a priori af Bevidsthedens
almindelige Natur, --- og desuden at forklare, hvorledes
Bevidsthedslivet under visse Forhold kan antage en sporadisk
Karakter. Paa den anden Side er der ogsaa mange Arter og
Grader af Isolering, --- og desuden skal det forklares, hvorledes
der kan fremtræde Enhed og Sammenhæng i Bevidstheden. Og
enhver Psykolog maa indrømme, at der saalænge Bevidsthedslivet
bestaar, finder en stadig Kamp Sted mellem den syntetiske
% --- p. 26 ---
\setcounter{footnote}{0}%
\opage{26}og den sporadiske Tendens. Hvor Helhedspræget overvejer,
bliver det da et Spørgsmaal, hvorvidt Helheden er et rent mekanisk
Resultat eller Aggregat, eller om den ikke hviler paa
dybereliggende Betingelser.

Den patologiske Psykologi synes afgjort at tale for den
syntetiske Opfattelse. Der kræves et stadigt Arbejde, uden hvilket
den »psykologiske Spænding« (tension psychologique, Pierre
Janet) ikke kan udholdes. Sjælssygdom viser sig i, at denne
Spænding, uden hvilken Bevidsthedens Enhed ikke kan bestaa
sammen med et rigt Indhold af kvalitativt forskellige Elementer,
kun med stor og pinlig Anstrengelse kan udholdes, og at
den kræver særligt stærke Paavirkninger og Foranledninger for
at Sprængning eller Slaphed ikke skulle indtræde. Det Isolerede
og det Sporadiske i Sjælelivet stiller stedse en Opgave, der skal
løses. Det fremtræder i Spireformerne og i Opløsningsformerne,
hvor Opgaven ofte ikke kan løses.

Begrebet psykisk Energi kunne vi bestemme saaledes. Der
er saa meget des mere psykisk Arbejde at gøre, jo flere Elementer
eller Tendenser der skulle forbindes i samme Tilstand,
--- jo mere forskelligartede de ere, --- jo stærkere de hver for
sig ere, --- jo inderligere de skulle forbindes, --- og jo fjernere
i Tiden de ligge paa hverandre og fra det nærværende Øjeblik.
Det vil i det enkelte Tilfælde ikke være let at anvende det
Begreb om psykisk Energi, der her viser sig muligt. Kun ved
omhyggelig Iagttagelse og alsidigt Kendskab til de specielle bestemte
og individuelle Forhold kan man give en Bestemmelse
af disse fem Faktorer. Elementernes eller Tendensernes Antal,
Graden af deres Forskellighed, Styrken af hvert enkelt af dem,
Forbindelsens Grad og Tidsafstandens Størrelse --- alt dette er
vanskeligt at konstatere, og det varierer fra Tilfælde til Tilfælde.
Og dertil kommer endnu den Hastighed, hvormed Forbindelsen af
Elementer eller Tendenser foregaar. Hine fem Faktorer svare til
Massen ved Bestemmelsen af fysisk Energi.

Der viser sig her en uudtømmelig Rigdom af Muligheder
paa Grund af de højst forskellige Forhold, der kunne være til
Stede mellem disse fem Omstændigheder. Der er her et stort
Omraade for Iagttagelse, Experiment og Sammenligning. En sam%
% --- p. 27 ---
\setcounter{footnote}{0}%
\opage{27}menlignende Individualitetspsykologi er endnu i sin første Begyndelse.
Kun paa Religionspsykologiens Omraade er der, særligt
her i Amerika, gjort Begyndelser i denne Retning. Tilsidst
vil ingen Teori kunne give en udtømmende Fremstilling af,
hvorledes de forskellige Elementer eller Tendenser virke sammen
i en enkelt Tilstand eller i en enkelt Personlighed. Som overalt
saaledes er ogsaa her den fuldendte Individualisering ikke
Videnskabens, men Kunstens Sag. Kun Kunsten kan tilsidst give
en Syntese, der nogenlunde lader Livets store og dybe Syntese
komme til sin Ret.

\section*{III}

Nye Problemer møde os, naar vi forsøge at karakterisere
Psykologiens Forhold til dens Nabovidenskaber. Psykologien
staar i et ejendommeligt Forhold til Naturvidenskaberne paa
den ene Side, til Aandsvidenskaberne paa den anden Side. Og Forholdet
kan korteligt angives saaledes, at den i Sammenligning
med Naturvidenskaberne har en afgjort syntetisk Karakter, medens
den i Sammenligning med Aandsvidenskaberne har en afgjort
analytisk Karakter. Derved bestemmes her de Problemer, som
overalt fremkomme, hvor forskellige Videnskaber mødes.

Som allerede omtalt, svare de simpleste psykiske Elementer,
vi kunne skelne, til komplicerede fysiologiske Processer.
Hvad der psykologisk staar som enkelt, er fysiologisk en Mangfoldighed.
Der maa altsaa i det enkelte psykiske Element være
sammenfattet, hvad der fysiologisk udfylder flere Øjeblikke og
breder sig over en hel Region. Dertil kommer en anden Omstændighed.
De psykiske Elementer ere kvalitativt forskellige
indbyrdes, medens der er Grund til at tro, at de tilsvarende
Hjerneprocesser kun frembyde Forskelle i Grad, Retning og
Kombination. Hvad der psykologisk fremtræder som kvalitative
Forskelligheder, maa da fra Naturvidenskabens Synspunkt opfattes
som kvantitativt forskelligt. Det bliver fra det naturvidenskabelige
Synspunkt lettere at paavise Kontinuitet, end fra det
psykologiske Synspunkt. Den gamle Sætning; at Naturen ingen
Spring gør, formaar Psykologien ikke at gennemføre saa fuldstændigt,
som Naturvidenskaben formaar det.

Af disse Omstændigheder have nogle Forskere ment at
% --- p. 28 ---
\setcounter{footnote}{0}%
\opage{28}maatte slutte, at en Videnskab om det Psykiske kun var mulig,
hvis man i Stedet for Forholdet mellem psykiske Tilstande
kunde sætte et Forhold mellem de tilsvarende Hjernetilstande.
Psykologien skulde da, for at blive Videnskab, forvandles til
Fysiologi. Ellers vilde det --- mener man --- være umuligt at
nærme sig til Idealet for videnskabelig Forstaaelse, nemlig Paavisning
af Kontinuitet og af Ækvivalensforhold mellem Fænomenerne.
Overfor denne Synsmaade maa jeg for det Første bemærke,
at vi paa alle Omraader begynde med at paavise Aarsagsforhold
mellem kvalitativt forskellige Fænomener; det bliver
saa senere Undersøgelsers Opgave at lade dette elementære Aarsagsforhold
afløse af et fuldkomnere Aarsagsforhold, hvis Led
staa i Kontinuitets- og Ækvivalensforhold. Nu er det vel saa,
at vi paa det psykologiske Omraade neppe kunne komme ud
over det elementære Aarsagsforhold, fordi vi ikke kunne finde
psykiske Enere, altsaa ingen kvantitativ Metode kunne gennemføre.
Men det udelukker ikke Berettigelsen af overhovedet at
antage et Aarsagsforhold mellem psykiske Tilstande. Og dette
bliver desuden nødvendigt: ti er der Aarsagsforhold mellem de
tilsvarende Hjerneprocesser, maa der ogsaa, idetmindste middelbart,
være et saadant mellem de psykiske Tilstande. For det
Andet er det foreløbigt kun rent skematiske Konstruktioner, man
formaar at give af de til psykiske Tilstande svarende Hjerneprocesser,
Konstruktioner, ved hvilke Analogien med de umiddelbart
iagttagne og analyserede psykiske Tilstande er ledende.
Vi slutte fra de psykiske Tilstandes Forskelligheder til Forskelligheder
mellem de tilsvarende fysiologiske Processer. Og denne
Slutning er ikke altid let at drage; desuden forudsætter den
nødvendigvis, at den psykologiske Iagttagelse og Analyse er fuldkommen.
Psykologiens Selvstændighed forudsættes altsaa.

Muligvis hænger baade Enkeltheden og den kvalitative Karakter
af de psykiske Elementer sammen med, at de selv ere
Frugter af en skjult Syntese, saa at hvis vi kunde gaa længere
ned paa de psykiske Differentialers Omraade, gaa ned til Differentialer
af anden eller tredie Grad, vilde hele Problemet staa
klarere for os. Men vi støde dog her tilsidst paa det store Problem
om Forholdet mellem det Psykiske og det Fysiske. Ogsaa
% --- p. 29 ---
\setcounter{footnote}{0}%
\opage{29}her bliver Forholdet mellem Analyse og Syntese af afgørende
Betydning. Ti den Forskel, som gøres mellem det Psykiske og
det Fysiske, skyldes netop en Analyse, en Sondren mellem Elementer,
Realiteten stedse frembyder i Sammenhæng. Vi foretage
her en Abstraktion; vi sønderbryde Realitetens Helhed og
undre os saa bagefter over, hvor vanskeligt det er igen at forene
de Stykker, i hvilke vi have delt den. Dette Punkt er især
med stor Styrke blevet belyst af Wilhelm Wundt og Roberto Ardigò. Realiteten staar stedse som den store Grundsyntese,
indenfor hvilken vi bevæge os med alle vore Analyser og Abstraktioner,
--- som en fuldtonende Helhedsklang overfor vor
møjsommelige Staven. Men der gives nu som sagt ingen anden
Vej til Erkendelse end den, der begynder med Analysen. Vor
Opmærksomhed gaar fra Punkt til Punkt og søger saa først bagefter
at sammenarbejde sine enkelte Resultater. Og ligesaa lidt
som vi tør tillægge vore Distinktioner absolut Gyldighed, ligesaa
lidt tør vi betragte det som en Tilfældighed, at vor Stræben
efter Erkendelse netop nøder os til at foretage disse bestemte
Distinktioner. Det er et Bidrag til Realitetens Karakteristik, at
den kun kan forstaas, naar der foretages omhyggelige Analyser
af dens Indhold.

Jeg skal ikke her komme ind paa en udførlig Drøftelse af
Hypoteserne om Forholdet mellem det Psykiske og det Fysiske,
men kun bemærke, at da Hjerneprocesser ikke kunne være andre
videnskabelige Synspunkter og Metoder underlagte end andre
materielle Fænomener, paahviler Bevisbyrden den, der vil hævde
en Indgriben af Sjælen i de fysiologiske Processers Gang. En
saadan Antagelse vilde ophæve Fysiologiens Selvstændighed, ligesom
Negtelsen af et Kausalitetsforhold mellem Sjæletilstande
vilde ophæve Psykologiens Selvstændighed. Der er ingen Grund
til at afvige fra den Grundsætning, som Naturvidenskaben har
fulgt i Aarhundreder, og som den skylder sine Triumfer: at
materielle Fænomener skulle forklares ved materielle Aarsager.
Spinozas Ord: »At sige, at Sjælen sætter Legemet i Bevægelse,
er i Grunden at sige, at man ikke véd, hvorledes denne Bevægelse
bliver til« --- staa ved fuld Magt den Dag idag. Den
eneste Arbejdshypotese, der gør et Samarbejde mellem Fysio%
% --- p. 30 ---
\setcounter{footnote}{0}%
\opage{30}logi og Psykologi uden Overgreb fra nogen af Siderne mulig, 
er den, som baade indenfor det Psykiske og indenfor det Fysiske
 fastholder saa stor Kontinuitet som muligt, og som opfatter
 selve Forholdet mellem de to Rækker af Fænomener som 
et Funktionsforhold i Ordets matematiske Betydning. En afsluttende
 metafysisk Tydning staar da stedse aaben, men vedkommer
 ikke Psykologien som saadan. Den Kritik, Parallelismen,
 eller som jeg hellere kalder den, Identitetshypotesen, 
har været Genstand for, har væsentligt angrebet den som metafysisk
 Hypotese. Men den er først og fremmest en Arbejdshypotese,
 og den eneste, der faktisk følges og kan følges paa Videnskabens
 nuværende Stadium. Det gælder ganske vist om
vore Arbejdshypoteser, hvad jeg ovenfor bemærkede om vore
Analyser og Distinktioner: vi have ingen Ret til at antage det
for en ren Tilfældighed, at vor Forstaaelse af Tilværelsen netop
gør disse bestemte Arbejdshypoteser nødvendige. En Metafysik,
der vilde konstruere en Verdensanskuelse uden nogetsomhelst
Hensyn til de Arbejdshypoteser, der have vist sig nødvendige,
vilde ikke have nogen filosofisk Værdi.

\section*{IV}

Medens Psykologien i Modsætning til Naturvidenskaben
staar med en afgjort syntetisk Karakter, staar den i Modsætning
til Aandsvidenskaberne med en afgjort analytisk Karakter.
Aandsvidenskaberne handle om menneskelige Værker eller menneskelige
Idealer, men Psykologien handler om det menneskelige
Sjælelivs Elementer og almindelige Love. Kalde vi Læren
om de menneskelige Værker Kulturvidenskab, Læren om de
menneskelige Idealer Etik, saa er Psykologiens Forhold til Kulturvidenskab
og Etik bestemt ved Forholdet mellem Elementer,
Værker og Idealer. Alle tre Undersøgelser skride frem jevnsides,
og de øse af samme dybe Kilde: det umiddelbare og uvilkaarlige
Aandsliv, det virkelige, konkrete Liv, der leves, og som
ingen Analyse vil kunne udtømme, saalidt som det nogensinde
fuldstændigt vil have udtrykt sig i noget Værk eller i noget
Ideal. Al Forsken staar her overfor det uendeligt Konkrete og
forsøger fra forskellige Sider at karakterisere dets Væsen, at
% --- p. 31 ---
\setcounter{footnote}{0}%
\opage{31}udtrykke dets Fylde i bestemte Former. Men Livets Toner ere
saa mangfoldige og klinge saa inderligt sammen, at ingen Videnskabs
Nodesystem fuldt ud kan udtrykke dem. Det gælder Kulturhistorikeren
og Etikeren saavel som Psykologen. Men indenfor
denne beslægtede Stilling er der igen et Vexelforhold mellem Psykologien og Aandsvidenskaberne. For at finde Sjælelivets
Elementer og Love er det ikke nok at studere det enkelte
Individ i dets enkelte Tilstande. Dertil kræves ogsaa Studium
af de menneskelige Værker og Idealer, i hvilke Sjælelivets Væsen
aabenbarer sig gennem Tiderne. Der existerer intet Sjæleliv i
al Almindelighed; det fremtræder i forskellige Former til forskellige
Tider og paa forskellige Steder, og i hver enkelt af disse
historiske Former ytrer det sig saavidt muligt efter sit fulde
Omfang, skønt Sammenspillet og dets Elementer faa forskellig
Klangfarve. Her har Psykologien en betydningsfuld Gruppe af
Fakta at analysere. Den sociologiske Metode træder her ved
Siden af den introspektive, den experimentale og den fysiologiske.
Dens Endemaal har en mere abstrakt Karakter end Kulturvidenskabens
og Etikens, idet Elementet er mere abstrakt end
Værket eller Idealet. Der bliver her Brug for den af Comte
anvendte og af Mill beskrevne omvendt deduktive Metode, idet
der først konstateres en Udviklingsgang, der har ført til visse
Værkers eller Idealers Opstaaen, og derefter forsøges en Udledning
og Forklaring af denne Udviklingsgang fra almindelige
psykologiske Elementer og Love. Ved ren Deduktion vil man
aldrig kunne udlede de specielle Aandsvidenskaber af den almindelige
Psykologi; en Reduktion, ikke en Deduktion er det
Højeste, vi kunne vente at naa. Vi kunne f. Ex. ikke deducere
Pædagogik, Æstetik og Etik ud fra den almindelige Psykologi;
men ved Iagttagelse af, hvorledes Livets uvilkaarlige Kunst udfolder
sig i Opdragelse, æstetisk Frembringen og etisk Handlen
kunne vi finde raadende Synspunkter og Ideer, der igen kunne
forstaas efter almindelige psykologiske Love. Saaledes er ogsaa
Psykologiens Forhold til Erkendelsesteorien og til Religionsvidenskaben.
Angaaende dette sidste Forhold har jeg i min Religionsfilosofi
udtalt mig saaledes: »Paa det religiøse Omraade have
Menneskene gjort nogle af deres dybeste, mest intensive sjæle%
% --- p. 32 ---
\setcounter{footnote}{0}%
\opage{32}lige Erfaringer. I Religionen (hvis den er ægte og oprindelig) virke alle Sjælelivets Elementer med en Energi og et Sammenspil som neppe paa noget andet Omraade. Derfor er Studiet af det religiøse Sjæleliv af saa stor Betydning for Psykologien i det Hele«. Et Vexelforhold vil der sikkert her mere og mere kunne komme i Stand, saaledes at Indsigten i de psykiske Elementers Natur og i Lovene for deres Vexelvirkning vil
virke vejledende og klarende ved de specielle Aandsvidenskabers Arbejde. En nærmere Betragtning af Kulturvidenskaberne og Etikens Historie vil ogsaa vise, at direkte eller indirekte
gaar der stadigt en saadan Indvirken. Hvis Psykologien har en Fremtid, vil denne Indvirken kunne blive endnu betydningsfuldere, end den hidtil har været. Psykologien staar i en stor Gæld til sine Nabovidenskaber, til Aandsvidenskaben saavel som til Naturvidenskaben. Lad os haabe, at den vil formaa at betale paa denne Gæld, selvom denne --- hvis de andre Videnskaber vedblive at skride fremad --- forhaabentligt stedse vil blive stiftet paany!
% --- p. 33 ---
\setcounter{footnote}{0}%
\opage{33}\essay[Begrebet Analogi i dets filosofiske Betydning. (1904)]{BEGREBET ANALOGI OG DETS FILOSOFISKE\\ BETYDNING\footnote[1]{Et Foredrag i Jowett-Society i Oxford. Den 26. November 1904.}.}
\streg

\section*{I}

Naar man følger den Retning, Drøftelsen af Videnskabens
første Forudsætninger har taget i de senere Aar, kan man ikke
undgaa at bemærke, at det sidste Kriterium for deres Gyldighed
mere og mere findes i deres Anvendelighed. Berettigelsen
af at opstille en Grundsætning begrundes ved, at den gør det
muligt at vinde Klarhed og finde Sammenhæng i Iagttagelser,
der ellers vilde staa dunkelt og spredt for os. Grundsætningernes
Sandhed bestaar da ikke i deres Overensstemmelse med en
absolut, i og for sig bestaaende Tingenes Orden. En Tingenes
Orden bliver først til for os, naar vi ved Grundsætningernes Hjælp
bringe Sammenhæng mellem Iagttagelserne. Vi frembringe altsaa
selv Sandheden, naar vi kunne finde de Grundsætninger, der i
største Omfang og i inderligste Grad kunne forbinde vore Iagttagelser.
Et kritisk eller dynamisk Sandhedsbegreb er saaledes
ifærd med at arbejde sig frem, i Modsætning til det dogmatiske
Sandhedsbegreb, der kan kaldes statisk, da det forudsætter en
hvilende, bestaaende Tingenes Orden, som Tanken blot havde at
gengive. Dette Sandhedsbegreb fordredes allerede af den kritiske
Filosofi, naar den lærte, at objektiv Gyldighed bestaar i den
fulde, lovmæssige Sammenhæng mellem Iagttagelserne. Men det
% --- p. 34 ---
\setcounter{footnote}{0}%
\opage{34}er af ikke ringe Betydning, at ikke blot Filosofer, men ogsaa
Naturforskere, som Clerk Maxwell, Ernst Mach og Heinrich
Hertz, mere og mere gaa i denne Retning. Det er ikke blot
Dogmatismen, men ogsaa Skepticismen, som derved fortrænges.
Ti Skepticismen opstaar ved, at man forudsætter det dogmatiske
eller statiske Sandhedsbegreb og saa viser, at det ikke kan ske
Fyldest.

Det gælder da om at opstille saadanne Grundsætninger, der
gøre det muligt at naa det Maal, man har sat sig, nemlig at
finde den størst mulige Sammenhæng indenfor det størst mulige
Antal Iagttagelser. Grundsætningerne behøve ikke at have andre
Egenskaber end netop dem, der i denne Henseende ere nødvendige.
Og navnligt behøve de ikke, hvad de Forestillinger,
der indeholdes i dem, angaar, at udtrykke eller gengive det, der
i sig selv eller i sidste Grund maatte være indeholdt i de givne
Iagttagelser. Deres Opgave er at afgive Synspunkter for disses
grundige Ordning, og denne Opgave kunne de løse, selvom de
benytte billedlige Udtryk, naar blot den indre Sammenhæng
mellem de Billeder, der anvendes, er en saadan, at Iagttagelsernes
Sammenhæng kan udtrykkes og forudsiges ved dens
Hjælp. Det er altsaa ikke et Identitetsforhold, men et Analogiforhold
mellem Grundsætninger (Principer, Ideer) og Iagttagelser
(Erfaringer, Fænomener), der kræves, for at Sandhed
kan vindes. Analogi er Lighed i Forholdene indenfor forskellige
Genstande, ikke Lighed i de enkelte Egenskaber. Det
kritiske eller dynamiske Sandhedsbegreb bliver altsaa tillige symbolsk.
Det kræver et System af Billeder eller Tegn, ved hvis
Hjælp vi kunne aflæse Forholdet mellem Iagttagelserne og ledes
til nye Iagttagelser. Vi betragte de Begivenheder, Erfaringen
fremstiller for os, \textit{som om} de blot vare Ændringer i vore billedlige
Forestillinger. Det vil da være naturligt at vælge den
simpleste og klareste Billedkreds, som kan fyldestgøre Opgaven.
Derved kommer man til, hvad Renaissancetidens Tænkere kaldte
Simpelhedens Princip, og hvad man i vor Tid har kaldt Tænken
af Verden efter det mindste Kraftmaals Princip. 

At vi ikke have Ret til at tillægge vor Erkendelse mere
% --- p. 35 ---
\setcounter{footnote}{0}%
\opage{35}end symbolsk Gyldighed, følger først og fremmest af, at vi ikke
have Ret til at tillægge den Egenskaber, som ikke ere nødvendige,
for at den skal kunne naa det Maal, den har sat sig. Naar
Analogi er tilstrækkelig dertil, har man ikke Lov at forudsætte
en Identitet, som desuden ikke kan bevises. Den Mulighed er
nemlig aldrig udelukket, at andre Grundsætninger (Principer,
Ideer) end de, vi anvende, kunne løse den samme Opgave. Overgangen
fra Grund til Følge er altid sikrere end den fra Følge
til Grund, fordi forskellige Forudsætninger kunne føre til samme
Slutning. Et og samme Maal kan naas ad forskellige Veje, og
at vor Videnskab har fundet den absolut eneste Vej til Ordning
og Sammenhæng i Erfaringens Verden, vil aldrig nogensinde
kunne bevises.

Men ogsaa af andre Grunde maa vor Erkendelse blive symbolsk,
naar den væsentligt opererer med Analogier.

Det ligger i vor Aands Natur, at vi ikke kunne tænke uden
Billeder. Hvad vi i speciel og streng Betydning kalde Tænkning,
er stedse den Virksomhed, ved hvilken Forholdet mellem
forskellige Billeder bliver saa klart og utvetydigt som muligt.
Enhver Tanke, den være saa abstrakt og ophøjet, den være vil,
stræbe vi uvilkaarligt at anskueliggøre ved at fremdrage Billeder,
i hvis indbyrdes Forhold den lægger sig for Dagen. Billedet
kan snart være fyldigt og konkret, snart et blot Skema eller
et Tegn. Det er ikke blot i stille Selvtænkning, vi gaa saaledes
frem, men ogsaa, og ganske særligt, hvor det gælder om at
fremstille det Tænkte for Andre. Vi tænke i Exempler, selvom
vi reducere disse Exempler til ganske simple Skemaer eller Tegn,
og endnu tydeligere er det, at vi maa benytte Exempler og
Tegn, naar det Tænkte skal meddeles Andre.

Vor Tænkning er stedse en Tolkning af noget Givet. Den
er en Virksomhed, der øves paa Grundlag af Modtagen. Modsætningen
mellem Virken og Modtagen, mellem Aktivitet og
Passivitet er ganske vist relativ, beror paa Synspunktet og paa
Udviklingens Trin. Hvad der fra et Synspunkt er passivt, kan
fra et andet Synspunkt og i Forhold til et lavere Udviklingstrin
være at betragte som aktivt. Men i hvert enkelt Tilfælde,
i hver enkelt Akt foretager Tænkningen en Omsætning af
% --- p. 36 ---
\setcounter{footnote}{0}%
\opage{36}et Modtaget til sin egen Form, og dette Omsætningsarbejde kan
ikke bevises at være af rent formel Betydning, saaledes at der
skulde være absolut Identitet mellem det Givne og selve Tankens
Elementer. Naar Tanken f. Ex. analyserer et materielt
eller et aandeligt Fænomen, kan der ikke føres noget Bevis for,
at de Elementer, den for at klare sig det givne Indhold skelner
imellem, ogsaa existere som selvstændige Dele i det Givne.
Det er tilstrækkeligt til Tænkningens Gyldighed, at Forholdet
mellem Tankeelementerne \textit{stemmer med} Forholdet mellem Erfaringens
Elementer.

Men ikke blot bestaar selve Tænkningen i en Omsætning
fra passiv til aktiv Form. I hvert enkelt Tilfælde vil det desuden
vise sig, at Tanken arbejder ved at sammenholde forskellige
Erfaringsomraader og lade det ene belyse det andet. Hele
vor Erkendelse, baade den uvilkaarligt udviklede og den videnskabelige,
er derfor gennemtrængt af Analogier. Ogsaa dette
hænger sammen \textit{med}, at Tankearbejdet er Løsning af en Opgave,
Forsøg paa at virkeliggøre et Formaal. Naar Tænkningen
gaar frem til nye Opgaver, vælger den ikke helt nye Midler og
Veje, men søger saavidt muligt at hjælpe sig med de tidligere
anvendte, især naar de have det Fortrin at være klare og tydelige.
Saaledes har Sproget ikke dannet helt nye Udtryk for
aandelige Fænomener, men hjulpet sig ved Overførelse af Udtryk,
der oprindeligt vare dannede for de materielle Fænomener,
der først havde tiltrukket sig Opmærksomheden. Og da der nu
idethele fremtræder afgørende Forskelligheder mellem Erfaringsomraaderne,
--- da Kendsgerningerne ikke ere ensartede, men
danne forskellige Grupper, hver med sine Ejendommeligheder,
saa er Tænkningen ogsaa nødt til at gaa saaledes frem, at den
benytter et Omraade, en Gruppe til at belyse en anden, og da
særligt saaledes, at det Omraade, der ordner sig for Tanken paa
den simpleste og tillige paa den frugtbareste Maade, lægges til
Grund ved Forstaaelsen af det andet. Dersom Tilværelsen ikke
frembød kvalitative Forskelligheder, vilde der ikke være nogen
Nødvendighed herfor. Men da Tilværelsens forskellige Dele, saaledes
som vi kende dem gennem Erfaringen, ikke ere ensartede,
bliver Analogien en nødvendig Vej til Forstaaelse.

% --- p. 37 ---
\setcounter{footnote}{0}%
\opage{37}Jeg skal nu forsøge at belyse Analogiens Karakter og Betydning
paa de forskellige Tankeomraader, hvor den anvendes.

\section*{II}

Enhver Dannelse af et Almenbegreb sker ved Hjælp af
Analogi. De Kendemærker, Almenbegrebet indeholder, variere i
de forskellige Fænomener, for hvilke \textit{det} skal gælde, og Almenbegrebet
bliver derfor kun muligt ved, at Egenskabernes Forhold
stemme overens, selvom de hver for sig ere forskellige.
Almenbegrebet »Hest« bliver f. Ex. kun muligt ved, at den
brune Farve staar i \textit{samme} Forhold til en Hest som den graa
Farve til en anden Hest. De menneskelige Egenskaber variere
fra Menneske til Menneske, men det Konstante er et vist Forhold
mellem disse Egenskaber, og derved bliver Almenbegrebet
»Menneske« muligt. Vi kunne ikke danne noget Begreb, hvis
Kendemærker paa aldeles identisk Maade forekomme i de enkelte
Fænomener. Begrebet dannes ved Muligheden af stadig Overgang
fra Exempel til Exempel gennem Analogi.

Dette har en særegen Interesse, da man fra Aristoteles
Tid af har stillet Analogi i Modsætning til Begreb. For Aristoteles
var Analogien det Baand mellem forskellige Omraader, der
endnu kunde gælde, efter at de mest omfattende Begreber vare
dannede. Slægtskabet ved Analogi var altsaa ikke noget begrebsmæssigt
Slægtskab. Denne Modsætning gælder strengt taget
ikke, naar selve Begrebsdannelsen beror paa Analogi. I Stedet
for denne Modsætning maa sættes en Modsætning mellem forskellige
Grader og Arter af Analogi.

Ved Begrebsdannelsen sammenholdes forskellige Exempler
og der sker en Overførelse fra det ene til det andet, saaledes
at det simpleste og klareste Exempel lægges til Grund og belyser
de andre. Naar det gælder om at finde Sætninger, der
udtrykke Forholdene i Erfaringen og kunne lede til at finde
dem, træder Analogidannelsen endnu tydeligere frem.

I sin strengeste Form er Analogien kvantitativ og kaldes
da ved sit latinske Navn: Proportion. Forholdet mellem to Tal
kan være \textit{det samme} som Forholdet mellem to andre Tal, skønt
alle fire Tal ere indbyrdes forskellige. Analogien er her For%
% --- p. 38 ---
\setcounter{footnote}{0}%
\opage{38}holdsidentitet. I saadanne Tilfælde kan man stedse finde et af
Leddene i Forholdet ved Hjælp af de andre. Den kvantitative
Analogi er konstruktiv. Den skaffer det Ubekendte frem ved
Hjælp af det kendte Forhold til det Bekendte, uden at behøve
at raadspørge Erfaringen. Hvor Talforhold kunne paavises, formaar
man ved deres Hjælp at aflæse Begivenhedernes Gang i
Erfaringens Verden. Men denne Anvendelighed beror selv paa
en Analogi --- en Analogi mellem Tallenes Plads i Talrækken
og Begivenhedernes Plads i Tidsrækken. Talrækken er den
simpleste og nøjagtigste Ordning, vi kunne danne, og det er da
naturligt, at man stræber at bringe alle andre Ordninger og
Rækker i saa nøje Forbindelse med den som muligt. Denne
Analogi er bleven fremdragen af Maxwell og Cournot.

Da nu de Forandringer, der fuldstændigst kunne bestemmes
ved Hjælp af Tal, ere Bevægelserne i Rummet, saa forstaas,
hvorfor den moderne Naturvidenskab siden Renaissancens Dage
er gaaet ud paa at forklare alle Forandringer i Naturen som
Bevægelser. Den skyder de kvalitative Forskelle og Forandringer
til Side eller skriver dem paa Subjektets Regning; undertiden
udtrykker den sig saaledes, \textit{som om} de blot vare Skin.
Det er dog tilstrækkeligt at forudsætte en Analogi mellem kvantitative
og kvalitative Forhold, saaledes at de kvalitative Forhold
kunne bestemmes og forudsiges ved de kvantitatives Hjælp. Videnskaben
opfatter da de kvalitative Forandringer, som om de
blot vare kvantitative Forandringer, og man tænker sig Forandringerne
foregaa i Noget, der kun frembyder kvantitative Forskelligheder.
Saaledes opstaa Begreber som Materie, Atom, Æter.
Det er Billeder, eller rettere Skemaer, som kunne bestemmes
helt igennem med Tal, --- som gøre en konsekvent Anvendelse
mulig, --- og som kunne føre til paalidelige Slutninger med Hensyn
til nye Iagttagelser. --- Denne Analogi har vist sig at være
saa frugtbar og saa betydningsfuld, at man har været tilbøjelig
til at forvexle den med Identitet, eller dog til at betragte den
som den eneste mulige Analogi. Men dertil foreligger der ingen
Berettigelse.

Af endnu mere fundamental Betydning end Talforholdet er
et Forhold, der gør sig gældende i al Tankevirksomhed, hvad
% --- p. 39 ---
\setcounter{footnote}{0}%
\opage{39}enten denne angaar det Kvantitative eller det Kvalitative, nemlig
det logiske Forhold mellem Grund og Følge. Det er den
inderligste Forbindelse, der kan være mellem vore Tanker, afset
fra det rene Identitets- eller Dækningsforhold. Det er da
naturligt, at vi søge saavidt muligt at nærme alle vore Iagttagelser,
alle Forhold mellem Begivenheder, som Erfaringen viser
os, til dette Forhold mellem vore Tanker. Vi søge at gennemføre
en saadan Opfattelse af Forholdet mellem Begivenhederne
i Tilværelsen, at den ene kommer til at forholde sig til den
anden som Grund til Følge. Vi ordne saavidt muligt Begivenhederne
i saadanne Rækker, at de foregaaende Led kunne føre
til Forstaaelse af de senere, maaske endog til forud at bestemme
disse. Men da Forholdet mellem Grund og Følge er et rent
logisk Forhold, medens Forholdet mellem Begivenhederne i Tilværelsen
er et Tidsforhold, saa er det ved Hjælp af en Analogi,
at vi anvende Begrundelsens Lov paa Erfaringens Begivenheder.
Dette er Grundtanken i Kants bekendte Lære om Erfaringens
Analogier, hvorved de logiske Grundsætningers afgørende
Betydning for den reale Forsken fastholdtes, medens
paa samme Tid ældre Systemers dogmatiske Identificeren af
Grund og Aarsag (ratio og causa) ophævedes. Hegel har senere
gjort et romantisk Forsøg paa igen at afløse Analogien ved Identitet;
men selv af dem, der i vore Dage tillægge Hegels Tænkning
størst Betydning, opfattes dette Forsøg dog neppe som
lykket. Et er, at Tankens Love kunne lede os paa vor Vej gennem
Erfaringens Verden, et Andet er, at Tilværelsen skulde
være udtømt og fuldeligt udtrykt i Tankens Love.

Talrækken, Bevægelsen og Begrundelsesforholdet staa som
betydningsfulde Exempler paa Analogiens videnskabelige Betydning.
Specielle Exempler kunde hentes fra de enkelte Videnskaber.
Saaledes er Ideen om en Bølgebevægelse overført fra
det ene fysiske Omraade til det andet i Kraft af Forholdsligheden.
Undertiden véd man end ikke, hvad det er, der bølger,
kun at der sker en Bølgebevægelse, som ved den saakaldte
»Æter«. Den Betydning, Analogien har indenfor Biologien, idet
forskellige Organer hos forskellige levende Væsener kunne blive
af samme Betydning for Livet, har allerede Aristoteles frem%
% --- p. 40 ---
\setcounter{footnote}{0}%
\opage{40}draget. Men jeg skal til at belyse Analogiens Betydning endnu
henvise til den Bestræbelse, der har gjort sig gældende i Sociologien
for ved Analogien mellem Samfund og Organisme at
faa en ledende Tanke, et Grundskema for denne Videnskab. En
Række af fremragende Forfattere have søgt at lægge denne Analogi
til Grund, men de have forgæves stræbt at gøre Billedet
til et konsekvent og frugtbart Skema for den sociologiske Forsken.
For det Første er Samfundet en Forbindelse af Organismer,
og Forholdene indenfor en og samme Organisme kunne
ikke overføres paa en Gruppe af organiske Væsener. Dernæst
er der indenfor Organismen en betydningsfuld Forskel mellem
Elementer, til hvis Funktion der er knyttet Bevidsthed, og saadanne,
ved hvilke dette ikke er Tilfældet; i Samfundet ere derimod
alle Væsener bevidste, og de staa derfor som Formaal i en
ganske anden Forstand, end Organismens enkelte Dele kunne
siges at gøre det. Men den Iver, hvormed man indenfor Sociologien
har søgt efter en Analogi, der kunde lægges til Grund,
vidner om den Betydning, Analogien har for Videnskabens Udvikling.
Tillige viser dette Exempel, hvilke Fordringer der stilles
til de Analogier, Videnskaben kan bruge. De maa ikke blot
passe paa enkelte Punkter, men paa alle Punkter, saa at man
kan drage den fulde Konsekvens af dem, uden at Anvendeligheden
paa det Omraade, der skal belyses, hører op. Symbolet
eller Billedet maa kunne blive til et Skema, man trygt kan
følge i det Enkelte, idet det stadigt fører til nye, fremskridende
Undersøgelser. Det er Naturvidenskabens Lykke, at den har et
saadant Skema i den mekaniske Bevægelse. Det har i tre Aarhundreder
ledet den fremad til nye Opgaver og nye Opdagelser.
Naar den spekulative Filosofi og den religiøse Tænkning
ikke have ført til saadanne Resultater, maa Grunden søges i,
at de ikke formaa at omforme deres Symboler til Arbejdsskemaer\footnote[1]{I videre Betydning ere ogsaa Skemaerne, endog Taltegnene og Talordningen Symboler. I denne Betydning brugte Leibniz Ordet Symbol, naar han stillede symbolsk Erkendelse i Modsætning til intuitiv Erkendelse. (\textit{Meditationes de cognitione.} 1684). I Distinktionen mellem Skema og Symbol følger jeg Kant (\textit{Kritik der Urtheilskraft.} § 59).}.
Og Grunden hertil maa igen søges i, at de ikke,
% --- p. 41 ---
\setcounter{footnote}{0}%
\opage{41}som den hidtil omtalte Erkendelse, søge at forklare en Del af
Tilværelsen ved Hjælp af en anden Del, men at kaste Lys over
hele Tilværelsen ved Hjælp af, hvad der kan vides om en enkelt
Del.

\section*{III}

Naar spekulativ eller metafysisk Erkendelse skal have nogen
videnskabelig Betydning, maa dens Opgave være at udforme
en Verdensanskuelse, en Anskuelse af Tilværelsen eller dennes
inderste Grund, saaledes at den stadigt bygger paa, hvad Erfaringsvidenskaben
lærer. Nu er Tilværelsen som Helhed aldrig
given, og kan ikke tænkes given for nogen Erkendelse. Den
Undersøgelse, der her kan være Tale om, kan da kun gaa ud
paa, om noget af vore ved Erfaring vundne Begreber kunde
egne sig til at lægges til Grund ved en Helhedsopfattelse af
Tilværelsen. Vore Begreber ere nu alle dannede paa Grundlag
af bestemte Erfaringer, og udtrykke altsaa hvert for sig en Del
eller en Side af Tilværelsen. Spørgsmaalet bliver da, om Tilværelsen
som Helhed kan forstaas og udtrykkes ved Hjælp af
en enkelt Del eller en enkelt Side. Det er ikke Delene, der
her skulle belyse hverandre, men den enkelte Del, der skal belyse
Helheden. Problemet om Erkendelsens Mulighed bliver
her betydeligt skarpere.

Hvor vi gaa over fra det ene Erfaringsomraade til det
andet, vil det være muligt at angive de Ændringer, Begreberne
maa undergaa, naar de overføres fra det ene Omraade til det
andet. Men den Ændring, et Begreb maa undergaa, naar det
skal overføres fra Del til Helhed, maa være ganske anderledes
indgribende, end naar Overførelsen sker fra Del til Del. Spørgsmaalet
er, om en Overførelse fra Del til Helhed overhovedet
er mulig, da Helheden kan have andre Love end de, der gælde
for nogensomhelst enkelt Del for sig alene betragtet. Helheden
kan desuden have Sider eller Dele, der slet ikke komme frem
i vor Erfaring. Maaske kende vi kun nogle af de Elementer,
af hvilke Tilværelsen som Helhed er sammensat. I hvert Tilfælde
er det her umuligt at verificere Brugen af Analogien, da
Helheden aldrig foreligger given. Derfor er det intet Under, at
Tænkningen her ikke kommer til nogen Afslutning. Hverken
% --- p. 42 ---
\setcounter{footnote}{0}%
\opage{42}Materie eller Aand, hverken Mangfoldighed eller Enhed, hverken
Væren eller Vorden egne sig til afsluttende Udtryk for
Tilværelsen; denne sprænger stedse paany gennem Erfaringernes
Rigdom den afsluttende Form, i hvilken man vil samle
deres Indhold. Og dog er der en stadig Trang i den menneskelige
Aand til at gaa til Grænsen, --- til at anvende sine Hovedtanker
som Udtryk for Tilværelsen som Helhed betragtet. Denne
Trang vil neppe kunne udryddes; men den maa forbindes med
den kritiske Bevidsthed om, at vi her stedse anvende Analogier,
der ikke kunne gennemføres, Symboler, der ikke kunne
blive til videnskabelige Skemaer.

Blandt de Tankeforsøg, der her kunne gøres, har den filosofiske
Idealisme en særegen Betydning. Den fremtræder under
en dobbelt Form, som det har sin Interesse at undersøge, da
den ene Form (som jeg vil kalde spekulativ Idealisme) ikke vil
indrømme, at den bygger paa Analogi. Denne Form for Idealisme
anser det for en Nødvendighed, at hvad der existerer,
hvad der har Realitet, paa én Gang maa være Enhed og Mangfoldighed;
jo større Fylde og jo inderligere en Enhed der i et
givet Tilfælde er til Stede, des højere en Grad af Realitet staa
vi overfor. I denne Maalestok for sand Realitet kundgør selve
det Absolutes Væsen sig. Paa en betydningsfuld og dybsindig
Maade er denne Ide i den nyeste Tid bleven udviklet af Francis
Herbert Bradley. Overfor dette Tankeforsøg har jeg to Bemærkninger
at gøre. For det Første er det ikke konstruktiv
eller egentligt spekulativ Idealisme, der her foreligger. Det er
ikke et System, men en Maalestok, et Ideal, der udvikles. Vi
mindes mere om Kant og Maimon end om Hegel. For det Andet
er selve denne Maalestok, dette Ideal hentet fra den menneskelige
Aand. I vor Aand finde vi stedse et givet Indhold, hvis
Mangfoldighed til en vis Grad er sammenarbejdet til Enhed, og
jo større Fylden er uden Skade for Enheden og Koncentrationen,
des højere en Udviklingsgrad tillægge vi Aandslivet. Det
er denne Maalestok, hin Idealisme fra Aandslivet udstrækker
til at gælde for al Realitet. Det er i Virkeligheden en Analogi,
der ligger til Grund for denne Udvidelse: der gaas ud fra, at
i silk hvad der gælder den os umiddelbart tilgængelige Del af Til%
% --- p. 43 ---
\setcounter{footnote}{0}%
\opage{43}værelsen, maa gælde for hele Tilværelsen. Vi føres derved over
til den anden Form for Idealisme, der fra første Færd vedkender
sig at bygge paa Analogi, og som jeg vil kalde metafysisk
Idealisme. De vigtigste Repræsentanter for dette Standpunkt ere
Leibniz og Lotze. Leibniz er den, der først med Klarhed og
Energi har udtalt, at al idealistisk Verdensanskuelse grunder
sig paa en Analogislutning. Fra den kritiske Filosofis Standpunkt
bliver nu Spørgsmaalet, om det er tvingende nødvendigt
at foretage denne Analogislutning. Det kunde synes, som om
man havde bundet Valg mellem denne Mulighed og Materialismen,
mellem at opfatte det Materielle i Analogi med det Aandelige
eller det Aandelige i Analogi med det Materielle. I en vis
Forstand er det ogsaa saa. Men denne Bundethed hænger
maaske sammen med, at vi erfaringsmæssigt kun kende to Tilværelsesformer,
den aandelige og den materielle. Der bestaar
dog her intet kontradiktorisk, men kun et kontrært Modsætningsforhold,
og den Mulighed er ikke udelukket, at Tilværelsen (saaledes
som Spinoza tænkte sig det) har langt flere Grundformer
(Attributer) end de to, vor Erfaring viser os. Allerede af denne
Grund maa Spørgsmaalet stadigt staa aabent, afset fra de Vanskeligheder,
der under alle Betingelser ligge i Overførelsen af et
Begreb fra Del til Helhed. Selvom Tilværelsen som Helhed
kunde udtrykkes ved en enkelt Del eller en enkelt Side, var
det jo ikke sagt, at det var en af de to, vi kende, der egnede
sig dertil.

Medens den anden Art af Idealisme har et Fortrin i den
klart bevidste Maade, paa hvilken den henviser til Analogien
som Grundlag, udmærker den første Art sig ved at angive en
Maalestok for Realitet og derved tillige en Metode for Forsken.
Idealismens videnskabelige Betydning bestaar i, at den ved at
finde Realitetens Maalestok i Forbindelsen af Mangfoldighed og
Enhed sætter et Arbejde i Gang for Udvidelse af Erfaringen
og for Sammenarbejden af det Fundne, uafbrudt Analyse og
Syntese i inderlig Vexelvirkning. Dens praktiske eller etiske Betydning
stemmer overens hermed; den kræver baade Livets Udvidelse
og Forhøjelse, forkynder et Plus ultra! og et Excelsior!
paa én Gang.

% --- p. 44 ---
\setcounter{footnote}{0}%
\opage{44}Metafysik er maaske mere en Kunst end en Videnskab.
Det er da Filosofiens Opgave ved Opstillingen af en Maalestok
som den omtalte at øve en Kritik af de uvilkaarligt udviklede
Verdensanskuelser og derigennem bringe dem i Forbindelse med
det egentlige videnskabelige Arbejde. I hvert Tilfælde viser der
sig her en Lighed imellem Filosofiens højeste Synspunkter og
den exakte Videnskabs Grundsætninger: begge Steder bestaar
Værdien og Berettigelsen i, at de opstillede Principer formaa
at indlede et frugtbart Tankearbejde, føre til nye Opgaver og
nye Løsninger.

\section*{IV}

Ogsaa Religionen arbejder i sin Forestillingskreds med Analogier.
Den hævder paa sin Vis, ligesom Metafysiken, en Helhedsopfattelse
af Tilværelsen. Men medens Metafysikeren ledes
af den intellektuelle Interesse at vinde en afsluttende Tanke og
overlader til Tænkningens og Erfaringens Konsekvens, hvad
denne Tankes Indhold skal være, er den væsentlige Interesse i
Religionen, hvilken Skæbne det Værdifulde (hvor man saa finder
det) faar under Tilværelsens store Kamp. Religion betinges
ved Erfaringen om, at Værdier opstaa og Værdier forgaa, og
det religiøse Problem bestaar i det Spørgsmaal, om det er muligt
at fastholde en Kontinuitet under denne stadige Skiften og ud
over den Begrænsning, vor Erfaring er underlagt. Hvor den
religiøse Bevidsthed er oprindelig og levende, er den uafbrudt
optagen af at samle sine Erfaringer og finde et fyldestgørende
Udtryk for dem. Den fastholder, at der i de Værdier, som erfares,
i Følelseskvaliteterne, ligesaa vel som i Sansekvaliteterne,
aabenbarer sig en stor Tingenes Orden, og den søger at forme
denne Overbevisning i bestemte Tanker. Værdierne kunne for
den religiøse Bevidsthed ikke staa som sporadiske Glimt, som
Biprodukter i Tilværelsens store Proces. Men under sin Udviklingsgang
bliver den religiøse Bevidsthed klar over, at hvad
der bestaar, ikke kan være de enkelte, empiriske Værdier;
deres Tid kan være forbi, og dog bestaar --- det er den religiøse
Tro --- Værdiernes Rige. Det koster en haard Kamp,
inden denne Erkendelse trænger igennem. Menneskene slippe
% --- p. 45 ---
\setcounter{footnote}{0}%
\opage{45}nødigt de nærmest liggende Erfaringer. Skal der bruges Symboler,
bruges helst de mest umiddelbare og nærliggende Kendsgerninger.
Allerhelst vil den religiøse Bevidsthed slippe for den
Indrømmelse, at alle dens Forestillinger ere symbolske. Dogmatismen
har her sit egentlige Hjemsted; men dens Træk ere
aldeles lignende, som hvor den ytrer sig paa det videnskabelige
Omraade.

Religionen er ligesom Metafysiken nødt til at udtrykke det,
der skal være en absolut Helhed, ved Forestillinger, der ere
hentede fra en enkelt Del eller Side. Dens Vanskeligheder ere
i denne Henseende af lignende Art som Metafysikens. Men
Kampen mellem Tanke og Billede bliver mere indgribende i
Religionen end i Metafysiken. Metafysikeren kæmper for at gøre
Billedet eller Symbolet til et videnskabeligt brugbart Skema.
Denne Interesse staar i Religionen ikke i første Linie. Det gælder
om store, gribende Billeder, der kunne udtrykke Stemningen
under Livskampen og holde Mod og Haab oppe. Muligheden
af at gennemføre Symbolet konsekvent træder tilbage, og
det er i Regelen en mere eller mindre fjendtlig Kritik, der
undersøger den. Naar Billedet blot udtrykker en væsentlig
Side ved Livserfaringen, fastholdes det, og Inkonsekvensen
holdes borte ved, at Sammenligning udelukkes saa længe som
muligt.

Det religiøse Problems Fremtid beror paa, om det er muligt
at fastholde Overbevisningen om Værdiernes Kontinuitet
under fremskridende Erfaring og under klar Bevidsthed om,
at enhver Formulering bliver symbolsk i udpræget Forstand.
Hvorledes det end vil gaa hermed, saa frembyder det religiøse
Problem en Analogi med de andre Hovedproblemer for den
menneskelige Aand. Og hvis man har Ret i, at det egentligt
religiøse Axiom er Sætningen om Værdiens Bestaaen, saa viser
det sig ogsaa her, at Grundsætningernes Betydning beror paa
det aandelige Arbejde, de fremkalde. Hint Axiom kan nemlig
kun hævdes gennem en stadig Stræben efter at opdage og
frembringe Værdier i Virkelighedens Verden. Det religiøse
Spørgsmaal bliver derfor tilsidst et etisk Spørgsmaal. Det er
% --- p. 46 ---
\setcounter{footnote}{0}%
\opage{46}i Religionen ligesaa lidt som i Videnskaben de enkelte Billeder eller Symboler, men den Brug, der gøres af dem, der er det Væsentlige. Der er en betydningsfuld Analogi mellem Religion og Videnskab, hvor de blive klare over deres Væsen, nemlig i den Brug, de hver for sig gøre af de Analogier, med hvilke de nu engang ere nødte til at arbejde.

% --- p. 47 ---
\setcounter{footnote}{0}%
\opage{47}\essay[Forord og Efterskrift til min Religionsfilosofi. (1905)]{FORORD OG EFTERSKRIFT TIL MIN\\ RELIGIONSFILOSOFI}
\begin{center}(1905).\end{center}
\streg

En Drøftelse af det religiøse Spørgsmaal, ikke blot i dets
Forhold til Videnskaben, men ogsaa i dets Forhold til det personlige
Liv, kan ikke undgaa at være farvet af den Personlighed,
der anstiller Drøftelsen. Dette lægger jeg i min Fremstilling
af Religionsfilosofien heller ikke Skjul paa. Det personlige
Element, der nødvendigvis gør sig gældende, kan virke baade
til Bedste for Klarheden og Grundigheden af Problembehandlingen
og til Skade for den. Hvilken af Delene der gælder med
Hensyn til mit Forsøg, kan jeg naturligvis ikke selv afgøre.
Men jeg kan maaske give et lille Bidrag til Afgørelsen ved
at se tilbage paa min personlige og videnskabelige Udviklingsgang.


Jeg har undertiden set den Bemærkning, at min nuværende
Stilling til det religiøse Spørgsmaal finder sin Forklaring ved,
at jeg i mine unge Aar studerede Teologi, og at der endnu
sidder mig Adskilligt i Blodet fra den Tid, der gør, at jeg i
dette Spørgsmaal staar noget anderledes end mange andre Fritænkere.
For mig selv staar det saaledes, at saa skarpt jeg end
føler Modsætningen mellem mine nuværende Anskuelser og
dem, i hvilke jeg fandt Næring i min Ungdom, saa er dog Sammenhængen
i min Udvikling mig mest slaaende. Jeg har den
% --- p. 48 ---
\setcounter{footnote}{0}%
\opage{48}Overbevisning, at jeg i Aarenes Løb har vundet fyldigere og
inderligere Tilfredsstillelse for en personlig Trang, jeg i mine
unge Aar søgte fyldestgjort ad helt andre Veje end nu. Hvad
jeg mener at have bevaret fra hine Dage, er Sansen for det
indre Livs Betydning og Trangen til at vinde Udtryk for Sindets
Higen og Stemning under Livets Gang --- til at faa Livets
Værdi og Livets Skæbne udtrykt i store og gribende Tanker
eller Billeder. Det var denne Trang, der i sin Tid førte mig
til saa længe som muligt at holde fast ved den positive Religions
Forudsætninger og kun Skridt for Skridt lade dem falde.
Jeg saa dengang i Kristendommen, saaledes som Kirken fremstillede
og repræsenterede den, den eneste Form, i hvilken en
saadan Trang kunde fyldestgøres.

Det var Læsningen af Søren Kierkegaards Skrifter, der
først rokkede ved denne Antagelse. Det blev mig klart, at det
var et stort Selvbedrag, naar Kirken og den »moderne Kristendom«
mente at have naat en Harmoni mellem det Personlighedsliv,
der kræves i den oprindelige Kristendom (»det nye
Testamentes Kristendom«), og Livet, som det føres i Familie og
Stat, i Kunst og i Videnskab. Teologien havde søgt at slaa en
Bro her, mendette Forsøg beroede paa en stor Misforstaaelse.

Det var tunge og mørke Aar, som nu fulgte. Jeg ser ikke
tilbage paa min Ungdom som en lykkelig Tid. Med Tvivlen
om, hvorhen jeg skulde vende mig med min Trang til at forstaa
mig selv, forbandt sig Tvivlen om, hvor jeg skulde finde en
Arbejdsplads. Mit Arbejde kunde ikke skaffe mig Tilfredsstillelse,
naar det ikke kunde blive af væsentlig Betydning for min
Personlighed. Men Bruddet med Teologien førte mig ind i en
Dobbelthed, en Dualisme mellem det indre og det ydre Liv,
som i Længden var uudholdelig. Det var Livet selv, der førte
mig ud over denne Dobbelthed ved at lære mig rent menneskelige
Forholds Værdi at kende af egen Erfaring. Som jeg under
Indflydelsen af Kierkegaards Problemer havde kastet mine teologiske
Bøger henad Gulvet i Utaalmodighed over deres uklare
Sammensmelten, saa blev det nu Kierkegaards Bøger, jeg kastede
fra mig, fordi jeg ikke kunde forlige deres Indhold med, hvad
Livet nu  lærte mig. Jeg kom i et levende Forhold til Livet
% --- p. 49 ---
\setcounter{footnote}{0}%
\opage{49}udenfor Lønkammeret. Og da Vejen først var aaben, førte det
ene Skridt det andet med sig. Snart blev ogsaa Videnskaben
for mig en Realitet, der krævede sin selvstændige Plads ved
enhver alvorlig Opgørelse af Livets Regnskab. Studiet af den
græske Tænkning i hine Aar var herved af stor Betydning for mig; den lærte mig at indse Muligheden og Værdien af en fri
og menneskelig Livsforstaaelse og Livsførelse. Jeg maatte nu i
Livet og i den virkelige Erfaring finde en Mulighed for Tilfredsstillelsen af den Trang, der havde ført mig først til Teologien, siden til den Kierkegaardske Dualisme. Jeg saa, at Vejen
 maatte være lang, og jeg har fra Begyndelsen af ikke havt stor
Tro til, at jeg vilde komme ret langt frem ad den; men jeg
fastholdt Forvisningen om, at her var et Problem, der ikke kunde
skydes til Side.

I Aaret 1870 disputerede jeg for den filosofiske Doktorgrad
(med en Afhandling om »Den antike Opfattelse af Menneskets
Villie«), og i den korte Selvbiografi, som i den Anledning
 blev trykt i et Universitetsprogram, udtalte jeg mig paa følgende
Maade om min Udviklingsgang og min Bestræbelse:

»Som mit Fremtidsstudium havde jeg i min Skoletid en
Tidlang tænkt mig Filologien paa Grund af min Interesse for
de gamle Sprog og Historien; dog blev Trangen til at udvikle
og klare den samtidigt vaagnende religiøse Sans den stærkeste, saa at jeg med Lyst og Iver og med det Formaal engang at
træde i Kirkens Tjeneste tog fat paa det teologiske Studium.
Men der opstod snart en Tvivl hos mig, som hindrede den Tilfredsstillelse,
jeg havde ventet at finde i dette Studium baade
i religiøs og i videnskabelig Henseende, saa at det virkede nedslaaende
og hæmmende i Stedet for klarende og befriende. Studiet
af S. Kierkegaards Skrifter og af Filosofien, saavidt jeg den
Gang efter Tid og Evne kunde drive det, aabnede efterhaanden
mine Øjne for det Berettigede i den Kritik, der med saa stort
Eftertryk rettedes mod Teologien fra Universitetets filosofiske
Lærestole\footnote[1]{Rasmus Nielsens og Hans Brøchners.}. Jeg afsluttede mit teologiske Kursus i Sommeren
1865 med den Overbevisning, at jeg ad andre Veje maatte søge
% --- p. 50 ---
\setcounter{footnote}{0}%
\opage{50}den Sandhed, jeg ikke havde fundet der. Tanken om at indtræde
i gejstlig Virksomhed havde jeg opgivet; jeg beskæftigede
mig nu især med Historie, Literatur og de gamle Sprog, i hvilke
Fag jeg ogsaa virkede som Lærer. Den antike Oldtid traadte
frem for mig i et friskere og klarere Lys, netop ved Modsætningen
til den skolastiske Periode, jeg forekom mig at have
gennemlevet; jeg glædede mig ret ved at læse dens Historieskrivere, dens Digtere og snart især dens Filosofer. Den dybe
Overbevisning, jeg herved vandt om Videnskabens og det humane
Livs selvstændige Betydning, i Forening med, at jeg endnu
saa det sande Udtryk for min religiøse Tro i de kirkelige Dogmer,
gjorde, at jeg troede at finde en Løsning for mine Tvivl
i en Anskuelse, som netop da fremtraadte med større Energi
for Offentligheden, men var traadt mig imøde fra mine første
Studenteraar, og som skarpt sondrede mellem Religion og Videnskab
som to selvstændige Riger\footnote[1]{Rasmus Nielsens Lære om Tro og Viden.}. Jeg har altid opfattet denne
Anskuelse mere fra den skeptiske, negative Side, mere som et
Middel til en personlig Løsning af Vanskelighederne, end som
en egen positiv Lære. Et lille Skrift, jeg dengang udgav, »Filosofi
og Teologi« (1866), søgte fra dette Synspunkt at bestemme
hin Anskuelses Forhold til dem, der tidligere vare fremtraadte
i vor Literatur. Hvad der senere har ført mig bort fra den,
er dels de tildels meget grundige Indlæg, der fremkom under
Sagens videre Forhandling, dels at den nye Anskuelse under
sin videre Udvikling stedse mere syntes at blive fra en skeptisk
eller kritisk til en positiv, dogmatisk. Hertil kom nu ogsaa
mit eget, efterhaanden mere selvstændige Studium af Filosofien.
Trangen til Enhed og Harmoni i Livs- og Verdensanskuelsen
voxede tilligemed Tilliden til den menneskelige Aands Kraft til
at opnaa en saadan. Herved er der vel endnu kun sat et Ideal,
men et Ideal, der netop maa anspore  til stedse fornyet Forsken.«

\stjerner

Af væsentlig Betydning blev  senere for mig Studiet af den
engelske Filosofi. De Poler, mellem hvilke min Udvikling har
% --- p. 51 ---
\setcounter{footnote}{0}%
\opage{51}svinget, er Kirkegaards Personlighedslære og den fransk-engelske
Positivisme. Brøchner, hvem jeg stod personligt nær i hans sidste
Leveaar, havde desuden henledet min Opmærksomhed paa Spinoza,
og jeg erindrer, hvorledes Læsningen af Spinozas Breve
paa en Fodtur i Jylland i Sommeren 1872 greb mig stærkt
og gav mig et centralt Synspunkt, jeg senere har fastholdt.

At det er lykkedes mig efter mine Læreaar at finde et nyt
Stade, tør jeg ikke sige. Men jeg har fundet et Stade, hvor jeg
selv kan staa, og hvorfra jeg maaske kan hjælpe Andre til at
finde, hvor de skulle staa. Jeg har vundet den Overbevisning,
at selv det Højeste og Ædleste i Aandens Verden underligger
en stor Naturordens Love, som det er Videnskabens Opgave at
udfinde, og at det virkeligt Værdifulde ikke krænkes, men vinder
ved at have sine faste Rødder i Virkelighedens Sammenhæng.
Selv de skønneste Toner ere ikke for gode til at udtrykkes
i Noder. Men mit Partitur er ikke færdigt, og bliver
det vel aldrig. Jeg kan høre flere og skønnere Toner, end jeg
kan udtrykke ved min Nodeskrift, og jeg finder i Erfaringen
mange Tegn, som jeg ikke kan tyde, men som muligvis
ere en Nodeskrift for Toner, jeg ikke kan høre. Det rokker
ikke ved min Overbevisning, og det rokker fremfor Alt ikke
ved min Metode. Problemerne ere blevne større og skarpere
for mig under min Vandring i Tankens Verden. Det kalde
Nogle min Skepticisme eller min Negativitet. Jeg selv finder
det naturligt, at der bliver højere til Himlen, jo højere man
kommer op paa sin Bjergbestigning.

\stjerner

For min Opfattelse af det religiøse Problem blev mit Forsøg
paa at fremstille Psykologien og Etiken paa Erfaringens
Grund, med Bortseen fra spekulative og religiøse Forudsætninger,
af væsentlig Betydning. Derved blev Erkendelsen skærpet
af Nødvendigheden af psykiske Ækvivalenter, d. v. s. af nye
aandelige Livsformer, der kunde træde i Stedet for dem, som forsvinde,
og bringe Erstatning for dem, saa at Livet ikke under sin
Udvikling mister sin Værdi (maaske ogsaa sin Energi). Jeg fandt
nu fremdeles --- ved de Studier, der forberedte Udgivelsen  af
% --- p. 52 ---
\setcounter{footnote}{0}%
\opage{52}mit Værk over den nyere Filosofis Historie --- at saa godt som
alle fremragende Filosofer siden det store kritiske og revolutionære
Gennembrud i Slutningen af det 18. Aarhundrede havde
været opmærksomme paa dette Punkt. I et Foredrag, jeg i
Foraaret 1895 holdt i Kristianias Studentersamfund, og som er
trykt i første Del af mine »Mindre Arbejder« under Titel:
»Kampen mellem gamle og nye Tanker. Et Tilbageblik«, søgte
jeg at vise, hvorledes Problemet om psykiske Ækvivalenter stadigt
havde rørt sig i det sidste Aarhundredes Tænkning. I en Universitetstale
(1897) om det religiøse Problem (der ligeledes findes
i første Del af »Mindre Arbejder«) betonede jeg, at det afgørende
Spørgsmaal for det religiøse Problems Fremtid er, om den
personlige Livsanskuelse skal og kan grundes alene paa, hvad
der kan udledes af de videnskabelige Resultater, eller om der
ikke stedse maa gøre sig Elementer og Tendenser gældende
ved det sidste Opgør om Livets Betydning, der ikke kunne retfærdiggøres
ad videnskabelig Vej, selvom de ikke stride mod
Videnskabens Resultater.

Disse Studier og Overvejelser førte mig til den Opfattelse,
som jeg har fremstillet i min, i Aaret 1901 udkomne Religionsfilosofi.


Religionsfilosofien raader ikke over andre Metoder end de
almindelige filosofiske. Dens Metode er en Forbindelse af erkendelsesteoretisk,
psykologisk og etisk Undersøgelse. --- For den
erkendelsesteoretiske Betragtning viser det sig mere og mere,
at Religion ikke fremgaar af rent intellektuelle Motiver, og at
religiøse Forestillinger ikke kunne hævdes ad rent intellektuel
Vej. Skulle religiøse Forestillinger fremdeles have Værdi, maa
det være som billedlige Udtryk for en Livserfaring og en Livstrang
af mere speciel og personlig Art end den Erfaring, Videnskaben
bygger paa. --- I det psykologiske Afsnit vises, hvorledes
den religiøse Følelse faar sin naturlige Plads i Bevidsthedslivet,
naar den opfattes som betinget ved de Erfaringer, der
gøres eller menes at være gjorte angaaende Forholdet mellem
Livsværdien og Virkeligheden. Det er Spændingsforholdet mellem
Værdi og Virkelighed, der atter og atter stiller Problemet
og vil vedblive at stille det, selvom de positive Religioners Tid
% --- p. 53 ---
\setcounter{footnote}{0}%
\opage{53}skulde være forbi, ja selvom man ikke længere vil bruge Ordet
Religion. En historisk Undersøgelse af de dybeste Motiver, som
den religiøse Bevidsthed appellerer til, godtgør den Sætning, at
den religiøse Trang er en Trang til at hævde Værdiens Bestaaen.
--- I et etisk Afsnit drøftes endeligt selve Religionens Værdi,
og det vises, at den religiøse Trang har den Betydning at føre
til Opdagelse af nye Værdier og til Hævdelse af de allerede
fundne. Skønt Religion ikke er en nødvendig Forudsætning
for etisk Livsførelse, kan den dog udfylde og fortsætte denne.
--- Resultatet bliver, at Religion er en Form for aandeligt Liv,
der --- hvis Livets Kraft ikke skal svækkes --- ikke kan forsvinde,
uden at en ny og ækvivalent Livsform udvikles. Det
store Spørgsmaal er, om dette er muligt. Min Betragtning vender
sig dels mod trygge Dogmatikere, dels mod selvglade Fritænkere.
»Prophete rechts --- Prophete links!« Problemet ligger
midt imellem dem, og de sky det begge, hver paa sin Vis.

\stjerner

Hos Tænkere med en teologisk Fortid kan denne virke
senere hen paa forskellige Maader. --- Der kan være en blot
Efterklang, vemodig, men svag, og stedse svagere, en Efterklang,
som kan hindre i at drage Konsekvenser, der ellers vilde staa
som uundgaaelige. --- Der kan gøre sig en Kontrastvirkning
gældende, som ytrer sig ved, at man senere hen ikke kan finde
Ord, der ere stærke nok til at udtrykke Modsætningen, Antipatien,
maaske Modbydeligheden. Denne Type minder undertiden
om Skoledrenge, der nyde deres Frihed ved at haane den
Skole, de --- ikke altid lige grundigt --- have gennemgaaet. Det
er en Hildethed, ligesaa vel som Efterklangen. --- En tredie
Type kan betegnes som den kritiske Kontinuitets. Den finder
en dybere liggende Sammenhæng bag den stærke Modsætning
ved at gaa tilbage til det, der under forskellige Former maa
danne og altid har dannet Grundlag for enhver alvorlig Forstaaelse
af Livet og dets Skæbner. Selvom de Ideer, i hvilke
Menneskene i tidligere Tider udtrykte denne Forstaaelse, for os
nu i det Højeste have poetisk Betydning (det er end ikke altid,
de have det), saa kan der dog være et Slægtskab til Stede, og
% --- p. 54 ---
\setcounter{footnote}{0}%
\opage{54}det er i hvert Tilfælde Religionspsykologiens Opgave at undersøge, hvorledes det i denne Henseende forholder sig.

Det er naturligvis vanskeligt for den Enkelte at afgøre, til
hvilken af disse Typer han selv hører. Men jeg har ingen Tvivl
om, hvilken af dem der er den frugtbareste og den mest berettigede.
Det maa være den, der uden at hildes af Efterklang
eller Kontrastvirkning benytter den Erfaring, der er vunden ved
at leve i en modsat Lejr, til Belysning af Problemet, saaledes
som dette stiller sig for os nu og foreløbigt maa stille sig. ---
Det er af ikke ringe Interesse, at Undersøgelser i denne Retning ere i Gang i hele den civiliserede Verden. (Smlgn. nedenfor Afhandlingen »Om nogle religionsfilosofiske Arbejder fra
den nyeste Tid«). Det er heller ikke uden Interesse at se, hvorledes det religiøse Problem, naar det opfattes paa den her antydede Maade, viser sig at være analogt med de andre store Problemer, der sysselsætte menneskelig Tænkning. Det er jo sandsynligt, at hvad der paa de forskellige Omraader bliver Problem for den menneskelige Tanke, maa frembyde analoge Træk
(Se herom mit Universitetsprogram Novbr. 1902: Filosofiske Problemer). ---

Om mit Arbejde kan tages til Indtægt fra den ene eller den
anden Side, er mig temmeligt ligegyldigt. Det vidner dog enten
om stor Nøjsomhed eller om naiv Beregning, naar man endog fra katolsk Side har fremstillet min Bog som et gunstigt Tidernes Tegn. Jeg selv finder Broen stedse mere afgjort kastet af mellem mig og enhver Art af dogmatisk Religion. Men maaske er man fra teologisk Side saa vant til at finde religiøse Spørgsmaal behandlede med Haan og Spot eller med Skoledrengeargumenter, at allerede det, at Problemet drøftes med en vis Alvor, kan fremkalde det Indtryk, at der er sket en Tilnærmelse. Det er dog ikke alle Teologer, der have opfattet min Fremstilling
 saaledes. I »Kristeligt Dagblad« stod der for nogle Aar tilbage en skarp, men ridderlig Kritik af mine Anskuelser. Kritikens Forfatter endte med at bevidne sin Modstander sin personlige Agtelse. Den daværende Formand for »Indre Mission« bebrejdede ham dette, idet han erklærede, at jeg var en Mand, ingen Kristen kunde have personlig Agtelse for, --- og den da%
% --- p. 55 ---
\setcounter{footnote}{0}%
\opage{55}værende Redaktion af »Kristeligt Dagblad« sluttede sig i en
Note til denne Erklæring. Sligt beroliger, naar man undertiden
føler sig ilde ved at blive tagen  til Indtægt fra en vis Side.

\stjerner

Jeg har forsøgt at sammenfatte  Religionens Væsen og Betydning
i en enkelt Sætning, som jeg har søgt at begrunde ad
psykologisk og ad historisk Vej: Sætningen om Værdiens Bestaaen.
Fra forskellige Sider har jeg erfaret Indvendinger mod
denne Hypotese. (For mere end en Hypotese har jeg ikke udgivet
den).

Det er blevet sagt, at en saadan almindelig Sætning ikke
kan udtrykke Religionens Væsen, da Religion er noget Levende
og Konkret. Denne Indvending er altid forekommet mig noget
barnlig. Enhver almindelig Sætning kan kun være et kortfattet
Udtryk for, hvad der i mangfoldige forskellige Former forekommer
i Virkeligheden. Ingen af disse mangfoldige Former kan
umiddelbart være udtrykt i den almindelige Sætning, men dennes
Betydning maa bestaa i at afgive en Ledetraad til at finde det,
hvori de enkelte Fænomener trods deres Forskellighed ligne
hverandre. Den almindelige Sætning, i hvilken vi definere et
Hvirveldyr, kan ikke indeholde Noget om Løvens Brølen eller
om Menneskets Tale, men den kan derfor meget godt gælde
baade for Løven og for Mennesket. Meningen med hin religionsfilosofiske
Sætning er naturligvis ikke den, at Religionerne
selv direkte tale om Værdier og om disses Bestaaen, men at
der i al Religion ytrer sig en Trang, som i vor psykologiske
Teori kan betegnes som en Trang til at antage en Bestaaen af
det Værdifulde, Mennesket paa det givne Udviklingstrin kender.
Jeg har ikke villet reducere Religion til abstrakt Metafysik, men
jeg har talt som Filosof om Noget, der er Andet og Mere end
Filosofi. Man kunde med samme Ret bebrejde Fysiologen, at
han sætter en abstrakt Teori i Stedet for det virkelige Aandedræt.


Man har dernæst fundet det Væsentlige i Religionen paa
andre Punkter end det, jeg lægger Vægten paa.

Fra et teologisk eller dog teistisk Standpunkt har man sagt,
% --- p. 56 ---
\setcounter{footnote}{0}%
\opage{56}at det Væsentlige i al Religion er, at Mennesket føler sig i Forhold
til et eller flere højere personlige Væsener: Religion bestaar
i et Forhold mellem et Jeg og et Du! --- Overfor denne
Paastand spørger jeg (og har jeg allerede i min Bog spurgt):
Hvad er det, der gør Guderne guddommelige, og hvilken Trang
er det, der fører til Troen paa dem? Dersom Guderne ikke antoges
at hævde det, Menneskene anerkende som godt og stort,
vilde Menneskene saa dyrke dem? Selv hvor Guderne blot
dyrkes som de mægtigste Væsener, der kendes, ligger der en
Trang til Grund til at fastholde en fortsat Bestaaen af det, som
Mennesket sætter højest (hvilket nu dette maatte være). Naar
Religion hidtil overvejende har været Tro paa personlige Væsener,
skyldes det for en stor Del Indflydelsen af Animismen,
d. v. s. Tendensen til at opfatte Begivenhederne i Verden som
Virkninger af personlige Væseners Indgriben, en Tendens, der
ytrer sig meget tidligt. De religiøse Forestillinger ere meget
ofte Omdannelse af Forestillinger, der have udviklet sig af andre
end religiøse Motiver. I Troen paa en Sjælevandring, der er
ejendommelig for mange Religioner, haves et Exempel paa det
Samme. Et andet Exempel frembyder den Maade, paa hvilken
den græske Filosofi benyttedes af den kristelige Teologi. Spørgsmaalet
er, om vi i saadanne Forestillinger have det Væsentlige
i Religionen, eller om det Væsentlige ikke maa være den Trang,
der føles til at optage og forme dem, en Trang, der under andre
Forhold kunde lede til at opføre helt andre aandelige Bygninger.

Fra et religionshistorisk Synspunkt har man fundet Religionens
Væsen i, at den er en Art Magi, en middelbar Magi,
der opstaar, naar Mennesket opgiver Haabet om umiddelbart at
kunne ændre Naturens Gang efter sine Ønsker. Ved Bønner
og Ofre mener han da nu at kunne foranledige mægtigere Væsener
(hvad enten disse forestilles som personlige eller ikke) til
at gribe ind og bøje Naturens Gang, saa at han kan opnaa sit
Ønske. Overfor denne Opfattelse spørger jeg paa lignende Maade
som for: \textit{hvorfor} søger Mennesket at vinde en Magt over Naturen,
han ikke kan vinde ved eget Arbejde? Mon ikke fordi
han erfarer Grænsen for sin egen Magt og dog føler Trang til
at vinde Tryghed og Sikkerhed for det, der ligger ham mest paa
% --- p. 57 ---
\setcounter{footnote}{0}%
\opage{57}Hjerte (hvad dette saa maatte være)? Atter her viser denne
Trang sig som det Afgørende, og til den maa man derfor gaa
tilbage, naar man vil efterspore de inderste Kræfter i det religiøse
Liv. Kun derved vil man kunne faa et Grundlag for
Drøftelsen af det religiøse Problem i dets hele Omfang. Trangen
kan bestaa, skønt Formerne skifte, og hverken Magi eller Animisme
behøve til evige Tider at være knyttede til Religionens
Væsen. Skulde Religionens Tid være forbi, fordi det erklæres,
at Kærlighed er bedre end Ofre, eller fordi intet Billede længere
staar som helt fyldestgørende? --- Modsætningen mellem
gamle og nye Former bliver derfor ikke betydningsløs. Den
kan til given Tid fremkalde Kampe paa Liv og Død i den ydre
Verden og de største Brydninger i Sindets indre Verden. Den
filosofiske Undersøgelse kan kun angaa \textit{Muligheden} af nye
aandelige Livsformer; den virkelige Frembringelse af dem er
Livets Sag.

% --- p. 58 ---
\setcounter{footnote}{0}%
\opage{58}\essay[Om nogle religionsfilosofiske Arbejder fra den nyeste Tid. (1903)]{OM\\ NOGLE RELIGIONSFILOSOFISKE ARBEJDER\\ FRA DEN NYESTE TID.}
\begin{center}(1903).\end{center}
\streg

\section*{INDLEDNING.}

Ved Religionsfilosofi forstaar jeg ikke et Forsøg paa at
sætte Spekulation i Religionens Sted, endnu mindre et Forsøg
paa ad Spekulationens Vej at begrunde religiøse Antagelser, men
et Forsøg paa at belyse Religionens Stilling i Tidens Aandsliv.
Religionen er her Genstand for Undersøgelse som en faktisk
given Form for aandeligt Liv, hvis Beskaffenhed, Betydning og
 Berettigelse det gælder at bestemme. Spørger man nu, hvilke
Synspunkter Filosofien raader over til Brug ved en saadan Undersøgelse, saa kan der sikkert kun være Tale om tre Synspunkter,
--- men det er egentligt ogsaa alle de Synspunkter, Filosofien
idethele raader over: det logiske (eller erkendelsesteoretiske), det
psykologiske og det etiske. Da enhver Religion --- hvor vid
eller hvor snever en Betydning man saa tager Ordet i --- faar
Udtryk i visse Forestillinger, opstaar naturligt Spørgsmaalet om
disse Forestillingers Gyldighed, deres Erkendelsesværdi. Indeholde
de en konsekvent og gyldig Erkendelse af Tilværelsen, og
kunne de bringes i Overensstemmelse med de Forestillinger, som
Videnskaben ad sine egne Veje er naat til at fastholde som
gyldige? Selvom nu en saadan \textit{erkendelsesteoretisk} Undersøgelse
skulde føre til det Resultat, at de religiøse Forestillingers Be%
% --- p. 59 ---
\setcounter{footnote}{0}%
\opage{59}tydning ikke kan bestaa i at skaffe os en videnskabelig Forstaaelse
af Tilværelsen, er dermed Spørgsmaalet om deres Værdi
og Berettigelse ikke afgjort. Den kan bero paa deres Forhold
til andre Sider af Sjælelivet end den rent intellektuelle. Og
direkte Iagttagelse viser jo ogsaa, at Følelse, Trang og Villie
spille nok saa stor en Rolle i Religionen som Erkendelse. Der
er da al Anledning til at foretage en \textit{psykologisk} Undersøgelse
af det religiøse Liv, saaledes som Iagttagelse lærer os det at
kende i Samtiden og historisk Overlevering i Fortiden. Denne
Undersøgelse gaar for det Første ud paa at finde de stadigt virkende
Motiver i Religionen og at udskille dem af de skiftende
Former og Forbindelser, i hvilke de under forskellige Forhold
fremtræde. Kun i disse Motiver kan det Væsentlige i Religionen
være at søge, --- det, der, saalænge det besidder Livskraft,
stedse vil kunne ytre sig i nye Skikkelser. Der existerer ingen
Religion i al Almindelighed; men Erfaringen viser os visse Grundmotiver
i en Rigdom af forskellige Forbindelser eller Variationer.
Derfor er det for det Andet Religionspsykologiens Opgave
at opsøge og beskrive de vigtigste Typer for Religion og religiøst
Liv. Dersom nu saaledes Motiver og Typer ere paaviste, er det
den \textit{etiske} Undersøgelses Sag at afgøre, hvilken Værdi Religionen
som ejendommelig Form for menneskeligt Aandsliv har for Menneskelivet
i dets hele, store Sammenhæng. Maalestokken er her
den samme, som vi ogsaa ellers anvende, naar vi drøfte en menneskelig
Egenskabs eller Handlings Værdi. Vi undersøge deres
Forhold til visse Grundværdier, i sidste Instans til en eneste
Grundværdi, som for os er den afgørende. Hvilken denne Grundværdi
er, er det Etikens Sag at udfinde. Hvis der er Uenighed
om den Grundværdi, der bestemmer, hvad vi i de enkelte Tilfælde
kalde værdifuldt, saa vil der ogsaa være Uenighed om Vurderingen
af Religionen. Spørgsmaalet vilde naturligvis være
afgjort paa Forhaand, naar man gik ud fra, at Religionen (d. v. s.
den Religion, man selv har eller mener at have) selv er Grundværdien,
efter hvilken Alt skal maales. Men saa vilde der heller
ikke bestaa noget religiøst Problem eller være Mulighed for
eller Trang til nogen Religionsfilosofi, ligesaa lidt som der vilde
være Mulighed for eller Trang til nogen Erkendelsesteori, naar
% --- p. 60 ---
\setcounter{footnote}{0}%
\opage{60}man gik ud fra, at en eller anden Religion indeholdt den absolute
Sandhed, efter hvilken Alt, hvad man ellers vilde kalde
sandt, skulde maales. Et Problem opstaar der derimod her, saasnart
man bliver sig bevidst, at man ved Anerkendelsen af Religionen
forudsætter en vis Maalestok for gyldig, en Maalestok,
som ikke er hentet fra selve Religionen alene, og især, naar
man opdager, at man paa det religiøse Omraade anvender en
Maalestok, der er forskellig fra den, man anvender paa andre
Omraader.

Af de tre Synspunkter, fra hvilke det religiøse Problem kan
og maa behandles, er det det psykologiske, der i den nyeste Tid
især tildrager sig Opmærksomheden. Det er naturligvis ikke
noget Nyt, at Religionen opfattes som et psykologisk Fænomen.
Enhver dybere gaaende religiøs Bevægelse maatte føre dertil,
især ved at gøre det tydeligt, at det kommer an paa ikke blot
at staa i Forestillingsforhold til det Guddommelige, men ogsaa
at bevæges af det i inderlig Følelse. Mere og mere henpeges
der fra religiøs Side til, at Religionens sidste Grundlag er menneskelig Livserfaring. Religionspsykologien anerkender fuldtud
dette, og den stiller sig da den Opgave at beskrive denne Erfaring.
Den gaar herved ud fra, at skal Beraabelsen paa Erfaring
have nogen Mening, maa man ved Erfaring forstaa umiddelbar
Oplevelse, og det kommer da an paa at bestemme, hvad
og hvor Meget der umiddelbart kan opleves. Det gælder at
skelne imellem, hvad der virkeligt er umiddelbart givet, og hvad
der føjes til som Forklaring. Forklaringen hentes fra en Forestillingskreds,
som ligger Mennesket nær. Og fra Forklaringen
maa igen skelnes Udtryksmaaden. Men disse tre Ting, Oplevelse,
Udtryksmaade og Forklaring, er det særligt ved den religiøse
Erfaring, der væsentligt er Følelseserfaring, overordentligt vanskeligt
at holde ude fra hverandre. Religionspsykologisk kommer
det da an paa at opsøge simple Exempler, hvor Forskellen
er tydelig, og ved deres Hjælp at forklare sig de mere indviklede
Tilfælde. Men saa Meget vil strax være klart, at der er en
lang Vej fra en umiddelbar religiøs Erfaring til de ofte meget
subtile Dogmer eller til de digterisk udformede Myter og Legender,
som de historisk foreliggende Religioner fastholde i deres
% --- p. 61 ---
\setcounter{footnote}{0}%
\opage{61}Overleveringer. Og der maa ganske særligt skelnes mellem
Oplevelser, som fremkaldes ved saadanne overleverede Forestillinger
under Forsøget paa at omsætte dem til selvvunden Erfaring,
og Oplevelser, som paa første Haand føre til visse bestemte
Forestillingers Dannelse. Kun denne sidste Art er oprindelig religiøs
Erfaring. Og selv ved den maa der endnu spørges, om
ikke den ene og samme Oplevelse kunde udtrykkes og forklares
ved andre Forestillinger og i andre Former, end det faktisk sker.
Kun hvis en saadan Kontraprøve kunde gøres, og hvis den saa
førte til et negativt Resultat, vilde Erfaringsbeviset være ført.

Jeg tror, at denne psykologiske Metode vil faa en meget
stor Betydning under den fortsatte religiøse Diskussion. Det er
med fuld Ret, man fra religiøs Side henviser til personlig Livserfaring
som sidste Kriterium overfor den ofte meget overfladiske
Beraabelse paa Fornuftens eller Videnskabens formentlige Resultater.
Men hvis man ikke behandler selve Begrebet Erfaring
med streng Kritik, sondrende bestemt mellem det Givne og dets
Fortolkning, saa indfører man en Dogmatisme, der er langt værre
end den gamle rationalistiske Dogmatisme, man fornemt ser ned
paa, saavist som det er værre at bedrage sig selv og Andre med
Hensyn til sit Livs inderligste Forhold end at begaa logiske
Fejl i et indviklet Ræsonnement. I hvert Tilfælde kan Psykologien
ikke lade sig hindre i at forsøge at analysere de ved
første Øjekast umiddelbare og usammensatte Fænomener, man
beraaber sig paa. Der ligger her et stort Materiale tilrede i
den religiøse Bevidstheds Udtalelser til forskellige Tider og under
højst forskellige Kulturforhold, tiltrækkende og forfriskende
især for dem, der have faaet nok af den hidtil gængse Maade at
behandle de religiøse Spørgsmaal paa, baade hos Angriberne og
hos Forsvarerne.

Det er dette, jeg efter bedste Evne har søgt at indskærpe og
gennemføre i min Religionsfilosofi, hvis mest omfattende Afsnit
netop er psykologisk. Jeg har da ogsaa havt den Tilfredsstillelse,
at en teologisk Kritiker (Pastor L. P. Larsen i Studenterhjemmets
Julebog 1901) har erklæret, at min Metode er den ene rette.
Men des større maa da min Forundring være, naar jeg ser, at
denne samme Kritiker bestrider mine Resultater --- ved Hjælp
% --- p. 62 ---
\setcounter{footnote}{0}%
\opage{62}af en ganske anden Metode end den, han saaledes har anerkendt!
Ti han bruger Begrebet Erfaring netop i den vage og vide Betydning,
i hvilken det ofte bruges i daglig Tale. Han stiller
gennemgaaende Spørgsmaalet saaledes: have visse overleverede
Dogmer Værdi for Livet? Men en gennemført religionspsykologisk
Undersøgelse gaar ogsaa den omvendte Vej og spørger:
forsvinder der virkeligt Livsværdier, fordi man ikke antager visse
Dogmer, eller fordi man sætter poetiske Symboler i Stedet for
Dogmerne? Jeg selv tør (uagtet jeg ikke lægger Skjul paa mit
eget personlige Standpunkt i denne Sag) ikke bestemt besvare
dette Spørgsmaal; jeg har udtrykkeligt fremhævet Muligheden
af, at der her kunde være et Værditab. Denne Mulighed kan
ingen Videnskab modbevise, og derfor grunder Antagelsen af
Værdiens Bestaaen sig paa en Tro, selvom man ikke udstrækker
denne Antagelse udover den menneskelige Verdens Grænser. ---

Paa et Omraade, hvor endnu --- og maaske til evig Tid ---
forskellige Retninger brydes, vil det være af stor Betydning at
undersøge, om den anvendte Fremgangsmaade er den eneste,
der kan være Tale om. Derfor har jeg ønsket at fremdrage en
Række religionsfilosofiske Forsøg, af hvilke de fleste først ere
fremkomne efter Afslutningen af min Bog, for at se, om der er
andre Synspunkter at anlægge end de tre nævnte (det logiske,
det psykologiske og det etiske), og specielt for at se, om den
forsøgte Bestemmelse af religiøs Erfaring og af Religionens
Grundtanke skulde trænge til Supplering eller Berigtigelse.

Ikke blot hvad Metoden, men ogsaa hvad Resultaterne angaar,
vil en saadan Sammenligning med andre Forskeres samtidige
Arbejder kunne være af Betydning. Jeg tillader mig da
at minde om, at de Resultater, til hvilke jeg mener at være
naat, hovedsageligt ere følgende tre. --- Det religiøse Problem
gaar ind under det almindelige Vurderingsproblem, og Metoden
for dets Behandling vil derfor mere og mere blive etisk Vurdering
ved Siden af psykologisk Beskrivelse. Ingen af disse to
Fremgangsmaader kan undværes, da man jo kun kan vurdere
nøjagtigt, hvad der foreligger beskrevet nøjagtigt. --- Den religiøse
Trang kan betegnes som en Trang til at hævde Værdiens
Bestaaen, og den religiøse Erfaring bestaar af, hvad Mennesket
% --- p. 63 ---
\setcounter{footnote}{0}%
\opage{63}har oplevet angaaende Forholdet mellem de Værdier, det kender,
og den Virkelighed, det kender. De specielle religiøse Antagelser
angive --- hvad enten de fremtræde i Mytens, Dogmets,
Legendens eller Symbolets Form --- de Maader, paa hvilke Menneskene
have søgt dels at udtrykke, dels at forklare, hvad de
saaledes have oplevet. --- Der er i Aandslivets Vilkaar i den
nyere Tid Intet til Hinder for, at en saadan Trang kan vedblive
at ytre sig og at finde Tilfredsstillelse, selvom det mere og
mere maa ske i andre Former end i tidligere Tider, --- og selvom
det ikke er nødvendigt, at den ytrer sig hos Alle. Det er en
religionshistorisk Hypotese, at Sætningen om Værdiens Bestaaen
er den religiøse Grundtanke, d. v. s. at det er den, der udtrykker
den gennemgaaende Tendens i de historisk foreliggende
Former for Religion. Men Sætningen om Værdiens Bestaaen
kan ogsaa opstilles som en filosofisk Hypotese. Den betegner en
Livsanskuelse, der kan gøre Fordring paa at være \textit{forenelig med}
en strengt gennemført videnskabelig Verdensanskuelse, selvom
den ikke kan \textit{udledes} af denne. Livsanskuelsen staar hos Alle
med en vis Selvstændighed overfor Verdensanskuelsen. Der sker
en stadig Indvirkning af denne paa hin; men ingensinde er
Livsanskuelsen et blot Produkt af den videnskabeligt begrundede
Verdensanskuelse. Ti Livsanskuelsen indeholder et Plus,
der skriver sig fra Personlighedens Vurdering af den Tilværelse,
den søger at forstaa ved Videnskabens Hjælp. Nu kan der vel
indenfor Vurderingen findes større eller mindre Konsekvens,
og de Synspunkter, der lægges til Grund ved Vurderingen,
kunne være mere eller mindre naturlige og væsentlige. Men
paa Vurderingens Omraade ere vi i langt højere Grad nødte til
at være Empirikere, end hvor det gælder Forstaaelse og Forklaring.
Det Irrationale (i matematisk Betydning), som vor
Tænkning stedse maa lade staa tilbage, har et større Omfang
paa Vurderingens end paa den rene Teoris Omraade. Ud over
\textit{Muligheden} af en Livsanskuelse, der bygger paa Grundtanken
om Værdiens Bestaaen, kan Filosofien ikke føre os. Grundmotivet
til enhver Livsanskuelse udspringer af den praktiske Erfaring
om Værdiernes Skæbne i Virkelighedens Verden; og dersom
denne Erfaring fører til Tro paa Værdiens Bestaaen, saa
% --- p. 64 ---
\setcounter{footnote}{0}%
\opage{64}vil denne Tro ikke behøve at staa som en mystisk Hypotese,
blot Genstand for passiv Kontemplation, men den vil naturligt
blive en Arbejdshypotese, idet den ansporer til at finde nye
Værdier, naar de gamles Tid er forbi, eller til maaske endog selv
at frembringe saadanne. Den giver da Mod til at arbejde baade
paa Tankens og paa Virkelighedens Omraade for at vinde fremskridende
Bekræftelse. Maaske ville ikke alle Individer af deres
Erfaringer føres til at nære en saadan Tro; det vil kunne være
tilstrækkeligt, om den holder sig i Slægten. Ti Slægtens Liv
behøver for at være rigt og kraftigt mange forskellige Individualiteter,
maaske endog en stadig Kamp mellem aandelige Retninger.
Mange Individer kunne være saa optagne af Tankens
og Virkelighedens specielle Arbejde, at de ikke faa Tid, Kraft
og Interesse tilovers til at danne sig nogen i bestemte Tanker
udtrykt Livsanskuelse. Der kan her, som i andre Henseender,
finde en Arbejdsdeling Sted. Hvorfor skulde Arbejdsdelingen
være indskrænket til Livets rent ydre Forhold? Er der ikke
Grund til at antage, at den griber ind i Menneskenes inderligste
Anliggender? Livet vil blive rigere derved, ligesom Kunstens
Verden bliver rigere, jo mere den enkelte Kunstner følger
sin individuelle Trang og Evne. Selve Sætningen om Værdiens
Bestaaen kræver maaske endog, at ikke Alle tro paa den eller
tænke paa den.

Disse sidste Bemærkninger lægger jeg særlig Vægt paa
overfor en Indvending, der er bleven rettet mod min Bog: at
den fremstiller en rent subjektiv Opfattelse. Jeg har ikke lagt
Skjul paa mit eget personlige Stade overfor de Spørgsmaal, jeg
drøfter; men Hovedsagen for mig har været at undersøge de
Muligheder, der under Nutidens Kulturforhold frembyde sig i
religiøs Henseende, og her var det ikke nok at holde sig til et
enkelt Individ. Det er en sammenlignende Undersøgelse, jeg
har villet anstille. Jeg har villet foretage en Sammenligning
mellem det Religiøses Stilling i Fortid og i Nutid, og mellem
dets Betydning for forskellige Personligheder og Standpunkter
i Nutiden. Det drejer sig i en Religionsfilosofi ikke om vedkommende
Filosofs egen Tro, men om de Muligheder for Tro,
han finder i Slægtens Aandsliv, saaledes som det nu kan føres
% --- p. 65 ---
\setcounter{footnote}{0}%
\opage{65}i Klarhed og i Dybde. Mindre og mindre har en Filosof Ret
til at antage, at Slægtens Psykologi er udtømt med hans egen
individuelle Psykologi.

En beslægtet Indvending er rejst mod min Fremstilling af
en energisk tysk Forsker, som jeg har den Ære at kunne kalde
min personlige Ven. I en udførlig Anmeldelse af min Bog i
Tidsskriftet »Das freie Wort« (1902) har FERDINAND TÖNNIES
kritiseret min Opfattelse af det religiøse Problem paa en Maade,
der forekommer mig at kunne belyse Forholdet mellem Religionshistorie
og Religionsfilosofi.

»For Guyau«, siger han, »og jeg tilstaar, at jeg deri slutter
mig til ham, er Religionen i første Linie Genstand for Socialpsykologi
eller Sociologi. For Høffding hører den næsten fuldstændigt
til Individualpsykologien, der jo falder sammen med
det, man plejer at forstaa ved Psykologi overhovedet. Han kalder
Buddhismen og Kristendommen »de højere Folkereligioner«:
i dem skal det indre psykologiske Drama fremtræde som det
egentlige Verdensdrama\dots{} Nu er det jo umiskendeligt, at i
begge disse højst vigtige religiøse Systemer har den individuelle
Trang og den individuelle Tænkning faaet en overvældende Betydning.
Men hvorfra stammer denne Sjælens Kamp? Er den
en oprindelig og væsentlig Bestanddel af Religionen som Folketro
og Folkesæd, eller dog af den religiøse Følelse, der »bærer«
Tro og Sæd? \textit{Eller} hidrører den, netop hvad dens renere og
højere Indhold angaar, fra en anden Kilde, nemlig fra en Tænkning,
en Filosoferen, der efter sit Væsen er fremmed for Folketro
og Folkesæd,\dots{} skønt den ganske vist i hine to Folkereligioner
har opnaaet en Udligning med populære og vildtvoxende
(naturwüchsige) Forestillinger og Følelser, ja, i Kristendommen
endog har sluttet et evigt Forbund med dem?«

Konsekvensen af denne Tankegang synes at maatte blive
den, at Buddhismen og Kristendommen ikke egentligt ere Religioner.
Ja, man maatte vel egentligt gaa videre. Gives der
nogensomhelst positiv Religion, i hvilken der ikke kan paavises
Elementer af Tænkning? De mest elementære, for ikke at sige
de mest primitive Religioner bygge paa Animismen, Menneskets
mest elementære Filosofi, og saaledes vil paa alle Stadier af Ud%
% --- p. 66 ---
\setcounter{footnote}{0}%
\opage{66}viklingen et Folks Religion bære Præget af dets forudgaaende
»Filosofi«. Historikerne have undertiden temmeligt romantiske
Forestillinger om den Umiddelbarhed og »Naturwüchsigkeit«,
hvormed Menneskene gaa op i deres Kultus, Myter og Dogmer.
Der foregaar, som Religionshistorien selv viser, stadige Kombinationer,
Suppleringer og Metamorfoser, hvorved et mere eller
mindre dybtgaaende Reflexionsarbejde medvirker. Men hvad der
for mig er Hovedsagen: det er misvisende her blot at tale om
Tænkning og Filosofi. Det er først og fremmest personlig Livserfaring,
der øver sin Indflydelse ved disse Ændringer, der ofte
ere Nydannelser. De Personligheder, som have været afgørende
for de højere Folkereligioners Tilblivelse og Udvikling, have
først og fremmest lagt deres egne Erfaringer om det Værdifulde
og dets Skæbne i Verden til Grund ved de Ændringer, de førtes
til at foretage i det Overleverede. Ingen Religion laver sig
selv. Religionen opstaar og existerer som Religion kun i Menneskers
Sind og Tanke. Det maa have været visse bestemte Mennesker,
der have formet Kultus, Myte og Dogme, ligesom det maa
have været visse bestemte Mennesker, der have digtet »Folkepoesien«,
og ligesom »Samfundet« til enhver Tid bestaar af individuelle
Mennesker. Forsaavidt maa al Psykologi være individualistisk,
selvom den med Iver anvender den sociologiske
Metode (ved Siden af andre psykologiske Metoder), hvad der
ganske særligt i Religionspsykologien er Anledning til at gøre. ---
I Slutningen af sin Anmeldelse udtaler Tönnies iøvrigt, efterat
have gentaget, at han i Religionen især ser en Genstand for
Socialpsykologi: »Ich bekenne meinerseits, dass mich Høffding
gelehrt hat, die Tragweite der individuellen Entscheidung besser
als bisher zu würdigen«. ---

Kun naar man lægger Vægt paa det subjektive Element i
Religionen, bliver det muligt at finde en indre Kontinuitet mellem
de skiftende Religioner og Livsanskuelser. Og paa dette
Element beror ogsaa den Betydning for Slægtens fortsatte Liv,
en Religion kan have, selvom dens Tid engang er forbi. Det
er de store menneskelige Skatte, den religiøse Aand har formaaet
at hæve i Meneskenes Indre, der udgøre den kosteligste
Arv, Fortidens Religioner efterlade. Et Digt som BUNYAN'S
% --- p. 67 ---
\setcounter{footnote}{0}%
\opage{67}»Pilgrims Progress«, der er blevet kaldet Protestantismens divina
comedia, har saaledes faaet en blivende Betydning (idetmindste
for den engelske Race) ud over den bestemte Konfessions Grænser,
indenfor hvilken det blev til. Og Grunden er den, at det
udtrykker, hvad der bevægede sig i et af Livets Modsætninger
dybt grebet Menneskes Sjæl. En udmærket Literaturhistoriker
har med Rette sagt om det: »Alt hvad der er bedst og mest
karakteristisk i Bunyans Skrifter, fremgaar af det indre Drama,
i hvilket de Optrædende vare Gud, Satan og en ensom Menneskesjæl.
Institutioner, Kirker, Ordinanser, Ritus og Ceremonier
hjælpe ham kun lidet eller slet ikke\dots{}. Lad Enhver med
Troskab fastholde sine inderligste Erfaringer i Livets store Anliggender,
og hans Ord ville vække mangefold Genlyd i andre
Sjæle; \textit{det dybeste Samfund findes ikke i Institutioner eller}
\textit{Korporationer eller Kirker, men i det ensomme} \textit{Hjertes} \textit{Hemmelighed}.«
(DOWDEN: \textit{Puritan and Anglican}. London 1900.
p. 234). I min Religionsfilosofi har jeg ogsaa peget hen paa,
hvorledes Sansen for de rent subjektive Elementers Betydning
paa det religiøse Omraade kan bane Vejen for en dybere Fællesskabsfølelse
end den, der betinges ved Sammenslutning om fælles
Dogmer. Til et og samme Dogme eller til en og samme Kultus
kan der svare højst forskellige personlige Stemninger og Tilstande.
Denne Forskelligartethed er et Bevis paa det subjektive
Elements Selvstændighed. Forholdet mellem de objektive og de
subjektive Elementer (eller maaske bedre: mellem de sociale og
de individuelle Elementer) i Religionen er stedse varierende,
saalænge der er Liv i Religionen. I de historisk udformede Religioner
er det de objektive, sociale Elementer, der stærkest falde
i Øjnene. Og hvis Udtrykket »positiv Religion« er at forstaa i
Analogi med Udtryk som »positiv Ret«, menes derved en Religion,
der optræder som Institution, med en udformet Kultus
og en fastslaaet Overlevering. Psykologisk set maa det dog
altid være det Væsentlige, hvad der motiverer Menneskets Tilslutning
til og Fastholden af det saaledes Fastslaaede, og hvad
der under visse Forhold fører til at foretage Ændringer deri.
Dette subjektive Elements Betydning kan stige saaledes, at det
objektive Element saa godt som elimineres. Vi komme da til
% --- p. 68 ---
\setcounter{footnote}{0}%
\opage{68}en Grænse, hvor der kan være Strid om, hvorvidt Begrebet Religion endnu kan anvendes. Det kan let blive en Strid om Ord.
Det kan føre til at opkaste dybsindige Spørgsmaal i Lighed med
det, en Mand stillede: »Mon nu den Stjerne, Astronomerne kalde
Uranus, ogsaa virkeligt er Uranus?« --- Dersom man lægger
Vægten paa Grundmotiverne, ikke paa de ydre Kendemærker,
vil det være berettiget at udstrække Begrebet Religion ud over
de saakaldte positive Religioners Sfære. Religion bliver da et
Slægtsbegreb, indenfor hvilket Religion som positiv Religion,
bygget paa Overlevering og forsynet med faste Kultusformer,
kun betegner en særegen Art. Derved undervurderes Modsætningerne
ikke; ti indenfor et og samme Slægtsbegreb kunne de
stærkeste Modsætninger fremtræde, ligesom der i Virkeligheden
kan føres en Kamp paa Liv og Død mellem Arter, der høre til
samme Slægt. ---

Efter disse indledende Bemærkninger gaar jeg nu over til
at omtale nogle religionsfilosofiske Arbejder fra den nyeste Tid
for derigennem at give et kritisk Supplement til mit eget Forsøg.
Jeg vil vælge saadanne Arbejder, der baade kunne siges
at være typiske for vor Tids Religionsfilosofi og at være lærerige
for dem, der interessere sig for det religiøse Problem. ---
Der er dog nogle Forfattere, som det i og for sig kunde være
naturligt at fremdrage her, men som jeg hellere vil omtale ved en anden Lejlighed, fordi de bedst ville forstaas i Sammenhæng med den nyeste Tids Filosofi idethele. Dertil høre især Guyau
og Nietzsche, som jeg om ikke lang Tid vil faa Lejlighed til at
omtale i et Skrift om moderne Filosofer\footnote[1]{Se \textit{Moderne Filosofer} (Kbhvn. 1904). p. 107---147.}, jeg har under Arbejde,
og til hvilket nærværende Afhandling ogsaa kan betragtes som
Supplement,

I Afhandlingens første Afsnit behandler jeg nogle systematiske
Forsøg; i dens andet Afsnit fremdrager jeg en Række
speciellere religionspsykologiske Undersøgelser, især fra den
amerikanske Literatur.

% --- p. 69 ---
\setcounter{footnote}{0}%
% CHECK p. 69: centred line: heading or verse
\opage{69}

\section*{I.\\ SYSTEMATISKE FORSØG.}
\streg

\subsection*{1. HERMANN SIEBECK.}

SIEBECK (f. 1842), der virker som Professor i Filosofi i
Giessen, har gjort sig bekendt ved et udmærket Arbejde over
Psykologiens Historie i Oldtid og Middelalder og ved sine æstetiske
Studier. Hans Hovedværk er den i Aaret 1893 udkomne
Religionsfilosofi. Oprindeligt udgaaet fra Herbarts Skole har
Siebeck senere modtaget Paavirkninger af Lotze; men især mærkes
Virkningen af ivrige Kantstudier i hans Hovedværk.

Religionsfilosofiens Opgave er for Siebeck stillet ved Forholdet
mellem Religion og Kultur. Den moderne Kultur er tilbøjelig
til at betragte sig selv som Formaal; Religionen kan
den, naar det kommer højt, kun anerkende som en af sine ufuldkomne
Forgængere. Fra religiøs Side paastaas derimod, at Kulturlivet
--- ligesom Livet og Arbejdet i Erfaringens Verden
idethele --- kun kan være et Middel for et højere, »oververdsligt«
Formaal. Den Strid, der saaledes opstaar, er det Religionsfilosofiens
Sag at afgøre. Den er altsaa ikke en religiøst
stemt Filosofi, som Nyplatonismen i Oldtiden eller som Schellings
senere Lære; den er heller ikke, som Middelalderens Skolastik,
en Fremstilling af en vis bestemt Religion i filosofisk Form.
Den betragter Religionen i dens forskellige Skikkelser som en
Kendsgerning i det menneskelige Liv og undersøger, hvilken Betydning
denne Kendsgerning har for det menneskelige Livs Udvikling.
Den søger at bestemme, paa hvilket Omraade af menneskeligt
Aandsliv Religionen særligt hører hjemme, og hvilken
Indflydelse den ud derfra øver paa andre Sider af Livet. Filosofien
undersøger her Religionen paa lignende Maade, som den
i andre af sine Discipliner undersøger Videnskaben, Kunsten,
Retten og Moralen.

Det maa ifølge Siebeck væsentligt være en indirekte Metode,
Religionsfilosofien kan anvende. Den maa prøve, om selve
Kulturlivets Bestaaen og Udvikling ikke blive værdiløse og selv%
% --- p. 70 ---
\setcounter{footnote}{0}%
\opage{70}modsigende, naar det afskærer sig fra de supplerende Synspunkter og Motiver, Religionen yder. En positiv Bestemmelse
af den højere Verden, til hvilken Religionen vil aabne Adgang,
kan Filosofien ikke give. Længere end til at paavise Nødvendigheden
af en Supplering kan den ikke naa; men den kan
prøve, om de Former, i hvilke den religiøse Bevidsthed har
søgt at udtrykke det Guddommelige, egne sig til at yde den
Supplering, en Kritik af Kulturudviklingen fører til at fordre.
Paa Forhaand udtaler Siebeck som sin Overbevisning, at der
i den angivne Betydning af Ordet kan gives en Paavisning af
»Normaliteten af den religiøse Bevidstheds væsentlige Indhold«.

Det vil for Religionsfilosofien, naar dens Opgave bestemmes
paa denne Maade, være naturligt at holde sig til de højeste
Former for Religion, da kun de blive af Betydning for Forholdet
mellem Religion og Kultur i det moderne Aandsliv. Siebeck
skelner mellem tre Hovedformer for Religion: Naturreligion,
etisk Religion (»Moralitätsreligion«) og Forløsningsreligion.
I den første Form opfattes Guderne som Naturmagter, og Forholdet
til dem bestemmes væsentligt ved Selvopholdelsestrangen.
I den anden Form have Guderne et Kald at udøve, og
deres Dyrkelse motiveres ved den Værdighed, dette Kald giver
dem. Den tredie Form, Forløsningsreligionen, der fremtræder
typisk i Buddhismen og i Kristendommen, er karakteriseret ved
den Betydning, profetiske Personligheder faa. Myte og Kultus,
som raade i de tidligere Former, træde her tilbage for de oprindelige
Personligheder. Der er i al selvstændig Personlighed,
især i den geniale Personlighed, noget Irrationelt, noget Nyt,
der ikke finder sin fulde Forklaring ved, hvad der gaar forud,
og som maaske kun langsomt finder Anerkendelse og Tilslutning.
Men allermest gælder dette om de store Personligheder,
fra hvem de nye Værdibegreber udgaa, hvis Gyldighed gennem
lange Tider skal proves i Slægtens Livserfaringer. Paa dem
passer \textit{Begrebet Aabenbaring}. Dette Begreb er, ifølge Siebeck,
selv et Værdibegreb, d. v. s. et Begreb, hvis Gyldighed ikke
beror paa den Maade, paa hvilken dets Indhold kan paavises i
Virkeligheden, men paa Indholdets praktiske Betydning for Livet.
Dette ses ved at sammenholde forskellige Udtalelser af
% --- p. 71 ---
\setcounter{footnote}{0}%
\opage{71}Siebeck. Et Sted (p. 192) hedder det, at det religiøse Indhold,
der kundgøres af de store Personligheder, kaldes Aabenbaring
formedelst dets »oververdslige« Karakter. Og hvad Begrebet
»det Oververdslige« angaar, hedder det et tidligere Sted (p. 18):
»Begrebet det Oververdslige kan først udspringe paa Grundlag
af en Verdensanskuelse, der ikke blot tager Verden med dens
Tilværelse og Indhold som given Kendsgerning, men i Kraft af
en af det Givne uafhængig, spontan aandelig Virken underkaster
denne Kendsgerning en \textit{Vurdering}, hvis Princip er hentet fra
selve Aandens Indre.«

Det er egentligt først indenfor denne tredie Form for Religion,
at Religion og Kultur træde i Modsætning til hinanden.
Forløsningen, der i Buddhismen forkyndes paa mere negativ, i
Kristendommen paa mere positiv Maade, bestaar i Personlighedens
Myndighed overfor Verden; den er given med Tilegnelsen
af den Tanke, at Verden kun har Værdi paa Grund af det
personlige Livs Existens. Personligheden er paa én Gang fremgaaet
af Verden og selvstændig overfor Verden; den vil ikke
blot forstaa Verden og hævde sig i den og overfor den, men
ogsaa arbejde i den for det, der fra det personlige Livs Standpunkt
staar som det Højeste. Deraf kommer det, at det falder
den religiøse Tro saa naturligt at opfatte selve Verdensgrunden
i Analogi med det Personlige. For Troen er Personligheden det
første Objekt, medens den er Videnskabens allersidste Objekt,
fordi den er den mest sammensatte og tillige mest centrale Genstand,
Erkendelsen kan beskæftige sig med.

Paa dette højeste Standpunkt fremkommer der en Modsætning
mellem Tro og Viden. I Kraft af visse Fænomeners Værdi
er den religiøse Bevidsthed tilbøjelig til at stille dem udenfor
den naturlige Aarsagssammenhæng, medens paa den anden Side
Videnskaben er tilbøjelig til at erklære sig selv for at være det
tilstrækkelige Grundlag for Livsanskuelse og Livsførelse. Siebeck
hævder overfor den første Paastand, at Undertro ikke er et nødvendigt
Element i den religiøse Tro. Jo klarere den religiøse
Bevidsthed bliver over sig selv, desmere vil den --- mener han ---
indse, at Undertroen ikke forherliger, men netop begrænser Gud,
idet den sætter Naturen som en Skranke, der skal gennembrydes
% --- p. 72 ---
\setcounter{footnote}{0}%
\opage{72}af Gud. Overfor Videnskabens Paastand om at være tilstrækkeligt
Grundlag for Livet hævder han, at selve denne Paastand
kun kan være en Trossag, da et videnskabeligt Bevis for den
ikke er muligt, og da Afvisningen af Religionens Nødvendighed
sker ved en Personlighedsakt, der motiveres ved den Forudsætning,
at med selve Videnskaben er det højeste Gode givet. Selve
Trangen til videnskabelig Forsken hviler paa denne ubeviselige
Forudsætning.

Der er i al Religion et subjektivt og et objektivt Element.
Det subjektive Element, som man kunde kalde Religiøsiteten ---
ytrer sig i Stemning, det objektive Element, Religion i Modsætning
til blot Religiøsitet, er Sag, Værk, Ydelse, synlig og
haandgribelig Skikkelse for den subjektive Stemning. Religiøsitetens
 Dybde er Værdimaaler for det objektive Religionsvæsen ---
Tilstrækkelighed og Livskraft; men den subjektive Stemning
vilde aldrig have udviklet sig uden under de objektive Elementers
 Indflydelse, saa at Religiøsitet aldrig i Tiden gaar forud
for Religion. Der er et stadigt Vexelvirkningsforhold mellem
dem. Men jo mere et Kultussystem organiseres, desmere udelukkes
nye Stemninger og nye Tanker fra det religiøse Liv-Det
er Kultus, der stedse gør den største Modstand mod Religionsfornyelse.
Jo højere et Udviklingstrin Religionen naar, des
mere udprægede individuelle Variationer vil det religiøse Liv
frembyde, saa at de store menneskelige Forskelligheder ikke
blot kunne gøre sig gældende i Forholdet mellem Religionerne,
men ogsaa indenfor en og samme Religion. ---

Dette er det væsentligste Indhold af den første Afdeling af
Siebecks Værk, den beskrivende Del, som afhandler Religionens
Væsen og Udvikling. Idet han nu gaar over til at drøfte Religionens
Sandhed, udtaler han, at det her ikke drejer sig om
nogen enkelt Religions Sandhed, men om Gyldigheden af det
religiøse Princip overhovedet. Denne Gyldighed kan kun begrundes
ved, at Aabenbaringen --- i den ovenfor angivne Betydning
af Ordet --- griber ind i den menneskelige Kultur og, skønt
den ikke selv udspringer af denne, dog yder et nødvendigt Supplement
til den.

Paa dette Punkt er det nu ikke ganske let at forstaa Sie%
% --- p. 73 ---
\setcounter{footnote}{0}%
\opage{73}becks Tankegang. I sin beskrivende Del (som iøvrigt ikke er
blot beskrivende) har han bestemt skelnet mellem Vurdering
og Aarsagsforklaring, saaledes at disse ikke behøve at træde i
Modsætning til hinanden, men betegne forskellige Synsmaader
for Tilværelsen. Men nu opstiller han den Paastand, at en
etisk Handling, en Handling, der gaar ud paa at hævde og gennemføre
det Værdifulde i Virkelighedens Verden, kun bliver
mulig, dersom denne i sidste Instans er underlagt en etisk
Magt, og han finder heri en Begrundelse for Gudsbegrebet. Enhver
etisk Handling hviler, mener han, paa den Forudsætning,
at det er muligt at virkeliggøre de etiske Ideers Indhold indenfor
Natursammenhængen, og denne Forudsætning er en religiøs
Forudsætning. Om den etiske Magt, der her forudsættes, kan
man kun ad Analogiens Vej gøre sig en Forestilling, men en
saadan Brug af Analogien anser Siebeck for berettiget, naar den
undgaar baade »Deismens« dualistiske Forhold mellem Gud og
Verden og »Panteismens« Ophævelse af enhver Forskel mellem
dem; derved naar man til en »Teisme«. --- Der savnes her hos
Siebeck en nærmere Undersøgelse baade af Analogiens Betydning
og Berettigelse og af Forholdet mellem Vurdering og Aarsagsforklaring.
Hvad det første Punkt angaar, er dette saa meget
des mærkeligere, som Siebeck selv tidligere i sin interessante
Afhandling \textit{Die metaphysischen Systeme in ihrem gemeinsamen}
\textit{Verhältnisse zur} \textit{Erfahrung} (i andet Bind af Vierteljahrsschrift
für wissenschaftliche Philosophie) har paavist Analogiens Betydning
for enhver Spekulation, der gaar ud over Erfaringen. Det
stod nu blot tilbage at anvende denne Betragtning paa de religiøse
Forestillinger. Hvad det andet Punkt angaar, fremhæver
han vel Vanskelighederne ved at forene den teleologiske og den
kausale Betragtningsmaade i fuld indbyrdes Harmoni, men han
mener, at den religiøse Bevidsthed ikke betragter Verden som
en afsluttet Totalitet, men som Del af en højere Tingenes Sammenhæng,
fra hvilken stedse nye skabende Kræfter kunne
strømme ind i den. Jeg er ikke sikker paa, at den religiøse Bevidsthed
her tydes rigtigt; den har i hvert Tilfælde meget ofte
en Tendens til netop at betragte »Verden« som et Hele for sig,
og Siebeck har selv fremhævet dette i sin Kritik af Mirakel%
% --- p. 74 ---
\setcounter{footnote}{0}%
\opage{74}forestillingen. Det er snarere for den videnskabelige Bevidsthed,
at Verden ikke danner nogen afsluttet Totalitet. »Verden« er
ikke noget videnskabeligt Begreb.

Idet Siebeck saaledes anerkender Berettigelsen af at tænke
sig den naturlige Aarsagssammenhæng underordnet et etisk Princip,
overser han ikke de Indvendinger, der stadigt rejse sig mod
denne Betragtningsmaade ved Existensen af Smerte og Ondskab
i en Tilværelse, der skulde have sin Grund i et saadant Princip.
Han ser ganske vist ikke, at denne gamle Indvending særligt
skarpt maa ramme »Teismen«, der tænker sig hint Princip
i personlig Form. Men han indrømmer, at Vanskeligheden
er uovervindelig for den religiøse Bevidsthed, saa meget end
denne trøster sig med, at det Onde maatte være muligt i en
Verden, hvor frie Væsener skulde udvikle sig. Han har givet
det Kapitel i sin Bog, i hvilket dette Spørgsmaal drøftes, Overskriften
»Theodice« og bemærker, at teoretisk er det i et Forsøg
paa Theodice Gud, der skal forsvares, men at hvad det for
den religiøse Bevidsthed gælder om, er at holde sit Haab oppe
trods Anfægtelserne. Man kunde hertil bemærke, at forud for
al »Theodice« gaar et Stadium, hvor det er Mennesket, der skal
retfærdiggøres for Guds Domstol. Saa kommer først senere et
Stadium, hvor Gud stedes for Mennesketankens Domstol; da
begynder man at skrive Theodiceer. Derpaa kommer det tredie
Stadium, hvor ingen af Delene sker, men hvor det tydeligt og
klart gælder en Kamp for at holde Haabet oppe, og hvor de
Forestillinger, man danner sig om, hvad der ligger udover Erfaringen,
udtrykkeligt staa som Symboler for dette Haab.

Men naar det saaledes viser sig, at de Problemer, Videnskaben
ikke løser, heller ikke løses i Religionen, ja netop fremtræde
i denne i den skarpeste Form, saa er den Supplering af
Kulturen, den Afslutning af Verdensanskuelsen, den religiøse
Ide skulde bringe, ikke lykkedes. Nødvendigheden af at optage
dette Supplement kan ikke paavises.

% --- p. 75 ---
\setcounter{footnote}{0}%
% CHECK p. 75: centred line: heading or verse
\opage{75}

\subsection*{2. AUGUSTE SABATIER.}

Det er et Tidens Tegn, at man ogsaa fra teologisk Side begynder
paany at interessere sig for Religionsfilosofi. Denne
Interesse har ellers været saa godt som forsvunden, efterat de
Forhaabninger, man havde næret til den romantiske Spekulation,
vare bristede. Den for faa Aar siden afdøde protestantisk-teologiske
Professor i Paris AUGUSTE SABATIER har i sit Værk
\textit{Esquisse d'une philosophie} de la \textit{religion} \textit{d'après} \textit{la psychologie}
\textit{et} l'histoire (5 éd. Paris 1898) sat sig til Opgave at
behandle det religiøse Problem paa psykologisk og historisk
Grundlag.

I psykologisk Henseende lægger Sabatier Vægten paa, at
hvad der tilstræbes i Religionen ikke er en Løsning af »Verdensgaaden«,
men af »Livsgaaden«. Det er det Gaadefulde og Modsigende
i vort eget Væsen, der afføder den religiøse Trang, og
Verdensforholdene interessere paa det religiøse Standpunkt kun,
forsaavidt de antages at have Indflydelse paa Løsningen af det
indre Livs Gaader. Den religiøse Trang er en Art højere Selvopholdelsestrang,
som vækkes ved Disharmonierne i vort Indre,
især ved Modsætningen mellem Afhængighed og Frihed, eller
mellem Passivitet og Aktivitet. Vort Livs Strøm, der udspringer
fra det Inderste i os, brydes og splittes som en afmægtig Bølge
mod de ydre Tings Skær og Klipper. Det er Grunden til al
Smerte, og da især til den højeste Smerte af alle, den, der betinges
ved Disharmonien mellem vort Ideal og hvad vi kunne
gennemføre i Virkeligheden. En Harmoni mellem disse Modsætninger
naas ifølge Sabatier i Bønnen, den egentlige religiøse
Akt. Her fremtræder nemlig paa én Gang Afhængighedsfølelsen
overfor den hemmelighedsfulde Magt, vi føle os i Forhold til,
og Frihedsfølelsen formedelst den Kraft, der fra denne Magt
tilstrømmer os. Det Væsentlige i Bønnen er ikke de Ord, den
iklædes, eller de Forestillinger, den knyttes til, men den indre
Forbindelse af Hengivelse og Selvvirksomhed, der her gensidigt
støtte hinanden og føre over i hinanden. Religionens Nødvendighed
beror paa, at den udspringer af vore inderste Livsvilkaar.
Derfor vil den vise sig selv der, hvor der er Antipati mod en%
% --- p. 76 ---
\setcounter{footnote}{0}%
\opage{76}hver ydre Religion. Der er et indre Alter, som kan bevares og
holdes i Hævd, selvom man omstyrter alle ydre Altre. Der gives
kun én sand Ateisme: Frivolitet og Blaserethed.

Naar Trangen er stærk og inderlig, er Tilfredsstillelsen allerede
given; den kommer ikke udefra. For Sabatier ligger hele
Religionens Mysterium i Pascals Ord: »Du vilde ikke søge mig
hvis du ikke havde fundet mig«. Svaret, Aabenbaringen er til
Stede i selve Bønnens Mulighed, i den indre Kraft og Higen ---
hvoraf Bønnen udspringer. I samme Øjeblik Mennesket søger
Gud, er det Gud, der søger Mennesket. Aabenbaringen er ikke
en Meddelelse én Gang for alle af uforanderlige Lærdomme, som
det saa blot gælder at holde fast, men den bestaar i den Maade -
paa hvilken Gudsbevidstheden hos den Enkelte og i Menneskeheden
opstaar og udvikler sig.

Dette Aabenbaringsbegreb betegner Sabatier som det psykologiske
eller det kritiske, i Modsætning dels til det mytologiske,
dels til det dogmatiske Aabenbaringsbegreb. Det har sine store
historiske Repræsentanter i de israelitiske Profeter og i Jesus,
men det finder Anvendelse i enhver religiøs Personlighed. Den
religiøse Inspiration er ligesaa hemmelighedsfuld som den poetiske,
men ikke mere mirakuløs end denne. Den er en indre
befrugtende Kraft, der virker i Harmoni med alle aandelige
Kræfter og Love. Den profetiske Inspiration er kun en højere
Grad af, hvad der rører sig i enhver stærkt bevæget religiøs
Følelsestilstand. Maalestokken for Aabenbaringens Værdi ligger
i den Berigelse, Frigørelse og Forhøjelse af det personlige Liv,
den bevirker.

De Forestillinger og Handlinger, i hvilke Religionen finder
Udtryk, variere med Tiderne. Det er Ortodoxiens, selv Reformatorernes
store Fejl, at de ikke have skelnet mellem den religiøse
Erfaring og dens Udtryk. Man har sammenblandet de
psykologiske Fænomener og de teologiske Forklaringer. Ved
ethvert Dogme maa der skelnes mellem den religiøse Erfaring,
der ligger til Grund, og det mere eller mindre ufuldkomne Udtryk,
Reflexionen har kunnet give den. Dette sidste Element er
det mest variable Element i Religionen. Der kan gives tre Arter
af dogmatiske Variationer. Dogmer kunne aftage og for%
% --- p. 77 ---
\setcounter{footnote}{0}%
\opage{77}svinde af Mangel paa Anvendelse (par désuétude), idet den religiøse
Erfaring ikke længere yder det Indhold, de skulle udtrykke.
Saaledes gaar det med Antagelsen af dæmonisk Besættelse og
af en personlig Djævel. Dogmer kunne dernæst faa et andet
Indhold end det oprindelige, idet de fortolkes efter nye Erfaringer.
En saadan Forskydning foregaar uafbrudt. Ved Inspiration,
Inkarnation og Treenighed forstaa senere Slægter ikke
det Samme som tidligere Slægter. Der heldes stadigt ny Vin
paa de gamle Læderflasker. Endeligt kan der opstaa nye Dogmer
som Udtryk for nye Erfaringer. Saaledes fremkom indenfor
Protestantismen Retfærdiggørelsen ved Tro og det almindelige
Præstedømme.

At proklamere Uforanderligheden af Kirkens Dogmer vilde
være at proklamere Uforanderligheden af den Filosofi, der anvendtes
til at udforme dem. Dogmatiken maa holde sig til de
religiøse Erfaringer, hvis Udlægning den er. Og naar der anvendes
visse Begreber for at kunne foretage denne Udlægning,
saa maa den ogsaa anerkende de Resultater, Erkendelsesteorien
kommer til angaaende vore Begrebers Betydning og Gyldighed.


Den religiøse Erkendelse er symbolsk Erkendelse. Til at
udtrykke Følelserne af Tillid, Frygt og Kærlighed benytter den
Billeder, der ere hentede fra menneskelige Forhold. Symbolets
Magt beror paa den nøje Sammenhæng mellem Følelse og Billede.
Den religiøse Følelse skelner fra først af ikke mellem Indhold
og Form; den er sig ikke bevidst, at den udtrykker sig
symbolsk. Det er den kritiske Filosofi, der indskærper Symbolismens
Uundgaaelighed, som baade Ortodoxien og Rationalismen
hver paa sin Maade miskende. --- Symbolerne kunne blive fælles
for store Kredse, og der er Symboler, som maaske aldrig ville
forsvinde fra det religiøse Liv.

Ikke blot Erkendelsesteorien fordrer Sabatier respekteret af
Dogmatiken, ogsaa Erfaringsvidenskabens almindelige Metode.
Han skelner, som vi have hørt, mellem Verdensgaaden og Livsgaaden.
Verdensgaaden er det Videnskabens Sag at drøfte. Verden
maa tages, som Videnskaben viser os, at den er, og Livsgaaden
stilles saa ved Spørgsmaalet om, hvorledes Videnskabens
% --- p. 78 ---
\setcounter{footnote}{0}%
\opage{78}Verdensbillede forholder sig til det aandelige Livs inderste Trang,
--- om det indeholder Midler eller Hindringer for Tilfredsstillelsen
af denne Trang.

Erfaringsvidenskabens Selvstændighed maa ogsaa komme
Psykologien til Gode, og her staa vi (som jeg allerede har gjort
opmærksom paa i min Religionsfilosofi §. 66) ved det vanskeligste
Punkt i Sabatiers Religionsfilosofi. Han har ført alle det
religiøse Livs Fænomener tilbage til religiøs Erfaring, saaledes
som han nærmere beskriver denne, og han har hævdet, at alle
religiøse Forestillinger ere Symboler, som udtrykke denne Erfaring.
Men alligevel antager han, at Gud er Aarsag til den
religiøse Erfaring, og her vilde da foreligge et Tilfælde, hvor
et Erfaringsfænomen fik en teologisk Forklaring, hvad Sabatier
selv udtrykkeligt har afvist Muligheden af, da Gud ikke selv
er Led i Fænomenernes Række og derfor ikke kan bruges til
videnskabelig Forklaring. Tilmed maa jo selve Gudsforestillingen
som alle andre religiøse Forestillinger have en symbolsk
Karakter.

Sabatiers teologiske Udgangspunkt viser sig dog ikke blot
i hans Forsøg paa at forklare den religiøse Erfaring, men ogsaa
i selve hans Beskrivelse af den. Han fremdrager kun en eneste
Type for saadan Erfaring og lægger ikke Vægt paa Variationerne,
skønt han i Principet anerkender deres Mulighed og deres
Berettigelse. Ligesom hos Siebeck føle vi ogsaa her Savnet af
et bredere psykologisk Materiale. Et saadant vilde maaske have ført til en mindre snever Bestemmelse af Religionens Natur.
Dog hindre disse Indvendinger ikke, at man skilles fra den
franske Teologs filosofiske Arbejde med Erindringen om en
ædel, fin og begejstret Mand.

\subsection*{3. RUDOLPH EUCKEN.}

Efter en Række Forarbejder (især \textit{Die Einheit des} \textit{Geisteslebens}
1888, og \textit{Der Kampf} um \textit{einen} \textit{geistigen} \textit{Lebensinhalt}
1896) har denne aandfulde Tænker, der virker i Idealismens
gamle Arnested Jena, udtalt sig om det religiøse Problem i sit
Skrift \textit{Der} Wahrheitsgehalt der Religion (1901).

% --- p. 79 ---
\setcounter{footnote}{0}%
\opage{79}Det er ikke et religionsfilosofisk System, Eucken her har
villet give. Et saadant er efter hans Opfattelse foreløbigt umuligt
paa Grund af de uklare og usikre Forhold paa det religiøse
Omraade; »Under vor Tids aandelige Anarki gives der
ikke noget fast og anerkendt Tilknytningspunkt; al dybere
gaaende Drøftelse maa gaa tilbage til det første Grundlag og
bygge op fra ny derpaa. Saaledes maatte ogsaa jeg gennem en
almindelig Undersøgelse af den menneskelige Tilværelse Skridt
for Skridt arbejde mig hen til det Sted, hvor Religionens Problem
bryder frem for saa strax at fremtræde som Midtpunktet
for al Søgen efter vor Tilværelses Sjæl og Betydning«. Forfatteren
erklærer sig selv for en Søgende, og han henvender sig
til Søgende. Han finder i Tiden en stærk Længsel efter Religion
forbunden med en tydelig Bevidsthed om Utilstrækkeligheden
af Religionens nuværende Form. Den nye Livstype, som
Renaissancen hævdede, og det nye Tankearbejde, som derved
indlededes, førte bort fra Kristendommen. Men den storartede
Kultur, i hvilken denne Livstype og dette Tankearbejde udformedes,
har medført indre Forviklinger, og Troen paa dens Tilstrækkelighed
er rokket. Derved kommer da Tanken frem, om
der ikke i Religionen skulde være Noget, der ikke kan og maa
gaa tabt. Der føles Savnet af et stort Maal, som ikke blot
kunde kaste Lys over Menneskelivets store Kreds, men ogsaa
hæve den Enkelte op over hans ofte smaalige Forhold. Og tillige
trænge de gamle Livsgaader, man betragtede som løste eller
dog som fortrængte, sig frem paany. Især føles der en skarp
Modsigelse mellem Menneskets aandelige Anlæg og dets virkelige
Tilstand. Alt dette fører til en Kamp mellem Religion og
Kultur.

Men hverken Religionen alene eller Kulturen alene raader
over Midler til at afgøre Kampen. Afgørelsen maa ske ved, at
der baade indenfor Religionen og indenfor Kulturen skelnes
mellem det Evige og det Forgængelige, og mellem det Uundværlige
og det Uudholdelige. Der maa foretages en Revision
af selve Grundbetingelserne for vort aandelige Liv, og at foretage
en saadan er Filosofiens Sag. Ikke i den Forstand, at denne
kan konstruere en Religion. Ti Religion er noget Faktisk, Ud%
% --- p. 80 ---
\setcounter{footnote}{0}%
\opage{80}tryk for en real Existens, som kan paavises, naar den er der,
men ikke kunstigt frembringes. Men da det her gælder Paavisning
af en aandelig Kendsgerning, er Paavisningen ikke en
ganske simpel Sag; der behøves aandeligt Arbejde for at finde
Religionens ejendommelige Plads i Livets hele Sammenhæng.
Filosofien kan ikke naa videre her end til at paavise Muligheder,
som først gennem store Personligheder og ofte under
store Omvæltninger i de menneskelige Forhold kunne omsættes
til Virkelighed. Men dersom det kan godtgøres, at Religionen
hænger sammen med vor Tilværelses inderste Grund, saa kan
ethvert Angreb paa den og enhver Kritik af den kun tjene til,
at dens inderste, evige Kærne mere og mere lægger sig for
Dagen.

Problemet stilles ved det aandelige Livs Stilling i Tilværelsen.
Dets Bestaaen og dets Udvikling forudsætter, at der bagved
og udover de aandelige Processer, Erfaringen viser os, ligger
en evig og sammenhængende aandelig Verden (eine geistige
Substanz, ein seelischer Bestand), der indeholder \textit{Mulighederne}
for nye Dannelser og \textit{Udbyttet} af hvad der allerede har udviklet
sig. Aandslivet maa være Mere end blot Fænomen, altsaa
Mere, end Psykologien kan paavise. Det maa have et evigt
Princip i sig. Ikke blot en psykologisk, men en metapsykisk
Betragtning maa anlægges. For en saadan Betragtning ere Væsen
og Værdi Et, idet dens højeste Genstand er selve Tilværelsens
værdifulde Kærne. Kun ved en saadan Betragtning bliver
et centralt Livsindhold muligt.

Den antike Kultur overvurderede Formen; den ved Renaissancen
grundlagte Kultur overvurderede Kraften. Den sande
Kultur, som har en religiøs Karakter, stiller \textit{Væsensdannelse} (die
Wesensbildung) i Modsætning baade til Formdannelse (die Formbildung)
og til Kraftdannelse (die Kraftbildung). Baade Formen
og Kraften skulle tjene et værdifuldt Indhold. Formdannelsen
alene fører til Stivnen i plastiske Kunstværker, i hvilke det
Evige antages at være udtrykt for alle Tider. Kraftdannelsen
fører til rastløs Bevægelse, saa at al Væren opløses. Væsensdannelsen
derimod hævder en stadig Vexelvirkning af Evighed
og Tid, idet de evige Muligheder gennem Arbejdet i Tiden om%
% --- p. 81 ---
\setcounter{footnote}{0}%
\opage{81}sættes til Virkelighed, og idet paa den anden Side dette Arbejdes
Udbytte bevares i Evighedens Form. Bevarelse af det
Givne og Frembringelse af nyt Indhold maa gaa Haand i Haand,
for at vi kunne trænge frem til den sande Virkelighed.

Den Tankegang, der hævder »Væsensdannelsens« Mulighed
og Nødvendighed, kalder Eucken den \textit{noologiske} (af noos, Fornuft,
Aand, der hos Nyplatonikerne stilles over psyche, Sjæl), og
han stiller den i Modsætning til den psykologiske Betragtning,
som blot beskæftiger sig med det i Erfaringen givne Bevidsthedsliv.
Det er efter hans Overbevisning ikke noget Nyt, han
herved proklamerer, men en gammel Sandhed, som især KANT
har indskærpet ved sin bestemte Skelnen mellem den logiske,
den etiske og den æstetiske Betragtningsmaade paa den ene Side,
den empirisk-psykologiske paa den anden Side. Her staar Gyldighed
eller Værdi i Modsætning til Virkelighed --- men saaledes,
at Gyldigheden tilsidst hævdes som den højeste Virkelighed.
--- Den noologiske Tankegang er forskellig fra en tidligere
Tids spekulative eller metafysiske Tankegang, som troede at
kunne vinde en Forstaaelse af Fænomenerne ved teoretiske Abstraktioner.
Den noologiske Metode bygger paa Erfaring, paa
Faktum. At forklare noologisk vil sige at indordne en aandelig
Livsytring i det aandelige Livs Helhed, paavise dens Stilling
og dens Opgave indenfor denne Helhed, og derved belyse og
befæste den. Denne Forklaring er forskellig fra den psykologiske,
som blot bestaar i Paavisningen af de Processer, gennem
hvilke Mennesket erhverver og tilegner sig et aandeligt Indhold.
Begge Metoder maa anvendes, forat Modsætningen mellem Idealisme
og Realisme kan overvindes. Idealismen har hævdet Aandsindholdets
Selvstændighed og Egenværdi; Realismens Berettigelse
beror paa dens Fremhæven af de naturlige Betingelser. Den
realistiske Betragtningsmaade kommer til sin Ret, hvor det indenfor
den aandelige Helhed, den noologiske Opfattelse hævder,
lykkes at finde de specielle Aarsager og Love; men den har
Uret i at mene, at Aarsagsforklaringen kan afgive den sidste
Maalestok for Sandhed. Idealismen har med Rette fordret et
værdifuldt Indhold for Tilværelsen, men har set bort fra de reale
Betingelser.

% --- p. 82 ---
\setcounter{footnote}{0}%
\opage{82}Den noologiske Betragtning fører nu ifølge Eucken netop
til den Kendsgerning, Religionen hævder: Nærværelsen af et
absolut Aandsliv, ophøjet over og dog virkende i de i Erfaringen
fremtrædende Livsfænomener. Der gives ingen Religion uden
levende Nærværelse af en højere Verden i vor empiriske Verden
--- men saaledes at disse to Verdener staa i et irrationelt
Forhold til hinanden. Religionen kan bestaa uden Gudetro; det
er den ægte Buddhisme et Exempel paa. Men uden Udblik ud
over Erfaringsverdenen er Religion et tomt Ord. Motivet til
Religion ligger i Modsætningen mellem Væsen og Existensform:
i den blotte Erfaring staar Aandslivet som et Tillæg til det
Materielle, som en Bisag; men dette strider mod de ubetingede
Krav, Aandslivet stiller, og de ophøjede Værdier, det hævder.

Ud fra Religionens Grundpaastand tager Virkeligheden sig
anderledes ud end før. En fra Religionen udgaaende Verdensbetragtning
vil rette sit Øje mod det Blivende (i Modsætning
til det Omskiftende), mod det Frie (i Modsætning til det blot
Naturlige) og mod det Fornuftige (i Modsætning til det Ufornuftige).
Den vil særligt mærke sig de Punkter, hvor en ny
Sfære begynder, saaledes Overgangen fra det Uorganiske til det
Organiske, fra Naturlivet til Sjælelivet, fra Sjælelivet til Aandslivet,
og den vil opfatte denne Trinrække som en Udfoldelse af
det Inderste i Tilværelsen. Men bevises kan en saadan religiøs
Verdensopfattelse ikke. Den har Karakteren af en Overbevisning,
ikke af en Videnskab. Det er Filosofiens Sag at bestemme
Forholdet og Grænsen mellem den videnskabelige og den religiøse Verdensbetragtning.

Motivet til Religion ligger i en aandelig Trang. Det drejer
sig ikke om Menneskets empiriske Lyksalighed og dens Bestaaen,
men det Afgørende er Aandslivets Helhed og Enhed,
Grundlaget for alle specielle og individuelle Livsytringer. Det
er dets Bestaaen, der staar paa Spil, idet den blotte Erfaring
ikke bekræfter den. Det kommer da an paa, om der kan vækkes
en stærk Trang til at fastholde en Overbevisning paa dette
Punkt. Empirisk set er der Grunde nok til Tvivl. Erfaringen
viser os ingenlunde Aandens Verden som en selvstændig Verden,
end mindre som Hovedsagen i Tilværelsen. Og under Kul%
% --- p. 83 ---
\setcounter{footnote}{0}%
\opage{83}turens Udvikling ofres Personlighederne for Samfundet og Fremskridtet.
Indenfor selve Aandslivet røre disharmoniske Tendenser
sig; Del staar mod Del, og Del staar mod Helhed; aandelige
Retninger bekæmpe hverandre, og det Onde viser sig at være
en virkelig Magt.

Overfor de Tvivl, der saaledes rejse sig, fremhæver Eucken
for det Første, at selve den stærke Følelse af Disharmonien er
en Følge af, at der stilles store Fordringer til Livet. Lidelsens
Bitterhed og Modstandens Haardhed vidne om Dybden af vor
Existens og om højere Kræfters Tilstedeværelse i vort Væsen.
Med Modsætningerne vilde ogsaa Smerten falde bort. Det var
aldrig i Tider med ulykkelige Oplevelser og tunge Forviklinger,
at Menneskeheden tvivlede om sin ideale Opgave, men derimod
trivedes denne Tvivl i saadanne Tider, hvor trægt Velbehag og
indbildt Mættelse raadede.

For det Andet er det Euckens Overbevisning, at hin Tvivl
maa bekæmpes ved en kraftigere og mere koncentreret Form
for \textit{Religion} end den ubestemte og universelle \textit{Religion}, som i
det Foregaaende er beskreven. Han henpeger til de store Personligheders
Betydning i Religionernes Historie. Hos dem rørte
det Guddommelige sig paa dybere og kraftigere Vis end hos
Andre, og de formaaede at udtrykke Livets mangfoldige Erfaringer
i klare og glødende Billeder. Et nyt Liv brød her igennem
paa afgørende Maade og greb bestemmende ind i Slægtens historiske
Udvikling. Ved at fremdrage deres Betydning gøre vi
Overgang fra \textit{den universale} Religion til \textit{den karakteristiske}
Religion.

»Karakteristisk Religion« betegner hos Eucken de historiske
eller positive Religioner, forsaavidt de --- hver paa sin ejendommelige
Maade --- give et nyt, i sine Grundtræk klart Totalbillede
af Livet og stille dette Billede i Modsætning til Erfaringens
Verden. Den adskiller sig fra »den universale Religion«
ved den bestemte Udformning, uden hvilken en saadan Modsætning
til Erfaringsverdenen ikke kan opretholdes. Den udspringer
af store Personligheder og hæver sig over al Folkekultur som
en umiddelbar Ytring af Aandslivets inderste Væsen. Uden
saadan karakteristisk Religion er selvstændig Inderlighed  ikke
% --- p. 84 ---
\setcounter{footnote}{0}%
\opage{84}mulig. Det drejer sig her ikke om en blot Udsmykning af Livet,
men om et nyt Livs, en ny Virkeligheds Gennembrud. Dog
ere de Gennembrud, de store positive Religioner vidne om, kun
Højde- og Vendepunkter indenfor en stor Proces. Det Guddommelige,
som her bryder igennem, er ikke først opstaaet paa
disse bestemte Steder.

Det højere Liv, der saaledes bryder frem, formaa vore Begreber
ikke at udtrykke. Fantasien maa komme os til Hjælp
med sine Symboler. Og her ligger Religionens saarlige Punkt.
Fantasien bestemmes stedse ved det givne Kulturtrin og dets
Overleveringer, og Symbolerne ville derfor faa en historisk Karakter.
Der vil engang komme en Tid, da det, der maaske for
Aartusinder var Udtryk for den højeste Sandhed, og som for
den menneskelige Bevidsthed havde indgaaet en inderlig Forbindelse
med det religiøse Indhold, rokkes i sin Betydning og
maa vige Pladsen for nye Former --- ofte først efter svare Rystelser
og haarde Kampe, under hvilke Tvivlen fra at være rettet
mod Symbolerne kan gaa over til at angribe Religionens inderste
Kærne.

Den universale Religion er stadigt nødvendig som Korrektiv
for eller Indledning til de karakteristiske Religioner, der overladte
til sig selv let indsnevres og blive til blotte Trøstemidler.
De karakteristiske Religioners Betydning beror ikke paa, at
nogen af dem skulde have fundet afsluttende Udtryk for den
højeste Sandhed, men paa deres levendegørende og forhøjende
Virken under Tidens Uro og Omskiftelser. Det maa stedse paany prøves, om det historisk Overleverede virkeligt udfører et saadant
Arbejde. Men Kirken har en Tendens til at stille sig slavisk
overfor et en Gang for alle givet Forbillede og opfatte al
Sandhed som bestaaende i Efterligning og Gentagelse. Den isolerer
det Højeste og giver det derved en passiv Stilling overfor
det virkelige Liv. Derved bliver Kamp stedse nødvendig
paa det religiøse Omraade. Det er de store religiøse Personligheders
Betydning, at de nøde os til at vælge. ---

Hvad der især bestemmer Euckens Plads i den moderne
Religionsfilosofi, er hans afgjorte Overbevisning om, at Værdierne
kun kunne hævdes, naar der bagved og ud over Erfarin%
% --- p. 85 ---
\setcounter{footnote}{0}%
\opage{85}gens Verden antages en højere Virkelighed, i hvilken det Aandsliv,
der i Erfaringen fremtræder spredt og i skiftende Former,
bestaar i Enhedens og Evighedens Form. Paa en ofte dybsindig
og stemningsfuld Maade hævder han denne Modsætnings
Nødvendighed, --- snart saaledes, at han minder om den gamle
Mystik, snart saaledes, at han minder om den romantiske Idealisme,
og snart saaledes, at det er den personlige Livskamps ophøjede
Poesi, der faar Udtryk i hans Ord. Dog betoner han
meget stærkt, at Filosofien her kun kan paavise Muligheder, og
at det er gennem profetiske Personligheder, saadanne Muligheder
vinde Virkelighed og bestemt Form. Naar han --- forat
undgaa Subjektivismen --- hævder Nødvendigheden af en Metafysik
(eller en »Metapsyki«), saa er det dog egentligt i Kraft af
det Sandes, Skønnes og Godes \textit{Værdi}, at denne Metafysik skal
være nødvendig. Hvad han kalder den »noologiske« Metode,
adskiller sig fra den psykologiske Metode derved, at den vil
paavise, at hin Værdi kun kan hævdes ved Antagelse af en
Helhedssammenhæng, der bærer de i vor Erfaring spredt og
flygtigt optrædende værdifulde Livsytringer. Især lægger den
noologiske Betragtningsmaade Vægt paa de store Gennembrud,
de betydningsfulde Begyndelser og de kvalitative Forskelligheder
indenfor Aandslivet, som vi kende dette i Erfaringen. Han anser
Erfaringsvidenskabens Forsøg paa at paavise Kontinuitet i
Fænomenernes Verden for haabløse, men fortolker Bruddene
paa den empiriske Kontinuitet som enkelte Udslag af en stor
transscendent Kontinuitet, et Mulighedernes, Endemaalenes og
Værdiernes Rige, der kun glimtvis aabenbarer sig i Endelighedens
Verden. Derved er det, han er beslægtet med den romantiske
Idealisme. Men Noologien adskiller sig fra den spekulative
Metode derved, at Antagelsen af en saadan Helhedssammenhæng
udtrykkeligt kun grundes paa Erfaringen om det
Værdifulde. Det synes da klart, at Udgangspunktet er empirisk
og personligt. Han faar Fart til at gaa over fra Erfaringen til
»Metapsykien« gennem den Spænding, som frembringes ved et
stort Dilemma. Enten --- paastaar han --- maa man betragte
Aandslivet som et flygtigt Fænomen, et Vedhæng ved den materielle
Natur, --- eller man maa antage, at det i sig selv ud%
% --- p. 86 ---
\setcounter{footnote}{0}%
\opage{86}gør en stor Totalitet og er Et med Tilværelsens inderste Kærne.
Dette Dilemma kommer frem hos ham, baade hvor han udvikler
Ideen om »den universale Religion«, altsaa vil gøre Overgangen
fra det empiriske Standpunkt til et religiøst Standpunkt af mere
ubestemt Art, og hvor han gaar over til den »karakteristiske
Religion«, altsaa vil paavise Nødvendigheden af en historisk udformet
Religion. Karakteristiske ere især følgende Udtalelser: »Ved
denne Afgørelse staar ikke dette eller hint, men Aandslivet
som Helhed paa Spil. I Menneskets Erfaring fremtræder
Aandslivet i en uholdbar Mellemstilling; fra denne maa Mennesket
enten vende tilbage til den blotte Natur og bekende, at al
ejendommelig Stræben kun var en tung Drøm, eller modigt gaa
fremad og sikre den nye Stræben en ny Verden. Kun en uklar
Tænkemaade kan søge en Mellemvej mellem disse Veje, og kun
et mat Sindelag kan roligt udholde Dilemmaet. I hvert Tilfælde
staar det fast, at uden Religion gives der ingen Sandhed
Aandslivet og ingen indre Storhed i Mennesket.« (Der Wahrheitsgehalt der Religion. p. 238.) »Er det menneskelige Liv et
blot Tillæg til Naturen, eller er det Begyndelsen til en ny Verden?
Af dette Spørgsmaals Besvarelse afhænger al Udformning
af vore Forhold, hele Retningen af vor Handlen. Religionen
vil hæve den menneskelige Tilværelse op paa en over Verden
hævet Højde og derved redde vort Liv fra den Intethed, som
det ellers sikkert forfalder til. Er dette dens Forsøg kun en
Ikarusflugt, saa falder alt Haab bort, saa staar det Ædleste og
Bedste som tom Indbildning, saa ender det Hele i Ufornuft.«
(ibid. p. 319). ---

Hvad der skal drive os til at vælge her, det er dog stedse
tilsidst den personlige Livserfaring, særligt den Grad, i hvilken
man betages af Værdiernes Skæbne i Virkeligheden. Der synes
ikke at være nogen Grund til at indføre en ny Terminus;
»Noologi« er hverken nogen nødvendig eller nogen heldig Berigelse
af den allerede i Forvejen tilstrækkeligt forsynede filosofiske
Terminologi. Af en privat Udtalelse af Eucken, som jeg
haaber at kunne anføre uden at begaa en Indiskretion, ses det
klart, at det er den personlige Erfaring om Livets Modsætninger,
særligt af den store Modsætning mellem Ideal og Virkelig%
% --- p. 87 ---
\setcounter{footnote}{0}%
\opage{87}hed, der har bestemt hans eget Valg og ført ham til »Noologien«.
I et Brev, han skrev til mig efter Læsningen af min
Religionsfilosofi, udtaler han, efter at have dvælet ved den overvejende
psykologiske Karakter af mit Forsøg paa at behandle
Problemet: »Mich treibt es, bei allem Gegensatz zur alten ontologischen
Metaphysik, immer wieder zur Metaphysik und so
auch zu einer metaphysischen Begründung der Religion; ich
fürchte ohne das dem Subjectivismus zu verfallen, und dagegen
sträube ich mich namentlich deswegen, weil ich stärker unter
dem Eindruck der grossen, in der Erfahrung unlösbaren Konflikte
unseres Daseins stehe, kurz gesagt, dualistischer denke
und denken muss.« I denne Udtalelse peger Eucken direkte
hen til individuel Erfaring som bestemt ved den enkelte Personligheds
Ejendommelighed. Det er den, der bestemmer hans
Standpunkt. Men jo klarere dette Grundlag fremtræder, des naturligere
bliver det dog at anvise Religionsfilosofien den Opgave
at undersøge den psykologiske Jordbund, af hvilken de forskellige
religiøse Standpunkter voxe frem, og dertil knytte en erkendelsesteoretisk
og en etisk Prøvelse af de Tanker og de Bestræbelser,
som disse Standpunkter give sig Udslag i. En ganske
særegen Metode behøver da Religionsfilosofien ikke, og den saakaldte
»Noologi« er jo ogsaa selv en »Mellemvej« mellem Spekulation
og praktisk Tro. Som Spekulation henfalder den til
Erkendelsesteoriens, som praktisk Tro til Psykologiens og Etikens
Behandling. Den Opfattelse af Religionsfilosofiens Opgave
og Fremgangsmaade, som heraf følger, synes at være mest i
den kritiske Filosofis Aand, i hvilken jo Eucken --- ligesom
Siebeck og Sabatier --- vil arbejde. Den »metafysiske« Begrundelse
af Religionen, som Eucken kræver af Filosofien, viser sig
efter hans egne Udtalelser tilsidst at støtte sig til Trangen til
Afgørelse af et stort praktisk Dilemma.

Jeg finder i Euckens Standpunkt --- i den »Noologi«, der
fører ham til hans religiøse Metafysik --- et Exempel paa, hvorledes
personlig Livserfaring kan føre til Tro paa Værdiens Bestaaen.
For mig er enhver Tro, selvom den fremtræder i de
mest filosofiske Former, stedse kun \textit{Genstand} for Filosofi, ikke
selv \textit{Produkt} af Filosofi. I hvert Tilfælde er den Filosofi, der
% --- p. 88 ---
\setcounter{footnote}{0}%
\opage{88}producerer den, selv Produkt af personlig Livserfaring. Filosofien
som Videnskab raader nu engang ikke over andre Metoder
end den logisk sluttende, den psykologisk forklarende og den
etisk vurderende.

Men fra mit Standpunkt faa saadanne Arbejder som Euckens,
og det Samme gælder Sabatiers, en overordentligt stor Interesse
ved den klare Maade, paa hvilken de afgørende Motiver
for religiøs Livsanskuelse træde frem. To idealt sindede, af Livet
grebne og med Tankens Energi udrustede Mænd have her udtalt
sig om et af de største menneskelige Problemer.

\section*{II.\\ RELIGIONSPSYKOLOGISKE UNDERSØGELSER.}
\streg

Et gennemført Forsøg paa at gaa tilbage til umiddelbar Erfaring
ved Behandlingen af det religiøse Problem maa særligt
henvende Opmærksomheden paa det Punkt, hvor Bevidstheden
fra umiddelbar Oplevelse gaar over til at forme en Dom, altsaa
paa den første Opstilling eller uvilkaarlige Tilblivelse af en religiøs
Sætning. Hvad der her gaar for sig, kan ikke være væsensforskelligt
fra, hvad der altid gaar for sig ved Overgangen
fra Iagttagelse til Dømmen. Den religiøse Bevidsthed har hverken
en Psykologi eller en Logik, der er forskellig fra den almindelige
Psykologi og Logik, og de religiøse Domme maa blive
til, som alle andre Domme blive til. Hvad der er ejendommeligt
for de religiøse Domme, er, at den Erfaring, de tilsidst grunde
sig paa, er en Følelseserfaring, en Erfaring om visse Oplevelsers
Værdi som fremhjælpende eller hæmmende, og at de Oplevelser,
der her er Tale om, angaa Forholdet mellem Ideal og Virkelighed,
mellem Værdiernes og Realitetens Verden. Religiøse Vær%
% --- p. 89 ---
\setcounter{footnote}{0}%
\opage{89}dier ere derfor Værdier i anden Potens eller sekundære Værdier:
de forudsætte, at Mennesket kender Noget, der giver Livet
Værdi (hvori dette saa bestaar), og at han saa erfarer, at der
er Magter i Verden, som fremme eller hæmme dette Værdifulde.
Paa den religiøse Følelses Højdepunkter falder naturligvis Forskellen
mellem primær og sekundær Værdi bort; der kendes
intet Hvorfor og anvendes ingen Prædikater. I mest udpræget
Form sker dette i Mystiken. Men hvor der kan finde en Dømmen
Sted, altsaa være Tale om Udsagn af religiøs Bevidsthed,
der vil --- paastaar jeg --- en nærmere Analyse godtgøre, at
hvad der udsiges Noget om, er Forholdet mellem Værdi og
Virkelighed, saaledes som Individet mener umiddelbart at have
erfaret dette. Religiøse Domme angaa Værdiernes Skæbne i
Virkelighedens Verden. --- Herved maa det dog vel bemærkes,
at religiøse Domme ikke altid, maaske endog saare sjeldent, dannes
paa første Haand.

Ved enhver Domsdannelse maa der (som jeg nærmere har
udviklet i min Afhandling Det psykologiske Grundlag \textit{for logiske}
\textit{Domme}) skelnes mellem Udgangsforestillingen og Slutningsforestilling.
Ved den sproglige Formulering af Domme optræde
disse to Forestillinger som logisk Subjekt og logisk Prædikat.
Ved Domsprocessens Begyndelse har Udgangsforestillingen
en vis Ubestemthed, som søges hævet ved Opdagelsen af et
Prædikat, en nærmere Bestemmelse, der sker gennem Analyse
af Udgangsforestillingen og af den Sammenhæng, i hvilken den
forekommer. Jo rigere Udgangsforestillingens Indhold er, des
vanskeligere vil det være at finde en definitiv Slutningsforestilling,
et afsluttende og adækvat Prædikat. I denne Henseende
frembyder den religiøse Bevidsthed, saaledes som den fremtræder
i Historien, de største Forskelligheder. Jo snevrere baade
Værdiernes og Virkelighedens Verden er, og jo mere elementært
Forholdet mellem dem fremtræder, des lettere er det for
Mennesket at slaa sig til Ro ved et Prædikat og dermed afslutte
den religiøse Tænkning. Her er alle mulige Grader lige fra
Fetischismens og Animismens begrænsede og barnlige Prædikater
til den højeste Mystiks sublime Grænsebegreber. Og den
Proces, ved hvilken Slutningsforestilling her knyttes til Udgangs%
% --- p. 90 ---
\setcounter{footnote}{0}%
\opage{90}
forestilling, er snart uvilkaarlig Ideassociation, snart Efterligning og Vane, snart spekulativ Grublen, snart
bekymret Søgen og Spørgen. Udgangsforestillingens Indhold er ved religiøs Erfaring
det harmoniske eller disharmoniske Forhold mellem Værdi og
Virkelighed, og Slutningsforestillingen skal angive, hvad det er, der lægger sig for Dagen i dette Forhold. Dersom en saadan Slutningsforestilling vindes eller findes, er Overgangen gjort fra Erfaring til Tro. Gudsbegrebet i dets forskellige Former er det Grundprædikat, i hvilket den religiøse Bevidsthed søger at faa Udtryk for og Forklaring af sine dybeste og væsentligste Erfaringer. Ordet Gud skal betyde »den, til hvem der ofres«, altsaa den, til hvem Mennesket føler sig i Gæld, til hvem det fører sine Erfaringer om Forholdet mellem Værdi og Virkelighed tilbage. Men hvem er det? Qvid est quod amo, quum te amo? spørger Augustinus. Her ligger Mulighed for stedse ny Grublen, især da Udgangsforestillingen ikke er konstant, eftersom Forholdet mellem Værdi og Virkelighed stedse ændrer sig. Der findes i Goethes Faust en Monolog, som egner sig godt til Exempel paa, hvad her menes med religiøs Erfaring. I Scenen »Wald und Höhle« foretager Faust et Opgør. Naturens og Tankens store Skatte ere aabnede for ham; han kan fordybe sig i den indre som i den ydre Verden i Forstaaelse og Stemning; men tillige er der medgivet ham en dæmonisk Ledsager, som han ikke kan undvære, og hvis Kulde og Frækhed forvandler hine
høje Gaver til et Intet --- samtidigt med, at han vækker det sanselige Begærs flammende Ild. Til en Slutningsforestilling kommer det her ikke; vil man søge en saadan i Digtet, maa den findes dels i Prologen til første Del, dels i Slutningen af anden
Del. --- Den store Modsætning, Goethes Faust her gør Rede for,
ligger --- kun i endnu skarpere Form --- til Grund for Schopenhauers
Livserfaring, og Schopenhauers Filosofi var i sin inderste
Grund et Forsøg paa at finde en Slutningsforestilling som Udtryk
for denne Erfaring. ---

Umiddelbar religiøs Erfaring er et Ideal, da ethvert Menneske, vi kende, begynder under Indflydelse af en Tradition og
strax tyder sine Oplevelser ved Overleveringens Hjælp, eller maaske nøjagtigere, strax ser dem i de Former og den Belys%
% --- p. 91 ---
\setcounter{footnote}{0}%
\opage{91}ning, Overleveringen har lagt tilrette. Allerede Indlevelsen i
Sproget bevirker, at der her ikke kan være Tale om absolut
oprindelig Erfaring. Enhver religiøs Erfaring er mere eller mindre
konfessionel. Baade Værdierne, Virkeligheden og Forholdet
mellem dem opfatter Mennesket fra først af under Indflydelse
af Forestillinger, som det ikke selv har dannet paa egen Haand,
endsige selv udledet af sine Erfaringer. Det vil da være af stor
Interesse at undersøge, hvorledes de overleverede Forestillinger
efterhaanden ændres ved Individets egne fremskridende Erfaringer,
og dernæst at anstille en Sammenligning mellem religiøse
Standpunkter og Udviklingsprocesser, der bestemmes ved væsentligt
forskellige Udgangspunkter.

Der foreligger nu om alle her berørte Emner --- baade om
religiøse Domme som Værdidomme, om en konfessionel Religiøsitets
Former og Modifikationer og om komparativ Religionspsykologi
--- en Række Arbejder, der for største Delen ere fremkomne
efter Afslutningen af min Religionsfilosofi, og som det
vil have sin Interesse at faa karakteriserede.

\subsection*{1. OM RELIGIØSE DOMME. \textit{(MAX REISCHLE).}}

De Forsøg paa Religionsfilosofi, som vi i det første Kapitel
beskæftigede os med, vare mere eller mindre konsekvent byggede
paa den kritiske Filosofi og dens Sondring mellem rent
teoretiske Domme og Værdidomme. Indenfor Teologien har denne
Distinktion i de senere Aar spillet en ikke ringe Rolle, især hos
ALBRECHT RITSCHL og hans Disciple. Auguste Sabatier staar
denne Retning nær, og han taler da ogsaa om Værdidomme
(jugements d'estimation ou de dignité). Den Vægt, der lægges
paa Forskellen mellem de to Arter af Domme, er karakteristisk
for den nyeste Tids Bestræbelser for at føre Religionen tilbage
til umiddelbar Erfaring og bestemme denne Erfarings egentlige
Beskaffenhed. Det er gennem Undersøgelsen af Værdidommene,
at der i den nyeste Tid er kommet en Berøring mellem Teologi
og Filosofi, som ellers i de sidste halvhundrede Aar --- især
paa Grund af Teologiens mere og mere udpræget konfessionelle
% --- p. 92 ---
\setcounter{footnote}{0}%
\opage{92}Karakter --- have staaet hinanden fjernt. --- Jeg skal her som
Exempel henlede Opmærksomheden paa en Afhandling af MAX
REISCHLE: \textit{Werturteile und Glaubensurteile}. (Halle 1900).

En Værdidom forudsætter ifølge Reischle en Værdifølelse,
der er forskellig fra den blotte Lystfølelse. Forudsætningen for
en Værdifølelses Opstaaen er, at Tilfredsstillelsen angaar hele
Jeget (das Gesammt-Ich), ikke en enkelt Del eller Tilstand af
Jeget. Værdi tillægger jeg en Ting, om hvilken jeg ad Eftertankens
Vej har indset, at den vil tilfredsstille mit Jeg som Helhed.
Dog indrømmer Reischle, at Grænsen mellem »hele Jeget«
og dets enkelte Dele eller Tilstande er flydende. Hvad han vil
fremhæve, er især det, at der forud for \textit{Værdidommen} gaar en
sammensat Proces, ved hvilken Værdifølelsen naar en Klarhed
og Sikkerhed, den blotte Lystfølelse mangler.

Værdidomme kunne ifølge Reischle være af to Arter: naturale
og ideale (eller normative).

\textit{Naturale Værdidomme} ere begrundede i Menneskets Natur
og Livsvilkaar, der bestemme den Tilfredsstillelse, der er mulig
for det. Her gøre sig mange Forskelligheder gældende mellem
de forskellige Individer, men da Interesserne dog til en vis Grad
ere solidariske, kan der gives kollektive Værdidomme, og da
Livsvilkaarene til en vis Grad ere fælles, kan der gives generelle
Værdidomme. Hvor en ydre Maalestok (f. Ex. Penge) er
mulig, kan Værdidommen faa en objektiv Karakter, saa at maaske
endog dens subjektive Oprindelse glemmes.

\textit{Ideale eller normative} Værdidomme forudsætte en Ide eller
en Norm, som vi anerkende og vente anerkendt af Andre. Herhen
høre æstetiske, etiske, intellektuelle og religiøse Værdidomme.
Disse bestemmes ikke blot ved den Maade, paa hvilken vi passivt
bestemmes udefra, men væsentligt ved Overensstemmelsen
med en Lov, der gælder for vor egen Aands Virksomhed. Ved
etiske Værdidomme f. Ex. følge vi en Norm eller en Ide, som
formaar at føre vor Villie til Harmoni og Frihed. Ved religiøse
Værdidomme er den tilgrundliggende Ide den Betydning, det
har at underkaste sig en højeste Magt for ved Enhed med den
at naa den højeste Frihed overfor Verden. De ideale Værdidomme
gøre Fordring paa Almengyldighed, fordi Ingen kan for%
% --- p. 93 ---
\setcounter{footnote}{0}%
\opage{93}negte de Ideer, der betinge dem, uden at opgive sin Personlighed,
sit Væsens Frihed og indre Sammenhæng. ---

Forskellen mellem naturale og ideale Værdidomme er ---
hvad Reischle ikke har bemærket --- aabenbart ligesaa flydende,
som han selv indrømmer, at Forskellen mellem Lystfølelse og
Værdifølelse er. Hint »Jeg som Helhed«, Værdifølelsen skal
forudsætte, indeholder nemlig en Norm, der allerede gælder for
de specielle og skiftende Elementer og Tilstande; disses Værdi
bestemmes ved Forholdet til det. Og desuden maa paa den
anden Side »Jeget som Helhed« selv godtgøre sin Værdi overfor
de enkelte Elementer og Tilstande, hvis Underordning kræves;
det er ingenlunde selvfølgeligt, at Helheden har Ret overfor
det enkelte Element eller den enkelte Tilstand. Endeligt beror
jo ifølge Reischle selv de æstetiske, etiske, intellektuelle og
religiøse Aandsfunktioners Værdi paa, at de skaffe vort Væsen
Harmoni og Frihed; --- men det vil ikke sige Andet, end at
de gøre et »Jeg som Helhed« muligt, og de ideale Værdidomme
blive fra denne Side set ikke forskellige fra de naturale. Det
er derfor ingenlunde let at opstille Værdimaalere paa det aandelige
Omraade, hvor der er den inderligste Vexelvirkning og gensidige
Bestemmelse af Elementerne indbyrdes og mellem Elementer
og Helhed. Indenfor den individuelle Bevidsthed, ligesom
indenfor Samfundet, kæmpe de forskellige Værdimaalere.
Enhver Ides eller Norms Værdi skal selv godtgøres --- og derfor
stilles Problemet stedse paany. Af alle Værdidomme frembyde
netop, som jeg et andet Sted har vist, de, der føre til Opstilling
og Anerkendelse af en Norm, de største Vanskeligheder
for Teorien. Overgangen fra det »Naturale« til det »Normative«
(Ideale) indbefatter store Problemer, som en blot Klassifikation
af de faktisk givne Værdidomme let tildækker. Reischle
har klart Øje for det analoge Problem, der opstaar, naar de
forskellige Religioner, hver med sine Grundværdier, staa overfor
hverandre; men han har ikke set, at Kampen allerede kan
begynde i den Enkeltes Indre, og at en ledende Ide eller Norm
stedse er Resultatet af en saadan, mere eller mindre heftig
Kamp. ---

I den sidste Del af sin Undersøgelse drøfter Reischle, hvor%
% --- p. 94 ---
\setcounter{footnote}{0}%
\opage{94}vidt \textit{Trossætninger} kunne kaldes Værdidomme. Han indrømmer,
at Autoritet og Tradition spille saa stor en Rolle baade i den
protestantiske og i den katolske Kirke, at de Trossætninger, de
Enkelte hylde, ikke kunne kaldes Værdidomme i den Betydning,
at disse Enkelte selv have udledet dem af deres Værdierfaringer.
Derimod anser Reischle det for utvivlsomt, at Trossætningerne
maa kaldes Værdidomme, naar der spørges, hvorpaa
deres Fordring om Gyldighed principielt beror. De ere Værdidomme
af ideal eller normativ Art; deres Betydning bestaar i,
at de angive Betingelser for Saligheden, d. v. s. for Personlighedens
fuldstændige Harmoni og Frihed. Det vilde være urigtigt
at kalde dem Postulater, som om det blot var vor egen
Trang og Higen, der i dem formulerede, hvad der kunde volde
Tilfredsstillelse. Der ligger ved de religiøse Domme bestemte
Kendsgerninger til Grund, for den Kristne saaledes Jesu Liv
og Død, og de religiøse Domme udtrykke Tilliden til saadanne
Kendsgerningers Magt og Betydning. Den Kristne tror paa Kristus,
fordi han kun i Tilslutning til ham kan finde det sande
Liv. Om alle historiske Religioner gælder det, at visse givne
Fænomener --- Personer og Begivenheder --- tydes som guddommelige
Aabenbaringer, og de religiøse Domme udtrykke Tilliden
til disse Aabenbaringer. Paafaldende, især velgørende eller
ødelæggende Fænomener i Naturen, i det bevidste eller det
ubevidste Sjæleliv, i Folket eller i Historien vække den uvilkaarlige
Forestilling om, at højere Magter her virke til Menneskers
Vel eller Ve. Ofte dikteres Dommen af Frygt, oftest dog
af Tillid. Og jo højere en Religion staar, des mere er det etiske
Ideer, der diktere Dommen. Selve det, at Noget kaldes »Aabenbaring«,
er altsaa efter Reischle en Værdidom; det beror paa,
om vedkommende Personligheder eller Begivenheder staa som
Former for store Værdiers Indtræden i Historien, og om de
kunne blive Genstand for tillidsfuld Tilslutning og Hengivelse.
En rent objektiv Afgørelse af Striden mellem forskellige Aabenbaringer
er derfor umulig. Dette er en nødvendig Konsekvens
af Empirismen i Religionsfilosofien, og Reischle drager ogsaa
denne Konsekvens, som han kun fornegter ved at sætte Ordet
% --- p. 95 ---
\setcounter{footnote}{0}%
\opage{95}»Aabenbaring« med Anførselstegn, hvor det ikke drejer sig om
hans egen Religion.

\subsection*{2. DEN RELIGIØSE FØLELSES PATOLOGI. \textit{(MURISIER).}}

For at kunne paavise Elementerne i den religiøse Følelse
kan det være hensigtsmæssigt at studere den i saadanne Former,
i hvilke den antager en sygelig Karakter. Det Sygelige bestaar
i, at Elementer, der ellers forekomme i harmonisk Sammenhæng med andre Elementer, udsondre sig og optræde med overvældende
Styrke, saa at de beherske hele Livets Gang. I sit
Skrift \textit{Les maladies du sentiment} religieux (Paris 1901) har E.
\textit{MURISIER} (Professor i Neufchâtel) sat sig den Opgave at undersøge
den religiøse Følelses sygelige Former, som han mener at
kunne føre tilbage til to, Extasen og Fanatismen, hin det extreme
Udtryk for den religiøse Følelse i dens individuelle Form,
denne dens extreme Udtryk i social Form. Disse to Former
staa i Modsætning til hinanden: den ene fører til absolut Isolering,
den anden til udadgaaende Virken og til Forbindelse
med Andre. Men der ligger, som Murisier viser, en og samme
Trang til Grund ved begge Typer. Mystikeren ser i Isolationen
og i Bortfjernelsen af al Mangfoldighed fra Bevidstheden det
eneste Middel til at bevare Enheden og Sikkerheden, til at fastholde
Livets totale Retning mod et eneste Maal. Hans Aand
mangler Energi til at magte Livets Mangfoldighed i ham og
udenfor ham, og han samler sig derfor paa et eneste Punkt.
Fanatikeren søger derimod at opnaa Enhed og Sikkerhed i Livets
Gang og Retning ved en praktisk Handlen, som fører saa
Mange som muligt til at dele hans Tro. Han søger at tilvejebringe
en social Enhed, indenfor hvilken ingen disharmoniske
og forstyrrende Tendenser mærkes. Ved social Ensformighed
støttes det egne Jegs Enhed og Identitet. Fanatikeren priser
kirkelig Enhed med ligesaa stor Iver som Mystikeren Tankens
Enhed. Og ligesaa ivrigt som Mystikeren holder alle Indtryk
borte, der kunde forstyrre hans Opgaaen i det Ene, der er det
Højeste for ham, ligesaa ivrigt undertrykker Fanatikeren enhver
% --- p. 96 ---
\setcounter{footnote}{0}%
\opage{96}Retning, der kunde bryde den sociale Enhed. Hans Forfølgen
svarer til Mystikerens Askese, og hans Proselytisme til Mystikerens
Koncentration. For mange Mennesker er Tilslutningen
til et religiøst Samfund, især et snevert og exklusivt, den eneste
Vej, ad hvilken de kunne finde Ligevægt og Fred i Sindet.

Gennem Studiet af disse modsatte Former for Religiøsitet,
støttet til et historisk, biografisk og medicinsk Materiale, finder
Murisier i de Karaktermærker, der ere fælles for dem, et Bidrag
til Karakteristik af Religionen. Religionen skyldes --- hvad
enten den fremtræder i individuel eller i social Form --- en
Enhedstrang, der hænger nøje sammen med Selvopholdelsestrangen.
Denne Enhedstrang fremtræder i Religionen særligt som
Trang til Ledelse (besoin de direction), og forsaavidt var Schleiermachers
Beskrivelse af Religionen som Afhængighedsfølelse et
heldigt Greb. I sygelig Form kan Enhedsretningen, den koncentrerede
Afhængighed kun naas gennem Udelukkelse af forskelligartede
Indtryk og Retninger; i sund Form naas den
paa mere harmonisk Maade gennem »Opbyggelse« og Organisation.


Religionen gaar i sin sociale Form forud for den individuelle
Religiøsitet; Religionshistorikerne beskæftige sig derfor mest
med hin og ere tilbøjelige til at betragte denne som et Opløsningstegn.
Naturfolkenes Religion bestaar mest i Trosforestillinger
og Ceremonier, som ere fælles for hele Samfundet. Det
primitive Menneskes Jeg har ikke udsondret sig af Familie, Klan
og Folk, ligesom Barnets Bevidsthed stadigt henter sit første
Indhold fra de nærmeste Omgivelser. Selve Extasen hos de
Vildes Troldmænd minder mere om Fanatikernes Aandstilstand
end om Mystikernes Extase; der hører en vis Udvikling til, før
den egentligt individuelle Form kan blive til. Men Religionens
Sociologi kan paa intet Punkt ganske skilles fra Religionens
Psykologi. Selve Trangen til at frembringe social Enhed eller
slutte sig til en saadan har en psykologisk Aarsag, der maa
studeres hos de enkelte Individer. Selv hvor Religionen nærmest
fyldestgør et socialt Krav, bæres den dog væsentligt af en indre
Higen i de Enkeltes Sind. ---

Der kan sikkert indvendes Adskilligt mod Murisiers Op%
% --- p. 97 ---
\setcounter{footnote}{0}%
% CHECK p. 97: indent after a broken word?
\opage{97}fattelse af Mystiken som et rent individualistisk Fænomen. Dels
begynder Mystikeren som Regel indenfor en konfessionel Religion,
og dels hører hans eget Jeg til det, han vil befries fra;
han vil ud af de snevre Forhold, ud over alle Grænser, og Mystiken
er ofte en til det Yderste dreven Altruisme. Mystiken
er et Expansionsfænomen, nok saa meget som et Koncentrationsfænomen,
eller rettere, der gives mangfoldige Nuancer indenfor
den mystiske Retning. --- Der gives ogsaa flere Former for
sygelig religiøs Følelse end de to Hovedformer, til hvilke Murisier
holder sig. Der er saaledes Tilstande af Angst og Fortvivelse,
hvor hverken Expansion eller Koncentration er mulig, og
der er i Modsætning dertil den overvættes Selvfølelse, hvor Æren
og Magten ved at være knyttet til det Højeste nydes paa egoistisk
Vis. Denne Modsætning mellem potenseret Afmagtsfølelse
og potenseret Magtfølelse er maaske nok saa betydningsfuld
for Religionens Patologi, som de af Murisier beskrevne Former.

Hvad der ikke kommer tilstrækkeligt frem hos Murisier, er
den religiøse Følelses Sammenhæng med Værdierfaringer. Man
vilde kunne have en Ledetraad til at finde de sygelige Former
for Religion ved at gaa ud fra de Forudsætninger, hvortil Studiet
af Religionen atter og atter vil føre tilbage: Erfaringer om
Værdier, om Virkeligheder og om Forholdet mellem Værdierne
og Virkeligheden. Men det mest videnskabelige Synspunkt vilde
dog være det, der betragter Religionens patologiske Former som
specielle Tilfælde af de almindelige Sindssygdomsformer. Anlæggelsen
af dette Synspunkt vilde ogsaa føre til at indse den Begrænsning,
indenfor hvilken man maa bevæge sig, naar man
fra det syge Bevidsthedslivs Fænomener vil drage Slutninger
angaaende det sunde Bevidsthedslivs Natur. --- Disse Indvendinger
hindre dog ikke i at anerkende, at der i Murisiers lille
Bog findes megen god og fin Psykologi, og at den --- indenfor
den Begrænsning, Forfatteren har valgt --- belyser det religiøse
Liv fra en karakteristisk Side.
% --- p. 98 ---
\setcounter{footnote}{0}%
\opage{98}

\subsection*{3. BIDRAG TIL DEN KONFESSIONELLE RELIGIØSITETS\\ PSYKOLOGI.}

Fra religiøs Side hævdes det mere og mere, at Grundlaget
for Religion er Erfaring, og ethvert Forsøg paa psykologisk Undersøgelse
af Religionen fører til samme Resultat. Men saa bliver
det af afgørende Betydning at undersøge Arten og Beskaffenheden
af den Erfaring, der her kan være Tale om. Ellers kan
man ikke vide, hvor Meget dette Grundlag formaar at bære,
og Dogmatismen holder igen sit Indtog. Men her indtræder nu
en Vanskelighed, som fremkommer ved, at enhver Følelsestilstand
har en Tilbøjelighed til hurtigt at indgaa Forbindelse med
bestemte Forestillinger og føre til bestemte Handlinger. Og de
Forestillinger og Handlinger, der stedse ligge nærmest, ere de
i det omgivende Samfund raadende. I dem lever Individet sig
ind, maaske endog før det har begyndt at gøre sine egne Livserfaringer.
Det arver Slægtens Vurderinger, før det kan begynde
at vurdere Livet paa egen Haand. Og naar det begynder
at gøre sine egne Erfaringer, vil det i hvert Tilfælde strax
lægge sig dem til Rette efter de overleverede Former. For en
stor Del, om ikke for største Delen, fremkaldes dets Erfaringer
netop ved det Indtryk, overleverede Forestillinger og Handlinger
gøre paa det, især gennem Omgivelsernes Betagethed af dem.
Kun med stor Vanskelighed og under særegne Forhold vil den
Kontraprøve blive mulig, som behøves til et virkeligt Erfaringsbevis,
--- nemlig Prøvelsen af, om den Erfaring og den Livsvurdering,
Individet selv efter sin Natur og sine Vilkaar kan
naa til, nu ogsaa virkeligt ene og alene kunne udtrykkes og
forklares gennem de overleverede Former, saa at Indhold og
Form ganske svare til hinanden. Muligvis kunde samme Indhold
udtrykkes gennem andre Former, eller andre Forestillinger
og Handlinger kunde endog gøre en rigere og værdifuldere Livserfaring
mulig.

Den religiøse Erfaring bliver derfor, historisk set, stedse
mere eller mindre konfessionel. Den bliver overvejende en Gentagelse
af, hvad Andre forud have erfaret og oplevet; den bevæger
sig i færdige Udtryk og Formler, og det, der fra først
% --- p. 99 ---
\setcounter{footnote}{0}%
\opage{99}af, nemlig da Udtrykket fandtes, og da Formlen skabtes, var Virkning
og Resultat, det virker i den konfessionelle Erfaring som
Aarsag, ofte Hovedaarsag. Dog er der indenfor den konfessionelle
Erfaring alle mulige Grader mellem den passive Indlevelse
i givne Former og den selvstændige Tilegnelse, der paa
sit Højdepunkt kan gaa over til at finde nye Veje og skabe
nye Former.

Religionspsykologien maa have denne Kendsgerning for Øje,
naar den indsamler sit Erfaringsstof. Et rigt Materiale have vi
her --- foruden i hvad egen Menneskekundskab yder --- i religiøse
Personligheders Udtalelser, især i Selvbiografier. Det er
især til et saadant Materiale, jeg har støttet mig i det psykologiske
Afsnit af min Religionsfilosofi, idet jeg saa tillige har
sammenlignet de Typer, jeg ad denne Vej fandt, med de Typer,
de store Folkereligioner, især deres Grundlæggere, frembyde.
Men det vil naturligvis være af Betydning at faa en bredere
Basis, og dette vilde kunne opnaas, dersom der kunde samles
Oplysninger og Udtalelser fra større Kredse, ogsaa fra ikkeliterære
Kredse. Forsøg paa at skaffe saadanne støde paa mange
Vanskeligheder, dels paa Grund af den Ubehjælpsomhed, som
vil hindre Mange i at udtale sig klart, dels paa Grund af den
Skyhed, der vil hindre Mange fra at udtale sig om det Inderligste
i deres Liv; det blotte Spørgsmaal vil maaske endog undertiden
blive betragtet som en Profanation. Men da psykologisk
Forskning som al anden Forskning øves for Sandhedens
og Klarhedens Skyld, vil det være fuldt berettiget at gaa denne
Vej og prøve, hvor langt den kan føre. Det er især i Nordamerika,
man har forsøgt det. I Nordamerika rører det religiøse
Liv sig fra den praktiske Side, og den herskende Metodisme
med den dertil knyttede Revivalbevægelse har naturligt henledet
Opmærksomheden paa de sjælelige Processer, der foregaa
under religiøse Oplevelser. Det er da ogsaa Omvendelsens Psykologi,
de amerikanske Psykologer især have beskæftiget sig
med. Allerede Revivalbevægelsens Ledere vare begyndte paa en
Statistik over Omvendelsernes Hyppighed, og de første Forsøg
fra rent psykologisk Side bleve tildels gjorte i Samvirken med
dem. Ved udsendte Spørgsmaalslister har man da søgt at faa
% --- p. 100 ---
\setcounter{footnote}{0}%
\opage{100}Oplysninger om Tilstandene før, under og efter det religiøse
Gennembrud, saavel som om andre Typer for religiøs Udvikling,
der ikke havde Gennembruddets Karakter, om Kriser og Tvivl
o. s. v. Gennem Analyse af de fremkomne Udtalelser har man
saa søgt at naa til Beskrivelse af visse Hovedformer for religiøs
Udvikling --- naturligvis indenfor de givne historiske Forudsætninger
(nemlig den amerikanske Metodisme). Om meget store
Tal kunde der ikke her være Tale, og derfor ikke om streng
statistisk Metode. Man har kastet Nettet ud paa et enkelt Sted
i de religiøse Fænomeners Hav og beskrevet de Arter, man fandt
ved sit Dræt. Der er, som det Følgende vil vise, virkeligt vundet
nyt Stof for sammenlignende Psykologi ad denne Vej.

Den Første i Rækken af disse Forskere er JAMES LEUBA
med sin Afhandling \textit{Psychology of religious} \textit{phenomena} (American
 Journal of Psychology. April 1896). Leuba havde fra sin
Barndom af bevæget sig i Kredse, hvor religiøse Gennembrud
vare hyppige, og senere har han af psykologisk Interesse udvidet
sin Erfaring ved Samtaler med Mennesker, der havde oplevet
saadanne, ved at besøge Møder og ved at udsende Spørgsmaalslister
angaaende Forholdene før, under og efter Omvendelsen.
»Spørgsmaalslisterne«, siger han, »bleve sendte til Personer,
som mente at have »erfaret Religion«, til Præster og til
Ledere af indre Mission, og gennem dem ogsaa til en stor Klasse,
det ikke havde været muligt at naa direkte til. De bekendtgjordes
desuden i »The Presbyterian« i Filadelfia, i »The Christian«
i London og i »The Outlook«\dots{} En lille Procentdel
af de modtagne Svar indeholdt de Talemaader, som ere gængse
ved revivals og i Missionshuse og smagte saa afgjort af religiøs
Affektation (religious cant) og saa lidet af personlig Erfaring, at
de maatte skydes til Side; men de fleste af Svarene bare indre
Vidnesbyrd om alvorlig Bestræbelse for at give nøjagtig Beskrivelse
af en Erfaring, man tillagde umaadelig Betydning\dots{}
Udbyttet af vore Spørgsmaal begrænsedes meget ved visse Folks
Mangel paa Tiltro til Nytten af at anvende Psykologi paa religiøse
Fænomener, ved de Vanskeligheder, sædvanlige Mennesker,
der ikke ere øvede i Selviagttagelse, føle ved at skulle
gøre Rede for deres egen Erfaring, ved nogle kirkelige Autori%
% --- p. 101 ---
\setcounter{footnote}{0}%
\opage{101}teters fjendtlige Holdning, og endeligt ved den Afsky, Opfordringen
til at granske religiøse Erfaringers Natur undertiden
vakte.« Angaaende dette sidste Punkt skriver Leuba i en senere
Afhandling Følgende: »Den Grundtone af dyb, undertiden endog
tragisk Alvor, der gaar igennem de fleste af mine Dokumenter,
vilde endog stemme en ironisk Skeptiker til Ærefrygt. Den
virkelige Værdi, man tilskriver Religionen, aabenbarer sig, naar
de intellektuelle Antagelser, af hvilke det religiøse Liv menes
at afhænge, trues\dots{} En Dame af mit Bekendtskab, hvem jeg
gav en Spørgsmaalsliste, fór tilbage i stor Rædsel, da hun havde
læst de første Spørgsmaal. Hendes uklare Begrundelse af sin
Vægring ved at svare godtgjorde, at hun ikke turde se sin Tro
lige i Ansigtet; blot Tanken derom indjog hende Skræk. En
anden Dame, der med paafaldende Heftighed havde afslaaet en
personlig Anmodning om Svar, skrev et Aar senere af egen
Drift for at forsvare sin Adfærd. Hendes Tro var paa den Tid,
da Spørgsmaalet kom, i en betænkelig Tilstand, og hun kæmpede
med Fortvivlelsens Energi for, at hendes Himmel ikke skulde
miste sin Gud. Hendes vigtigste Forsvarsmiddel var at forjage
fra sin Bevidsthed enhversomhelst Ting, der kunde true hendes
Tro. Den kritiske Undersøgelse, jeg foreslog, saarede hende paa
det Dybeste. Disse gribende Exempler paa en Klasse af vore
Medmennesker, der foretrække at stikke Hovedet i en Busk
fremfor at lade de haarde Fakta berøve deres Verden den Sjæl,
der giver den Ynde, vækker til betydningsfuld Eftertanke om
Sammenhængen mellem Tro og Trang.«

I en Række senere Afhandlinger har Leuba fortsat sine Undersøgelser
(\textit{The} \textit{personifying} passion in youth. The Monist.
July 1900. --- \textit{Introduction} \textit{to a} \textit{psychological study of religion}.
ibid. January 1901. --- \textit{The content of} religious consciousness.
ibid. July 1901. --- Les tendances \textit{religieuses chez} \textit{les mystiques}
\textit{chrétiens}. Revue philosophique. Juillet—Novembre 1902). Hvad
der karakteriserer Leuba er hans Interesse for de individuelle
Tilfælde i deres Mangfoldighed og Ejendommelighed. Han har
tilvejebragt en Samling af religiøse Selvkarakteristiker, der vise,
hvormange Nuancer og Overgange der findes paa dette Omraade,
og hvor lidt vore sædvanlige Rubriker slaa til. ---

% --- p. 102 ---
\setcounter{footnote}{0}%
\opage{102}I Anlæg og i Betydning beslægtet med Leubas Undersøgelser
er en nyligt udgiven Afhandling af LUCIEN ARRÉAT:
\textit{Religion} in \textit{France}. (\textit{The} Monist. January 1903). Ved sin Undersøgelse
af fransk Katolicisme giver han et interessant Modstykke
til Leubas Undersøgelse af amerikansk Metodisme. ---

EDWIN STARBUCK (The \textit{psychology of} religion. An empirical study \textit{of} the growth of religious consciousness. London 1899) raader over et meget omfattende statistisk Materiale, som han har underkastet en gennemgaaende Analyse. Han opfatter Omvendelsen som en Proces, gennem hvilken den Unge skal gøre
det socialt Overleverede til fuld Alvor og Virkelighed i sig, og derfor som et sociologisk Fænomen, der kan paavises hos alle Folk og paa alle aandelige Omraader. Af særlig Interesse er
hans Paavisning af forskellige Arter af Omvendelse og hans Sammenligning af Omvendelsesprocessen med andre Former for religiøs
 Udvikling.

GEORGE COE (\textit{The spiritual life.} Studies in the science \textit{of}
\textit{religion}. New-York 1900) har ad lignende Vej som Leuba og
Starbuck særligt undersøgt Temperamentets Indflydelse paa den religiøse Udviklings Karakter og dets Forhold til andre medvirkende
Elementer.

GEORGE SANTAYANA har i sit Skrift \textit{Interpretations of}
\textit{Poetry} and Religion (London 1900) særligt fremdraget det poetiske
Element i Religionen og er derved bleven ført til at stille
sig i Modsætning baade til Ortodoxien og til den gængse Kritik.
For ham er Religionen den højeste Poesi, medens han tillige
søger at vise, hvorledes Poesien paa sit Højdepunkt gaar
over til Religion. ---

Gennemgaaende er de amerikanske Psykologers Standpunkt
overfor det religiøse Liv et sympatisk. De hævde energisk Psykologiens
Selvstændighed overfor al Teologi og Spekulation, men
de fordybe sig med Begejstring og Interesse i de religiøse
Fænomener. Deres eget religiøse Standpunkt synes ikke at
være udpræget antikonfessionelt, selvom de mene at kunne paavise
Nødvendigheden af, at den konfessionelle Begrænsning viger
for den værdifulde Mangfoldighed af Veje og Typer, uhildet
Iagttagelse opdager. Saavidt muligt vil jeg i den følgende Drøf%
% --- p. 103 ---
\setcounter{footnote}{0}%
% CHECK p. 103: indent after a broken word?
\opage{103}telse af deres Undersøgelser ogsaa søge at bestemme, til hvilken
Type de selv høre, for derved at bestemme den personlige Ligning og angive en Begrænsning for deres Undersøgelsers Betydning.

\subsubsection*{\textit{a. Gennembrud og Kontinuitet i den religiøse Udvikling.}}

STARBUCK fandt i det af ham indsamlede Materiale to Typer
af religiøs Udvikling. Den ene bevægede sig gennem en
mere eller mindre pludselig Afgørelse frem til en ny Livsretning;
den anden foregik mere kontinuerligt.

\textit{Gennembrudstypen} frembyder igen Forskelligheder af Interesse.

Alderen, i hvilken Gennembruddet finder Sted, kan være
forskellig. Det er hyppigst i de unge Aar og indtræder noget
tidligere for Kvinder end for Mænd; senere end i 40 Aars Alder
forekom det sjeldent (indenfor de 1265 Tilfælde, hvor Alderen kunde konstateres). --- Anledningen til Gennembruddet var forskellig
En stor Del af Tilfældene skyldtes Vækkelsesmøder
(revivals). Nogle fandt Sted under sædvanlig Gudstjeneste. Ved
andre igen kunde ingen særlig ydre Anledning paavises. ---
 Hvad Motiverne angaar, forekom Frygt langt hyppigere end
Haab, og altruistiske Motiver (Idealitet og Menneskekærlighed)
vare sjeldne. Karakteristisk er det, at Syndsfølelse var sjeldnere
ved revivals end ved andre Anledninger; Starbuck mener, at
ved revivals faar man ikke Tid til at fordybe sig i sig selv og
gruble over sin Færd og sin Fortid.

Vigtigst ere Forskellighederne med Hensyn til selve Gennembruddets
Karakter. Der fremtræder her to Former, som dog
kun tilnærmelsesvis kunne paavises i deres Renhed. I den ene
Form er det en Vane, der brydes, en lastefuld Tilstand, fra
hvilken man undslipper; forud gaar da en dyb Følelse af Synd.
I den anden Form er det et nyt Lys, der bryder ind; der udfyldes
et Savn og naas en positiv Fuldendelse; forud gaar da
en Følelse af Ufuldstændighed. Hin Form har en negativ, denne
en positiv Karakter. Den negative Form behøver fuldstændig
Overvindelse og Fortrængen af Fortiden for at blive til; den
% --- p. 104 ---
\setcounter{footnote}{0}%
\opage{104}positive Form behøver kun en lille Impuls for at træde i Kraft.
Starbuck bemærker, at indenfor Grænsen af hans Materiale er
den positive Type (spiritual illumination, forud for hvilken der
findes a sense of incompleteness) hyppigere end den negative
(escape from sin, forud for hvilken der gaar a feeling of wretchedness),
og han fremhæver det Betænkelige ved den i Kirken
herskende Tendens til at gøre den negative Type til den ene
rette og normale. De, der kæmpe for fuldere Lys, maa behandles
anderledes end de, der kæmpe for at befries fra Mørke.
Ogsaa i det af Leuba indsamlede Materiale forekommer den
positive Type hyppigst. Starbuck og Leuba drøfte ikke Aarsagen
hertil. Det er ikke ganske sikkert, at den positive Type i Virkeligheden
er den hyppigste; maaske er den negative Type mere
sky og utilbøjelig til at udtale sig, saa at deri ligger Grunden
til, at der kan konstateres færre Tilfælde af den.

Af meget stor Interesse er Vexelvirkningen af det Bevidste
og det Ubevidste ved Gennembruddet. Der er her mange forskellige
Grader af Bevidsthed, hvilket sandsynligvis hænger sammen
med, at der er forskellige Grader i de indre Hindringer, der skulle overvindes. En Forudsætning for Gennembruddet
er det --- indenfor begge Former --- at der er foregaaet en Væxt
under Bevidsthedens Tærskel (within the sub-liminal consciousness).
Nu er det ikke sagt, at den bevidste Stræbens Retning
falder sammen med den ubevidste Væxts Retning, --- og da saa alligevel Villiens Tilslutning til det ubevidst Anlagte er nødvendig,
for at Gennembruddet kan indtræde, er det intet Under
at der her foreligger saa mange psykologiske Gaader. Hvad der
skal komme i Stand, er jo, at det, der hidtil kun har staaet
som Genstand, som fjernt og ophøjet Ideal, skal blive selve det
der lever og rører sig i Individet. Hvad der hidtil kun var
Objekt, skal blive Subjekt.

Er Gennembruddet indtraadt, opstaar der en Følelse af Nyhed, i Regelen ogsaa af Glæde og Fred, og hyppigt af Selvforglemmelse
og Menneskekærlighed. Hvad der hidtil holdt Selvet tilbage indenfor snevre Skranker, er fortrængt, og en Udvidelse
kan ske.

Psykologisk set er Omvendelsen i begge dens Former ikke
% --- p. 105 ---
\setcounter{footnote}{0}%
\opage{105}noget enestaaende Fænomen. Der gives en Række analoge Erfaringer,
saaledes naar Uvished og Nedstemthed pludseligt afløses af frejdig Afgørelse, eller naar en lang Anstrengelse for
at tilegne sig en Færdighed som med ét Slag afløses af Herredømme
over Vanskelighederne, eller naar en Følelsestilstand
melder sig pludseligt ved Kontrast eller forberedt ved indre organiske
Processer. Især forekomme saadanne Fænomener i den Alder, da den Unge skal indleve sig i Slægtens Vaner og Traditioner
og gøre dem til sine egne --- skal udvide sit Selv til
at omfatte større Interesser, nedbryde Skranken mellem den individuelle
og den sociale Villie. Fysiologisk set stilles denne
Opgave paa en Tid, hvor Væxten nærmer sig sin Fuldendelse,
og hvor nye Opgaver og Funktioner begynde at kræve deres
Ret. Det er en Bevægelsens og Modsætningernes Tid; baade i
Organismen og i Bevidstheden træder Modsætningen mellem det
Gamle og det Nye skarpt frem. Men ved religiøs Paavirkning
forstærkes denne Modsætning ofte paa unaturlig og usund Maade.
Man anerkender, som allerede bemærket, ikke altid Forskellen
mellem de to Typer, men behandler dem ens. Starbuck anfører
et Exempel paa, hvorledes det kan føre til Fortvivlelse, naar
Alle skulle tvinges ind i ét bestemt Udviklingsskema og nødes til at gøre Erfaringer, som for dem ere umulige. ---

I Modsætning til Gennembrudstypen i dens to Former stiller
Starbuck \textit{Kontinuitetstypen}. Der er her ingen Bevidsthed
om afgørende Forandring, om Modsætning mellem Gammelt og
Nyt. Der sker en uvilkaarlig Opvaagnen, som vel har en vis
Lighed med den positive Gennembrudsform, men mangler Erfaringen
om pludselig Forandring. Udviklingen foregaar kontinuerligt
fra Barndommen gennem de unge Aar. Dog kan der
komme en Periode af »Storm og Trængsel« (storm and stress),
der hænger sammen med Overgangsaarenes organiske Forandringer
og kan medføre Angst og Spænding. Det er den Periode,
hvor den Unge skal indleve sig i en større Verden og gøre
Slægtens Erhvervelser til sine egne. Mangfoldige Fordringer og
nye Muligheder fremstille sig. I intellektuel Form (og noget
senere) fremtræder den saaledes vakte Uro som Tvivl. Det gælder
om med klarere Bevidsthed at fastholde, hvad man hidtil
% --- p. 106 ---
\setcounter{footnote}{0}%
\opage{106}havde antaget paa Autoritet og med barnlig Troskyldighed. En
saadan Periode var meget hyppig hos dem, Starbuck lærte nærmere
at kende; den kunde begynde, forløbe og afslutte uden
bratte Overgange. Ofte førte den over til en Tid, hvor Individet
stod fremmed overfor det Overleverede og betragtede dette som
aldeles værdiløst. Denne Fremmedhedsperiode (period of alienation)
skyldtes ofte Slappelse efter Uroen eller Tvivlen; men den
kunde ogsaa have sin Grund i, at Individet nu vilde hævde sig
selv overfor de Former, Samfundet vilde føre det ind i, eller
Lysten ved ubunden Brug af Evner og Kræfter. I Stedet for den religiøse Tradition træde i en saadan Periode ofte æstetiske, intellektuelle eller andre Interesser i Forgrunden. Ved Slutningen
af Ungdomsperioden indtræder der som Regel en Rekonstruktion:
Individet har vundet en bestemt Tro, som ikke altid stemmer
med den overleverede, men i Regelen er en Blanding
af selvstændig Produktion og Genoptagelse af det Gamle. Starbuck
er forundret over at finde denne Periode saa bestemt udpræget.
Han forklarer sig det ved, at Livsforholdene kræve et
bestemt Standpunkt i praktisk Henseende, saa at der maaske
endog i Slægten har udviklet sig en Tendens til saadan Rekonstruktion
efter de bevægede Ungdomsaar. --- Gennem alle disse
Perioder forløber Udviklingen indenfor denne Type saa kontinuerligt, at det ofte kun er gennem Sammenligning af fjernt
fra hinanden liggende Punkter, at man bliver opmærksom paa
dem.

Karakteristisk for Udviklingens Afslutning er --- og dette
gælder begge Hovedtyper --- en vis Ændring i de religiøse Forestillinger
og i hele Opfattelsen af Religionen. Selvom de konfessionelle
Forestillinger væsentligt fastholdes, kan der dog i
Regelen paavises en gennemgaaende Forskydning fra det Ydre
til det Indre, fra det Dogmatiske til det Etiske, fra det isoleret
Personlige til Bevidstheden om Delagtighed i et mere omfattende
Liv. Religionen opfattes mere og mere som et indre Liv, der
er Del af en universel Tingenes Orden.

\textit{Sammenlignes nu til Slutning} Gennembrudstypen med
\textit{Kontinuitetstypen}, fremtræder hin selvfølgeligt med langt større
Uro og med »Storm og Trængsel« i mere udpræget Grad end
% --- p. 107 ---
\setcounter{footnote}{0}%
\opage{107}denne. Gennembrudstypen er mere modtagelig for ydre Paavirkning;
Modsætningen mellem det Gamle og det Nye gør sig
stærkere gældende i den, og Følelseslivet er heftigere. Gennembrudstypen holder sig desuden langt mere end Kontinuitetstypen
til de bestemte konfessionelle Former, til de organiserede Institutioner.
Sine indre Erfaringer udtrykker den helst i bibelske
og konfessionelle Udtryk og har ingen Trang til at danne selvstændige
Former og Udtryk eller finde nye Analogier. En saadan
Trang er derimod meget hyppig indenfor Kontinuitetstypen,
hvor ogsaa Opfattelsen af Religionen som et indre Liv forekommer
hyppigere. --- Starbucks almindelige Resultat er, at det
religiøse Livs Udvikling under konfessionelle Former kan belyses
fra almindelige psykologiske og fysiologiske Synspunkter.
Der udfoldes i Religionen Anlæg og Tilskyndelser af almindelig
menneskelig Natur, og ogsaa de forskellige Former, i hvilke den
religiøse Udvikling foregaar, kendes fra andre sjælelige Omraader.
Selv betragter han Religionen »i dens højeste Form« som
en Udstraaling, et Sammenspil og en Idealisering af Tendenser,
der allerede forud ere til Stede i den menneskelige Natur. De
Former og Udtryk, Slægten har fundet herfor, føres den Unge
i Overgangsaarene for Alvor ind i, og Udviklingens Gang og
Kriser betinges ved de Vanskeligheder, her ere at overvinde.
Han er tilbøjelig til at mene, at de Ulemper, som komme frem
under denne Udvikling, langt mere skyldes de Ældres Mangel
paa Forstaaelse af Menneskenaturens Krav end Særheder og
Mangler hos de Unge. --- Med denne sidste Bemærkning kan
sammenstilles en Udtalelse af Professor PEABODY i hans Afhandling
\textit{The religion} of a college \textit{student}. (Forum. June 1901.
p. 447): »For én ung Mand, der fjernes fra Religionen gennem
Paavirkning af Videnskab, Filosofi eller Kunst, er der ganske
sikkert ti, det gaar saaledes paa Grund af uforstandig eller upraktisk
Religionsundervisning\dots{}. Det er ikke Universitetet, der
skal forny Overensstemmelsen med Kirken; det er Kirken, der
maa aabne sine Øjne for den vidunderlige Udvidelse af den intellektuelle
Horizont, som enhver Studerende i vore Dage har
Adgang til at erfare.«

% --- p. 108 ---
\setcounter{footnote}{0}%
% CHECK p. 108: centred line: heading or verse
\opage{108}

\subsubsection*{\textit{b. Temperamentets Indflydelse.}}

Ogsaa GEORGE COE fremhæver Betydningen af et psykologisk
Studium af Religionen i dens nuværende Tilstand. Og
det er ogsaa Ungdomstiden, han særligt har henvendt sin Opmærksomhed
paa. Denne Tid er, paastaar han, langt mindre
kendt og forstaaet end Barndommen og den modne Alder. Hvad
der foregaar i denne Tid, er en Overgang fra at være optagen
af det Enkelte til at fatte det Universelle, fra at overvældes
af det Ydre til at fordybe sig i det Indre, fra at lade sig
lede af Autoritet til at bestemmes ved Grunde og fra at stanse
ved det Givne til at stræbe mod Idealet. Der spørges mere
indtrængende om Livets Betydning, og de sociale Interesser
begynde at vaagne. Alt dette gør, at der føles Trang til
en mere personlig Tilegnelse af Religionens Indhold, end der
er Mulighed for i Barndommen. Omvendelse eller en tilsvarende
Inderliggørelsesproces er derfor et naturligt Led i Ungdomstidens
Udvikling, og Coe finder ogsaa faktisk, at en saadan
Forandring er overordentligt hyppig (især ved 16 Aars
Alder). Han minder om, at det hos Folkeslag rundt om paa
Jorden er almindeligt at henlægge en religiøs Indvielse til dette
Tidspunkt. Ofte sendes den unge Mand ud i Ensomheden for
at vente paa en Aabenbaring fra Stammens Gud og for ved
dens Hjælp at blive klar over sit Kald. I Tilslutning til Starbuck
hævder han, at det hos unge Amerikanere indenfor Kristendommen
ingenlunde blot er Frygt eller Anger, der ere Motiver
ved den religiøse Udvikling, men ogsaa mere ubestemte
Tilstande: Uro, Tomhed, Følelse af Ufuldstændighed, Higen mod
det Uendelige, --- selvom de Fleste udtrykke deres Erfaringer
i det konfessionelle Sprog. Det viser sig fremdeles gennemgaaende,
at rent intellektuelle Problemer og Motiver ere underordnede;
Grundene findes bagefter; man søger først at leve
det Liv, der udtrykker det dybeste Selv, og siden ser man
maaske til, hvorledes man ved Fornuftens Hjælp kan vise Berettigelsen
af det. Og det er ikke nogen enkelt Erfaring, der
fører til Maalet. Det Meste af, hvad Ungdommen fordrer afgjort
paa én Gang, kan man kun leve sig ind i efterhaanden, alt
% --- p. 109 ---
\setcounter{footnote}{0}%
\opage{109}som Livet udfolder sig med sit Arbejde, sine Skuffelser og sine
Sejre.

Coe vil studere disse Fænomener uden teologiske og metafysiske Fordomme. Fristelsens Psykologi f. Ex. er og bliver den
samme, hvorledes man saa forklarer sig Oprindelsen til de Tilskyndelser,
der medføre den, --- om man mener, at de skrive
sig fra Djævelen, eller at de ere Rester af den Dyriskhed, hvorfra
Menneskeheden har udviklet sig. Det Samme gælder om
Omvendelsens Psykologi, som Coe særligt har beskæftiget sig
med. Det er hans Opgave at paavise de vigtigste Faktorer, som
medvirke ved dette psykiske Fænomen. Han gaar her den Vej
at anstille en nøjagtig Undersøgelse af en forholdsvis lille Række
af individuelle Tilfælde. Foruden Udsendelse af Spørgsmaalslister
anvendte han indgaaende Samtale med Enkelte, indhentede
Andres Vidnesbyrd om dem, anvendte Hypnose, søgte ad alle mulige Veje at naa til en individuel Karakteristik. Antallet af
Individer, han undersøgte, var 77 (52 Mænd, 25 Kvinder); de
vare næsten alle Studerende, Mændene gennemgaaende 24 Aar
gamle, Kvinderne 22 Aar; de Fleste vare opdragne i den metodistiske
Kirke, der jo lægger saa stor Vægt paa personlig religiøs
Udvikling.

Det viste sig da, at hos saadanne unge Mennesker fremtræder
der store Forskelligheder med Hensyn til Omvendelsen. Der
kan i denne Henseende skelnes mellem tre Klasser af Individer:
A) saadanne, som ventede en Forvandling (transformation) og
virkeligt erfarede den; B) saadanne, som ventede en Forvandling,
men ikke erfarede den; C) saadanne, som ikke tilfredsstilledes
ved, hvad de erfarede, eller som dog stræbte at opnaa en
fyldigere Erfaring.

Ved nærmere Undersøgelse viser det sig nu, at Klassen A
overvejende var Følelsesmennesker (sensibility distinctly predominant),
medens Klassen B overvejende var Forstandsmennesker
(intellect distinctly predominant). Og medens Klassen A
nærmest mindede om de saakaldte melankolske og sangvinske
Temperamenter (the slow-intense and the prompt-weak varieties),
mindede Klassen B mere om det saakaldte koleriske Temperament
(the prompt-intense variety). Klassen C frembød Forskel%
% --- p. 110 ---
\setcounter{footnote}{0}%
\opage{110}ligheder, der snart mindede om A og snart om B. Temperamentet
viser sig saaledes at være en væsentlig Faktor.

Hermed forbinder sig nu en anden Iagttagelse. I Klassen
B forekom saa godt som ikke automatiske Fænomener i Sammenhæng
med de religiøse Tilstande, hverken paa Forestillingslivets Omraade (stærke Drømme, Hallucinationer), eller paa Bevægelsesorganernes Omraade (ubetvingelig Latter, Sammenfaren,
Klappen i Hænderne, Højttalen med sig selv o. s. v.), medens
saadanne Fænomener langt hyppigere forekom i Klasse A (og
i C). --- Naar Klasserne desuden undersøgtes med Hensyn til
Modtagelighed for Suggestion og Hypnose, viste det sig, at A's
Suggestibilitet overvejende er passiv, idet Individerne her ganske
rette sig efter den givne Indskydelse og følge den i det Enkelte,
medens Klassen B's Suggestibilitet kan kaldes spontan, 
idet Individerne enten give sig selv deres Indskydelser eller
selvstændigt føje Noget til den oprindelige Indskydelse.

Foruden den vakte Forventning spiller altsaa Temperament,
Tendens til Automatisme og Suggestibilitet en Rolle ved Omvendelsen. Karakteren af den religiøse Erfaring bestemmes ved
disse Faktorer. ---

Ligesom Starbuck fremhæver Coe meget stærkt Faren ved
at anse en eneste Type for ene berettiget. Vækkes der en Forventning om at opleve en mere eller mindre pludselig Sindsomvæltning, og denne Forventnings Opfyldelse efter Individets
hele psykologiske Natur er umulig, saa kan det føre til Fortvivlelse. Coe mener, at det egentligt kun er det melankolske
og det sangvinske Temperament, Kirken har anerkendt i sin
Kultus, sine Hymner og sin Moral. I den Kristustype, som
raader i Kirken, lægger dette sig særligt for Dagen; det er
Sorgernes Mand, det lidende Offerlam. Den virkelige Jesus led,
fordi han trods al Modstand fastholdt sin Tanke til det Yderste;
han var en aktiv, ikke en passiv Individualitet; der var et
kolerisk Element til Stede, der er forsvundet fra det traditionelle
Billede. --- Kirken maa være ligesaa bred, som Menneskenaturen
er det. Men den har vænnet sig til at se paa Tingene
med kvindelige Øjne: derfor er det især Mændene, der vende
sig fra den.

% --- p. 111 ---
\setcounter{footnote}{0}%
\opage{111}Coe finder, trods sin øjensynlige Sympati for Religionen,
ingen Betænkelighed ved en gennemført psykologisk Drøftelse
af de religiøse Tilstande. Hvis vi --- siger han --- ville forstaa
religiøse Erfaringer, maa vi ikke forklare dem ved Guds uransagelige
Veje, men ved paaviselige Ejendommeligheder i Menneskets
Natur; vi skulle ikke spekulere over den guddommelige
Naade, men studere den menneskelige Aands skjulte Liv. Ligesom
andre Religionsfilosofer i vore Dage finder han den sidste
Maalestok for religiøse Værdier i det Etiske. Det kommer an
paa, hvad der ved religiøs Udvikling vindes for et dybere, indholdsrigere
aktivt Liv.

\subsubsection*{\textit{c. Forskelligheder i Forholdet til Konfessionen.}}

Der foreligger her, som allerede omtalt, to Undersøgelser,
en af Leuba om amerikansk Metodisme, en anden af Arréat om
fransk Katolicisme.

LEUBA har dels studeret det religiøse Liv, som det fremtræder
hos Samtidige (i hans Hjemland Nordamerika), dels Mystikernes
Religiøsitet. Begge Studier føre ham til at gøre en
bestemt Forskel mellem Religionshistorikerens og Religionspsykologens
Synspunkt. Historikerens Emne er ikke saa meget selve
de religiøse Erfaringer, som de i Kultus, Myte og Dogme foreliggende
Produkter af disse Erfaringer, saaledes som de forme
sig indenfor et Samfund. For at naa til den umiddelbare religiøse
Erfaring maa man henvende sig til den individuelle religiøse
Bevidsthed. Af den bæres og betinges tilsidst al Religion,
og ved de højere Folkereligioner bliver dette Forhold særligt
tydeligt. Det er ikke Guderne, men Forholdet til Guderne, der
interesserer Religionspsykologien. Dennes Opgave er at give en
religiøs Dynamik, d. v. s. en Lære om det religiøse Livs Motiver.
(Udtrykket »religiøs Dynamik« havde Coe brugt før Leuba).
For at faa Materiale hertil fra Samtidige har Leuba gennem
udsendte Spørgsmaalslister indhentet Udtalelser fra en Række
Personer, der stode paa meget forskellige Standpunkter, men
alle indenfor den angloprotestantiske Verden. Spørgsmaalene
angik Grundene til deres Deltagelse eller Ikke-Deltagelse i Kul%
% --- p. 112 ---
\setcounter{footnote}{0}%
\opage{112}tus (offentlig eller privat Kultus), --- de Omstændigheder og Lo
kaliteter, der øvede særlig religiøs Indflydelse paa dem, --- de
mest karakteristiske religiøse Erfaringer, de kendte, --- de Forskelle, de gjorde mellem religiøse og ikke religiøse Tilstande og
Handlinger. Det er lykkedes ham at samle en Række karakteristiske
 Udtalelser. Om Statistik er der ikke Tale; men Leuba
mener at kunne paavise visse Ejendommeligheder ved det reli
giøse Liv i den Kreds, han har med at gøre.

Et Hovedpunkt er for ham det Resultat, at hvad der ytrer sig i Religionen, ikke er en Stræben efter Erkendelse, men en Stræben efter Væren. Religionen udspringer, efter hans Opfat
telse, af en Følelse af Brudstykkeagtighed (unwholeness) i For
bindelse med en Følelse af indre Delthed og Splittethed. Hin
Følelse medfører Trang til at blive delagtig i et større, mere
omfattende Liv, medens Delthedsfølelsen medfører Trang til at ledes af et eneste Formaal (unification of motives). Syndsfølelsen er en Form for denne sidste Følelse. Til de to nævnte Motiver slutte sig endnu Altruisme og Afhængighedsfølelse: den første
er beslægtet med Trangen til Enhed, den anden er betinget ved
den indre Modstand, der er at overvinde under Bestræbelser
for at naa til Enhed. De store Religionsstiftere agtede ikke paa
teoretiske Problemer, men stræbte efter ideal Fuldkommenhed og indre Enhed. Og der ytrer sig nu til Dags mere og mere
en Trang til at gaa tilbage til de umiddelbare religiøse Erfaringer, af hvilke Religionen stedse fra først af er voxet frem. Skønt man endnu ikke har haft Mod eller Forstand nok til at gøre dette fuldt ud, saa ere dog de første Skridt i denne Retning
(to return to plain religious empiricism) gjorte, især indenfor
Protestantismen. I Stilhed føres der --- ved Siden af Kampen om Tilværelsen mellem Individerne indbyrdes --- en rent indre
Kamp, der i Strenghed og tragisk Alvor ofte overgaar den ydre. Den føres paa forskellig Maade indenfor de forskellige Religioner og finder Udtryk i forskellige Former: men stedse hænger den sammen med Trangen til et større og tillige mere udelt Liv. Den religiøse Trang er egentligt ingen ganske speciel Trang, men har sit dybe og varige psykologiske og fysiologiske Grundlag.

Ved Sammenligning mellem Religionen, som den fremtræder
% --- p. 113 ---
\setcounter{footnote}{0}%
\opage{113}i nyere Tid i Modsætning til tidligere Former for den, finder
Leuba en Række af Forskelligheder.

I primitive Samfund er Religionen et sammensat Produkt
af alle Menneskets fundamentale Fornødenheder og Bestræbelser,
baade teoretiske og praktiske. De teoretiske Motiver have
nu udskilt sig og finde deres Tilfredsstillelse i Videnskaben; det
er ikke paa dem, at Religionens Betydning i nyere Tid beror,
og allerede i Oldtiden staar Buddhismen som Vidnesbyrd om,
at disse Motiver ikke ere de afgørende. Selvom alle hidtil til
Religionen knyttede intellektuelle Elementer falde bort, vil Religionens
psykologiske Væsen, saaledes som det ovenfor er beskrevet, kunne bestaa.

En anden betydningsfuld Forskel mellem mere primitiv og
mere moderne Religion finder Leuba deri, at medens forhen
Følelserne af Lydighed og Ærefrygt vare fremherskende, ere
disse nu --- i hvert Tilfælde i den angelsaksiske Verden --- ikke
saa hyppige, som det antages. Ikke den Lydighed, Katekismen
indskærper, men den klare Forstaaelse af, hvad der er ret og
godt, er nu Motiv til Handlen, og Enheds- og Frihedsfølelsen
er ved at fortrænge Underkastelsen. Hermed staar det Træk i
Forbindelse, at religiøs Sandhed ikke længere betragtes som
Noget, der kan og skal begrundes ad objektiv Vej; det Afgørende
er, hvad der tilfredsstiller den personlige Trang og Higen. Mange af de Udspurgte udtalte dette med største Uforbeholdenhed.
Følelsesmotiverne havde afgjort Overvægten over
alle intellektuelle Interesser, og der er maaske endog fuldstændig
Ligegyldighed overfor den rent objektive Betydning af de
religiøse Forestillinger, selv hos Individer, som ellers paa speciellere
Omraader viste sig i Besiddelse af kritisk Sans og Evne.
Man kender ikke Gud; man \textit{forstaar} ham ikke; men man \textit{bruger}
ham. Ikke Gud, men Liv, mere Liv, højere, rigere, mere fyldestgørende
Liv er Religionens Formaal. I Fortiden viser dette sig klart hos de saakaldte Mystikere, og her slutter Leubas Studier
over Mystiken sig til hans Studier over samtidige Fænomener. Det mystiske Liv maa ikke fortolkes og bedømmes efter
 de Begreber og Dogmer, Mystikerne anvende. Dets Ejendommelighed er kun en mere extrem og koncentreret Form for
% --- p. 114 ---
\setcounter{footnote}{0}%
\opage{114}Ejendommeligheden ved al Religion: en Higen efter højere og
fuldstændigere Tilfredsstillelse, efter Fred og Enhed i Villien.
Mystikernes Higen er mere lidenskabelig end den sædvanlige
Religiøsitets, og Koncentrationen mod et eneste Maal er stærkere.
Maal og Retning ere ens, men de gaa frem med langt
større Energi.

Endnu et karakteristisk Træk i den nyere Tids Religiøsitet
finder Leuba deri, at der fremtræder saa mange og betydelige individuelle Forskelligheder, selv indenfor en og samme Civilisation. Under Indflydelse af den uendeligt større Horizont
Livet har faaet siden Konciliernes Tid, er ogsaa Religionen
undergaaet Forandringer, og der viser sig Forskelligheder ikke
blot i de ydre Ting ved Religionen, men i Individernes inderste
Sjæleliv. Fornødenhederne ere forskellige, og Forhaabningerne
ere forskellige. Og dog gaar man i Teologien ud fra, at en og
samme Føde passer for Alle. ---

Indenfor sin Samling af Selvbekendelser paaviser Leuba en hel Række af Nuancer ligefra et Standpunkt, hvor det Religiøse
og det Konfessionelle gaa i den Grad i Et, at Individet ikke formaar
 at udtrykke sig om sine indre Oplevelser uden i bibelsk
og konfessionelle Udtryk, til saadanne Standpunkter, hvor de
aldeles ikke kan forbindes nogensomhelst Mening med den Art Udtryk. Mellem begge Yderpunkter ligge de interessanteste
Tilfælde; ti hverken af dogmatisk Position eller af dogmatisk
Negation er der nogen synderlig Psykologi at lære. Kun nogle
enkelte Træk kunne her anføres fra de interessante Selvkarakteristiker,
Leuba har samlet.

En videnskabeligt dannet ung Kvinde, Tilhænger af den
presbyterianske Kirke, skelner bestemt mellem, hvad der skyldes
hendes konfessionelle Opdragelse, og hvad der skyldes Trangen
til Noget, som har evig Værdi (a desire for things of permanent
value). Disse to Ting udelukke ikke hinanden, men de falde
heller ikke sammen. Naar hun deltager i sin Kirkes Kultus, er
det dels af Vane, men dels ogsaa fordi hun trænger til at træde
i Forhold til et højere Væsen og befries for selviske Tanker.
Forskellen mellem religiøse og ikke-religiøse Tanker og Følelser
er for hende en Gradsforskel og beror paa, at hine fremtræde
% --- p. 115 ---
\setcounter{footnote}{0}%
\opage{115}
som betydningsfuldere og stærkere. Men den ideale Trang tilfredsstilles
aldrig fuldt. »Medens de gængse religiøse Paastande
erklære Religionen for den eneste Ting i Verden, der kan stille
tilfreds, beror«, siger hun, »for mig dens Værdi paa, at den
stadigt driver mig frem til nye Anstrengelser.« --- Andre Korrespondenter
fremhæve Menneskekærlighed, barnligt Sind eller
Taknemmelighed som det Væsentlige i, hvad de kalde Religion.
--- En Læge, som deltager i Kirkens Kultus, skønt hans egen
private Andagt for ham har en mere afgjort religiøs Karakter,
udtaler bl. A. Følgende: »Jeg er langt fra at have den antropomorfistiske
Forestilling om Gud, som jeg havde i mine unge
Aar. Ved Gud forstaar jeg den Magt eller det Princip, der bevirker
al Udvikling i Universet henimod en højere Tingenes
Orden, og som ved sine Love hæver Livets Stade fra det Ubevidste
til det Bevidste, fra det Dyriske til det Fornuftige, fra
det Forgængelige til det Uforgængelige. Hos mig selv kalder
jeg alle Drifter, Tendenser og Anstrengelser guddommelige, som
føre til at udvikle hos mig større Modtagelighed for Skønhed,
Sandhed og Retfærdighed, højere Fornuft og Evne til moralsk
Fremgang«. --- En beslægtet Udtalelse fra en anden Korrespondent er følgende: »Jeg har opgivet ethvert Forsøg paa at forestille
mig Gud og føler ikke mere nogen Trang til et Billede
eller en Iklædning af det guddommelige Princip, som det er en
Nødvendighed for mig at tro paa. Mit religiøse Forhold til det,
man kalder Gud, kan jeg beskrive fra to Synspunkter: 1) Jeg
føler i min inderste Bevidsthed Trang til at være Et med et stort Princip, Trang til Enhed med Noget, som jeg ikke kan
definere, men som omgiver mig, som jeg saa at sige er dykket
ned i, og som jeg ønsker at være i Harmoni med, et altomfattende
Nærværende\dots{}. 2) Jeg har næsten uafbrudt en dyb
 Bevidsthed om at staa langt under et indre Ideal, en Type, der
er anlagt i mig, og som jeg stedse har for Øje. Jeg føler, at
jeg \textit{maa} virkeliggøre denne Type, der stadigt bor i mig og søger
at lede mig. Erfaring har lært mig, at jeg bliver lykkelig ved
at følge den, ulykkelig ved at fjerne mig fra den. Til disse
to fundamentale Instinkter eller Institutioner har jeg reduceret
alle mine religiøse Forestillinger. Jeg er en Discipel  af Kristus
% --- p. 116 ---
\setcounter{footnote}{0}%
\opage{116}og stræber at efterligne ham; mit Ideal er at blive Et med det
Guddommelige, ligesom han var det. Men der er mange af de
store kristelige Samfunds Antagelser, jeg ikke deler.« --- En
Kvinde, der lever med Mand og Barn paa Landet i stadig Kamp
med Fattigdom, erklærer, at for hende er Religion Sandhed, og
ved Sandhed forstaar hun den Lov, der gælder for vort Væsen.
»Naar jeg siger Gud, mener jeg det Gode, mit Væsens Lov\dots{}
Jeg tror, at Liv er Udvikling, Væxt\dots{} og at vi have saa megen Himmel og saa meget Helvede her i Livet, som vi selv
ville\dots{} Det Gode er vort Ideal; det Onde er vor Lærer, som
skal faa os til at forstaa det Gode\dots{} Mennesker kunne ikke
deles i gode og onde, men i udviklede og uudviklede\dots{} Jeg
kender ingen religiøs Handling uden selve det at leve\dots{}
Maaske var jeg nærmest ved religiøs Erfaring, da det første
Gang blev klart for mig, at jeg var fri, fri for Tro, Lære, Form og Sentimentalitet, fri til at voxe, til at blive Alt, hvad der
ligger i min Natur\dots{} Hvis De ikke kan finde nogen Religion
i dette, vil De neppe finde nogen Religion hos mig\dots{}. Jeg
er imod Læren om Kristi Efterfølgelse, fordi den strider mod
selvstændig Tænkning, og fordi sjælelig og aandelig Udvikling
beror paa selvstændig Tænkning. Vi kunde ligesaa godt lade
en Anden spise for os som tænke for os. Begejstring kan en
Aand indgive en anden, men ikke Existens\dots{}. Kristendommen har været en stor Faktor i Udviklingen af den menneskelige
Aand, men Dogmernes Kristendom er man nu næsten voxet
fra\dots{}. Vi ere ikke ved at miste vor Religion, men vi ere
ved at finde os selv. Vi betvinge Frygten og lære at have Tillid
til den guddommelige Lov, der raader for Livet«.

Alle disse Exempler høre til den Gruppe, der mener at
være positivt religiøs, skønt det ofte, især for det sidste Exempels
Vedkommende, er tvivlsomt, med hvilken Ret. Men de ere
karakteristiske Udtryk for Trangen til personligt at gøre sig
Rede for, hvad man --- den Enkelte --- kan finde af Værdi
for sit Liv i den givne positive Religion. Gennem nogle Overgangstilfælde
af mere ubestemt Form komme vi i Leubas Række
over til afgjort ikke-religiøse Former, som --- hvor langt man
saa strækker Begrebet Religion --- ikke kunne komme ind under
% --- p. 117 ---
\setcounter{footnote}{0}%
\opage{117}det. Dels mangler der her religiøs Erfaring i enhver Betydning
af Ordet, dels er der bestemte intellektuelle og moralske Grunde,
som føre bort fra Religionen --- ganske vist fordi Religion og
Kirkevæsen antages at falde sammen. Saaledes bemærker en
Korrespondent: »Jeg benegter Kirkernes moralske Overlegenhed.
De tage aldrig Ledelsen af en moralsk Reform. De vare paa
den gale Side i Slavespørgsmaalet. De ere fejge, smaatskaarne
og trællesindede i vor Tids store industrielle Spørgsmaal.« ---
Leuba mener, at Antallet af ikke-religiøse Personer er langt
større, end man i Almindelighed tror. Naar der ikke er mange
Udtalelser fra saadanne i hans Samling, er det efter hans Mening
dels, fordi hans Spørgsmaalslister ikke have naat til dem,
dels, fordi de ikke altid ere klare over deres virkelige Standpunkt.
---

Leuba har ikke selv afsluttet Analysen af sit Materiale;
hans hidtil foreliggende Afhandlinger ere kun Resultater af en
første Behandling, og han ender egentligt med et ubesvaret
Spørgsmaal. I det religiøse Grundmotiv, de religiøse Tendenser
finder han Intet, der er særligt for Religionen, men særegne
Former for Selvopholdelses- og Selvudviklingstrangen. Denne
Opfattelse finder han afgjort bekræftet ved sine Korrespondenters
Udtalelser. Han afviser enhver Opfattelse af Religionen,
hvorefter denne skal gaa ud paa Erkendelse, hvad enten denne
Erkendelse skal give en Aabenbaring, eller konstatere et Mysterium.
(The Monist. July 1901. p. 571---573). Men samtidigt
med at han i sin sidste Afhandling (Revue philos. Novbr. 1902. p. 487) viser tilbage til disse Udtalelser, udtaler han: »Hvad
er da det Element, som adskiller det religiøse Liv fra det verdslige
(ti der er dog en Forskel imellem dem)? Det er det Middel eller de Midler, Mennesket anvender for at opnaa sine Ønsker,
og ikke selve disse Ønsker, der adskille Religionen fra det
øvrige Liv. Psykologisk udtrykt: det er ikke Villien (i Schopenhauers Betydning af Ordet), men Erkendelsen, der skiller dem.
 Naar man indser, at Forskellen mellem de religiøse Mennesker
og dem, som ikke ere det (forsaavidt der er saadanne), ikke bestaar
i Stræben, Ønske og Tendens, altsaa ikke i selve Livets
 Kilder, men kun i de Midler, der anvendes for at naa samme
% --- p. 118 ---
\setcounter{footnote}{0}%
\opage{118}Maal, fremtræder Religionen i et andet Lys end det, i hvilket
den sædvanligvis staar i Folks Øjne.« --- En nærmere Forklaring
af disse Antydninger ville forhaabentligt Leubas følgende
Arbejder bringe. Men naar han bestemt afviser den intellektualistiske
Opfattelse, at det skulde være Religionens \textit{Formaal} at
give Erkendelser af nogensomhelst Art, og dog mener, at Forskellen
mellem Religionen og andre aandelige Livsformer \textit{beror}
\textit{paa} Erkendelse, saa kan Meningen vel kun være den, at Forskellen
 beror paa den Maade, paa hvilken Erkendelsen i Religionen bruges som \textit{Middel}. Og Middel maa her betyde Fortolkningsmiddel, altsaa Symbol. Jeg tror ogsaa, at Leuba af sit Materiale vilde kunne uddrage et Resultat i denne Retning. Der er i al Religion en Trang til at forme det Oplevede i Tanker.
I den konfessionelle Religion tilfredsstilles denne Trang
umiddelbart ved de overleverede Forestillinger; i Overgangsformerne søges der efter andre Udtryk til at fortolke selve disse
overleverede Forestillingers Betydning for Livet. Til Grund ligger stedse en Erfaring om Forholdet mellem Ideal og Virkelighed, en Erfaring, der --- hvor Religion idethele opstaar ---
gaar over til Tro og Tillid overfor Magter, der antages at virke
i Retning af et harmonisk Forhold mellem Ideal og Virkelighed. Jeg finder min Hypotese om Religionens Væsen som Tro
paa Værdiernes Bestaaen bekræftet ved Leubas og hans Kollegers
Undersøgelser.

Derimod forekommer det mig, at Leuba lægger for lidet Vægt paa herskende eller overleverede Forestillingers Indflydelse. Det er skæbnesvangert for en Følelses hele Karakter at blive
knyttet til en bestemt Kreds af Forestillinger; den faar herved
en Retning og et Indhold, der afgør dens Kvalitet, og naar
Forbindelsen mellem den og de Forestillinger, den hidtil har
været knyttet til, skal sønderrives, kan det ske, eller i hvert
Tilfælde true med, at Følelsen forsvinder med det Samme. Man
kan ganske vist med Leuba skelne mellem Formaal, som kunne
siges at være fælles for Religionen og andre aandelige Livsformer,
og Midler, der ere forskellige for hin og for disse; men
det maa da vel fastholdes, at selve det, at et Maal kan naas
ved forskellige Midler, tyder paa en Forskellighed i Maalet selv.
% --- p. 119 ---
\setcounter{footnote}{0}%
\opage{119}En ganske udvortes Sammenhæng mellem Maal og Midler er
der paa saadanne Omraader, som det her drejer sig om, ikke.
Det religiøse Problems Brod ligger netop paa dette Punkt. Historien
lærer jo ogsaa, hvor ondt Mennesker have ved at finde
sig i Kvalitetsforskelligheder indenfor Religionen, og Striden om
Grænserne for selve Begrebet Religion beror ligeledes paa det
Spørgsmaal, om vi ved en Kvalitetsforskel komme udenfor Religionens
hele Modalitet (for at bruge Sansepsykologiens Terminologi).
Det bliver herved af særlig Betydning, at der gives
Naturer, som kun formaa at finde Udtryk for deres indre Liv
i \textit{fælles} Symboler, og som, hvis de nødes til at opgive dem,
føle sig aldeles hjemløse.

Dog maa det gentages, at Leuba ikke endnu har fuldendt
Bearbejdelsen af sit Materiale. Det vil være af Interesse at
høre de Ord, hvormed han afslutter sine foreløbige Meddelelser.
» Om Forskellighederne mellem Religion og Ikke-Religion kun
ere overfladiske, eller om de ere fundamentale, og da i hvilken
Grad, --- om Opdragelsen i Hjemmet og andre Omstændigheder
i og for sig kunne forklare disse Forskelligheder, eller om de
have dybere Rødder i særegne, medfødte Anlæg, --- om den ene
er bestemt til at fortrænge den anden eller blot er et forberedende
Stadium for den eller ikke, --- det er et Spørgsmaal, som
maa forbigaas i Taushed i denne foreløbige Afhandling. Kun
det ønsker jeg at betone her: selv mellem Individer af samme
Samfund og omtrent af samme Intelligens og Kultur gives der
Forskelligheder, som ere tilstrækkelige til at gøre den Ene til
et Menneske, for hvem Religion er Livets Alt, og lade den Anden
staa aldeles ligegyldig og utilgængelig overfor den, selvom den
fortolkes i videste Forstand.« ---

LUCIEN ARRÉAT har --- ad lignende Veje som de amerikanske Psykologer --- søgt at finde forskellige Typer indenfor
den franske Katolicisme. Det mærkes strax, at Forholdene ere
andre. Forskellen mellem Katolicisme og Protestantisme hænger
utvivlsomt sammen med Raceforskelligheder. Medens man fra
anden Side (Charles Renouvier) har betragtet det som det franske
Folks Syndefald, at det ikke sluttede sig til Reformationen, hævder Arréat, at denne stred mod den franske Races Ejendomme%
% --- p. 120 ---
\setcounter{footnote}{0}%
\opage{120}lighed. Protestantismen krænkede Skønhedssansen ved sin Ødelæggelse
af Billeder og Bygninger, og hvad den berøvede Fantasien,
gav den ikke Erstatning for paa Tankens Omraade. Den
protestantiske Individualisme fordrer af den Enkelte en selvstændig
Drøftelse og Forstaaelse, der strider mod den latinske
Races Utilbøjelighed til teologiske Diskussioner. Katoliken ser
Læren fastsat en Gang for alle, og hans Inderlighed bryder frem
i det personlige Følelsesforhold, den begejstrede Tilslutning til
den faststaaende evige Sandhed, som Kirken garanterer. Arréat
har derfor heller ikke i sit Materiale fundet nogen Antydning
af den Overgang fra Lydighedsfølelse til Frihedsfølelse i Gudsforholdet, som Leuba mente at finde i Angloprotestantismen; den katolske Religiøsitet er væsentligt Underkastelse. --- Efter sine
Iagttagelser beskriver Arréat fire forskellige Klasser inden
for den franske Katolicisme.

Den første Klasse, der især findes i udpræget Form i Sydfrankrig,
indbefatter den største Mængde af Katoliker og har
Repræsentanter i alle Klasser af Samfundet. Fromheden bestaar
her i Overholdelse af de religiøse Skikke, der opfattes paa en
Maade, som nærmer sig til Afgudsdyrkelse, hvad der især viser
sig i den stærke Tilstrømning til de Steder, hvor der sker Mirakler.
Her virker den Racekarakter, hvis Tilbøjelighed til Overtro
allerede Cæsar omtalte.

De sande Katoliker udgøre den anden Klasse. Gennem
Deltagelse i Kirkens Naademidler styrke de Tryghedsfølelsen
ved at være ombord paa St. Peters Skib, der midt i Tidernes
Uro føres af en ufejlbar Styrmand. Den katolske Religion prises
for sin Ophøjethed og sin betagende Skønhed, for sit Fuldkommenhedsideal
og for sin virksomme Menneskekærlighed. Uden
Trang til Mystik eller til selvstændig Tænkning ytrer der sig
hor en aandelig Ligevægt, som dog undertiden er forbunden
med en udtrykkelig Forvisning om, at det kun er Kirkens vældige
Myndighed, der holder Tvivl og Fare borte.

Til en tredie Klasse regnes de, der støtte deres kirkelige
Tro ved mere eller mindre bevidst Tænkning. Mange Gejstlige
høre herhen, og til denne Klasse mener Arréat ogsaa at kunne
henregne Mange, som vægrede sig ved at udtale sig. Man pro%
% --- p. 121 ---
\setcounter{footnote}{0}%
\opage{121}testerer her imod, at Religionen skulde være en Følelsessag:
den grunder sig paa rationel Erkendelse, skønt Naaden bringer
den sidste Afgørelse. Den katolske Lære staar som en Helhed,
af hvilken Intet kan udskilles, og som umiddelbart betager Sjælen.
Dog udtales det undertiden, at Kirken kun er Troens
ydre og forgængelige Klædebon, der vil kunne omdannes til
nye Former; det egentligt Religiøse sættes i Trangen til det
Uendelige. Ved saadanne Udtalelser gøres Overgangen til den
sidste Klasse.

Hos mange Mænd og Kvinder, der erklære sig selv for
gode Katoliker, kan en vis Opløsningstilstand spores. De vælge
f. Ex. mellem Kirkens Lærdomme og Skikke og holde sig til,
hvad de have Brug for. Ofte er det kun en udtrykkelig Villiesakt,
der holder dem tilbage i Kirken; man betoner de sociale
Motiver, Lejligheden til at udøve Menneskekærlighed f. Ex.,
som det, der knytter til Kirken. Eller man opfatter det Religiøse
som ideal Stræben, hylder en praktisk Idealisme uden
Metafysik og benytter kun Kirkens Kultus som Stemningsmiddel.
Man udtaler den inderligste Sympati med de Troende, fordi
de ikke, som mange Fritænkere, give Afkald paa Religionens
ophøjede Poesi eller ere blaserede overfor de vældige Problemer,
Ordet Gud peger hen paa. Et sidste Skridt gøres, naar
man gaar over til Positivismen eller ad anden Vej søger et
helt nyt Udtryk for sin Livsanskuelse. Dette er dog usædvanligt.
Enten vender man tilbage til Barndommens Tro efter en
Uroens Tid, eller man synker i aandelig Henseende ned paa et
lavere Trin.

Det er --- udtaler Arréat til Slutning --- umuligt for Øjeblikket
at sige, hvorhen den nuværende Krise paa det religiøse
Omraade vil føre, eller hvilken Energi den religiøse Følelse vil
kunne besidde, naar den ikke støttes til Dogmetro eller poetiske
Myter: »Jeg ser, hvad man vil udrydde af det moderne Sjæleliv;
jeg ser ikke, hvad der ydes til dets Berigelse.«

% --- p. 122 ---
\setcounter{footnote}{0}%
% CHECK p. 122: centred line: heading or verse
\opage{122}

\subsubsection*{\textit{d. Poesi og Religion.}}

GEORGE SANTAYANA har i sit ovenfor nævnte Skrift drøftet
Forholdet mellem Poesi, Religion og Videnskab paa en Maade, der giver ham en Plads i Religionspsykologernes Række. Det er for ham Fantasien, den billeddannende Evne, der ligger til Grund for alle hine tre aandelige Livsytringer. De Billeder, vi kunne bruge i Praksis, og ved hvis Hjælp vi kunne forudsige fremtidige Erfaringer, kalde vi Forstandsforestillinger, medens
Billeder, der ikke kunne udnyttes paa denne Maade, vedblive at hedde Fantasiforestillinger. Forstand er en Fantasi, der gør
Erfaringsvidenskab mulig. Vi vælge under Erfaringens Indflydelse
 mellem de uvilkaarligt opdukkende Billeder. Men ved de
Fantasibilleder, der ikke kunne gaa over til Forstandsforestil
linger, maa igen skelnes mellem to Arter. Naar Fantasibillederne, skønt de ikke udtrykke empirisk Realitet, dog staa som Udtryk
for Livsværdier, Idealer og Formaal, faa de Religionens Karakter, medens de staa som rent poetiske Former, naar de ingen saadan Livsbetydning have. Religion og Poesi ere i deres inderste
 Væsen beslægtede, ja identiske, og blot forskellige ved
deres forskellige Forhold til Livet. Poesien gaar paa sit Højdepunkt over til Religion, naar den udtrykker vore Erfaringers
Værdi og giver store, ledende og begejstrende Antecipationer. Omvendt gaar Religionen over til Poesi, naar den opgiver Forudsætningen om at indeholde metafysiske Sandheder. Det gælder, for ret at forstaa dette, at skelne mellem Værdiernes og
Kendsgerningernes Verden: Værdi beror ikke paa selve det Givne, men paa dets Betydning.

Vor Villie er uendelig, men vor Erfaring er begrænset. Poesien fører os ud over denne Begrænsning og gør derved en højere Udvikling og Øvelse af Villien mulig, idet dennes højeste
Maal nu ved Poesiens Hjælp kan fremtræde med større Klarhed
end den blotte Erfaring tillader. En Stræben, en Lidenskab
forstaa vi kun ret, naar vi se det Ophøjede i dens hele og fulde
Udfoldelse. Men fordi Poesien bliver Religion, hører den ikke
op at være Poesi. Paa dette Punkt forsynde ikke blot de Ortodoxe,
men ogsaa deres Kritikere sig. Allerede den græske

% --- p. 123 ---
\setcounter{footnote}{0}%
\opage{123}Filosofis Kritik af Polyteismen var uretfærdig, idet den ikke tog
Hensyn til Mytologiens nøje Sammenhæng med Livet. En stor
Religion er et Produkt af Menneskeaanden i dens Helhed, og i
Forhold til det staar Kritiken, der ledes af den specielle, intellektuelle
Interesse, som begrænset og smaatskaaren. I Kristendommen
traadte historisk Idealisme i Stedet for den antike Religions Naturalisme; den Harmoni, den gamle Religion havde
søgt at fremstille i Rummet og i Naturen, søgte den nye at
fremstille i Tiden og i Historien. Sit Fodfæste i Historien fik
Kristendommen derved, at den havde sin Rod i en national Tro,
der havde formet sig under Jødefolkets Prøvelser; men ud herfra
formede den et stort poetisk Udtryk for Erfaringer, Menneskene
kendte fra deres Sjæls Inderste. Det var Menneskehedens
store Selvbiografi, der her blev gjort til Verdenslivets Midtpunkt.
Den Betydning, en Religion har og bevarer, er ikke nødvendigvis
den, som den havde for dens Grundlæggere; disse betages
saaledes af deres Syner, at de betragte disses Indhold som ydre
Kendsgerninger. Religion er ikke forskellig fra Overtro i psykologisk Henseende; ikke paa Oprindelsen, men paa Værdien
beror Forskellen. Selvom den etiske Værdi ikke var Grunden
til Religionernes Tilblivelse, saa er den dog den sidste Maalestok
for deres Betydning. ---

Hele denne Tankegang, som Santayana udfører i det Enkelte
gennem interessante Karakteristiker af religiøse og poetiske Værker, faar en særegen Interesse, naar man sammenholder den
med nogle af Arréats, Leubas og Coes Iagttagelser, og naar man
tillige mindes de idealistiske eller mystiske Tendenser, der røre sig i den nyeste Tids Literatur. Hvad jeg savner hos Santayana
er en mere indtrængende psykologisk Analyse af Forholdet
mellem Erfaring og Tro; men iøvrigt har det for mig været af
stor Interesse at finde en med min egen beslægtet Betragtningsmaade
i et Skrift, der er udkommet omtrent samtidigt med
mit eget religionsfilosofiske Forsøg. Jeg tillader mig til Sammenligning
at anføre en enkelt Udtalelse fra min Bog (Religionsfilosofi
p. 334 f.): »Maaske forholder det sig saaledes, at Poesi
er et fuldkomnere Udtryk for det Højeste, end noget videnskabeligt
Begreb vilde kunne være. Ved Poesi maa da her ikke
% --- p. 124 ---
\setcounter{footnote}{0}%
\opage{124}tænkes paa løse Stemninger og Fantasier, men paa den uvilkaarlige
og livfulde Form, det virkeligt Oplevede iklæder sig i
stærkt bevægede Øjeblikke\dots{} Det er en Poesi, der hverken
staar i Modsætning til Villie eller til Tanke, selvom den særligt
bliver til, hvor Tanke og Villie støde paa Skranker\dots{}. Al høj
Kunst forudsætter en saadan Livspoesi og er dens Omsætning
i klare Former.«

\subsubsection*{\textit{e. Sammenligning mellem forskellige Typerækker.}}

I min Religionsfilosofi har jeg, især støttet til et biografisk
og historisk Materiale, fremstillet en Række typiske Forskelligheder
i religiøs Tro (p. 111—120) og, i en anden Sammenhæng,
en Række Forskelligheder i Henseende til de oprindelige Dispositioner,
der betinge det religiøse Livs Ejendommelighed hos
de Enkelte (p. 251—256). Sammenligner man disse Rækker med
de Typer, som de amerikanske Religionspsykologer have opstillet,
vil man især finde én Modsætning. som i alle Fremstillinger
spiller en Hovedrolle, nemlig den, jeg har beskrevet som
en Modsætning mellem Tro som Koncentration og Hvile efter
Udligning af Disharmonier og Tro som højeste Form for Livsudfoldelse.
Hertil svarer indenfor Dispositionerne Modsætningen
mellem disharmoniske og expansive Naturer. Baade hos Leuba,
Starbuck og Coe er dette den mest fremtrædende Modsætning.
Men det viste sig tillige, at de lægge ganske særlig Vægt paa
den expansive Type og ere tilbøjelige til at betragte den disharmoniske
eller diskontinuerlige Type som et specielt Tilfælde
af den. Der er naturligvis altid nogen Modstand at overvinde,
nogle Kriser at gennemgaa, og det er kun, naar Modstand og
Kriser blive særligt stærke, at vi efter disse Forskeres Mening
komme over fra den expansive til den disharmoniske Type.
Om Grunden til, at den expansive Type saa hyppigt forekommer
i deres Materiale, har jeg ovenfor udtalt en Formodning.
Man kunde maaske føje til, at de nævnte Forfattere (og det
Samme gælder aabenbart Santayana) selv personligt afgjort
synes at høre til den expansive Type, forsaavidt man kan lære
deres psykiske Habitus at kende af den Maade, hvorpaa de op%
% --- p. 125 ---
\setcounter{footnote}{0}%
\opage{125}fatte og vurdere Andres. Jeg for mit Vedkommende vilde være
tilbøjelig til at anse Forskellen mellem de to Typer for større
og mere fundamental, end hine Forfattere antage. Maaske ligger
der i denne Modsætning Spiren til en aandelig Kamp, der atter
og atter vil komme til at staa mellem aandelige Livsretninger.
Dybest inde forstaa de tilsidst ikke hinanden. I det danske
Aandsliv have vi et klassisk Exempel paa denne Modsætning i
Forholdet mellem Grundtvigianismen og den indre Mission.

Skulde jeg reducere min Typerække, vilde denne Modsætning
dog blive staaende. Jeg har idethele syv Typer, nemlig
foruden de to allerede nævnte [Troen som 1) Hvile og som
2) Udfoldelse]: 3) Kontemplation, 4) tillidsfuld Voven, 5) Resignation,
6) krampagtig Religiøsitet og 7) Kulsviertro (fides implicita).
Af disse vilde 3, 4 og 7 gaa ind under 2, medens 5 og
6 vilde høre ind under 1. Hvad Forskellighederne i Dispositioner
angaar, opstiller jeg paa Følelses- og Villieslivets Omraade
følgende Række af Modsætninger: mellem de Disharmoniske
og de Expansive, mellem de Passive og de Aktive, mellem
de Affektive og de Kontinuerlige og mellem de Idiopatiske og
de Sympatiske, --- og paa Erkendelsens Omraade Modsætningen
mellem analytiske og syntetiske Naturer og mellem diskursive
og intuitive Naturer. Det er let at se, at der er en Analogi
mellem alle disse Modsætningsforhold, og at de alle kunne indordnes
under det første, saa at det bliver Modsætningen mellem
de Disharmoniske og de Expansive, der ogsaa fra denne Side
set kommer til at staa i første Linie.

I den Overbevisning, at der for Fremtiden i langt højere
Grad end hidtil ved Forstaaelse og Bedømmelse af religiøse
Fænomener maa lægges Vægt paa de individuelle Forskelligheder,
stemmer jeg ganske overens med mine amerikanske
Kolleger.

\subsection*{4. KOMPARATIV RELIGIONSPSYKOLOGI.}

Et sammenfattende Arbejde paa Religionspsykologiens Omraade
foreligger i den amerikanske Filosof WILLIAM JAMES
% --- p. 126 ---
\setcounter{footnote}{0}%
\opage{126}Bog \textit{The varieties of} religious experience. A study in human
nature. London, New-York and Bombay. (1902). James hører
til vor Tids mest fremragende Tænkere. Han forener omfattende
Kundskab med stor Iagttagelsesevne, skarp Kritik med
idealistisk Begejstring, Fordomsfrihed med alvorlig Overbevisning. Det betydeligste af hans Skrifter er \textit{Principles of psychology} (1890), der paa flere Punkter af Psykologien gør nye Synspunkter gældende og overalt drøfter Spørgsmaalene paa en aandfuld
og vækkende Maade. I en Samling af Afhandlinger, som
er udkommen under Titelen \textit{The} will \textit{to believe} \textit{and other essays}
in popular philosophy (1897), har han i noget friere Form udtalt
sin ejendommelige Forening af Erfaringsfilosofi og udpræget
idealistisk Livsanskuelse, og allerede her kommer han ind
paa \textit{religionsfilosofiske} Spørgsmaal. Det nye Værk er ganske
helliget det religiøse Problem, særligt den rent psykologiske Be
lysning af det, dog saaledes, at ogsaa de andre Synspunkter, fra hvilke Problemet maa behandles, komme frem. James benytter
de yngre amerikanske Forfatteres Specialarbejder, men støtter
sig desuden til en omfattende Literatur af biografisk, især selvbiografisk
 Art. Han ser klart, at Emnet er for mangfoldigt til
at kunne udtømmes i en enkelt Fremstilling, og derfor --- tillige
for at slippe for en Ordstrid om, hvad Religion er --- erklærer
han, at han til Brug ved denne sin Undersøgelse vil
holde sig til en bestemt Betydning af Ordet Religion eller til
en bestemt Side ved Religionen, og han definerer saa Religion
som: individuelle Menneskers Følelse, Handling og Erfaring,
forsaavidt disse Mennesker i deres Ensomhed opfatte sig selv
som staaende i Forhold til det Guddommelige, hvorledes de
end opfatte dette. Denne personlige Religion stiller han i Modsætning
til den institutionelle Religion med dens mere eller
mindre udviklede Teologi eller Mytologi og dens kirkelige Organisation.
Han hævder, at \textit{institutional religion} (med Alt, hvad
dertil hører) stedse er afledet i Forhold til \textit{personal} religion; det
følger allerede af, at den bygger paa Overlevering. Hvis man
vilde holde sig til den institutionelle Religion alene, vilde Undersøgelsen
blot komme til at angaa Gudsdyrkelse, Ofring, Ceremoniel
og kirkelig Organisation, og man kr Re
% --- p. 127 ---
\setcounter{footnote}{0}%
\opage{127}ligion som den Kunst at vinde Gudernes Gunst. I den personlige
Religion er det derimod Menneskets egne indre Tilstande,
dets Samvittighed, dets Værdi, dets Hjælpeløshed eller dets
Ufuldkommenhed, Vægten ligger paa. Holder man sig til denne
Side af Religionen, faar Begrebet Religion et større Omfang,
end naar man anser det Institutionelle for Hovedsagen.

Den personlige Religion er den Maade, hvorpaa Mennesket
reagerer mod Livet som Helhed (a man's total reaction upon
life), i Modsætning til de øjeblikkelige og specielle Reaktioner.
De Erfaringer, vi have gjort om Livet og om Verden, virke paa
vort individuelle Temperament og gøre os enten energiske eller
ligegyldige, enten fromme eller spottende, enten nedstemte eller
begejstrede overfor Livet som Helhed. Vor »Reaktion«, der ofte
foregaar uvilkaarligt og halvt ubevidst, udtrykker sig i det fuldstændigste
Svar, vi kunne give paa det Spørgsmaal: Hvad er
Karakteren af denne Tilværelse, vi høre til? --- Dog vil James
kun kalde en saadan Reaktion religiøs, naar den bæres af en
alvorlig og ophøjet Stemning, en Stemning, der paa sit Højdepunkt
fører til Opgivelse af egne Ønsker, til villig Underkastelse
under Virkelighedens Smerte, fordi man føler sig baaren af en
Kraft, der er større end den, man selv raader over. Der er i
al udpræget religiøs Følelse et tragisk Element, som beror paa,
at en højere Lykkefølelse har fortrængt en lavere. Den religiøse
Følelse indeholder eller forudsætter saa mange Modsætninger, at
den er en af de rigeste menneskelige Følelser. Verden er rigere
ved, at der er en Djævel, naar der er en St. Michael, som kan
sætte sin Fod paa hans Nakke. Dette Særkende træder naturligvis
mest frem i de ekstreme Former for Religion, Former, der
for den jevne, dagligdags Opfattelse staa som abnorme; men i
alle Former for religiøs Bevidsthed vil det til en vis Grad forekomme.


Denne Beskrivelse angiver den almindelige Ramme, indenfor
hvilken James bevæger sig i sin Undersøgelse. Til nærmere
Karakteristik af hans Arbejde skal jeg først vise, hvorledes de
tre Synspunkter, der --- som omtalt i Indledningen --- maa
komme frem ved enhver omfattende Drøftelse af det religiøse
Problem, ogsaa fremtræde hos James, skønt han hovedsageligt
% --- p. 128 ---
\setcounter{footnote}{0}%
\opage{128}har holdt sig til den rent psykologiske Betragtningsmaade. Dernæst
skal jeg vise, hvorledes han opfatter den religiøse Erfarings
Væsen, hvilke Typer for Religion han finder, og til
hvilken Grundtanke Religionen efter hans Opfattelse kan føres
tilbage.

\subsubsection*{\textit{a. De tre Synspunkter.}}

James betegner sit Standpunkt, baade i almindelig filosofisk
Henseende og i religionsfilosofisk Henseende, som Erfaringens
Standpunkt. Han vil være gennemført Empiriker. Alle intellektuelle
Operationer --- hvad enten de bestaa i Konstruktion,
Sammenligning eller i Kritik --- forudsætte umiddelbar Erfaring. Han slutter sig til den gamle engelske Skole, den, der begyndte
med John Locke, og som stedse krævede bestemt Regnskab for
de Erfaringer, paa hvilke vi bygge vore Begreber. Denne Skole
bør efter hans Mening langt snarere end Kant have Æren for
den kritiske Metode i Filosofien, den eneste Metode, der gør Filosofi til et Studium, som kan drives af alvorlige Mennesker.
Men det drejer sig ikke blot om de Erfaringer, vi have havt,
ogsaa om dem, vi maa vente eller konsekvent selv fremkalde,
hvis vore Tanker skulle være gyldige. Vore Tankers Værdi beror
paa de Slutninger, vi kunne drage af dem. Hvis alle Sætninger,
vi opstille, vare praktisk ligegyldige, d. v. s. hvis der
slet Intet fulgte af dem, kunde vi saa skelne mellem sande og
falske Tanker? Forgæves har Dogmatismen søgt efter Kriterier
for Sandhed, der kunde fritage for at appellere til Fremtiden. I
Modsætning til Dogmatismen staar her Pragmatismen, som prøver
hver Sætning ved at forfølge dens Konsekvenser. Denne
Opfattelse (i hvilken James slutter sig til den amerikanske Tænker
CHARLES PEIRCE, fra hvem ogsaa Navnet Pragmatisme
stammer) minder om den biologisk-økonomiske Erkendelsesteori,
der i Tyskland er bleven udviklet af Richard Avenarius og
Ernst Mach.

James' empiriske Standpunkt fører ham til Skepsis overfor
spekulative og teologiske Synsmaader. Han hævder særligt Mangfoldigheden
i Tilværelsen og er kritisk stemt overfor Monismen
% --- p. 129 ---
\setcounter{footnote}{0}%
\opage{129}i dennes forskellige Former. Pluralismen, Mangfoldighedslæren
vil, tror han, mere og mere komme til Ære --- baade i Videnskaben
og i Religionen. Ud fra Erfaringen kan man intet Universelt,
ingen absolut og altomfattende Enhed naa. Og forskellige
Tydninger kunde gælde for forskellige Egne af Tilværelsen.
Religiøst set blive her især Erfaringerne om det Onde af Betydning:
har Jobs Bog ikke en Gang for alle godtgjort Umuligheden
af at løse det Problem, disse Erfaringer føre til at stille?
Her vil den monistiske Hypotese stedse frembyde store Vanskeligheder,
og vilde det saa ikke være den simpleste Udvej at antage,
at Verden fra først af har indbefattet Elementer, der ikke
vare i Harmoni med Helheden? Sædvanlige Mennesker have
stedse været mere eller mindre polyteistiske, og det er da ogsaa
klart, at man ikke af nogen Erfaring kan udlede Guds Uendelighed.
Hvad vi kunne erfare, kan i det Højeste være det, at
vi staa i indre Sammenhæng med Noget, der er større end vi
selv; men at dette er et Uendeligt, kan ikke godtgøres.

Hvad Tænkningen kan udrette paa det religiøse Omraade,
kan kun være det: at ordne, bestemme og fortolke Kendsgerningerne,
ikke at frembringe dem. Filosofien er her stedse sekundær.
Dog kan den derfor godt være til Nytte. Ti den religiøse
Bevidsthed udtrykker stedse sine Erfaringer i Former,
der stamme fra den intellektuelle Overlevering, under hvis Indflydelse
den staar. Derfor kan kun kritisk Sammenligning finde
det virkeligt Umiddelbare og Væsentlige og udsondre det fra
det Lokale og Tilfældige. Tillige er det Filosofiens Opgave at
udskille Antagelser, der ere inkonsekvente eller afgjort stride
mod videnskabelig Erfaring. Hvad der bestaar gennem denne
Skærsild, kan saa betragtes som Hypoteser, hvis Berettigelse og
Betydning kunne blive Genstand for nærmere Drøftelse, efterat
de ere blevne reducerede til deres simpleste Former.

Saavidt den erkendelsesteoretiske Betragtning. Saa energisk
James gør den gældende, saa er det dog den psykologiske Betragtningsmaade,
han har størst Interesse for, og som giver hans
Arbejde dets største Betydning.

Religionspsykologien yder den almindelige Psykologi vigtige
Bidrag. Der er mange psykologiske Fænomener, der sær%
% --- p. 130 ---
\setcounter{footnote}{0}%
\opage{130}ligt godt kunne studeres paa det religiøse Livs Omraade. James
fremhæver saaledes, at de religiøse Fænomener godtgøre, hvor
ringe en Del af vort aandelige Liv der kan forklares klart og
tydeligt. Bevidstheden fortoner sig her gennem mange Grader
over i det Ubevidste, eller, som James hellere vil sige, i det
Underbevidste. Der rører sig umiddelbare Tilskyndelser og Anelser,
som ofte, uden at vi selv mærke det, give os de første
Præmisser, paa hvilke klart bevidste Tanker grunde sig. Bevidste
Argumenter bevæge sig ofte kun paa Overfladen af vort Væsen,
og en uvilkaarlig og umiddelbar Forvisning er det Dybeste i os
(the deep thing in us). Fremdeles lærer Studiet af det religiøse
Liv os store Forskelligheder mellem Mennesker at kende, baade 
i Henseende til Arten af Følelseslivet og i Henseende til Graden
af Modtagelighed, --- Forskelligheder, vi ellers neppe vilde
blive saa opmærksomme paa. Vi lære, hvorledes Udviklingen
henimod et og samme Maal kan foregaa højst forskelligt, snart
kontinuerligt, snart springende. Allerede rent teoretisk eller naturhistorisk
set har et psykologisk Studium af det religiøse Liv
sin Betydning.

Den psykologiske Betragtningsmaade bliver saa igen af Betydning
for Forstaaelsen af det religiøse Problem. De religiøse
Processer ses som specielle Former for Processer, som ogsaa
kendes paa andre psykiske Omraader, og adlyde Love, der ogsaa
ellers kendes. James viser her hen til Starbucks Undersøgelser
af Omvendelsen i dens Sammenhæng med den organiske og psykiske
Udvikling i Overgangsaarene. Ganske særlig Vægt lægger
James paa det, der bevæger sig under Bevidsthedens »Tærskel«
eller udenfor dens »Margin«. Han finder det meget vanskeligt
her at drage en bestemt Grænse og forkaster den atomistiske
Psykologi, for hvem Bevidstheden er sammensat af enkelte, hvert
for sig faste og bestemte Elementer. I hvert Tilfælde gøre de
subliminale eller ultramarginale Indflydelser sig stadigt gældende,
og James er for sit Vedkommende tilbøjelig til at betragte
disse Indflydelser som de Veje, ad hvilke en højere Tingenes
Orden virker i vort Indre. Disse indre Tilkendegivelser
udgøre vor nærmeste, vor egentlige Verden; saa inderligt tilhøre
vi ikke den ydre Verden. Her finder James et Enheds%
% --- p. 131 ---
\setcounter{footnote}{0}%
% CHECK p. 131: indent after a broken word?
\opage{131}punkt mellem Religion og Psykologi, skønt han erklærer enhver
nærmere Bestemmelse deraf for en Fortolkning, der ikke længere
er rent psykologisk. Fortolkningen er heller ikke det Afgørende;
den er stedse afledet, og en og samme Oplevelse kan
være Genstand for forskellige religiøse Fortolkninger. Hvad
de fleste Mennesker paa dette Punkt især mangle, er Kritik
og Forsigtighed, ikke Tro (faith). De ere for tilbøjelige til at
lade dogmatisk Antagelse (belief) følge hovedkulds efter enhver
livfuld Forestilling, især naar denne støttes af uvilkaarlig Interesse. ---

Det tredie Synspunkt, det etiske, kommer frem ved Værdsættelsen
af de religiøse Fænomener. Ligesaa bestemt, som
James skelner mellem Oplevelse og Fortolkning, ligesaa bestemt
skelner han mellem Beskrivelse eller Forklaring af det Oplevede og Vurdering af det. Selv det, der fra et rent medicinsk
Synspunkt maa kaldes abnormt, kan ved sit Indhold have stor
Værdi. En Tilstands Værdi beror ikke paa, hvorledes den er
bleven til. Naar vi kalde visse Tilstande højere end andre, er
det ikke paa Grund af, hvad vi mene at vide om deres organiske
Betingelser. Hvis der er noget Værdifuldt i S\textsuperscript{ta} Teresas
Tanker, saa er det ligegyldigt, om hun var hysterisk eller ikke.
Vi vurdere en Kraft efter dens Virkning, ikke efter dens Udspring.


Ved Vurderingen kan intet spekulativt eller teologisk System
lægges til Grund. Ogsaa her maa vi være Empirikere. Vi
bedømme de religiøse Fænomener efter deres Frugter, og det
have Menneskene egentligt bestandigt gjort. En Gudsdyrkelse
hører op, naar den ingen Virkning øver paa Sindet, og naar
den i sin hele Karakter kommer i Strid med Noget, hvis Værdi
vi have erfaret saa grundigt, at vi ikke kunne fornegte det. De
Guder, Menneskene holde fast ved, ere de, som de kunne bruge,
og hvis Bud bekræfte de Fordringer, de stille til sig selv og
til Andre. Vi bruge stedse menneskelige Maalestokke (human
standards). Vi spørge, hvorvidt det religiøse Liv godtgør sig
som en ideal Form for menneskeligt Liv. Det er naturligvis
en Maalestok, der varierer med Tiderne, men anden Maalestok
have vi nu engang ikke. Kun maa Prøvelsen  atter og
% --- p. 132 ---
\setcounter{footnote}{0}%
\opage{132}atter gentages, da den maaske under nye Forhold kunde føre
til nye Resultater. Og Vurderingen vil ofte falde forskelligt ud for forskellige Individer. Naar psykologisk Undersøgelse godtgør,
at der er store Forskelligheder mellem Mennesker i deres
inderste Natur, maa man ogsaa antage, at de trænge til forskellig aandelig Næring, saa at de komme til at staa forskelligt i religiøs Henseende.

Overfor Religionen idethele staar James afgjort sympatisk. De bedste Frugter af religiøs Erfaring ere efter hans Overbe
visning de bedste Ting, Historien har at fremvise. Han føler her det indre Liv røre sig med Alvor og Energi, virke med
en Inderlighed og en Koncentration, i Modsætning til Alt, hvad
der kan hæmme og splitte. Det er, som man var hævet op i en højere og renere Atmosfære. --- James henviser paa dette Punkt
til Sainte Beuve's Vurdering af de religiøse Porsonligheder,
han har studeret i sit berømte Værk om Port-Royal. ---

Vi finde saaledes hos James de tre Hovedsynspunkter ved
filosofisk Drøftelse af religiøse Spørgsmaal. Han bebuder, at
han i et senere Værk vil imødegaa de Forsøg, der endnu i vore
Dage ere gjorte paa en mere spekulativ Behandling af Religionsfilosofien,
--- Forsøg, der mærkværdigt nok forekomme hyppigere
i England og Amerika end paa Kontinentet. Foreløbigt staar hans
Arbejde som et betydningsfuldt Indlæg til Gunst for, at det
netop er hine tre Synspunkter, der ere at anlægge i Religionsfilosofien.
--- Desværre kan jeg i min Karakteristik af James
Standpunkt ikke lade den Aabenhed, Friskhed og sproglige
Kunst, han raader over, komme til deres fulde Ret.

\subsubsection*{\textit{b. Den religiøse Erfarings Natur.}}

Religiøs Følelse er ifølge James ikke nogen speciel eller
elementær Følelse. Enhver Følelse kan under visse Betingelser
faa en religiøs Karakter, nemlig naar den medvirker ved det
Opgør af de menneskelige Livserfaringer, James kalder »den totale
Reaktion overfor Livet«. En noget nærmere Bestemmelse
af, hvad Religion er, giver den Udtalelse, at Religion er en af
de Veje, ad hvilke der kan komme Enhed i vort Liv. Med En%
% --- p. 133 ---
\setcounter{footnote}{0}%
\opage{133}heden følger en Følelse af Fred og Lykke, som overgaar Alt,
hvad der ellers kan erfares af den Art. \textit{Denne} Enhed kan naas
ved pludseligt Gennembrud eller gennem en lang Udvikling;
her lægge de individuelle Forskelligheder sig for Dagen. Men
Hovedsagen er, at Individet, der før har følt sig delt og hæmmet,
lavt og ulykkeligt, nu har Følelsen af Harmoni, Frihed og
Opløftelse. Tanker, der før stode som ideale Muligheder ude i
Sjælens Periferi, ere nu blevne Sjælens centrale Ejendom. Det
er, som om der er tilstrømmet Individet ny Energi fra Kilder,
der ligge udenfor dets Bevidsthed, og som det ikke vilde kunne
udløse direkte ved bevidst Anstrengelse. Denne underordnede
Betydning af det Bevidste og Vilkaarlige (the hopeless inferiority
of voluntary to instinctive action) i Forhold til det Ubevidste
og Uvilkaarlige er et karakteristisk Træk i al Religiøsitet.
Sygelige Former opstaa, naar der mangler intellektuel Kraft og
Interesse til at forme og tolke saadanne Erfaringer.

Religiøs Erfaring har en mystisk Karakter, dels paa Grund
af Uudsigeligheden, idet Enhver her kun indirekte kan forstaa
de Andre, dels paa Grund af Umiddelbarheden, idet den diskursive
Tænken ingen Rolle spiller, dels paa Grund af Tilstandenes
korte Varighed og Vanskeligheden ved at genkalde
sig dem, dels endeligt paa Grund af det passive Forhold, Individet
staar i til disse Erfaringer. Det ikke-religiøse Bevidsthedsliv
kender lignende Tilstande. Den religiøse Karakter fremkommer,
som allerede bemærket, derved, at Menneskets Forhold til
Livet som Helhed gøres op i saadanne Tilstande; derved faa
disse en Enhed og Energi, der ellers mangler. Den særligt saakaldte
Mystik er kun en bestemt og ofte extrem Form for, hvad
der fremtræder i al religiøs Erfaring.

\subsubsection*{\textit{c. De vigtigste Typer.}}

James erklærer, at det vigtigste Resultat, han er kommen
til ved sine Undersøgelser af den religiøse Bevidsthed, er de
store Forskelligheder, den frembyder. En betydningsfuld Del
af hans Arbejde gaar derfor naturligt ud paa at føre disse Forskelligheder
tilbage til visse \textit{Typer}.

% --- p. 134 ---
\setcounter{footnote}{0}%
\opage{134}Der er særligt én Modsætning, han meget stærkt fremhæver,
nemlig den mellem det sunde og det syge Sind, eller ---
for at bruge et Udtryk, han har fra Francis Newman --- mellem
de én Gang Fødte og de to Gange Fødte (the once-born and the
twice-born). De syge Sjæle behøve en Forandring, en Krise, en
Omvendelse for at naa til Enhed og Fred, medens der er sunde
Sjæle, som ved umiddelbar Udfoldelse og Hengivelse naa til
Maalet. De Sundes Religion betegner James som Naturalisme,
de Syges som Salvationisme. Men indenfor hver af disse Grupper
kunne igen højst forskellige Personligheder og Retninger
fremtræde.

De sunde Sjæles Religiøsitet er en Entusiasme, der fastholder
den lyse Side ved Tilværelsen og maaske endog erklærer
de mørke Sider for Illusioner. Selvom de ikke mangle Blik
for det Hæmmende og Formørkende, ere de dog forvissede om,
at det kan overvindes; saa stor er deres umiddelbare Tillid til
den indre Energi og til den Harmoni, der raader i Universet.
Denne Retning fremtræder mere i Græciteten end i Kristendommen,
mere i Katolicismen end i Protestantismen, mere hos
de latinske Racer end hos de germanske. I nyere Tid fremtræder
den hos Spinoza og Rousseau, i den liberale Retning
indenfor Kristendommen, især i Unitarismen (Emerson, Parker).
I Amerika træder den i skarp Modsætning til Metodismen og
Revivalbevægelsen. I nyeste Tid er den saakaldte mind-cure
movement et karakteristisk Exempel. Den helbreder baade Sjæl
og Legeme ved stadigt at indskyde og fastholde den Tanke, at
Smerten er en Illusion, og at Alt i Virkeligheden er lyst og
herligt. Frygt og Bekymring banlyses. Den virker gennem
Suggestioner, ikke mindst Selvsuggestioner. Ved stadigt at raabe
Hurra for Universet bestyrker man sig i Troen paa, at Universet
fortjener et Hurra.

Den syge Sjæl ligger under for Smerten og Disharmonien.
Bevidsthedens Tærskel overfor den Art Erfaringer er en Gang
for alle sænket, medens den hos den sunde Sjæl er hævet. Verden
staar som tom og forfængelig; der næres Frygt for, hvad
Livets Gang kan bringe; Sjælen føler sig syndig og fordærvet.
Tomhedsfølelse, Frygtfølelse og Syndfølelse ere de tre Hoved%
% --- p. 135 ---
\setcounter{footnote}{0}%
\opage{135}former for sjælelig Sygdom. Hos Tolstoi raader Tomhedsfølelsen,
hos Bunyan Frygtfølelsen og Syndfølelsen. Der er her en oprindelig
Disharmoni i Temperamentet, som kræver en anden
Fødsel. Og selvom denne naas, ville Naturer af denne Art
neppe naa til den overstrømmende Livsfølelse, som de Sunde
kunne have.

Der findes naturligvis i Virkeligheden en Mængde Overgangsformer
mellem de to Typer. Men i deres extreme Former
er Modsætningen mellem dem saa stor, at de neppe forstaa hinanden.
De ville hver for sig betragte deres Livsanskuelse som
den ene rette. Saaledes hævdede Wesley Omvendelse som eneste
Vej til Fred, medens Emerson tillagde den Natur, der kun behøvede
én Fødsel, højere Værdi end den, der maatte kæmpe
haardt med Djævlene i sit Indre. James selv mener, at de Livserfaringer,
de to Gange Fødte kunne bygge paa, ere dybere og
fuldstændigere end den første Types. Ingen Profet kan bringe
et afgørende Budskab til Menneskene, naar han ikke siger Ting,
der kunne klinge som Realiteter for Mænd som Bunyan og
Tolstoi. De fuldstændigste Religioner maa være de, der baade
kende og overvinde Livets mørkeste Elementer. Derpaa beror
Buddhismens og Kristendommens Overlegenhed. ---

I Forhold til denne store Modsætning ere de andre Modsætninger,
der kunne paavises, for James af mindre Betydning.
Saaledes Modsætningen mellem springende og kontinuerlig Udvikling,
--- Modsætningen mellem Naturer, hos hvem de Elementer,
der skulle overvindes, først maa have mistet deres Livskraft,
og andre Naturer, hos hvem det nye Liv kan trænge igennem,
før det gamle er helt udtømt, --- og Modsætningen mellem
en Udvikling, der kræver bevidst Anstrengelse (selvom denne
ikke direkte fører til Maalet), og en Udvikling, der gennemgaaende
forløber uvilkaarligt. ---

Lejlighedsvis giver James foruden denne almindelige Typelære
ogsaa Antydninger af sin egen rent individuelle Tro. Han
har i en længe raadende nervøst slappet Tilstand søgt Hjælp
hos mind-cure Bevægelsen og mener at være bleven helbredet
ved dens Hjælp, hvilket ikke hindrer ham i at stille sig kritisk
overfor den. Naar han udtaler, at Mystiken kun giver i extrem
% --- p. 136 ---
\setcounter{footnote}{0}%
\opage{136}Form, hvad al religiøs Erfaring indeholder, føjer han til, at han
selv ikke personligt kender mystiske Tilstande. Sin Tro karakteriserer
han ved, at skønt han ikke kan slutte sig til den gængse
Kristendom eller til den spekulative Teisme, er han overbevist
om, at der i de Tilstande, i hvilke Mennesket føler sig i den
inderligste Enhed med det Højeste, det kender, kommer ny Kraft
ind i Verden og sættes nye Udgangspunkter; forsaavidt er hans
Tro supranaturalistisk, antager guddommelige Indvirkninger indenfor
den naturlige Erfarings Omraade. Han føjer til, at han meget
vel véd, at den herskende Mening i akademiske Kredse gaar i
modsat Retning. --- Udødelighedstroen spiller kun en underordnet
Rolle for ham (skønt han i et særligt lille Skrift har søgt
at vise, at den ikke kan gendrives). Han siger herom: Naar
der blot »i al Evighed« sørges for vore Idealer, ser jeg ikke,
hvorfor vi ikke skulde være villige til at overgive Omsorgen
for dem i andre Hænder end vore.

\subsubsection*{\textit{d. Den religiøse Grundtanke.}}

Skønt Religionen efter James har sin Grund i Følelseslivet
og dets Erfaringer, og skønt den fremtræder i højst forskellige
Former, er der dog en Grundtanke eller en Hypotese, som den
stedse hviler paa. Det er den, at der kan naas en Tilstand af
Harmoni, Enhed og Fred, selvom det maa ske gennem Kampe
og Kriser, og at dette Maal naas ved Tilstrømning af Energi
fra en mere omfattende Tingenes Orden, der psykologisk set
ytrer sig i den ubevidste Natur i os. Hvorledes vi i det Enkelte
fortolke, hvad der her sker, er ikke det Væsentlige; her
ville stedse rent individuelle Suppleringer og Udførelser (overbelief)
komme til. Hovedsagen er, at for den religiøse Bevidsthed
er Tragedien, saa ofte den end melder sig paa Livets Vej,
stedse kun foreløbig og partiel. Skibbrud og Opløsning ere ikke
Tilværelsens sidste Ord. Al Religion hviler saaledes tilsidst paa
en kosmisk Hypotese. Derimod er Religionens Bestaaen ikke
afhængig af nogensomhelst speciel dogmatisk Antagelse. Naturligvis
--- hvis Religion gøres til Et med eller ganske afhængig
af en animistisk Verdensopfattelse, saa vil dens Tid være forbi,
% --- p. 137 ---
\setcounter{footnote}{0}%
\opage{137}naar Animismen uddør; ligesaa, naar man gør den til Magi eller
Fetischisme. Det bliver for James let en Strid om Ord, her
føres. Det Afgørende for ham er, at de indre Erfaringer, hvorpaa
Religionen væsentligt bygges, ere de mest konkrete og
umiddelbare Erfaringer, vi have. I Forhold til dem har Alt,
hvad vi kunne erfare om de ydre Verdensforhold, en abstrakt
og generel Karakter. Der er ingen Grund til at antage, at saadanne
Erfaringers Tid skulde være forbi, selvom de maa udskilles
fra Forbindelse med Elementer, de hidtil have været tilbøjelige
til at smelte sammen med.

\subsubsection*{\textit{e. Bemærkninger om James' Religionsfilosofi.}}

Studiet af James' Bog, der udkom et Aar efter mit eget
religionsfilosofiske Forsøg, har været af stor Interesse for mig.
Ikke blot er det længe siden, jeg har læst en Bog, der saaledes
har lært mig at se med nye og friske Øjne paa Livet og paa
Menneskene, som denne Bog af den amerikanske Filosof har
formaaet. Men Bogen fik en særegen Betydning for mig derved,
at hele Retningen og Behandlingsmaaden, trods al Forskellighed
paa ikke uvæsentlige Punkter (f. Ex. i det, James kalder
sin »Supranaturalisme«), i saa høj Grad er beslægtet. Jeg har
derfor en personlig Grund til at henlede Opmærksomheden paa
den. Ved sin aandfulde Form og sine mange interessante Enkeltheder
vil den iøvrigt lønne et Studium, selvom man staar
Bogens Retning mere fjernt end jeg. Kun nogle faa \textit{Bemærkninger}
vil jeg her fremføre for at betegne mit Standpunkt overfor
den.

Det er i mine Øjne med fuld Ret, at James skelner mellem
personlig og institutionel Religion. Det er ganske vist en
Forskel, som er lettere at gøre ved de højere end ved de lavere
Religioner, og det er maaske derfor, Religionshistorikere undertiden
krympe sig ved at indrømme den. Men den psykologiske
Interesse maa nu engang efter Sagens Natur nærmest angaa det
personlige Element i Religionen, medens den objektive eller
»institutionelle« Side kun har Betydning som Symptom. Og
hvad der er Hovedsagen: det religiøse Problem angaar egentligt
% --- p. 138 ---
\setcounter{footnote}{0}%
\opage{138}kun den subjektive Side. Spørgsmaalet om Religionens Stilling ---
nu og i Fremtiden --- i det menneskelige Aandsliv gaar jo dog
nærmest ud paa, om Menneskets Følelse og Trang vil kunne
vedblive at fastholde de hidtil givne Former og Symboler, og
om der, hvis disse skulde vise sig ufyldestgørende, vil kunne
findes eller skabes nye. Den, der fæster sit Blik paa Kontinuiteten
i den aandelige Udvikling, maa afgjort lægge Hovedvægten
paa det subjektive, personlige Omraade, saa meget dette end
tilsløres ved de ydre Formers historiske Overmagt.

Dog forekommer det mig, at James gaar for let henover
de Vanskeligheder, Kontinuitetsproblemet her frembyder. Det er
ingenlunde selvfølgeligt, at det personlige Element i Religionen
kan bevares, dersom de historiske Former falde bort. Som det
vil ses af min Religionsfilosofi, staar jeg her saaledes, at jeg ser
en \textit{Mulighed} for en Bevarelse af den Inderlighed og Koncentration
i Livsopfattelse, der har havt en saa væsentlig Støtte i
de historiske Former for Religion, skønt jeg ikke kan føre noget
tvingende Bevis for Nødvendigheden deraf, --- end ikke for, at
der kan gives Ækvivalenter derfor. Vi staa overfor den Mulighed,
at der her vil gaa Værdi tabt. Historien viser os et stadigt
Vexelforhold mellem det Personlige og det Institutionelle. Skal
dette Vexelforhold bestaa i Fremtiden, saa er det kun muligt
dersom den menneskelige Aand under Udviklingens Gang bevarer
Evnen til at finde eller skabe store Symboler for de
dybeste Livserfaringer, Symboler, der kunne faa fælles Betydning
for større eller mindre Grupper af Mennesker. Det kan
haabes, men ikke bevises. Og dog er det et af de vigtigste
Spørgsmaal med Hensyn til vor Slægts Fremtid. Man maa se
paa Livets største Konflikter med meget let Sind og meget uklart
Blik, naar man ikke føler Brodden i dette Problem. --- Hvorvidt
man under de nye Forhold vil kunne vedblive at bruge
\textit{Ordet} Religion, er et ligegyldigt Spørgsmaal, som James behandler
med den Ironi, det fortjener. ---

Sammenlignes James' Beskrivelse af religiøs Erfaring med
den, jeg har givet, fremtræder der en Overensstemmelse, der
viser, at vi have de samme Fænomener for Øje. Religiøs Erfaring
er for ham karakteriseret ved at være Menneskets »totale
% --- p. 139 ---
\setcounter{footnote}{0}%
\opage{139}Reaktion« overfor Livet; den er en Erfaring af Fred, Enhed
og Energi, ofte vunden gennem en Krise, der førte til Fortrængen
af Mørke og Disharmoni. Hvad jeg her savner, er en
bestemt Fremhæven af, at den Disharmoni, den religiøse Bevidsthed
under alle Former kæmper for at opløse, er en Disharmoni mellem Værdi og Virkelighed. Ved Forholdet mellem
de Værdier, Mennesket kender, og den Virkelighed, han kender,
er hans Religion bestemt. Det Værdifulde kan for ham være
selve det fysiske Liv i dets Bestaaen og Udfoldelse, eller det
kan være ideelle Goder. Stedse bestemmes hans Totalfølelse
overfor Livet ved, om Værdiernes Himmel er formørket eller
ikke. Det er ud fra dette Synspunkt, at jeg er kommen til min
Hypotese om, at Værdiens Bestaaen er den religiøse Grundtanke
eller det religiøse Axiom. James lægger Vægten paa den
indre Enhedsfølelse og paa Erfaringen om Tilstrømning af indre
Energi, uden at han tager tilstrækkeligt Hensyn til, at denne
Enhed og Energiforhøjelse for Mennesket betyder, at Disharmonien
mellem Værdi og Virkelighed er hævet. Dog tror jeg ikke,
at James vil være uenig med mig i min Hypotese. Som allerede
ovenfor omtalt, antyder han jo, at Religion vel først og
fremmest bygges paa subjektiv Erfaring, men at den mere eller
mindre bevidst hviler paa den Forudsætning, at det Tragiske i
Verden kun kan være partielt og forbigaaende, --- at Skibbrud
og Opløsning ikke kunne være Tilværelsens sidste Ord. Den
religiøse Bevidsthed kommer altsaa ifølge sin Natur til at forudsætte
Noget om kosmiske Forhold, om Tilværelsens inderste
Natur. Det er paa dette Punkt, Religionen stedse har havt Tilbøjelighed
til at blive Metafysik, ligesom omvendt Metafysiken,
uden altid at være sig det bevidst, er gaaet ud fra religiøse
Forudsætninger i sine Konstruktioner. Ifølge James er ethvert
Forsøg paa bestemtere og konkretere Formulering af hine sidste
Forudsætninger »over-belief«, d. v. s. fører fra Tænkningens og
Erfaringens Omraade over i Fantasiens og Mysticismens. James'
egen »Supranaturalisme« er aabenbart et Exempel herpaa. Om
de nye Begyndelser, han antager, har han ikke Ret til at sige
Mere, end at deres Aarsag foreløbigt ikke kan angives; at de
% --- p. 140 ---
\setcounter{footnote}{0}%
\opage{140}skulde have en overnaturlig Aarsag, er ikke motiveret efter hans
egen Opfattelse. ---

Der er endnu et væsentligt Punkt, i hvilket jeg stemmer
overens med James, nemlig i den store Vægt, han lægger paa
de individuelle Forskelligheder i Henseende til religiøs Erfaring
og Tro. Hans mesterlige Skildring af Modsætningen mellem de
sunde og de syge Sjæles Maade at tage Tingene paa, er et stort
Indlæg til Fordel for, hvad jeg kalder Personlighedsprincipet.
Denne Modsætning svarer til den allerede ovenfor omtalte Modsætning
mellem expansive og disharmoniske Naturer. Men James
har ført Karakteristiken igennem med en saadan Kraft og Klarhed,
at Spørgsmaalet om, hvorvidt forskellige Personligheder
virkeligt kunne have »samme« Religion, melder sig med endnu
større Nødvendighed end før. ---

% --- p. 141 ---
\setcounter{footnote}{0}%
% CHECK p. 141: centred line: heading or verse
\opage{141}\essay[\textit{Parmenides} fra Elea. (1901)]{PARMENIDES FRA ELEA.}
\streg

A. Patin, \textit{Parmenides im Kampfe gegen} Heraklit. (Bes.
Abdruck aus dem 25. Supplementband der Jahrbücher für klassische
Philologie.) Leipzig 1899, Teubner. P. 491—660.

Parmenides' Stilling indenfor den græske Filosofis Historie
har i de senere Aar tildraget sig særlig Opmærksomhed. Hvad
der foreligger fra hans Haand, er som bekendt Fragmenter af et
Læredigt, bestaaende af ca. 150 Vers. I en sammentrængt Fremstilling,
i hvilken den poetiske Form spiller en meget underordnet
Rolle, giver han sin strenge, radikale Lære om Enheden
som det eneste Reale, med Polemik mod al Antagelse af reale
Modsætninger, af Vorden, Bevægelse og Forandring. Der foreligger
her for første Gang i den evropæiske Tænknings Historie
en Verdensanskuelse, som i forskellige Former senere gentagne
Gange er traadt frem. Platon og Spinoza ere aandelige Frænder
af Parmenides. I Indien fremtræder en beslægtet Opfattelse
flere Aarhundreder før Parmenides i Vedantalæren, der er fremstillet
i Upanishad, og omtrent samtidigt med Parmenides gør
Sig --- mere som Livsanskuelse end som Verdensanskuelse ---
en analog Anskuelse gældende i Buddha's Lære. Dog er det
hos Parmenides ikke blot en metafysisk Verdensanskuelse, der
er formet som Læredigt. Hans Tænkning fremdrager paa en tydelig
Maade den skarpe Modsætning, som vi nu betegne som Modsætningen
mellem formal (logisk og matematisk) Videnskab og
real (empirisk) Videnskab, og denne Modsætning fremtræder her
% --- p. 142 ---
\setcounter{footnote}{0}%
\opage{142}for første Gang i principiel Form. Bag de dogmatiske Formuleringer
hos Parmenides spores som Grundlag den Hovedsætning,
at kun de Sætninger, der kunne formes som rene Identiteter,
kunne være Genstand for fast Overbevisning. Hvad der,
i Modsætning til streng logisk og matematisk Slutning, bygges
paa Iagttagelse og Erfaring, kan kun være »Dødeliges Meninger«,
--- kan kun være Noget, der synes at være som det er,
--- viser os kun Tingene fra »Meningens« Synspunkt. Alligevel
giver Parmenides ogsaa en Fremstilling af, hvorledes Tingene
maa tage sig ud fra dette Synspunkt. Her optræde kvalitative
og kvantitative Forskelligheder og stridende Modsætninger, og
der viser sig at foregaa en stor Verdensproces med skiftende
Faser, en hel Kosmogoni. Parmenides vil skildre, »hvorledes
Jord og Sol og Maane og den fælles Æter og den himmelske
Mælkevej og den yderste Olymp og Stjernernes varme Kraft
stræbte frem til Vorden«.

To Hovedspørgsmaal stille sig naturligt med Hensyn til
denne hele Lære: 1) Hvorledes er den bleven til, i hvilket Forhold
staar den til den forudgaaende filosofiske Tænkning? 2) Hvorledes
er Forholdet mellem Læredigtets to Dele, der synes at
staa i uforsonlig Modstrid med hinanden?

Det foreliggende indholdsrige og vel gennemtænkte Arbejde
giver et betydningsfuldt Indlæg i Diskussionen om disse to Spørgsmaal,
der begge have mere end rent historisk og filologisk Interesse.
Det forekommer mig, at Forfatteren har levet sig mere
ind i Parmenides' Grundtanke end tidligere Forskere. Derved
har han formaaet at finde en Løsning, hvor hans Forgængere
enten maatte erklære at staa overfor uforsonede Modsætninger
eller maatte tage deres Tilflugt til tvungne Forklaringer. Undertiden
fører maaske hans Stræben efter at finde karakteristiske
Hentydninger og Sammenhænge ham for vidt; men hans Energi
har ladet ham genfinde den parmenideiske Grundtanke i Enkeltheder,
der hidtil have staaet som Tilfældigheder eller Urimeligheder.


Hvad det første Punkt, Forholdet til Forgængerne, angaar,
er her det Spørgsmaal af Betydning, hvorvidt Parmenides er paavirket
af Herakleitos, hvis Lære om den absolute Vorden og den
% --- p. 143 ---
\setcounter{footnote}{0}%
\opage{143}stadige Strømmen staar i saa skarp Modsætning til Parmenides'
Hævden af den evige og uforanderlige Væren. --- Efter den
ældre Opfattelse, der hævdes af en Autoritet som \textit{Eduard} Zeller
(endnu i de seneste Udgaver af hans Værker), staar Parmenides
som en Konsekvens af den ældre joniske Skole, der søgte
efter Urstoffet i Alt. Denne Enhedssøgen fuldendes efter denne
Opfattelse hos Parmenides, idet han gav Enhedsprincipet den
mest konsekvente og abstrakte Form, som paa det givne Udviklingstrin
var muligt. Han opererer med selve det abstrakte
Begreb Væren og drager en Række af Slutninger ud fra den
Forudsætning, at hvad der er, maa være identisk med sig selv.
Han hævder, hvad man nu vilde kalde Identitetsprincipet og
Kontinuitetsprincipet i deres strengeste Former, der gør, at al
Forskel, Vorden og Forandring komme til at staa som selvmodsigende
og illusoriske. Det positive, anskuelige Udtryk for den
absolute, identiske og kontinuerlige Væren er for ham Rummet,
og Tilværelsen staar for ham som en helt igennem ensartet
Kugle, afsluttet og hvilende i sig selv. Af Joniernes Søgen efter
Urstoffet er saaledes kun det Stofliges abstrakte Form, Rumformen,
bleven tilbage. Der skal da her foreligge en direkte,
positiv Udvikling fra Thales gennem Anaximandros og Anaximenes
til Parmenides, saaledes at Xenofanes' Polemik mod Polyteismen
bliver Overgangsleddet. --- Efter en Opfattelse, der først
udtales af \textit{Bernays} og senere er hævdet af \textit{Gomperz og Diels},
stiller Sagen sig mere indviklet, idet Parmenides' Tænkning ikke
blot er en Konsekvens og Fuldendelse af den ældre joniske
Filosofi, men ganske særligt er bestemt ved et Modsætningsforhold
til Herakleitos. En Polemik mod den stadige Vorden antages
her at ligge til Grund for den energiske Hævden af den
evige Væren, og gennem det parmenideiske Læredigt finder man
en stadig Polemik mod den efesiske Filosof. Man søger at vise,
at de kronologiske Vanskeligheder, Zeller finder i denne Opfattelse,
falde bort ved nærmere Undersøgelse. Til denne Opfattelse
slutter \textit{Patin} sig, og han begrunder den dybere og alsidigere,
end det hidtil er sket.

Parmenides' Læredigt er kortfattet. Med kort, men fyndig
Begrundelse opstiller det en Række Prædikater, som skulle gælde
% --- p. 144 ---
\setcounter{footnote}{0}%
\opage{144}for det i Sandhed Værende. Det er et i poetisk Form sammentrængt
Resumé af et langt Livs energiske Eftertanke. Det vilde
nu, fremhæver Patin, være vanskeligt at tænke sig en saadan
Række af \textit{Domme} og Postulater, naar de ikke betegne Reaktioner
mod en modsat Tankegang og altsaa ere Afvisninger af modstridende
Paastande\footnote[1]{Det kan tjene til Belysning heraf at mindes, at der er to Slags positive Domme: 1) en umiddelbar og uvilkaarlig Hævden, og 2) en forsikrende Hævden, der forudsætter Tvivl eller Negation, og som udtrykkes ved særegne Partikler eller ved stærk Betoning af Kopula. Se min Afhandling \textit{Det psykologiske Grundlag for logiske Domme.} p. 30.}. Og da især, naar det ved nærmere Betragtning
viser sig, at de Prædikater, der staa fast, væsentligt
ere af negativ Art: ufødt, uforgængelig, urokkelig uforanderlig,
ubevægelig, udelelig, ikke ufærdig, uden Fortid og uden Fremtid!
Her afvises en Række Tanker, og skulde Parmenides have
fundet disse Tanker ved egen Grublen, eller er det ikke sandsynligere,
at de ere traadte ham i Møde som Andres Paastande?
Selv Udtryk som »hel« og »samtidigt værende helt igennem«
blive kun ret forstaaelige som Afvisninger. Og de Prædikater,
der saaledes afvises, ere ikke af nogen Tænker blevne hævdede
med saadan Energi som af Herakleitos. Det viser sig fremdeles
ved en nærmere Undersøgelse, som før Patin især Diels (\textit{Parmenides}
\textit{Lehrgedicht}, Berlin 1897) har foretaget, at de enkelte
sproglige Udtryk og Vendinger i Parmenides' Lære og Polemik
pege hen til Herakleitos. --- Foruden Herakleitos er det især
Anaximandros' Lære om den kosmiske Udvikling som en Differentiation
ud fra det ensartede Urstof, der er Genstand for Bekæmpelse
i det parmenideiske Digt.

Medens Patin og Diels ere enige i at finde stadige Hentydninger
til Herakleitos i dette Værk, afvige de dog i et meget
væsentligt Punkt fra hinanden. Efter Diels staar Parmenides
blot i et polemisk Forhold til Herakleitos, idet han kun tager
Hensyn til dennes Udsagn om den stadige Vorden, Udsagn, der
ifølge Parmenides bero paa den falske Forudsætning, at baade
Ikke-Væren og Væren have Realitet. Nu har jo dog faktisk
Herakleitos, foruden Paastanden om den absolute Strømmen og
% --- p. 145 ---
\setcounter{footnote}{0}%
\opage{145}de stridende Modsætninger, ogsaa Ideen om et Konstant, et
Uforanderligt, en Enhed, nemlig Loven, Tanken, Urvæsenet
(»Ilden«), der ikke forandres, men netop bestaar under Forandringernes
Bølgegang som »det, der altid var, er og vil være«
(Fragm. 20 ed. Bywater). Han stansede ikke ved de blotte Forskelligheder
og Forandringer. Hvad han bebrejdede Menneskenes
Mængde, var netop, at de ikke saa, hvorledes der trods den
stadige Opstaaen af Forskelligheder dog stedse bestod en evig
Harmoni (Fragm. 45). Diels mener nu, at Parmenides --- paa
en Maade, der ikke kan siges at være ualmindelig i filosofisk
Polemik --- har holdt sig til den ene Side af Herakleitos' Lære,
rettet sine Angreb mod den og overset den anden Side, den
positive Logoslære. »Freilich«, siger Diels (p. 69), »dies ist die
Regel in der historischen Entwickelung. So einseitig polemisiert
in der Historie wie in der Philosophie eigentlich Jeder gegen
seinen Vorgänger. Der Nachfolger stellt immer die schroffste
Reaktion dar. So war es stets, und so muss es stets sein nach
dem Gesetze der historischen Spirale«. I Modsætning til Diels
lægger Patin stor Vægt paa det, der er fælles for Herakleitos
og Parmenides. Han mener, at Parmenides kun polemiserer
mod Herakleitos, forsaavidt denne antager Forandringer i Tilværelsen.
Naar Herakleitos tillige antager en Enhed, der netop
bestaar og hævdes i og gennem disse Forandringer, saa er her
et Berøringspunkt mellem de to Tænkere. Parmenides protesterede
her kun imod, at Enheden skulde bestaa i og gennem
Forandringerne, og hævdede, at kun den havde Realitet. Men
Enhedstanken er fælles, og den lod de to Tænkere føle sig
polemisk stemte overfor Menneskenes Mængde, disse »Tvehoveder«,
der kun se Forskellighederne og Forandringerne. Hvad
denne fælles Tanke angaar, staar nu Herakleitos afgjort som
Forgængeren: »Heraklit gehört die Erfindung, Parmenides folgt
seiner Einheitslehre, bekämpft seine Vielheit und wandelt nur
in der Begründung neue Pfade« (p. 575). Herakleitos' Enhedslære
var efter Patin en Forgænger for Parmenides': »Konnte
die Vorstellung: Das Eine ist unveränderlich und unteilbar,
obwohl eine scheinbare Bewegung den Schein einer Vielheit erzeugt,
gefasst werden, ehe gedacht und ausgesprochen war, dass
% --- p. 146 ---
\setcounter{footnote}{0}%
\opage{146}die von je geglaubte Vielheit der Dinge nicht eigentlich sei,
sondern nur als das Scheinerzeugnis eines ewig bewegten Einen?
Ist denn nicht die Leugnung der Vielheit Vorbedingung der
strengen Einheitslehre? Musste nicht der Glaube an sie durch
die Annahme einer der gewöhnlichen Vorstellung immer noch
näher stehenden ewigen Bewegung im Einen erschüttert sein,
ehe man so weit gehen konnte, auch diese Bewegung selbst
noch für Täuschung und Trug zu erklären? Und musste nicht
erst schon die Einheit des die Individuen trennenden Gegensatzes
behauptet sein, ehe man wagen konnte, die Annahme
von Gegensätzen selbst schon als den Anfang des Irrtums zu
bezeichnen? (p. 511 f.)« --- Man maa sikkert indrømme, at en
Udvikling som den her antydede er naturlig, selvom det ikke
var nødvendigt, at de to Tankeskridt, det drejer sig om, bleve
foretagne af to forskellige Tænkere. Men da nu Herakleitos'
Filosofi foreligger historisk, og da Kronologien (efter hvad Diels
og Patin have søgt at vise) tillader at opfatte den som historisk
Forudsætning for den parmenideiske Lære, har Patins Opfattelse
Sandsynlighed for sig, og Herakleitos har da paavirket Parmenides,
ikke blot ved at ægge hans Modsigelse frem, men ogsaa
ved at udtale en Tanke, der kunde udvikles videre og blev udviklet
videre, endog i extrem og paradoxal Form.

Den Opfattelse, Patin saaledes er kommen til angaaende det
første af de to ovenfor fremhævede Punkter, bliver nu ogsaa af
Betydning for hans Opfattelse af det andet Punkt. Naar der
mellem Herakleitos og Parmenides ikke blot har bestaaet et
negativt, men ogsaa et positivt Forhold, vil det ogsaa vise sig
lettere at finde en indre Sammenhæng mellem det parmenideiske
Læredigts to Dele, Fremstillingen af den uforanderlige Væren
og Fremstillingen af den vordende og modsætningsfulde Verden.

I Regelen har man antaget, at der bestod et aldeles udvortes
Forhold mellem Læredigtets to Dele. Man mente, at
Parmenides i anden Del fremstillede en Lære, han ikke selv
hyldede, men at han gjorde »de Dødelige« den Indrømmelse for
en Stund at stille sig paa deres Standpunkt og vise, hvorledes
Tingene her tage sig ud. Ud fra hans Teori om den uforanderlige
og udelte Væren var hverken en Naturfilosofi eller en Kos%
% --- p. 147 ---
\setcounter{footnote}{0}%
\opage{147}mogoni mulige. En saadan Indrømmelse til den sanselige Anskuelse
var naturligvis en Inkonsekvens fra Parmenides' Standpunkt,
og Gomperz har vel Ret i, at Parmenides ikke kunde
have gjort denne Indrømmelse, hvis der ikke havde været en
Vaklen i hans eget Sind, --- hvis han ikke havde følt, at den
deduktive Tankegang ikke kunde udtømme Tilværelsen, og at
i hvert Tilfælde Skinnet af kvalitativ Flerhed og Forandring
stod tilbage og æskede Forklaring. Det er jo klart, at Problemet
om Enhed og Mangfoldighed ikke blot maa blive uløseligt,
men føre til et skrigende Paradox, naar man begynder med at
erklære Enheden for det eneste Gyldige og Reale. I den nyeste Tid
har Diels dog hævdet, at den anden Del af Læredigtet hverken beror
paa en Akkommodation eller tyder paa en Vaklen hos Parmenides.
Diels finder ogsaa i denne anden Del en stadig Polemik
mod Herakleitos og antager, at den Skizze af en empirisk
Verdensopfattelse, der her gives, kun skal vise, hvorledes Tænkere,
der ikke kunne hæve sig over de sanselige Forskelligheder
og Forandringer, maa opfatte Sagen. Ved »de Dødelige«, hvis
Meninger den anden Del af Digtet fremstiller, skal der altsaa
efter Diels tænkes paa visse Filosofer, især paa Herakleitos.
»Den anden Del af Digtet«, siger Diels (p. 63), »er ikke andet
end en kritisk Oversigt over de forudgaaende Tænkeres stridende
Anskuelser, en Doxografi, der --- ligesom hos Peripatetikerne
blot skal tjene Skolens propædeutiske Formaal«. Allerede i en
Afhandling i et Festskrift til Zeller (1887) for nogle Aar siden
havde Diels udtalt denne Opfattelse, efter hvilken Forholdet mellem
Digtets to Dele jo egentligt bliver endnu mere løst og udvortes end efter de tidligere Opfattelser.

Patin har nu her, forekommer det mig, den Fortjeneste at
have gjort et Forsøg paa at bringe den anden Del af Digtet i
indre positiv Sammenhæng med den første Del. Han gaar ud
fra, at en saa indtrængende og speciel Naturforklaring, som Parmenides
efter Fragmenterne at dømme her har forsøgt, nødvendigvis
maa være udsprungen af en alvorlig og positiv Interesse
for Emnet. Det kan ikke blot have været andre Dødeliges Meninger,
han har villet fremstille. Til de Dødelige  regner han
% --- p. 148 ---
\setcounter{footnote}{0}%
\opage{148}ogsaa sig selv, og han har altsaa indset, at der foruden den
sikre Viden, der bevæger sig paa den rene Logiks Grund, gives
en Opfattelse af Tingenes sanselige Forandringer og Forskelle,
en Opfattelse, der kan være mere eller mindre rigtig. Det er
ved hele Digtets Begyndelse en Gudinde (Sandhedens eller
Retfærdighedens Gudinde), der anviser Tænkeren den rene Logiks
ensomme Vej; og det er den samme Gudinde, der ved
Overgangen til Digtets anden Del erklærer at ville lære ham
de Dødeliges Meninger og vise ham Verdens Indretning eller
Udvikling paa en saa grundig Maade, »at ingen menneskelig
Anskuelse skal kunne gøre ham Rangen stridig«. Der er altsaa
dog ogsaa paa det sanselige, empiriske Omraade en Rangforskel
mellem Anskuelser, og det kan da ikke være det blotte,
tilfældige Skin, der raader, men der maa være Principer, som
gøre det muligt at skelne mellem det Rigtige og det Falske.
Og ved nærmere Betragtning viser det sig nu, at disse Principer
ere de samme som de, der ligge til Grund ved Identitetslæren
i Digtets første Del. Ogsaa i Digtets anden Del opererer
Parmenides med Identitets- og Kontinuitetsprincipet (for
at bruge et moderne Udtryk), hvor han antyder en Kritik af
de forskellige Meninger. Bag Fremstillingen af den sanselige
Verdens Modsætninger og Udviklingsprocessen ligger stedse som
Baggrund Enhedstanken som den, hvortil Alt viser tilbage. Saa
store Problemer, en fuldstændig Forbindelse af Digtets to Synspunkter
end volder, saa kan det dog, ifølge Patin, ikke siges,
at Parmenides selv har opgivet det første Synspunkt, da han
foresatte sig at benytte det andet. Hint danner stadigt Grundlaget
for dette.

Selvom Patin i det Hele skulde have Ret i sin Opfattelse,
falde naturligvis de filosofiske Vanskeligheder ved Forholdet
mellem Læredigtets to Tankerækker ikke bort. De have siden
den Tid fulgt den filosofiske Tænkning paa dens Vej under forskellige
Former. Men der kommer unegteligt større Enhed ind
i Parmenides' Fremstilling og Tankegang, naar man følger Patin's
Fortolkning, som vidner om en energisk Fordybelse i den
gamle Grækers Ideer.

% --- p. 149 ---
\setcounter{footnote}{0}%
\opage{149}Man kan vente en interessant Fortsættelse af Diskussionen
om det parmenideiske Digt, da Diels (Parmenides Lehrgedicht
p. 4) udtrykkeligt har betegnet sin 1897 udkomne Bearbejdelse
 som foreløbig og forbereder en ny Udgave, der særligt
skal have den filosofiske Sammenhæng for Øje, medens han
i det foreløbige Arbejde blot rettede sin Opmærksomhed mod
Parmenides som Digter.

% --- p. 150 ---
\setcounter{footnote}{0}%
\opage{150}\essay[Holberg som Filosof. (1903)]{LUDVIG HOLBERG SOM FILOSOF.}
\streg

Holberg er kendt af Alle som Digter, af Færre som Historiker, og af endnu Færre som Filosof. Og dog, hvis man
kunde spørge ham, om der ikke til Grund for hans hele storartede og mangfoldige Produktion, der omfattede baade Komediedigtning, Historie og Filosofi, laa en fælles Interesse, vilde han
sikkert svare, at han først og fremmest var Filosof.

Det var, siger han, hans Forsæt at moralisere paa alle mulige Maader: ved historiske Skrifter, ved Satirer, ved Epigrammer, ved den underjordiske Rejse og ved særlige filosofiske Afhandlinger (i Epistlerne og i »Moralske Tanker«). Overalt har
det været hans Hensigt »at bestride almindelige Vildfarelser og
at distingvere Dyders og Lasters Realitet fra deres Apparence, ---
at vise, hvorledes Skyggen tages for Legemet, og hvorledes 
Lyder konfunderes med Dyder« (Fortalen til »Moralske Tanker«).
Ganske særligt betragter han sine Komedier fra dette
Synspunkt: »Hvad Comoedier angaar, da kand man sige, at
ingen kraftigere Skrivemaade er opfunden at moralisere paa, og
hvorudi Dyders og Lasters Character livagtigen kand gives«.
»Det er et Arbejde, som er en Philosopho anstændigt«. »Det
er i den Henseende, at jeg selv ikke haver holdet det min
Stand og Alder uværdigt at moralisere paa den Maade«. Og
Molière, som har genoptaget »den Konst at forfærdige Moralske
Skuespil«, der havde ligget nede siden Plautus' og Terentius' 
Dage, kan »regnes blandt vore Tiders store Philosophos«. (»Mo-
% --- p. 151 ---
\setcounter{footnote}{0}%
\opage{151}ralske Tanker« 1744 p. 15—21). »Der udfordres af en Comoedie-Skriver,
først at han er Philosophus og har nøje udstuderet
det, som kaldes Ridiculum udi det menneskelige Kiøn; for det
andet, at han har det Pund at gennemhegle Lyder saaledes,
at han diverterer tillige«. (Just Justesens Betænkning over Comoedie).


Holberg glæder sig da i sit Levned over, at han har bragt det
dertil, »at man nu paa Dansk kand læse om Dyders og Lasters
Lærdom«, og at selv Tjenestepiger nu (ved Hjælp af hans Komedier)
kunne filosofere over det menneskelige Køns Daarligheder.


For os staar det klart, at Holberg ikke vilde have været saa stor en Digter, Historiker og Tænker, som han var, dersom
der virkeligt saaledes havde ligget et bevidst Forsæt om at opnaa
en moralsk Virkning til Grund for hans Produktion. Forsætter
komme jo meget ofte bagefter, som Forsøg paa at retfærdiggøre
det, der er gjort af uvilkaarlig Trang og umiddelbar Indskydelse.
Og at Holberg har haft den rene Kunstner- og Forskerglæde
ved sine Arbejder, fremgaar tydeligt nok af hans Komedier.
Interessen for Karakterskildring, den rige Fylde af komiske Enkeltheder
og den overlegne Humor har i hvert Tilfælde i Produktionens
Øjeblikke taget Magten fra det moraliserende Forsæt,
om det har været til Stede. Moralen i Stykkernes Slutning
er tydeligt nok paasat; den svarer ikke altid til Stykkernes Indhold
eller deres Grundstemning. Det er Ædrueligheden efter
Produktionens Rus, der nu fører Ordet, og Holberg ofrer paa
sin pedantiske Tids Alter, idet han selv er et Barn af sin Tid.
Og noget Tilsvarende gælder, som jeg skal forsøge at vise, ogsaa
for hans filosofiske Arbejder.

Der forekommer da ogsaa enkelte Udtalelser af ham, i hvilke
den geniale Evne og den kunstneriske Interesse betones. Saaledes
naar han (i »Moralske Tanker«) siger om Plautus' Komedier:
»Derudi fremskinner over alt den Festivitet, som er
Siælen udi en Comoedie eller Anima Theatri, som ingen kand
tilvejebringe, uden den, som Naturen haver dannet dertil, da
mange derimod med Fliid, Arbejde og en Slags Skiønsomhed
kan udarbejde saadanne Skue Spil som Terentii«. Eller naar
% --- p. 152 ---
\setcounter{footnote}{0}%
\opage{152}han (i samme Bog) erklærer: »Sigtet af mine Skrifter haver
ikke saa meget været at censurere Fejl og Laster, som at give Specimina udi visse ziirlige Videnskaber«.

Man manglede paa Holbergs Tid et klart Begreb om det,
vi nu kalde Geni, og derfor fremstillede man gerne som Frugt
af bevidst Plan og af et nøgterne Forsæt, hvad der i Virkeligheden
skyldtes en indre, skabende Trang, der ikke véd, hvad
den gør, og hvis Værk bagefter vækker Undren hos selve det
Frembringende. Det moderne Begreb om Geni i denne Betydning
af Ordet arbejdede sig frem i det 18. Aarhundredes engelske
Æstetik (under Indflydelse af det fornyede Shakespearestudium)
og udformedes bestemtere af Rousseau, Kant og Goethe ---

En saadan uvilkaarlig Trang har ført Holberg til at digte ---
forske og tænke. Og søger man et Mellemled eller et fælles
Grundlag for hans Digten og Forsken, saa er den langt snarere
at finde i hans psykologiske Interesse end i en moralsk Hensigt.
Den psykologiske Interesse forbinder Digteren, Historieforskeren
og Tænkeren. Digteren kan kun virke ved det enkeltes
Billede, den enkelte Skildring, uden at han kan afgøre, eller behøver
at afgøre, om det er et Exempel eller en Undtagelse,
han fremstiller. Det er Glæden og Interessen ved selve Billedet,
ved den plastisk formede Fylde, hvorpaa Kunstværkets Interesse
beror. Historikeren viser Menneskekaraktererne i deres
hele historiske Sammenhæng, som Virkninger og som Aarsager,
og han formaar at skaffe Stof til Sammenligning. Tænkerens
Opgave er det da at opstille Typer og finde almindelige Love
eller Problemer. I alle disse Retninger har Holbergs psykologiske
Interesse virket. Der fremtræder, som Georg Brandes i
sin Bog om Holberg (p. 135) har paapeget, en vis Modsætning
mellem hans psykologiske Fremstillinger i Komedierne og
hans psykologiske Iagttagelser og Reflexioner i de filosofiske
Arbejder. Som Filosof interesserer han sig især for psykologiske
Gaader og for sammensatte og tilsyneladende paradoxe
Karakterer. Han nævner (i »Moralske Tanker«) Gnieren,
den fattige Ødeland, Fromheds og Ugudeligheds samtidige Tiltagen,
Tilbøjeligheden til at gøre det Forbudte. Hans Forskerinteresse
kommer i Bevægelse for at finde Løsningen af disse

% --- p. 153 ---
\setcounter{footnote}{0}%
\opage{153}Gaader. Men denne Interesse for det psykologisk Gaadefulde
mærkes ikke i hans dramatiske Værker; Karaktererne i hans
Komedier ere ud af ét Stykke. Sammensatte og indviklede Karakterer
har Holberg ikke formaaet at forme i sin Kunst, skønt
han havde Blik for dem i Livet, og skønt de interesserede ham
som Tænker. Maaske har han her været ledet af Tidens Smag:
han siger udtrykkeligt (Epist. 66), at »et Skuespil uden outrerede
Karakterer ingen Virkning har«. Men en »outreret« Karakter
er netop en Karakter, i hvilken et enkelt Træk behersker det
Hele, uden at komme i Modsætning eller Vexelvirkningsforhold
til andre Træk.

Holbergs Blik paa Menneskelivet var ikke lyst. Han fandt
megen Tomhed og Forfængelighed om sig, megen Skal og liden
Kærne. »Den heroiske Dyd«, der trodser Skade og Fortræd, er
sjelden. Og de fleste Mennesker ere maskerede og paatage sig
ligesom Komedianter fremmede Skikkelser, saalænge som Spillet
vaarer. Oprigtigt Venskab er sjeldnere end oprigtigt Had. Desuden
ere moralske Sygdomme fast inkurable; den, der lider af
Laster, tror ikke selv, at han fejler noget. Selve Moralisterne
ere ikke bedre end de Andre, kun mere irritable: det er lige-saa let at gøre en Helgen vred, som det er vanskeligt at gøre
ham god igen, og Holberg vil hellere have et helt Regiment
Soldater paa Halsen end en eneste Helgen! De deklamere
mere end de belære; »vore Tiders sminkede Moral bestaar mere
af Ord og Talemaader end af Realitet«.

Om det altid har været saaledes, vil Holberg ikke afgøre.
Han vil ikke tage Parti i den Strid, om Verden gaar frem eller
tilbage. Hvis han blev tagen til Dommer, vilde han, efter at
have hørt paa begge Parter, afvise Sagen og dømme, at Omkostningerne
skulle ophæves paa begge Sider.

Hvad der kan gøres, og hvad Holberg betegner som sin
egentlige Opgave, er at afsløre Selvbedraget og vise, hvor ofte
man tager Skyggen for Legemet. Dette betegner han som sit
»fornemste Sigte«. Atter og atter kommer Udtrykket om at tage
Skyggen for Legemet igen hos ham. Han har --- som Hamlet
forlanger af Skuespilkunsten --- baade i sin Kunst og i sin
Forsken villet holde et Spejl op for Naturen. Ved Arbejdet paa
% --- p. 154 ---
\setcounter{footnote}{0}%
\opage{154}denne Opgave kunde hans Kunst og hans Tænkning, hans Poesi
og hans Prosa, hans Genialitet og hans Reflexion mødes. Hvad
han tilsidst og paa alle Omraader søgte, var --- for at bruge et
moderne Udtryk --- den personlige Sandhed. Det var et meget
udvortes Udtryk herfor, han fandt, naar han betegnede det som
sit egentlige Forsæt »at moralisere«. ---

Fra Holbergs Filosofi i den videste Betydning, i hvilken
han bruger dette Ord, gaar jeg nu over til hans Tænkning over
Emner, der i snevrere Betydning kunne kaldes filosofiske.

Af to udenlandske Tænkere er Holberg især paavirket, af
John Locke, »den store engelske Philosophus Locke«, der saa
energisk tog Ordet for at føre al Erkendelse tilbage til Erfaring,
og af Pierre Bayle, i hvis Skrifter han fandt »en dybsindig
Philosophie, en grundig Critique og en ugemeen Accuratesse«; og som han især priste for den Skarphed, hvormed han fast
holdt Problemernes Brod overfor forhastede Løsningsforsøg. ---
Det er karakteristisk for Holberg, at netop disse to Mænd have
været saa meget for ham. Overfor Skolastiken og Dogmatismen,
som endnu raadede i den lærde Verden, stillede han Erfaringen
 og de store Gaader. Han hævder paa én Gang Ædrueligheden og Nødvendigheden af stadigt fortsat Søgen. Forsigtigt
nærmer han sig til Videnskabens Grænser, og end ikke der
taber han Besindelsen, men fastholder Retten til og Nødvendigheden
af uhildet Undersøgelse. Den nærmere Samtids Forfattere
satte han ikke stor Pris paa. Han kommer flere Gange
(baade i sit Levned, i de »moralske Tanker« og i Epistlerne)
tilbage til den Paastand, at der er sket en Tilbagegang i den
videnskabelige og den literære Verden, en Paastand, som Rousseau
nogen Tid senere udtalte paany, og som laa til Grund for
hans Fordring om en ny Kultur.

Ved at gennemgaa Holbergs Udtalelser om filosofiske Emner
finder jeg, at de kunne henføres til tre Grupper. Den første
Gruppe indeholder Udtalelser om speciellere etiske og religiøse
Spørgsmaal; den anden Gruppe giver Antydninger i Retning af
en almindelig Verdensanskuelse; den tredje Gruppe indeholder
Udtalelser, som indskærpe Nødvendigheden af stadig Søgen og
Stræben overfor Tilværelsens og Tankens højeste Spørgsmaal.

% --- p. 155 ---
\setcounter{footnote}{0}%
\opage{155}1. En Sætning af Bayle, som Holberg især med Forkærlighed
kommer tilbage til, er den, at Menneskenes Handlinger
ikke saa meget bestemmes ved deres Tro, deres Principer og
deres Fornuft, som ved deres Temperament, eller ved deres
»Affecter og Passioner«. »De fleste Menneskers Idrætter styres
ikke af Fornuften, men Enhvers Temperament er Hiulet, som
driver det frem\dots{} Lærdom, Underviisning og god Optugtelse
er vel ikke unyttig, men giver dog Mennesket ikke et andet
Hierte og en anden Complexion: Thi den rører ikkun ligesom
den udvortes Hud«. (Ep. 424 kfr. Ep. 210). Hvis man ret
lagde Vind paa at kende sig selv, vilde man --- som Holberg
mener (»Moralske Tanker« 1744, p. 91) --- finde, at, hvad man
er tilbøjelig til at opfatte som hellig Nidkærhed, skyldes rent
organiske Tilstande. Det viser sig især, naar Nidkærheden har
sin Ebbe og Flod efter det legemlige Befindende og ikke er
den samme i Medgang og i Modgang, i Sundhed og i Sygdom.
Man skal derfor vogte sig for at hente Beviser for en Læres
Sandhed af dens Tilhængeres Liv; thi Mennesker have levet ædelt
og skønt med vantro Meninger, og Rettroenhed er ikke altid
forbunden med Dyd. (Ep. 464). --- Denne Modsætning mellem
Lære og Liv hørte til de psykologiske Iagttagelser, der særligt
interesserede Holberg, ikke mindst fordi Motivet til Religionsforfølgelse
maa falde bort, naar man ikke fra visse bestemte
religiøse Meninger kan slutte til nogen bestemt praktisk Adfærd.
En dyberegaaende Undersøgelse af Forholdet imellem Sjælelivets
forskellige Elementer, særligt af Tankelivets og Viljelivets
indbyrdes Forhold, havde Tidens Psykologi ikke Betingelser for.
Holberg stansede ved det Paradoks, han her fandt, og lod sig
nøje med at drage en gavnlig praktisk Konsekvens deraf. Og
det var af stor Betydning, at det blev paavist, at Forholdet
mellem Livsanskuelse og Livsførelse var mere indviklet, end
man hidtil havde antaget.

Alligevel lægger Holberg stor Vægt paa fornuftmæssig Begrundelse
af Moralen. Moralfilosofi hørte efter hans egen Erklæring
til hans Yndlingsstudier i hans ældre Dage. Moralen
bør, mener han, gaa forud for Teologien, og Børn burde undervises
i Moral, før de lære Katekismus. Han vil --- med Fare
% --- p. 156 ---
\setcounter{footnote}{0}%
\opage{156}for at blive kaldt Hedning --- vise Menneskets Pligt af Naturens
Lys. En Sætning som den, at vi ikke maa gøre mod
Andre, hvad vi ville, at de ikke maa gøre mod os, kan vel
ikke siges at være medfødt; deri giver Holberg Locke Ret; men
»dette Naturens store Bud« vil Enhver dog uvilkaarligt anerkende,
saa snart det forklares for ham. (Ep. 144). Meget dybt
gaaende er iøvrigt Holbergs Etik ikke. Han har et skarpt Blik
for Yderligheder og indskærper: »Alt hvad som er for meget
skader; thi all Dyd bestaaer udi Mediocritet, og, saa snart der
gaar ud over de Grændser, bliver den metamorphoseret til 
Last.« (»Moralske Tanker« p. 433). Dog viser det sig ved nærmere Undersøgelse, at Holberg ved sin Lovprisning af Middelvejen nærmest tænker paa, at hver Enkelt og ethvert enkelt
Samfund efter sin særegne, individuelle Natur kræver en vis
bestemt Udvikling, og at hvad der er for meget for den Enkelte,
kan være for lidet for den Anden, »ligesom [i Sygdomstilfælde
et Legeme haver nok udi een Pille, men et andet vil have to
eller tre af samme Vægt« (»Moralske Tanker« p. 429). Dette
stemmer ogsaa ganske med hans stærke Fremhæven af Temperamentets
 og Naturgrundlagets Betydning. Ogsaa paa Bedømmelsen
 af Karakterer og Handlinger anvender han denne Betragtningsmaade, idet han fordrer, at der skal tages Hensyn til de
Betingelser og Hindringer, den Enkelte har i sin egen Natur.
Han antager, at Gud fælder en mildere Dom ved samme Misgerning
»over en Cholericum end over En, der er født med
temperered Complexion«, og fordrer mere af Den, der er udstyret
med store Naturgaver, end af en tykhudet Person. (»Moralske
Tanker« p. 363). Man undskylder en Podagrist, at han
ikke kan danse, men ikke en Hypokondrist, at han ikke kan være
lystig og tækkelig i Selskab. (Mor. T. p. 408; 412). Den, der ingen
Galde har, kan ikke berommes for Koldsindighed; at styre en
vild og modig Hest koster dog større Kunst end at ride paa et
spagt Æsel (Mor. T. p. 538). »Endskiønt en Misgierning er hos
tvende Personer lige stor, saa er den dog ikke lige strafværdig,
efterdi En har hæftige Affecter at bestride, men den Anden
synder af fri Vilje.« (Ep. 467). --- Hvad Holberg her, støttende
sig til psykologisk Erfaring, indvender imod den sædvanlige
% --- p. 157 ---
\setcounter{footnote}{0}%
\opage{157} Morals Fordring om, at alle skulle skæres over én Kam, det har
endnu ikke faaet sin fulde Anerkendelse og Gennemførelse i
Etiken og i Retslæren. Der ligger i denne Fordring Nødvendigheden
af en Individualisering, som baade i Praxis og i Teori
stiller os de største Opgaver.

2. Ligesom Holberg i sin Moral søger at støtte sig til Erfaringen,
saaledes stræber han at bygge sin Verdensanskuelse
paa et Erfaringsgrundlag. Han polemiserer stadigt mod skolastisk
og spekulativ Metafysik. Forsigtig, som han er, taler han
 ikke altid rent ud; men det er tydeligt nok, til hvilken Side
han hælder.

Den nyere Tids astronomiske Opdagelser have ifølge Holberg
givet Menneskene en større Idé om Gud, end de havde forhen,
da de betragtede Jorden som en stor Del af Verden og som
den, for hvis Skyld de andre Himmellegemer vare til. Nu er
man klar over, at Jorden kun er et Sandskorn i Forhold til
hele Verden, og Guds Storhed og Ubegribelighed bliver da indlysende,
om den ikke var det forud. Man maa nu tilstaa, at
der intet daarligere Foretagende kan være, end at ville beskrive
Gud! (Mor. T. p. 134---137.) --- Ja, Holberg er endog aabenbart
tiltalt af den Mening, der i Bibelen ved »Verden« kun
forstaar vort Solsystem. Derved vil det, mener han, være lettere
at forstaa, hvorledes Gud kan have tøvet saa længe med at
 skabe Verden. Han siger herom (Ep. 1): »Mange fornuftige
Mænd holde for, at man med fuldkommen Vished ikke kand
sige, om den nu værende Verdens Skabning, som Moses beskriver,
er Guds første Bygning, eller om det med Verden i Almindelighed
ikke er af samme Beskaffenhed, som med dens Parter
i Særdeleshed, nemlig at af eens Undergang flyder en andens
Opkomst og Generation. Nogle ere endog af den Mening, at
under Verdens Dannelse, som udi Skabelsens Historie beskrives,
befattes aleene den Hvirvel, hvorudi Jorden med de andre Planeter
velter om Solen; thi den vide Verden, hvoraf vi end med
Perspectiver, maa ske, ikkun se den mindste Deel, er større end
den meenige Mand indbilder sig: Herudi deciderer jeg dog intet«\footnote[1]{Til den Hypotese, at Skabelseshistorien i første Mosebog kun angaar \opage{158}vort Solsystem, sigter aabenbart en Udtalelse i det omtrent 20 Aar forud for Epistlerne skrevne første Brev om hans Levned (dateret 31. Decbr. 1727). Efter at have omtalt, hvorledes han var kommen ud over visse religiøse Tvivlsmaal, siger han: »Der var vel endda nogle Dubia tilbage, som Studium Physicæ coelestis, hvilket jeg havde synderlig Smag udi, foraarsagede mig; mig syntes, at Vortices Cartesiani \& Mundorum Pluralitas vare imod den Underretning, som Moses giver om Skabelsens Historie; men nogle sunde Fortolkere over Moses Bøger betoge mig og disse Scrupler, i det de fornuftig og lykkelig foreenede vore Tiders Philosophie med Skabelsens Historie.«}.


% --- p. 158 ---
\setcounter{footnote}{0}%
\opage{158}Denne Udtalelse er karakteristisk for Holberg, idet den viser hans Stræben efter at forfølge de Tanker, Erfaringsvidenskaben,
fremkalder, i deres sidste Konsekvenser, selvom han derved kommer
i Strid med gængse og overleverede Antagelser. Verden
har udvidet sig for ham, og det bliver vanskeligere for ham at
naa til sikre og afsluttende Ideer.

Paa lignende Maade stillede han sig overfor Striden mellem
Bayle og Leibniz om, hvorvidt det Onde i Verden kunde forenes
med Verdens Udspring af et almægtigt og algodt Væsen.
mener vel, at Antagelsen af et Syndefald først i Englenes
og siden i Menneskenes Verden hjælper til at holde sig Vanskelighederne
fra Livet. Men han er klar over, »at all Tvistighed ikke
derved kand ophæves«. Heller ikke her »drister han sig til at skride til nogen Decision«. (Ep. 320). Med en vis Sympati
omtaler han (Ep. 225) en af hans Ven B. S. (som formodes
at være Rektor Bernhard Schnabel i Roskilde) fremsat Mening,
som løste Bayles Problem ved at antage en oprindelig Forskel
mellem et aktivt Princip, Guddommen, og et passivt Princip,
Materien. Guddommens Planer kunne kun ufuldkomment realiseres
paa Grund af Materiens Modstand. Denne Lære, som i
Oldtiden Platon og Stoikerne hyldede, staar i Modsætning baade
til den kirkelige Teologi og til den manikæiske Dualisme, hvilken
sidste ifølge Bayles Mening var den nødvendige Konsekvens,
naar man ikke tog Fornuften fangen i Troens Lydighed. »For
at sige mine ringe Tanker derom,« siger Holberg, »da synes
mig, at dette Argument kunde komme Monsr. Bayle til at svede,
hvis han havde med en Hedenske Philosopho at bestille«. Men
--- føjer han til --- Bayles Modstandere vare jo Kristne, der af

% --- p. 159 ---
\setcounter{footnote}{0}%
\opage{159}Aabenbaringen vide, at Sætningen »Af Intet kommer Intet« ikke
gælder! --- Den Løsning, Holberg her omtaler, var den, Rousseau
og Voltaire, og i det 19. Aarhundrede Stuart Mill, holdt
sig til. Holberg har altsaa anset den for holdbar, naar man
udelukkende gik ud fra den menneskelige Fornuft.

3. Holberg finder overhovedet, at naar man ser bort fra
almindeligt antagne Meninger og grunder sine Domme alene paa
Erfaringen, der er den tilforladeligste Lærerinde, saa er der
mange Grundsætninger, hvis Gyldighed kan omtvistes. (Mor. T.
p. 493; 505). Filosofien bliver for ham --- hvad den senere var
for Mænd som Carlyle og Stuart Mill --- en Kamp mod Vanen.
Naar han saa ofte erklærer ikke at ville »decidere«, er det ikke
blot af Hensyn til de herskende Meninger og Autoriteter; det
er ogsaa, fordi han har større Blik for Problemer og Vanskeligheder
end hans Samtid, og fordi selve det at søge, at »ransage«,
at forske for ham var noget Værdifuldt. Reformationens
Grundvold, siger han, var den kristelige Frihed i at ransage,
forend man tror; men denne Grundvold holdt man ikke fast
Ved, og Følgen var, at man fik mange Paver i Stedet for én.
(Mor. T. p. 335)\footnote[1]{De samme Klager findes, og i lignende Udtryk, hos Giordano Bruno og Thomas Hobbes.}. Det staar i et Menneskes Magt at ransage,
men ikke at tro. »Jeg kan,« erklærer Holberg, »for mit Vedkommende
sige, at hvis jeg farer vild, saa haver det kostet mig
mere at vildfare, end det haver kostet mange Andre at tro ret.«
(Mor. T. p. 349). I den Strid, der paa hans Tid førtes mellem
Fritænkere og Ortodoxe, forarges han oftere over disse end over
hine: »Jeg haver læset de fleeste nye Deistiske Skrifter uden
nogen merkelig Anfegtning: men jeg haver forarget mig over
adskillige Besvaringer.« »Jeg haver fundet udi Deisternes Skrifter
Objectioner, hvorved jeg haver stønnet, og Knuder, som jeg
ikke saa let haver kunnet løse. Saadant haver dog ikke betaget
mig Modet, men heller opmuntret mig til at undersøge alting
nøye, hvilket jeg endnu dagligen giør.« (Ep. 2). At beraabe sig
paa Fornuftens Ufuldkommenhed og Guds uransagelige Veje
tjener enhver Religion lige godt og fører let til at begunstige
Dovenskab (Ep. 16—17). »De, som sige sig ingen Vanskelig%
% --- p. 160 ---
\setcounter{footnote}{0}%
\opage{160}hed at se, tractere alting cavalierement, og tilkendegive, at de
aldrig have umaget sig at examinere. De blive let til magelige
fede og veltrivende Conformister«, der »med Frækhed fordømme
alle dem, som ikke blindt følge deres Meeninger, hvilke de selv
aldrig i Grunden have ransaget.« (Ep. 25).

Holberg havde et aabent Blik for, at der er Spørgsmaal,
som ikke kunne løses. Ved at behandle dem, saaledes som de
fremstille sig i Erfaringens Verden, f. Ex. ved at undersøge
Sjælelivet, saaledes som det fremtræder i sine Ytringer, lære
vi dog altid noget, hvorimod abstrakte Spekulationer, f. Ex. om
Sjælens Natur og Tilblivelse, kun ere » Eftersøgninger, med hvilke
Sindet forgieves plages.« (Mor. T. p. 152 f.). Holberg vil lade
Problemerne udfolde sig paa Erfaringens Grund og følge dem
saa langt, som den menneskelige Tænkeevne formaar; men heller
ikke længer. Han erklærer, at han ikke er nogen Hader af det
metafysiske Studium, uden hvis man ved Metafysik forstaar »de
mangfoldige Termini og subtile Distinctioner, som man udi
Disputatzer misbruger til at indvikle sin Contra-Part udi«. Men
der gives en anden Betydning af Ordet, i hvilken Metafysik er
»et Studium, som er en Philosopho anstændigt«, selvom man
her stadigt bevæger sig paa »de uvisse Tings« Omraade. (Tredje
Brev om sit Levned). Man skal blot bevare sig fri for falsk
Sikkerhed. Undersøgelsen har sin Værdi i sig selv. Holberg
udtaler sig undertiden herom saaledes, at man næsten tror at
høre Lessing eller Stuart Mill. »Jeg skriver,« siger han (Ep. 52)
»ikke saa meget for at dogmatisere, som for at tilkiendegive
visse Meeninger, som jeg ved Meditation er falden paa, og som
ikke maa ansees uden som Problemata«. Skeptiker vil han ikke
være; men han anser med Bayle den skeptiske Filosofi for en
gavnlig Modvægt mod Teologers og Filosofers Paastaaelighed
og Selvtillid (Ep. 333). Paa det aandelige Omraade er Strid
gavnlig; Fred betyder i Længden Sløvhed og Uvidenhed. »Det
er Mennesket tienligere at trette og tviste, end ligesom umælende
Creature at leve udi en brutal Samdrægtighed. Derfor have
adskillige ikke uden Grund lignet de Uleiligheder, som fødes af
Usamdrægtighed, med Stormvinde, hvilke nedslaae Huse og
% --- p. 161 ---
\setcounter{footnote}{0}%
\opage{161}Træer, men tilligemed rense Luften og hindre, at den ikke af
for megen Stilhed raadner«. (Mor. T. p. 501). ---

I Holbergs Filosofi finde vi som Hovedtræk Interesse for
Erfaring, dristigt Udblik til større Horisonter end de sædvanlige,
udpræget etisk Interesse og inderlig Overbevisning om den
stadige Forskens og den aandelige Kamps Betydning. Paa enkelte
Punkter mærkes hos den ædruelige Forfatter en stille Begejstring
og en stærk Undersøgelsestrang. Fremfor Alt vil han
undgaa at tage Skyggen for Legemet, og hans Filosoferen hænger
saaledes sammen med Motivet for hans digteriske Produktion.


Hos den danske Litteraturs Grundlægger finde vi Træk, der
ogsaa have været karakteristiske for den danske Tænkning i den
følgende Tid, overalt hvor den har udfoldet sig naturligt. Maaske
tør man i disse Træk se en Ytring af selve Folkekarakteren.
En Sammenligning med den svenske Filosofis Historie bekræfter
dette. (Se min Afhandling »Svensk og dansk Filosofi« i
»Norden 1902«).

% --- p. 162 ---
\setcounter{footnote}{0}%
% CHECK p. 162: centred line: heading or verse
\opage{162}\essay[Henrik Steffens (1902)]{HENRIK STEFFENS.}
\begin{center}{\small TALE VED EN FEST PAA ELERS KOLLEGIUM}\\
{\small DEN 11. NOVEMBER 1902.}\end{center}
\streg

Da Henrik Steffens som gammel Mand skrev paa den sidste
Del af sit Levned og saa' tilbage paa sin Skæbne og sit Værk,
kom det Spørgsmaal frem hos ham: hvad har Du egentligt ud
rettet?

»I min stille, ensomme Ungdom blev jeg greben af en
Længsel, der førte mig hen imod Religionen og Naturen i deres
hele Fylde. Et uroligt, ja vildt Temperament drog mig i modsat
Retning, og ofte nok mit hele Liv igennem laa jeg under i 
den heftige Kamp; men den stille, indre Stræben forlod mig
ikke. Jeg førtes i et betydningsfuldt Øjeblik ind paa en større
Skueplads; en stor aandelig Bevægelse greb mig. Jeg blev dragen
ind i en filosofisk Kamps gærende Midtpunkt; men den Mester,
jeg selv anerkender, vilde ikke kalde mig Filosof. Jeg tog Del
i en af de mest storartede Krige, Aartusinder have oplevet, men kom tilbage som den Grubler, jeg var, da jeg drog ud, og Soldat
blev jeg ikke. Det tyske Folks Tilstand nødte mig til Betragtninger
af forskellig Art; men Statsmand blev jeg ikke. Skal
jeg nu sige: Dit Liv var aldeles uden Betydning; Alt, hvad Du
vilde udføre, var kun en Række mislykkede Forsøg?«

Hans Svar er: »Mit Liv var et helt indre Liv. Jeg dukkede
uden Betænkning under i Tidens vilde Bevægelser; jeg
svømmede i forskellige Retninger paa Forskningens og Handlin%
% --- p. 163 ---
\setcounter{footnote}{0}%
\opage{163}gens svulmende Vover: men stadigt kun for at vende tilbage
til mig selv og hengive mig til indre Betragtning i større Sammenhæng.
Deraf udsprang, hvad jeg har gjort. Og har min Gerning
da været saa ganske uden Frugt? \textit{Har} \textit{den ikke til} meget
\textit{forskellige Tider sat de Unges Sind} i Bevægelse?\dots{} Har ikke
den Sæd, jeg har saat i fremmed Jord, ført til betydningsfuldere
Vækst, end min egen Jord kunde frembringe?«

Steffens har her selv udtalt det, hvorpaa hans Betydning i
Aandslivets Historie beror, og hvorpaa ganske særligt hans Betydning
for det danske Aandsliv beror. Han kendes her hos os
kun gennem den øjeblikkelige, men vidt rækkende Bevægelse,
han vakte hos dem, der for hundrede Aar siden vare unge. Deres
Sind blev bevæget, saa de aldrig glemte det, og saa det gennem
deres Taknemmelighed og gennem deres Værk holdtes i levende
Minde hos vort Folk og staar som en af de navnkundigste Episoder
i dets Liv. Vækkelsen blev dyb og varig, og det er maaske
den, vi skylder, at vi endnu er et Folk. Men den, der bragte
Vækkelsen, dvælede her kun flygtigt. Og vor Litteraturhistorie
omtaler ham ogsaa kun for dette flygtige Besøgs Skyld. De
indholdsrige Dage, da vor største Digter ved hans Hjælp opdagede
sin egen indre Rigdom, have vi faaet klart og levende
beskrevne, og de staa for os i en Glans, der stadigt paany tiltrækker
os. Men naar vi nu fejre Hundredaarsdagen for denne
Begivenhed, ville vi vise vor Taknemmelighed overfor den, der
fremkaldte den, ved at vise, at det ikke var nogen Tilfældighed
for ham selv, at han kom til at gribe ind paa den Maade, han
gjorde. Den Vækkelse, han bragte, hang --- som vi allerede have
hørt ham selv udtale --- nøje sammen med hans Karakter og
med hans Livsgang.

Da Steffens begyndte sine Forelæsninger her den 11. November
1802, var han 29 Aar gammel. Han er født i Stavanger
den 2. Maj 1773. Hans Fader var en Holstener, hans Moder
var dansk. Saaledes tilhørte han fra først af tre Lande, og
han elskede dem hele sit Liv igennem inderligt. Danmark kaldte
han dog stedse sit Fædreland. Og hans Barndoms- og Ungdomsminder
ere knyttede til det. Han gik i Skole i Helsingør
og Roskilde og blev Student i København. Det var Naturstudiet, 
% --- p. 164 ---
\setcounter{footnote}{0}%
\opage{164}der især tiltrak ham, men han havde tidligt fordybet sig i Goethe og Spinoza, og den nye Filosofi, der udviklede sig i Tyskland
under Indflydelse af Romantikken, vakte snart hans Opmærksomhed. Fra en naturvidenskabelig Studierejse i Norge, som ikke bragte det Udbytte, der ventedes, drog han pludselig til Tyskland, og efter en haard Kamp med Fattigdom og
Modgang blev han Docent i Naturvidenskab i Kiel. Snart sendte den danske Regering ham paa en Rejse til Udlandet i naturvidenskabeligt Øjemed. Det var paa denne Rejse, at han helt blev optagen af Romantikens Verden. Der hviler en herlig og eventyrlig Glans over det Bind af Selvbiografien, i hvilket Steffens skildrer sit Ophold i Jena, Weimar og Freiberg. Snart fordybede han sig i Naturen paa ensomme Bjergvandringer, snart studerede han Geologi hos Werner i Freiberg, snart filosoferede han med Schelling i Jena, snart var han Goethes Gæst i Weimar, snart fik han dybe, religiøst stemte Indtryk af Novalis' Personlighed og Poesi. Iagttagelse, Tænkning, Fantasi og Stemningsliv virkede sammen til en overstrømmende Helhedsfylde, hvis enkelte Elementer han vanskeligt kunde holde ud fra hverandre. Der var frembragt en Klangfarve i hans Sind, og naar han gav det, der rørte sig hos ham, Ord i begejstret Tale, maatte hos Tilhørerne mange forskellige Strenge sættes i Bevægelse. Det var ved denne samlede Kraft, han virkede, paa én Gang Tænkerens, Digterens, Naturforskerens og Profetens Kraft. Han har selv erkendt, at det ikke var et barbarisk Land, han kom tilbage til. Den Begejstring, Kampen den 2. April Aaret forud havde vakt, havde ogsaa grebet ham i hans Ensomhed i Udlandet og vakt hans Stolthed ved at være Dansk. Og han fandt i Ørstedernes og Mynsternes Kreds en selvstændig og forstaaende Interesse for, hvad der bevægede sig i Fædrelands Aandsliv. Men til at faa den unge Slægt vakt til Aandsarbejde krævedes der en koncentreret Energi som han mødte med, omstraalet som han var af Glansen fra Romantikens Hjemland.

Kun svagt fornemmes denne Koncentration, Glans og Begejstring i de trykte Forelæsninger. Det halvtyske Sprog og mange Kunstudtryk skade Virkningen. Vi maa stadigt
% --- p. 165 ---
\setcounter{footnote}{0}%
\opage{165}holde, at Steffens hørte til dem, hos hvem det mundtlige Ord staar
langt over det skrevne, og Personligheden langt over Værket.

Steffens har selv et Sted betegnet Længsel efter det Uendelige
som det Ejendommelige for den romantiske Aand. Intet
Enkelt og Endeligt kunde tilfredsstille den. Den iagttagende og
kritiske Forsken stanser ved Enkelthederne, sammenligner og beregner
dem. Men deraf kommer intet Hele ud, Intet, som helt
kan fylde og tilfredsstille. En Helhedsanskuelse er ikke Resultat,
men forudsætter et oprindeligt Organ og kommer som en
lykkelig Gave.

Dette, at det Højeste i Aandens Verden stod aaben for Enhver,
hvis Sans var aaben, har især været det, der gjorde Indtryk
paa de to betydeligste unge Mænd, der hørte Steffens' Forelæsninger,
Oehlenschläger og Grundtvig. Det maatte vække deres
Haab og Tillid at høre dette udtalt af en Mand, der selv havde
set sig om i Aandens Verden. Han vidnede om, at Poesiens og Religionens Kilder ikke vare hentørrede, men kun hæmmede af
Tidens Prosa og Udvorteshed. Og han hævdede, at Poesi og
Religion udsprang af samme indre Væld som Tanken og Forskningen,
saa snart disse ikke begrænses paa smaalig Vis. Thi det
skyldes --- lærte han fremdeles --- kun den kritiske Forstands
Overherredømme i nyere Tid, at de forskellige aandelige Magter --- Poesi, Religion, Videnskab, Handling --- staa som forskellige,
og som indbyrdes stridende. Der har været lykkeligere
Tider, hvor disse Magter virkede inderligt forenede. Det var den
store Fyldes Tid, som den endnu fremtræder for os i Oldtidens
udelte Værker. »Livet selv, frisk som det udsprang fra Guddommen,
havde umiddelbart den Betydning, som vi nu kun ane. Intet
Spor til den adsplittende Død, som i nyere Tider har gjort Videnskaber
og Poesi, og ved dem selv Religionen, monstrøse«.
Nu kan denne herlige Tid fornyes --- selvom foreløbigt det
Skønneste og Herligste, der gennemtrænger vor Tid, er Sansen
for hin gamle, forsvundne Tid!

En saadan Helhedsanskuelse, der viser ud over alt Enkelt
Enkelt og Isoleret, rører sig --- det anerkender Steffens --- i enhver
virkelig Forsker. Det er den, der giver Naturforskeren Mod til
at efterspore dunkle Sammenhænge og Historikeren til at for%
% --- p. 166 ---
\setcounter{footnote}{0}%
\opage{166}dybe sig i fjerne Tider. Men den maa gennemføres endnu
videre end de gøre. Vi maa betragte Naturen som Aandens Forhistorie.
--- Og her benytter Steffens saa sine Aaret forud udkomne
»Beiträge zur inneren Naturgeschichte der Erde« til at
vise, hvorledes der i hele Naturen ytrer sig en individualiserende
Tendens, idet allerede den uorganiske Natur viser Antydninger,
som først i Plante-, Dyre- og især i Menneskelivet komme i
fuld Udvikling. Hvad der fremtræder paa det aandelige Omraade,
ligger allerede som dunkelt Anlæg i den ubevidste Natur.
Derfor maa vi i vor egen Aands Inderste søge Nøglen til
Naturens Hemmeligheder. Gaaderne løses ikke ved blot Ophoben
af Fakta. --- Og som vi betragte Naturen som Aandens
Forhistorie, saaledes finde vi i Menneskets Historie Skæbnens
store, naturnødvendige Gang. Ingen enkelt Tid maa gøres til
Maalestok for andre Tider; det gælder at fordybe sig i den
indre Sammenhæng igennem Tiderne, den, der ikke blot træder
frem i de store Genier, i hvem Udviklingen paa visse Punkter
koncentrerer sig, som om de vare det Maal, der skulde naas,
--- men aabenbarer sig selv i det ringeste, mest begrænsede
Menneske. der ligesaa vel som hine er bestemt ved den historiske
Udvikling. ---

Steffens' Filosofi er --- som hele den Tids saakaldte »Naturfilosofi«
--- bygget paa en dristig Analogisering, hvor en mere
eller mindre dyb Poesi som oftest fører Ordet paa Tankens Vegne.
Hvad der adskiller ham fra denne Retnings egentlige Hoved,
ham, Steffens kaldte sin Mester, Schelling, er især den Vægt,
han lægger paa det Individuelle i Natur og Menneskeliv, og paa
den historiske Udviklingsgang. Begge Dele hænge igen sammen
med, at virkelig Naturiagttagelse spiller en større Rolle hos ham
end hos de andre spekulative Filosofer. Paa den anden Side
var det denne Sans for det Organiske, Individuelle og Historiske,
der gjorde ham fremmed for den exakte, realistiske Naturforsken;
denne forekom ham at opfatte Alt som dødt, alle Processer som
udvortes mekaniske. Naar han lagde saa stor Vægt paa Udviklingen,
betyder det derfor heller ikke, at han var Evolutionist
i moderne Forstand; thi han antog ikke, at Udviklingen formedelst
bestemte reale Aarsager arbejdede sig frem fra Trin til
% --- p. 167 ---
\setcounter{footnote}{0}%
\opage{167}Trin; det var for ham Naturens Aand, der gik fra Forsøg til
Forsøg, stedse stræbende op imod en højeste Individualitetsform.

Det er karakteristisk, at han, hvis Personlighed staar over
hans Værk, stræbte at yde en Personlighedsfilosofi ved at opfatte
Tilværelsens inderste Stræben som en Stræben efter Individualitet.
I sit Levned har han en i denne Henseende interessant Udtalelse:
»Det var Anskuelsen af det altomfattende Liv, som beherskede
min Ungdom, ja allerede min Barndom. Hvad en dunkel
Anelse og Længsel søgte, var ikke et rent ydre Forhold til Naturen,
men en indre, aandelig Enhed i Alt; derfor var det Ringeste
mig saa kært; derfor var det mindste Græsstraa i dets bestemte
Form mig saa værdifuldt og stod i et indre, saa at sige personligt
Forhold til det ubetydeligste Insekt. Dunkelt foresvævede
mig dette ved mine Studier, og idet jeg optog og delte fremmede
Anskuelser, vendte jeg dog stedse tilbage til det, der
ganske vist længe var fri Fantasi og havde mere digterisk end
videnskabelig Betydning. Hvad Schelling blev for mig, er bekendt;
ja, den Drift, der førte mig til Schelling, og som gav
mit Liv dets Betydning, er et Hovedtema for mine Livserindringer.
Men for mig er det, jeg kalder Naturfilosofi, ikke Andet
end den Overbevisning, at der vil udvikle sig en konsekvent
organisk Opfattelse, som i enhver Genstand for Forsken erkender
et ejendommeligt Noget, der udvikler sig paa sin egen Vis,
--- en Opfattelse, som først ved saadan Anerkendelse formaar
at fatte det Enkeltes Betydning for det Hele.« ---

Det faldt ikke i Steffens' Lod at udvikle disse Tanker
videre paa dansk Grund. Han gik tilbage til Tyskland og blev
Professor i Filosofi og Naturvidenskab (særligt Mineralogi) i
Halle. Her har Samlivet med Schleiermacher været frugtbart
for ham og har bidraget til at udvikle det individualistiske Element
i hans Filosofi. Indledningen til »Grundzüge der Naturwissenschaft«
(1806) slutter med en begejstret Lovprisning af Schleiermacher,
og dette Skrift er ogsaa præget af hans Indflydelse.

I sit nye Fædreland levede han sig dybt ind. Han tog
begejstret Del i dets nye Tanke- og Kunstliv; han led med det
i Nederlagenes og Ydmygelsernes Dage. Og det faldt i hans
Lod --- ved den vidunderlige Evne, han havde til at bevæge de
% --- p. 168 ---
\setcounter{footnote}{0}%
\opage{168}Unges Sind --- at tale det første Ord, der kaldte Tysklands
Ungdom til Frihedskampen. Han var Professor i Breslau, da
de tyske Hære maatte deltage i Napoleons Felttog i Rusland ---
og han følte Nedværdigelsen herved. Nu kom Rygterne om Katastrofen
i Rusland; Napoleon var flygtet gennem Schlesien
et preussisk Armekorps havde paa egen Haand sluttet sig til
Rusland, og det preussiske Hof kom nu til Breslau, hvor de
følte sig mere sikkert end i Berlin, der havde fransk Besætning. Kongen
var længe vaklende, men man havde bevæget ham til
at underskrive et Dekret, der kaldte Ungdommen til Fanerne---
Aftenen før det skulde offentliggøres, var Steffens sammen med
nogle unge Officerer af Scharnhorst's Skole. Man frygtede for,
at Dekretet paa Grund af sine ubestemte Udtryk ikke vilde
gøre den rette Virkning. Da fattede Steffens i Løbet af Natten
den Beslutning at gribe ind. Han bekendtgjorde paa sin
Morgenforelæsning, at han senere paa Dagen vilde holde en
Tale om det kongelige Opraab, der vilde blive udstedt. Rygtet
herom udbredtes, og en stor Skare fyldte Høresalen, Trappen
og Gaden udenfor. Steffens erklærer, at han ikke kan gøre
Rede for, hvad han sagde; det var, hvad der bevægede Alles
Tanker, men nu første Gang offentligt fik Luft. Han endte med
at opfordre de Unge til at melde sig og erklærede, at han
ogsaa selv vilde være Soldat. Dagen efter maatte han gentage
Talen for en endnu større Kreds. Der var nu vakt en Bevægelse,
som opfyldte de tyske Patrioters Ønske og Længsel.
Scharnhorst omfavnede ham og raabte: »Steffens, jeg lykønsker
Dem; De véd ikke, hvad De har gjort!« --- Steffens tog som
Løjtnant og Adjutant Del i Krigen.

Da han efter Felttogets Uro igen var vendt tilbage til sin
stille Gerning i Breslau, varede det ikke længe, før hans Ungdoms
Ideer kom til at staa i et noget andet Lys for ham end
før. Han opgav dem ikke, men han saa, at Vejen til deres
Gennemførelse var længere, end han i Ungdommens Begejstring
og Overmod havde tænkt. Det stod nu for ham (siger han i
et Brev til Tieck), som om det var et Babelstaarn, han havde
været med at bygge paa. Arbejdsdelingen paa det videnskabelige,
og idethele paa det aandelige Omraade, gjorde sig stær%
% --- p. 169 ---
\setcounter{footnote}{0}%
% CHECK p. 169: indent after a broken word?
\opage{169}kere og stærkere gældende; de forskellige Fakulteter og Fag
isolerede sig og gik hvert sine Veje, saa at den Enhed i Tankens
og Aandens Verden, han havde ment var ved at komme,
stod som et fjernt Ideal. »Jeg saa' kun altfor klart,« siger
han, »at ikke Begejstringen ved Aarhundredets Begyndelse,
men kun Historiens langsomme Gang efterhaanden skulde hidføre,
hvad det raske, flammende Liv allerede tror at besidde i
fuld Modenhed. Men dette forsvundne Haab havde forvandlet
sig til en religiøs Tro, der ikke vaklede, thi den støttede sig
ikke til egen Kraft, men til en guddommelig Kraft. Saaledes
kunde jeg opleve, at Alt mislykkedes mellem Hænderne paa mig,
uden at jeg dog blev tvivlraadig«. »Al Udvikling er stille. Du
ser ikke Græsset voxe, men det kan vel lykkes den roligt betragtende
Aand at komme i stedse større Nærhed af Udviklingens
tavse Gang.«

Denne milde Resignation er saa meget desto mærkeligere,
som Steffens’ Livsanskuelse ellers i de senere Aar fik et strengere
Præg og skarpt betonede Elementernes Modstrid i den
Enkeltes og i Slægtens Liv. I »Caricaturen des Heiligsten«
finder. han Syndigheden at bestaa i det enkelte Elements Isolation
fra den Helhed, til hvilken det hører, og en saadan Isolation
fremtræder for ham paa betænkelig Maade i den moderne
Individualisme og Liberalisme. Men Karikatur er for ham
ogsaa enhver ensidig videnskabelig Anskuelse, og det forekom
ham at være Menneskets Skæbne snart at falde under den ene,
snart under den anden Tankes Herredømme og derved miste
Blikket for den sande Ejendommelighed, der kun kan hævdes
indenfor en omfattende Helhed. Han tog endnu bestemtere og
skarpere Afstand fra Rationalismen, end han havde gjort i sin
Ungdoms romantiske Dage. Hans religiøse Tro fik en udpræget
konfessionel Karakter, idet han (i Skriftet »Von der falschen
Theologie und dem wahren Glauben«) sluttede sig til den strenge
Lutheranisme og dens Modstand mod den Union med de Reformerte,
Regeringen med Magt vilde føre igennem.

Men sin Filosofi ændrede han ikke væsentligt. I hans
»Anthropologie« og hans »Religionsphilosophie« findes de samme
Grundtanker som i hans Ungdomsskrifter (kun i udpræget teo%
% --- p. 170 ---
\setcounter{footnote}{0}%
\opage{170}logisk Form), de samme Forsøg paa at finde Antydninger af Menneskelivets Kræfter og Kampe paa de lavere Naturtrin, at læse Naturen, saa at sige ovenfra nedad. Romantikkens Tanker vare snarere poetiske end filosofiske, og Steffens har selv ogsaa udtalt, at et Digt i Lighed med Dantes stadigt arbejdede sig frem i hans Indre, uden nogensinde at vinde Form. Som en indre Melodi bevaredes op i den høje Alderdom den Anskuelse, han i sin Ungdom havde troet at kunne forme til System. Og naar han saa' Videnskaben, som han havde troet at kunne opløse i Poesi, vende tilbage til det strenge Arbejdes Prosa, fulgte han ogsaa denne Udvikling med Interesse (som det f. Ex. ogsaa ses af hans Deltagelse i det første nordiske
Naturforskermøde 1840), maaske fordi han --- ligesom Hans Christian Ørsted --- var vis paa, at denne ærlige Prosa vilde kunne føre til en dybere Poesi end den, man havde troet at vinde ved at sky den.

Sine sidste Aar levede han som Professor ved Universitetet i Berlin. Da han en Aften hos en Ven havde hørt en ny Danteoversættelse oplæse, fik han pludseligt en Blodstyrtning og døde faa Dage efter (13. Februar 1845), 72 Aar gammel.

I sin Mindetale over ham sagde Schelling: »Er starb in seiner Jugend!« --- Et sandt Ord, thi han bevarede sin Ungdom gennem Troen paa, at Ungdomsidealerne vilde arbejde sig frem ad andre Veje, end han havde tænkt sig fra først af. Man siger, at den, Guderne elske, dør ung. Hvis Steffens var død lige efter sin Optræden i København, vilde han maaske
staa i et endnu mere romantisk Lys for os. Men mon det ikke Guderne elske den endnu mere, hvem det forundes at bevare Ungdommens Sind usvækket gennem et langt Livs Kampe og Skuffelser? ---

Vi kunne da, efter at have fulgt ham til hans Død, atter vende tilbage til hans Optræden i hans Fædreland som den unge romantiske Budbringer. Hvad han har virket her, har han virket ud af sin hele Personlighed, og det var kun et Led af en Livsgang, der bares af en ideal Tro fra først til sidst. Det var da intet Tilfælde, at han virkede her, som han gjorde. Og det var heller intet Tilfælde, at han gav \textit{os} af det Bedste, han havde%
% --- p. 171 ---
\setcounter{footnote}{0}%
\opage{171}at yde. Thi Kærligheden til hans Fædreland og Længslen efter
det strakte sig, som hans Livserindringer vise, gennem hele hans
Liv. Saa ville vi da ogsaa paany sande, hvad Grundtvig sang
ved hans Død, at han har

% CHECK p. 171: centred line: heading or verse

\begin{verse}
talt et Ord, som aldrig glemmes, \\
vakt en Aand, hvis Kraft fornemmes, \\
mens vi har et Modersmaal!
\end{verse}

% --- p. 172 ---
\setcounter{footnote}{0}%
% CHECK p. 172: centred line: heading or verse
\opage{172}\essay[Henrik Steffens i det danske Aandsliv. (1902)]{HENRIK STEFFENS I DET DANSKE\\ AANDSLIV.}
\begin{center}{\small TALE VED KØBENHAVNS UNIVERSITETS REFORMATIONSFEST}\\
{\small DEN 20. NOVEMBER 1902.}\end{center}
\streg

Det er i denne Maaned hundrede Aar siden Steffens begyndte
 den Række af Forelæsninger her ved Universitetet, som
skulde gribe saa stærkt ind i det danske Aandslivs Udvikling.
Mindet herom er allerede blevet fejret af en Kreds af Studerende, der særligt ere knyttede til det Lokale, hvor Steffens
holdt sine Forelæsninger. Jeg havde den Ære at føre Ordet
ved denne Lejlighed og søgte at vise, at den pludselige og vækkende
Indflydelse, Steffens øvede herhjemme under sit kortvarige
Ophold, hang nøje sammen med hele hans Personlighed,
saaledes som denne fremtræder for os i hans Livsgang. Paa
Askov Folkehøjskole har man ogsaa fejret dette Minde; der er
en historisk Sammenhæng mellem Steffens' Optræden og Folkehøjskolebevægelsen.
 Jeg skal nu forsøge noget nærmere at paa-
pege, hvorledes hans Optræden har haft Virkninger paa Aandslivets
forskellige Omraader herhjemme.

Henrik Steffens kom hjem efter et toaarigt Ophold i Tyskland,
hvor han havde færdedes mellem de største Aander i Ti
den og var bleven indviet i den romantiske Skoles Tankekreds.
Naturforsken og spekulativ Filosofi, Poesi og Religion --- alle
disse forskellige Strenge klang sammen i hans Forkyndelse. Han
% --- p. 173 ---
\setcounter{footnote}{0}%
\opage{173}de netop vække en Helhedsanskuelse, i hvilken det, den Tids
Oplysning og Prosa var vant til at sondre, virkede sammen.
Her vilde han vise en indre Sammenhæng i Naturen, saaledes
at Naturudviklingen kom til at staa som Aandslivets symbolske
rhistorie. Nøglen til Forstaaelse af Naturen skulde søges i
Aandens inderste Dyb. Tillige pegede han hen paa de ideale
Magter, der aabenbarede sig i Historien, i Kunsten og i Religionen.
Om Tilhørerne end ikke forstode Alt i hans Fremstilling,
saa følte de dog, at han var begejstret for det Store, og
det gik op for dem, at der var aandelige Skatte i dem og om
dem, for hvilke de hidtil havde været blinde.

Paa tre Mænd herhjemme har Steffens især virket dels
strax ved disse Forelæsninger, dels senere.

Mest bekendt er hans Indflydelse paa Oehlenschläger, hvem
han --- som han selv har sagt --- vakte til Bevidsthed om sin
egen aandelige Rigdom. Det var den danske Digtnings Guldalder,
som derved indviedes, og det var ikke blot den store
Digter alene, det var selve vort Folk, som derved kom til Bevidsthed
om sin aandelige Rigdom. Forholdet mellem Steffens
og Oehlenschläger er saa udførligt, tildels saa genialt skildret af
andre Literaturhistorikere, at jeg tør forudsætte det tilstrækkeligt
kendt.

Hos en anden ung Mand, der hørte Steffens, var det vel
saa Poesien, der først og fremmest blev vakt --- men den betragtedes
dog tilsidst som Indledning til en endnu dyberegaaende
Bevægelse. Grundtvig var ikke lidet betænkelig ved den romantiske
Filosofi, Steffens forkyndte. Han saa i den »en skæv
en stærk Modvirkning af det Poetiske i Mennesket mod den
t opløsende Forstandsbrug«. Han glædede sig over, at Steffens
talte »med en forbausende Ærbødighed« om Kristus og
»å haant om den Oplysning, der hidtil var udbasuneret«. Men
han var dog tillige glad over, at »kun Faa forstode noget Sammenhængende
af hans med mange Kunstord udstafferede Tale«,
--- ti »der stak megen Ugudelighed i hans Lære«. Selv senere,
da Steffens i religiøs Henseende indtog et afgjort konfessionelt
Standpunkt, udtalte Grundtvig efter en Sammenkomst med ham
(under hans Besøg her 1824) sin Sorg over, at de vare »himmel%
% --- p. 174 ---
\setcounter{footnote}{0}%
\opage{174}vidt adskilte«. Men trods Alt holdt han fast ved, at hans egen
»historisk-poetiske Virksomhed« i Grunden kun var en Fortsættelse
af Steffens' Optræden her ved Aarhundredets Begyndelse.


Naar man tænker paa, hvilken Betydning det, Grundtvig
her kalder sin historisk-poetiske Virksomhed, har havt og har
for store Kredse i vort Folk --- hvorledes hele Folkehøjskolebevægelsen
er en Frugt af den --- er denne Ytring af Grundtvig
(som forekommer i et Brev til Ingemann af 26. November
1824) af stor Interesse. Det er under Romantikens Fane, Folke
højskolen har arbejdet, selvom den med stor Dygtighed paa sin
Vej har optaget Elementer af anden Art. Den største folke
opdragende Bevægelse i Danmark i det 19. Aarhundrede kan
saaledes føres tilbage til Steffens' Forelæsninger som første Kilde.

Det kan ikke negtes, og jeg ønsker særligt at udtale det fra
dette Sted, at man i Universitetskredse for længe har miskendt
denne Bevægelses Betydning. Det er naturligt, at Videnskabens
Mænd ere betænkelige og kritiske overfor en Popularisering af
Videnskaben, da derved de enkelte Kendsgerninger og Lærdomme
let rives ud af hele den videnskabelige Sammenhæng,
uden hvilken de ikke kunne ses i deres Betydning og i deres
Begrænsning. Men Livet har sin Ret, og Livet tilegner sig Videnskabens
Tanker paa sin egen Vis, ud fra Synspunkter og
Formaal, som Livet selv har ført til at anlægge. Kun saa kan
Tilegnelsen blive levende og frugtbar. Og her har den med Romantiken
indledede »historisk-poetiske« Opfattelse havt sin store
Betydning. Men saa maa det heller ikke glemmes, at Livet
under sin fortsatte Udvikling kan og maa hidføre nye Synspunkter
og Formaal, om hvilke Aandens og Tankens Skatte
kunne samles, og i hvis Lys de kunne tydes. Nye Tider, nye
Tanker! Ogsaa Steffens kom jo som Budbringer om en ny Tid
og om nye Tanker. Sligt kan ikke gives en Gang for alle.
Kommer der Strid mellem de forskellige Synsmaader, saa er det
kun til Gavn, naar Striden føres paa en Maade, der er Emnet
værdig uden Mistænkeliggørelse og uden Fanatisme. Det tjener
til vort Folks Ære, at det er vakt til en saadan Strid, saa at
% --- p. 175 ---
\setcounter{footnote}{0}%
\opage{175}det ikke blot er i de Lærdes Verden, at der kæmpes om Livets
store Spørgsmaal.

Universitetet har ikke den Opgave, umiddelbart at bringe
Sandhederne ud i Folkets store Kredse. Men derfor har man 
ikke Ret til at stille den videnskabelige Højskole i Modsætning
til Folket. Den videnskabelige Forskning øves paa hele Folkets
Vegne. Saa ensom den enkelte Forsker end ofte kan føle sig,
og saa uforstaaeligt det, han finder, ofte foreløbigt kan være for
alle Andre, saa vil han dog kunne betragte sig som en Udkigspost,
der skal iagttage og undersøge, hvad der kan erfares paa
hans Stade, og melde, hvad de Andre ikke kunne se paa første
Haand. Der er ingen absolut Modsætning mellem Tankens Liv
og de andre Former for menneskeligt Liv. ---

Medens det var ved det Romantiske i sin Personlighed og
sin Lære, at Steffens virkede paa Grundtvig, bestemtes hans Indflydelse
paa F. C. Sibbern mest ved de Sider af hans Personlighed
og hans Lære, ved hvilke han afveg fra de andre Ordførere
for Romantiken. Og gennem Sibbern øver Steffens en
betydningsfuld Indflydelse paa den danske Filosofi i det 19.
Aarhundrede.

Dansk Tænkning har ligefra Holbergs Tid havt en empirisk
og kritisk Karakter. Psykologisk Iagttagelse og Analyse, Drøftelse
af etiske Spørgsmaal og af Erkendelsens Betingelser og
Grænser have ligget de danske Filosofer nærmere end rene
Spekulationer. Hvad der udmærkede Steffens i Sammenligning
med de andre Filosofer af den romantiske Skole, var nu netop
hans stærke Fremhæven af det Individuelles Betydning i Naturen
og Menneskelivet, hans Fordring om levende Naturanskuelse
som Grundlag for den filosofiske Tænkning og hans
Utilbøjelighed til abstrakte Spekulationer som dem, Schelling
og Hegel svælgede i. I denne Aand talte han allerede 1802,
og i denne Aand talte han, da han sidste Gang offentligt havde
Ordet herhjemme, nemlig paa Naturforskermødet 1840. Og
det var ogsaa i denne Aand, Sibbern virkede hele sit Liv igennem.
Det er i Æstetiken og i Teologien, hos Heiberg og hos
Martensen, at den spekulative Filosofi har havt sit egentlige
Tilhold herhjemme. Den Udviklingslinie, der gaar fra Sibbern
% --- p. 176 ---
\setcounter{footnote}{0}%
\opage{176}gennem Poul Møller til Kierkegaard og derfra videre, betegner
en anden Retning, og utvivlsomt en mere dansk Retning.

Sibbern var bleven Student en Maaned før Steffens begyndte
sine Forelæsninger. Han skal have forsøgt at høre dem,
men fik ikke Noget ud af dem. Senere hen er Steffens kommen
til at staa i et straalende Lys for ham, uden at det vides,
hvad der har fremkaldt det. Han higede efter at komme i personligt
Forhold til ham. Det var (som han skriver til A. S.
Ørsted den 1. Oktober 1811) kun derfor, han strax i Begyndelsen af sin lange Udenlandsrejse drog til Halle, hvor han
saa ogsaa, efter sin egen Udtalelse, fik sit Ønske opfyldt, og
det »paa en skønnere Maade, end han nogensinde turde haabe«.
Sibbern var paa denne Tid i stærk Gæring, dels ved sin Stræ
ben efter at vinde fast Standpunkt i Tankens Verden, dels ved
en haabløs Lidenskab (den, han har skildret i første Del af
Gabrielis' Breve). Ved selve sin Personlighed har Steffens været
overordentligt Meget for ham: i et Brev taler han om det
Hjertekvægende og Lærerige i Omgangen med Steffens. Han har
ved Omgangen med ham ikke blot følt sig fremmet i sit Tanke
arbejde, men ogsaa beroliget og forsonet. Sit sidste Besøg hos
Steffens sammenligner han med Orests Besøg i Gudindens hel
lige Lund. Hans Sorg og Lidenskab er her bleven forvandlet
til den frejdige Entusiasme, der var den fortræffelige Mands
fremherskende Stemning Livet igennem, selvom han atter og 
atter maatte tilkæmpe sig den paa Trods af Mismod og Ilterhed. Begejstringen for Steffens bevarede han bestandigt senere
den havde, efter Poul Møllers Udsagn, Karakteren af en Forelskelse, og han var ellevild af Glæde, naar Steffens senere be
søgte Danmark. 

Paa sin Rejse besøgte Sibbern ogsaa Schelling og var
meget begejstret over ham. Men hans eget filosofiske Arbejde
minder mere om Steffens (saavel som om den i filosofisk Henseende
Steffens nærstaaende Schleiermacher). Han havde for
stor Respekt for Erfaringsvidenskaben til, at han kunde tiltales af den romantiske Naturfilosofi. Saaledes skriver han i et Brev 
1824: »Det er raadeligt for enhver Naturbetragter at holde sig i Erfaringen saa nær som muligt og at være varsom med ethvert
% --- p. 177 ---
\setcounter{footnote}{0}%
\opage{177}Forsøg paa at ville oplyse og udtyde Naturtingene ved en Betragtning
\textit{fra oven}; ti dertil ere Naturbetragtningerne endnu
langt fra at være modnede.« Hvad Sibbern mener med en
Udtydning »fra oven«, kan kun være Romantikens Forsøg paa
at læse det aandelige Livs Forhistorie ud af Naturen i Stedet
for at støtte sig til selve Lovene for Naturfænomenernes Opstaaen.
Sibbern selv manglede ikke Interesse for omfattende
Naturanskuelse. I hans Filosofi spiller en saadan en afgørende
Rolle. Men det er karakteristisk, at han endog i en vis Forstand
anlægger et Synspunkt, der er Romantikens modsat. Han
søger netop at belyse den aandelige Udvikling ved Analogier,
der ere hentede fra Naturens Udvikling. Og hvad selve Naturudviklingen
angaar, lægger han stor Vægt paa, at den foregaar
ud fra spredte Udgangspunkter; i Krystallisationen og i
den organiske Væxt finder han Exempler herpaa. Men hvad
der saaledes begynder fra sondrede Udgangspunkter, forbindes
ikke altid saa let til et harmonisk Hele. Det viser sig især,
hvor disse Udgangspunkter ere særskilte Individer. Der staar
i en Kamp, før Maalet naas. Sibbern har saaledes i denne
sin Lære om den sporadiske Udvikling givet Steffens' Lære
om Naturens individualiserende Tendens en realistisk Videreførelse.\footnote[1]{Det kunde have ligget Sibbern nær at henvise til Statsøkonomen Chr. Olufsens Afhandling »Om Menneskets Rolle i den fysiske Verden« (1822), som Professor Joh. Steenstrup for nogle Aar siden fremdrog af Forglemmelse, og som indeholder en Ide, der paa en interessant Maade staar som Forgænger for Darwins Lære om Udviklingen gennem Kampen for Tilværelsen. Olufsens Lære gaar ud paa, at »Livskraften« i Naturen paa Grundlag af de Lag og Former, hvis Tid er forbi, arbejder sig frem til nye Former, som bedre passe til de forandrede Betingelser. Det synes efter Olufsens Fremstilling ikke at være nødvendigt, at der er organisk Kontinuitet mellem Formerne; Kontinuiteten er i selve »Livskraften«. Hans Lære er en vitalistisk Udviklingslære, der staar imellem den romantiske Naturfilosof og Darwins Hypotese.}
Endnu 1843 citerer han (i sit Skrift om Filosofiens
Begreb) Steffens' Sætning om, at al videnskabelig Spekulation
maa føre til »Anskuelse af Individualitet« (Steffens' Grundzüge
der philosophischen Naturwissenschaft, 1806, p. 36, hvor 
Be%
% --- p. 178 ---
\setcounter{footnote}{0}%
\opage{178}grundelsen er, at i Individualiteten ere alle enkelte Elementer
sammenfattede ved en Helhedslov).

Paa en aandfuld Maade anvender Sibbern sin Lære om
Udviklingens sporadiske Karakter paa Erkendelsesteorien. Vi
Mennesker befinde os paa et enkelt af de mange Udviklingspunkter i Tilværelsen og kunne derfor ikke overskue denne i
dens Helhed. Og desuden befinder hver enkelt Tænker sig
paa et bestemt Udgangspunkt indenfor den menneskelige Udvikling.
Intet Under, at der da finder en stadig Debat og Kamp
Sted mellem de forskellige Verdensanskuelser. Denne Kamp
er et Tegn paa, at vort aandelige Liv endnu er i sin Vorden
og stadigt maa søge sin Vej og sin Retning. Dette Synspunkt
maatte gøre ham betænkelig ved forhastet Systemdannelse. Og
det var hans Overbevisning, at en Filosofi virker bedst medens
den selv er i Gæring og i Vorden; naar den er begyndt at
sætte sig, naar dens »Bevægelsesperiode« er forbi, vil dens Be
tydning ogsaa være forbi. Saaledes stiller Sibbern »den rige
Schellingske Foraarsperiode« (hvis Barn jo ogsaa Steffens var)
i Modsætning til Hegels Systematisering.

Det er især paa dette Punkt, at man forstaar hans Indflydelse paa Poul Møller og Kierkegaard. For Poul Møller
blev ethvert System let »Affektation«, fordi der til dets Gennemførelse
 i Regelen maa bruges Begreber, der opstilles for
Symmetriens Skyld, uden at være vundne ved virkelig Er
faring og selvstændigt Tankearbejde. Og Kierkegaards Livsopgave blev det at bekæmpe »Systemet« i alle Former, paa
det praktiske saavel som paa det teoretiske Omraade --- Kirkens Systematisering af det religiøse Livs Inderlighed saavel
som den spekulative Filosofis Systematisering af Tænkningen.

Steffens var nu dog trods Alt, hvad der gav ham en
særegen Betydning, en fuldblods Romantiker i sin Filosofi, og
det ligger da nær til Slutning at paapege Forholdet mellem
Romantikens Naturopfattelse og den, der fastholdes i Nutidens 
Tænkning.

Romantiken kom i Filosofien ikke blot som en Reaktion
mod Oplysningsperiodens Prosa, men tillige som en Reaktion
mod den dogmatiske Opfattelse af den moderne Naturviden%
% --- p. 179 ---
\setcounter{footnote}{0}%
\opage{179}skab, i Følge hvilken Naturens inderste Væsen skulde være udtømt
ved de mekaniske Synspunkter. I Modsætning til den
mekaniske Opfattelse, der fortolker Alt i Naturen som Bevægelse,
vilde Romantikerne fortolke de forskellige Naturprocesser og Naturtrin ved Analogier fra Aandslivet, læse Naturen »ovenfra«
i Stedet for »nedenfra«. Paa Forhaand kan det ene Forsøg være ligesaa berettiget som det andet; men Romantikernes Forsøg lykkedes ikke. Man kan ikke ved Hjælp af psykiske Analogier give en sammenhængende Tydning af Naturfænomenerne,
en Tydning, der gør det muligt at slutte fra det Forbigangne til det Tilkommende. Derimod har den moderne Naturvidenskabs Historie vist, at man ved at opfatte det, der
er i Naturen, som Bevægelser i Rummet, kan formulere strenge og klare Love, som atter og atter staa deres Prøve i Erfaringen. Derfor maatte Romantiken i det 19. Aarhundredes sidste Halvdel vige Pladsen for Realismen.

Dog maa det fastholdes, at Romantikens Filosofi gik ud fra et virkeligt Problem, der staar ved Magt den Dag i Dag,
Problemet om det Aandeliges Plads i Naturen og Sammenhæng
med Naturen. Hvad den ikke saa, var, at dette Problem videnskabeligt kun kan behandles fra to Synspunkter, det erkendelsesteoretiske
og det psykofysiske. Vi kende kun Naturen
gennem Iagttagelse og Tænkning, altsaa gennem aandelige Virkeligheder,
og selve Naturvidenskabens Grundbegreber og Grundforudsætninger
ere vundne og formede ved vor Bevidstheds Aktivitet og efter dens Love. Dette er det erkendelsesteoretiske
Synspunkt. Det psykofysiske Synspunkt anlægges, naar
Fysiologien og Psykologien hver paa sin Side undersøge, hvorvidt og hvorledes aandelige Tilstande hænge sammen med organiske
Tilstande. Ud fra begge Synspunkter drøftes Problemet i vore Dage, og forskellige Videnskaber stræbe hver fra sin Side at belyse det. Men det ligger langt dybere, end Romantiken antog.

Det Mærkelige er nu, at man i de senere Aar kan spore en ny Reaktion mod den mekaniske Naturopfattelse. Dels nye Fænomener, der ikke uden videre kunne forstaas ud fra den, dels Studier over Naturvidenskabens Historie, dels erkendelses%
% --- p. 180 ---
\setcounter{footnote}{0}%
\opage{180}teoretiske Undersøgelser have ført til den klare Indsigt, at Naturvidenskabens
Grundbegreber og Grundsætninger ikke ere
evige Ideer og Sandheder, men Tankeredskaber, Hjælpeforestillinger
og Symboler, ved hvilke det bliver os muligt at bringe Sammenhæng
mellem Iagttagelserne og slutte fra Fortid til Fremtid.
Fordi vi fra visse Grundantagelser kunne drage Slutninger,
der bekræftes ved Iagttagelse, have vi ikke Ret til at paastaa,
at hine Grundantagelser ere absolute Sandheder; den Mulighed
er ikke udelukket, at samme Resultat kunde naas ud
fra andre Grundantagelser. Hvis man ved Realisme forstaar
den dogmatiske Antagelse af de naturvidenskabelige Forudsætninger
som absolute Sandheder, kan man sige, at Realismen er
kommen tilkort ligesom Romantiken. Sandheden ligger dybere,
end Nogen af dem saa. Den sidste Virkelighed er for os det,
der gør, at vi kun forstaa vore Iagttagelser ved Hjælp af visse
bestemte Forudsætninger. Et afsluttende Begreb om denne Virkelighed
ere vi efter Sagens Natur afskaarne fra at danne. Hvad
det kommer an paa, er at danne de rette Arbejdstanker. De
ere ikke Sandheden; men Sandheden er i Nødvendigheden af
at danne dem.

Saa fjernt end Videnskaben ofte kan synes at staa Livet,
saa er der dog et indre Slægtskab mellem dem --- idetmindste
deri, at det begge Steder gælder Arbejdsværdier. Livet behøver
Tanker, der godtgøre deres Værdi ved at føre fremad. Hverken
Livet eller Videnskaben har Brug for Tanker, der kun staa som
Genstand for passiv Skuen, for mystisk Kontemplation.

Denne Opfattelse af Tanken er mindre romantisk end den,
man havde for hundrede Aar siden. Men den indeholder Spiren
til et frejdigt Haab for Tanken, for Livet og for Forholdet
mellem Tanke og Liv. Og den giver os et Synspunkt, hvorfra
vi --- trods Tidernes Forskel --- kunne forstaa Betydningen af
de Tanker, hvormed Henrik Steffens for hundrede Aar siden
gjorde saa dybt Indtryk paa nogle af de Ypperste af dem, der
dengang vare unge.

% --- p. 181 ---
\setcounter{footnote}{0}%
% CHECK p. 181: centred line: heading or verse
\opage{181}\essay[Søren Kierkegaard. En Karakteristik. (1903)]{SØREN KIERKEGAARD.}
\begin{center}{\small EN KARAKTERISTIK.}\\
(1902).\end{center}
\streg

Kierkegaards Betydning for det danske Aandsliv bestaar i
hans Indskærpen af Koncentrationen om det personlige Liv og
dettes Behov i Modsætning til den Opgaaen i Fantasi og Spekulation,
der var ejendommelig for Romantikens Tidsalder. Fra Tyskland
var den store romantiske Bevægelse kommen ind i det
danske Aandsliv og havde vakt stor poetisk og religiøs Begejstring,
kaldt geniale Kræfter frem til Virksomhed, men ogsaa
begunstiget et Hang til at svælge i Fantasiens Verden og i abstrakte
Spekulationer, der ofte førte bort fra den virkelige Verdens
Opgaver og Kampe. Det var mere paa det æstetiske og
det teologiske Omraade end paa det rent filosofiske Omraade,
at Romantiken øvede sin Indflydelse i Danmark. Den danske
Filosofi har lige fra Holbergs Tid havt en afgjort kritisk, etisk
og psykologisk Karakter, der ikke tillod de abstrakte Spekulationer
at brede sig.

Den, der skal modvirke en Tendens, maa helst selv have
erfaret dens Indflydelse. Søren Kierkegaard betegner ganske
vist en Reaktion imod Romantiken, men han staar selv paa Romantikens
Grund, og det Tragiske i hans Stilling ved Slutningen
af hans store Livskamp bestemmes for en væsentlig Del
derved. En stor digterisk Evne, en Hang til at vugge sig i
Stemninger og til at lade Erindring og Fantasi brede sig som
% --- p. 182 ---
\setcounter{footnote}{0}%
\opage{182}Modsætning til Virkelighedens Prosa gjorde ham dybt modtagelig
for den romantiske Poesis Værker, og hans utrættelige Reflexion
tumlede med Tankerne og arbejdede i Problemernes
Tjeneste med en Lidenskab, der ingen Grænser kendte. Hans
ydre Forhold begunstigede i hans Ungdom Interessen for Literatur
og Tænkning. Han var Søn af en velhavende Købmand
der fra Vestjylland var kommen til København og havde arbejdet
sig frem der. Fra ham arvede han en Formue, der stillede 
ham uafhængig, og uden hvilken han ikke vilde have kunnet
udgive sine Skrifter. Han henlevede sine Studenteraar i literære
Sysler uden praktiske Hensyn, vandt 28 Aar gammel Doktorgraden
ved en interessant Afhandling om »Begrebet Ironi
med stadigt Hensyn til Sokrates« (1841) og begyndte snart efter
sin omfattende Forfattervirksomhed.

Men bag Evnen og Trangen til Fantaseren og Reflekteren
laa der en Magt, der holdt igen og bandt Billeder og Tanker
til Jorden. Det var en Sjæl i Vaande, et Tungsind, der syslede
med de mørkeste Muligheder og stedse gjorde ny Kamp og
Selvanstrengelse nødvendig. Det er en Racearv. Den vestjyske
Race har Hang til mørkt Sind, saavel som Evne til Vid og Reflexion.
Det er ikke for Intet Hamlets Egn. Og Søren Kierkegaards
Fader har ikke blot givet sine Sønner denne Hang i Arv,
men ogsaa næret den gennem Opdragelsen og ved at indvie
dem i sine egne Oplevelser. Da han som lille Dreng gik og
vogtede Faar paa den jydske Hede og led ondt af Sult, Kulde
og Ensomhed, steg han en Dag op paa en Høj og forbandede
den Gud, der havde givet ham dette usle Liv. Dette kunde han
selv i sin høje Alder ikke glemme. Det stod for ham, at han
havde begaaet den Synd, der ikke kan tilgives, og at den strenge
Guds Straf vilde falde paa hans Slægt. Under disse gammeltestamentlige
Forestillinger voxede Søren Kierkegaard op; man
ser af hans efterladte Papirer, at de stadigt have sysselsat ham.

Vil man forstaa Søren Kierkegaards Forfattervirksomhed,
maa man fastholde begge de to Kræfter, der rørte sig hos ham:
Trangen til at lade Tanke og Fantasi røre sig frit --- og Trangen
til at kredse om sit eget Selv, tabe sig i det Mørke og

% --- p. 183 ---
\setcounter{footnote}{0}%
\opage{183}Tunge, som var en medfødt Arv og \textit{næredes} ved Slægtens
Tradition.

Kierkegaard har selv paa et Tidspunkt af sit Liv, nemlig
efter Afslutningen af den første, storartede Periode af hans
Forfattervirksomhed (1843---46), den fra hvilken Rækken af de
\textit{æstetiske} og filosofiske \textit{Skrifter} foruden de første rent religiøse
Skrifter stamme, følt Trang til at finde et Enhedssynspunkt for
sin mangfoldigartede Produktion. Han overraskedes ved, at der
i den brogede Række af Skrifter, han havde forfattet i lidenskabelig
og aandeløs Hast, saa at Produktionen ofte var ved at
sprænge hans Kraft, tydeligt fremtraadte et »Totalanlæg«, idet
det Hele kunde synes at være rettet med Bevidsthed mod et
bestemt Maal. Det var for ham en stor Fristelse at antage, at
han oprindeligt havde haft en bevidst Hensigt med den hele
Produktion --- men han havde store Betænkeligheder ved at
tillægge sig en saadan Tanke. I hans Dagbøger fra 1849 kunne
vi se, hvorledes han her har svinget frem og tilbage. --- Hans
Dagbøger ere de Monologer, han holdt, naar han gik op og ned
i sine Værelser, af og til stansende ved Pulten for at nedskrive
en Reflexion. Til Andre talte han aldrig om sine inderste Tanker. I Samtaler forholdt han sig væsentligt spøgende og ironiserende,
--- naar det ikke gjaldt om at trøste, thi deri var han
efter Fleres Vidnesbyrd en Mester; som han var en Mester i
at opdage Sorgens eller den trykkede Stemnings uvilkaarlige
Tegn, forstod han ogsaa som Faa at bringe Lindring og Opmuntring.


Førend jeg gaar nærmere ind paa, hvad der afholdt Kierkegaard
fra at tillægge sig en bevidst Hensigt med sit Forfatterskab,
vil jeg først give en kort Oversigt over det.

Det var \textit{de æstetiske} Skrifter, som først vakte Opsigt. De
vigtigste af dem ere »Enten-Eller« og »Stadier paa Livets Vej«.
De skildre i poetisk Form forskellige Livsanskuelser i deres
Modsætning til hverandre. Ved et »Stadie« forstaar Kierkegaard
ikke en Livsperiode, der i Kraft af en naturlig Udviklingslov
afløses af en anden Livsperiode. Nej, hvert enkelt Standpunkt
skildres som saa bestemt udpræget og afsluttet, at Overgangen
til et andet Standpunkt kun kan ske ved et ubegribeligt Spring,
% --- p. 184 ---
\setcounter{footnote}{0}%
\opage{184}ved et Valg, \textit{der} ikke kan psykologisk begrundes. Kierkegaard
beskriver især tre saadanne Livsanskuelser. --- For den »æstetiske« Livsanskuelse er Nydelsen, Øjeblikkets Nydelse, Alt. Hvad
det gælder om her, er at besidde en saadan Virtuositet eller
Elasticitet, at man kan gribe det enkelte Øjeblik og dets Lyst
i Flugten, uden at lade sig binde og hilde og uden at berøve
sig Evnen til med samme Frihed at gaa op i de følgende Øjeblikke.
Det er en Vilkaarlighedens lunefulde Svæven over Livet
og Livets Forhold, som her er det Karakteristiske. --- For
den »etiske« Livsanskuelse gælder det derimod stadigt Arbejde,
Kamp og Sejr, Trofasthed gennem Tiden, saa at der bliver en
indre Sammenhæng mellem de enkelte Øjeblikke. Ægteskabet
træder her i Modsætning til den momentane Erotik, Kaldsvirksomheden
i Modsætning til den flygtige Beskæftigelse. --- Lever
»Æstetikeren« i Øjeblikket, »Etikeren« i Trofasthed gennem Tiden,
saa lever »den Religiøse« i et disharmonisk Forhold mellem
Tid og Evighed. Hans Maal er Evigheden, i hvilken han
efter sit sande Væsen har hjemme, og han lever derfor i Tiden
som i et med hans Væsen uensartet Medium, --- som en Fisk
paa Landjorden. Jo skarpere Modsætningen er mellem Tid og
Evighed, des mere udpræget er det religiøse Standpunkt. I den
virkeligt kristelige Religiøsitet naas Højdepunktet.

Det vigtigste af de \textit{filosofiske Skrifter} er det, Kierkegaard
har kaldet »Uvidenskabelig Efterskrift«. Her gives en mere
begrebsmæssig Udvikling af Forholdet mellem de forskellige
Livsanskuelser, som i de æstetiske \textit{Skrifter} mere skildredes paa
poetisk-anskuelig Maade og gennem typiske Repræsentanter. Desuden
kritiserer Kierkegaard her skarpt og spottende den spekulative
Teologi og udvikler en Erkendelsesteori, der afgjort er i
den kritiske Filosofis Aand. Vanskelighederne og Grænserne
for en afsluttet Erkendelse paavises. Erfaringen frembyder Modsætninger,
der ingenlunde, som den romantiske Spekulation
mente, let kunne bringes ind under »en højere Enhed«. Og
fremfor Alt: Livet leves i Tiden, midt i Tiden og dens Spænding!
Al Forstaaelse vindes ved at se tilbage, men Livets Opgaver
pege fremad. Vi forstaa baglængs, men vi leve forlængs.
Livet kræver derfor langt mere Tro og Haab, end Videnskaben
% --- p. 185 ---
\setcounter{footnote}{0}%
\opage{185}kan begrunde. Saaledes er det Forholdet til Tiden, der baade
stiller Livsanskuelsen og Tænkningen de største Problemer.

\textit{De} \textit{religiøse Skrifter} begyndte --- hvad Kierkegaard lægger meget stor Vægt paa --- samtidigt med den æstetiske Produktion.
Det er en Række af »opbyggelige Taler«, rettede til
»den Enkelte«, rige paa Stemning, undertiden for rige paa Reflexion. Kierkegaards religiøse Standpunkt betoner fra først af
Inderligheden som det Afgørende; den personlige Følelse og
Lidenskab i det religiøse Forhold er vigtigere end den objektive
Rigtighed af Troens Indhold. En Hedning, siger han, der
beder i Inderlighed til en Afgud, beder i Virkeligheden til den
sande Gud, medens den, der paa udvortes og upersonlig Maade
beder til den sande Gud, i Virkeligheden beder til en Afgud.
Gud kommer man kun i Forhold til gennem den lidenskabelige
Trang, der klamrer sig til ham for ikke at forgaa. Gud er et
Nødværge --- ikke Genstand før Spekulation. Overfor Kristendommens Angribere (især Ludwig Feuerbach) staar Kierkegaard saaledes, at han giver dem Ret overfor Teologien og Kirken,
men ikke overfor »det nye Testamentes Kristendom«, skønt han
mener, at de ofte give et rigtigere Billede af denne end de ortodoxe Teologer. ---

De tre Grupper af \textit{Skrifter} kunne ikke holdes strengt ude fra hverandre. Det er tre Strømme, der her flyde samtidigt
fremad, snart skilte, snart forenede. I Kierkegaards senere Produktion
(efter 1846) høre de to første op at rinde, og der sker
en fuldstændig Koncentration i det Religiøse; det er ogsaa nu
først, at Kierkegaard selv gennemlever en saadan indre Bevægelse,
at han tør sige om sig selv personligt, at nu staar han
»ved Troen i dybeste Forstand«. Men naar han nu kastede et
Blik tilbage paa den overvældende og mangesidige Produktion
fra de foregaaende Aar, stod det for ham, som om den fra først
til sidst var anlagt paa at føre til et eneste Maal: til at føre
Menneskene, Samtidens æstetiserende, spekulerende og verdslige
Mennesker, til Selvforstaaelse, til Klarhed over, hvilket »Stadium«
de tilhørte, og til Beslutning om gennem et kraftigt »Enten-Eller«
at tilhøre det Højeste eller dog klart at erkende deres
Afstand fra det. Men saa var det, at Betænkelighederne og de
% --- p. 186 ---
\setcounter{footnote}{0}%
\opage{186}mange Overvejelser for og imod vaagnede. Skulde han være
fuldt ærlig, kunde han ikke tillægge sig en saadan oprindelig
Hensigt.

Betænkelighederne vare af forskellig Art, og det lønner sig
at dvæle ved dem, da de kaste Lys over Kierkegaards inderste
Væsen.

Kierkegaard var Psykolog nok til at vide, at en æstetisk
og intellektuel Produktion i større Stil ikke kan fremkaldes med
fuld Bevidsthed og som Middel til et udenfor liggende Formaal.
Og naar han besindede sig, blev han klar over det Uvilkaarlige
i hans Produktion, det, han ikke frygtede for at kalde det Geniale.
»Ethvert Geni er overvejende Umiddelbarhed, kender
intet Hvorfor«. Først bagefter kan det uvilkaarligt frembragte
Værks Betydning blive klar. Her passer Kierkegaards egen, i
en anden Sammenhæng udtalte Tanke, at vi forstaa baglængs,
men leve forlængs; dette Sidste gælder særligt om det skabende
Geni. Og ofte var Produktionstrangen saa stærk, at den truede
med at sprænge ham, ødelægge ham paa Sjæl og Legeme. Han
maatte med Selvbeherskelse lægge en Dæmper paa sin Lyst til
at skrive, saa at han kun arbejdede nogle Timer om Dagen. ---
Paa den anden Side turde han heller ikke betragte sig som et
udvalgt Guds Redskab, der uden at være sig det bevidst havde
arbejdet i det Religiøses Tjeneste, uagtet han var klar over, at
det Religiøse laa langt dybere i hans Sjæl fra først af end de
æstetiske og filosofiske Interesser.

Dertil kom, at han fra først af havde tænkt sig sin Forfattervirksomhed
som noget blot Midlertidigt, et Afløb for den
store Trang til literær Produktion, før hans praktiske Virksomhed
som Præst begyndte. Og senere var han flere Gange ved
at holde op; hvad der førte ham videre, var dels den voldsomme
Produktionstrang, dels ydre Forhold af forskellig Art. Om disse
ydre Forhold bruger han det Udtryk: »Virkelighedens Spænding
satte en ny Streng paa mit Instrument«. Denne Spænding
blev først frembragt ved Angrebene paa ham i Vittighedsbladet
»Korsaren«. Det saarede ham dybt, at baade hans legemlige
og aandelige Personlighed her blev karikeret. Paa den Tid var
man endnu ikke vant til Sligt, og Kierkegaard var særligt meget
% --- p. 187 ---
\setcounter{footnote}{0}%
\opage{187}lidet haardhudet. Det gør et tragikomisk Indtryk at se, hvorledes
han i sine Dagbøger, altsaa i sin dybe Ensomhed, beskæftiger
sig med denne Sag og føler Trang til at stemple det som
»Løgn«, at hans ene Buxeben skulde være kortere end det
andet, saaledes som det var paa Karikaturbillederne i »Korsaren«.
Men Kierkegaard saa i disse Angreb ikke blot en personlig
Fornærmelse, men en Nationalforbrydelse, fordi de Indflydelsesrige
i Samfundet taalte Sligt uden Protest. Nu kunde og vilde han ikke stanse sin Virksomhed og drage over som
Præst til den jyske Hede, hvad der havde været hans Ønske;
det vilde være at vige for Nederdrægtigheden! --- En anden
Streng blev der sat paa hans Instrument ved den officielle Kirkes
Stilling til hans Virksomhed, især til hans bestandigt strengere
Betoning af Forskellen mellem den oprindelige Kristendom,
»det nye Testamentes Kristendom«, og den moderne, i Nutidens
Kirke repræsenterede Kristendom. Hvad han paa den Tid (ved
Aar 1849) fordrede, var blot, at den officielle Kirkes Repræsentanter
skulde vedkende sig hin Forskel, vedkende sig, at de vare
langt fra Idealet, i Stedet for at give det Udseende af, at der
var en ubrudt Sammenhæng mellem Nutidskristendom og Urkristendom,
ikke blot hvad Dogmatiken angik (thi her nærede
Kierkegaard ingen Tvivl), men hvad den etiske Livsførelse, særligt
Forholdet til »Verden« (d. e. Kulturen) angik. Men denne
Indrømmelse vilde man ikke gøre. Man stod fornemt ignorerende
overfor hans Virksomhed. Dette gjorde --- i Forbindelse
med det stærkere personlige Gennembrud hos ham i disse Aar
--- at hans religiøse Virksomhed tog ny Fart, med strengere
Betoning af de ideelle Fordringer og med stedse skarpere Polemik
mod Kirken, indtil han i sit sidste Leveaar (1854--- 1855)
førte den revolutionære Kamp mod den bestaaende Kristendom,
der tilsidst lagde ham det Ord paa Tungen: det nye Testamentes
Kristendom er ikke til! ---

Det er karakteristisk for Kierkegaard, at han har tænkt sig,
at hans Forfatterskab kunde blive en isoleret Episode i hans
Liv. Det tyder paa en skarp Modsætning mellem hans literære
Begavelse og hans Temperament. Han har følt, at Fantasiens
og Tankens Liv, saa tillokkende det var for ham, ikke vilde
% --- p. 188 ---
\setcounter{footnote}{0}%
\opage{188}kunne optage ham helt, fordi der var en dunkel Magt i ham, som
det ikke vilde kunne overvinde. I hans Personligheds inderste Væsen have Produktionstrangen og Tungsindet mødt hinanden, snart
i Strid, snart under gensidig Indflydelse. Kun ved Vexelvirkningen
mellem dem forstaar man Kierkegaard. Uden Tungsindet, og
hvad det førte med sig, vilde han maaske være bleven en ren
Digter eller en ren Tænker. Nu hindredes han i at hengive sig
til Kunsten eller Videnskaben for deres egen Skyld alene. Han
fik ikke Lov at gaa bort efter Tangenten, men maatte lade alle
Billeder og Tanker kredse om det dybe Centrum, der stedse
drog dem til sig. --- Det er Dagbøgerne fra 1848—1849, der
have aabnet os denne Forstaaelse af Kierkegaards inderste Natur.

Tungsindet har virket koncentrerende. I Tungsindet trækker
Sjælen sig tilbage i sig selv eller bøjer sig mod sig selv,
som den Saarede bøjer sig over Saaret. Og i Tungsindet tabe
alle Tanker deres Værdi --- uden den Tanke, der for den Tungsindige
kan afgive et fast Tilhold, et Nødanker i det indre Livs
Storme. --- Det har virket isolerende. Kierkegaard siger, at det
tidligt stod klart for ham, at for ham var der ingen Hjælp eller
Trøst at finde hos Andre. Samliv med Andre var ham en Umulighed
paa Grund af den dybe Forskel mellem hans og Andres
Stemning overfor Livet. Det var Grunden til, at han hævede
sin Forlovelse. Ved Siden af den unge Piges lyse og lette Sind
følte han først ret, hvor ulykkelig han var --- og han følte tillige,
at hun ikke vilde kunne forstaa ham, og at han vilde gøre
hende ulykkelig ved sit tunge Sind. At han ikke overtog noget
Embede, udleder han af samme Grund: han var ikke Herre
over sit Sind --- og kunde derfor ikke gaa ind i en Stilling,
hvor han var afhængig af Andre! Det havde maaske været et
Gode for ham, hvis han af økonomiske Aarsager havde været
nødt til at paatage sig bestemte ydre Pligter. Det at \textit{skulle}
udføre Noget til bestemt Tid kan hjælpe over mangen Forstemning
og kan være en Støtte i Kampen mod dens Mørke og
Lammelse. --- Nu havde han kun den Afledning af Tungsindets
Magt, som det literære Arbejde skaffede. Den stærke Produktionstrang
skærpedes ved Kontrasten til Tungsindets Uro, og
desuden svækkedes Tungsindets Magt, naar saa megen Kraft
% --- p. 189 ---
\setcounter{footnote}{0}%
\opage{189}som muligt blev sat ind paa at producere. »Lig hin Prinsesse
i 1001 Nat frelste jeg Livet ved at fortælle, d. v. s. at producere.
Et uhyre Tungsind, indre Lidelser af sympatetisk Art,
Alt kunde jeg magte, naar jeg fik Lov til at producere.« (Dagbøgerne
1849). --- Endeligt maatte Tungsindet føre til stedse
ny Grublen ved at stille et særligt personligt Problem: er Tungsindet
Brøde eller Sygdom? Er det ikke formasteligt at frakende
Alt Værdi? Kierkegaard sammenlignede selv sit Tungsind
med acedia, det middelalderlige Udtryk for aandeligt Mørke
og Slaphed, en af de syv Dødssynder (hvorfor ogsaa acediosi
af Dante sættes i Helvede --- fordi de have været »lidet glade,
hvor Solen skinned sødt paa Enge«). Her var Stof til stadig
Grublen og til Konstruktion af mørke Muligheder. ---

Under Indflydelse af de to Kræfter --- den stærke Produktionstrang
og det dybe Tungsind --- er Kierkegaards literære
Værk blevet til. Udsprunget som det var af hans inderste Natur
og som Forsøg paa at naa til fuld Klarhed over sig selv,
maatte det blive af stor vækkende og opdragende Betydning for
Andre. Det udøvede en Mission, ikke fordi det var oprindeligt
anlagt til det, men fordi det selv var Frugten af det indre Livs
Kampe hos et af Livet stærkt grebet Menneske.

Som Tænker førtes Kierkegaard ifølge sin Natur til at lægge
Vægt paa Erfaringen og have Mistillid til Spekulationen, der
førte bort fra det personlige Menneskelivs Kriser og Opgaver.
Hans psykologiske Interesse holdt ham bunden og førte ham ad
Iagttagelsens Veje, især da han i høj Grad besad Beskrivelsens
Kunst. Hans etiske Interesse førte ham til at betone Betydningen
af Arbejdet, Kampen og Lidelsen i Tiden: vi leve jo
forlængs --- derfor ere Kamp og smertelig Spænding uundgaaelige
for Enhver, der ikke tankeløst fortaber sig i fantastisk Spekulation.
Og hans skarpe kritiske Evne førte ham endeligt til
at indse det Uvidenskabelige i den spekulative Teologi, i hvilken
Mange i Datiden søgte Løsningen af Tilværelsens Gaader. --- I
denne sin Tænknings realistiske og kritiske Karakter staar han
ikke ene i Danmark. I hele den Maade, paa hvilken Romantikens
Filosofi blev modtagen i Danmark, ytrer sig en Tilbøjelighed
til at sætte psykologisk Erfaring og kritisk Prøvelse i
% --- p. 190 ---
\setcounter{footnote}{0}%
\opage{190}Stedet for abstrakt og mystisk Spekulation og en Hævden af
Tidsudviklingens og det etiske Arbejdes Betydning overfor Romantikens
Hang til at svælge i det Evige, --- en Tilbøjelighed,
i hvilken vi maaske kunne se en Ejendommelighed ved vor
Nationalkarakter. Hos Henrik Steffens, Brødrene Ørsted, Frederik
Christian Sibbern og Poul Møller fremtræder denne hele
Tendens mere eller mindre udpræget, saa meget de end maa
siges at tilhøre Romantikens Tid.

Kierkegaard var nu, trods sin Reaktion mod Romantiken,
selv en udpræget Romantiker. Ja, han har --- forunderligt nok
hos en Mand, der hævder Livet og Existensen i deres Modsætning
til Tanken, --- den Fjernhed fra Livet, der snart minder
om Romantikens Poesi, snart om Middelalderens Askese,
eller rettere: som kun i disse to Former finder Livet værdt at
leve. Jeg har allerede fremhævet, at han mangler Blik for Livets
sammenhængende Udvikling. Overgangen mellem hans
»Stadier« skal stedse ske ved »Spring«. Han negter udtrykkeligt,
at der bestaar en Analogi mellem personlig Udvikling og
organisk Udvikling. Og medens han fra først af havde skelnet
mellem tre væsentlige Stadier eller Livsstandpunkter, skyder
han senere mere og mere det andet Stadium, den etiske Livsanskuelse
ud: den bliver for ham Et med Aandløshed, Filistrøsitet.
»De Aldre, man i Henseende til Idealitet kan lære af,
ere: Barnet, Ynglingen, den unge Pige, Oldingen --- af den
virksomme Mand, af den travle Husmoder lærer man i denne
Henseende ikke, og hvorfor ikke? fordi de væsentligen ere beskæftigede
med Endelighedens Formaal«. Altsaa kun Ungdom
og Alderdom skulle have Idealitet i sig, men ikke den modne
og kraftige Alders Stræben og Virken! Det kan man kun mene,
naar Begrebet om Idealitet enten er romantisk eller asketisk,
eller begge Dele. Realismen er her mere idealistisk end den
overspændte Idealisme. I selve Virkelighedens Sammenhæng og
i selve Virkelighedens Arbejde maa den poetiske og den religiøse
Idealitet være at finde, om den idethele skal kunne findes
mere. Og den maa findes der, hvis Slægten ikke skal opløses;
thi paa Haab og Erindring uden mellemliggende værdifuld Virken
og Virkelighed kan der i Længden ikke leves. Af selve
% --- p. 191 ---
\setcounter{footnote}{0}%
\opage{191}Virkeligheden maa de nye Idealer stedse frembringes, som Gnisten
af Flinten; forkaster man Flinten, bliver Gnisten ogsaa
borte. Her ligger Kierkegaards kolossale Begrænsning.

Og her ligger ogsaa det Punkt, hvor den, der har været
under hans stærke Paavirkning, maa kæmpe for at frigøre sig.
Mangen ung Mand har trods sin Beundring for ham kastet hans
Bøger paa Gulvet og vendt ham Ryggen, naar Livets Ret og
naturlige Sammenhæng gik op for Tanken, og naar der for Erkendelsen
meldte sig andre Problemer end de, der stilledes i
den religiøse Traditions Paradoxer. Og Kierkegaard selv vilde
og kunde ikke have Tilhængere. Han vilde ikke! »Jeg brugte«,
siger han, »næsten Halvdelen af min Kraft til at modarbejde
mig selv; jeg vilde ingen Tilhængere have« (Dagbøgerne 1849).
Han vilde vække en personlig Sandhed, som kun kan vindes
ved Selvvirksomhed, ikke ved direkte Tilslutning og Efterligning.
Og han kunde ikke have Tilhængere: det sørgede han
for ved sit sidste Ord, som var det, at hvad han havde kæmpet
for sit hele Liv, --- at det ikke var til! Under den voldsomme
Kamp med Kirken i sit sidste Aar udtalte han: »Det nye Testamentes
Kristendom er ikke til!« Da han i den mest lidenskabelige
Polemik, det danske Aandsliv har kendt, havde hævdet
denne Sætning, sank han sammen; hans Kraft var udtømt,
og Døden lukkede hans klare Øjne. Hvad saa? Hvor vilde
han have staaet, hvis han ikke var død saa tidligt (42 Aar gammel)?
Paa dette Spørgsmaal kan intet Svar gives.

Han var en Martyr for Tanker, der maa føres videre ad
andre Veje, end han saa. Men ved energisk og konsekvent at
gaa frem ad den Vej, han saa, har han gjort os, der kom efter
ham, den bedste Tjeneste. Vi kunne fortjene Navn af hans
Efterfølgere og vise vor Taknemmelighed imod ham, selvom vi
gaa frem ad helt andre Veje end han.
% --- p. 192 ---
\setcounter{footnote}{0}%
\opage{192}\essay[Fichtes »Taler til den tyske Nation« og de danske Nordslesvigere (1899)]{FICHTE'S »TALER TIL DEN TYSKE NATION«\\ {\small OG}\\ NORDSLESVIGERNES NATIONALE KAMP.}
\begin{center}(1898).\end{center}
\streg

\section*{I.}
\streg

Et Folk bliver til, naar fælles Skæbne og fælles Arbejde
frembringer Fællesskab i Sindelag og Tankegang. At man hører
til samme Folk, kendes paa, at man deler Fare og Tryghed,
Spænding og Haab, Harme og Beundring med hverandre under
Skæbnens Gang gennem Historien. Dermed forbinder sig Tilbøjeligheden
til at vende tilbage til de samme Minder og til at
forme beslægtede Idealer. Schweiz's Exempel viser, at et fælles
Sprog ikke er absolut nødvendigt.

Under lykkelige Forhold er Staten den ydre, faste Organisation,
indenfor hvilken Folkets Liv udfolder sig. Men Historien
viser Exempler paa en Modsætning mellem Folk og Stat,
især naar en Befolkning ved Erobring bliver løsreven fra den
Stat, den i lange Tider har tilhørt. Folket maa da føre Kampen
for sin Bestaaen uden en Statsmagts Midler og ofte i Modsætning
til en bestaaende Stats Magt.

I Begyndelsen af dette Aarhundrede var der en Tid, i
hvilken det nu saa mægtige tyske Folk fik Brug for denne
Skelnen mellem Folk og Stat, og en af dets mest energiske og
patriotiske Sønner trøstede sine Landsmænd med denne Forskel.
Det har for os Danske --- med Henblik paa den Kamp, vore
% --- p. 193 ---
\setcounter{footnote}{0}%
\opage{193}nordslesvigske Landsmænd føre for Bevarelsen af deres Sprog
og deres Nationalitet --- sin Interesse at fremdrage denne Episode.
Det er dette, der har bragt mig til paany at læse J. G.
Fichte's »Reden an die deutsche Nation« (1808).

Lad mig først minde om den historiske Situation, i hvilken
disse Taler bleve holdte.

Under Revolutionens og Kejsertidens Krige, der for en stor
Del gik ud over Tyskland, havde Preussen, efter en kort Deltagelse,
holdt sig tilbage. Det havde roligt set paa Østerrigs
Nederlag og det tyske Riges Opløsning. Men da kom Turen
til det. I Slaget ved Jena (1806) blev den preussiske Krigsmagt
slaaet til Jorden. Den ene Fæstning maatte overgive sig
efter den anden. Napoleon holdt sit Indtog i Berlin. Kongen og Regeringen flygtede til Kønigsberg, og da Fjenden ogsaa kom
her, længere øster paa. Ved Freden i Tilsit (1807) afstod Preussen
det Halve af sit Landomraade og kom for Resten af Riget
i en fuldstændig Afhængighed af Sejrherren. Dets Politik i ydre,
tildels ogsaa i indre Anliggender bestemtes af Napoleon. Denne
vaagede mistænksomt over enhver Bevægelse i det overvundne
Land, han mente kunde blive ham farlig. Saaledes maatte den
energiske Statsmand von Stein afskediges; ja, han blev erklæret
fredløs og maatte drage til Udlandet.

Under disse Forhold opstod hos de Bedste i Tyskland den
Overbevisning, at der maatte vindes i det Indre, hvad der var
tabt i det Ydre.\footnote[1]{I en Betænkning af Overpræsident von Vincke blev det udtrykkeligt udtalt, at det gjaldt »at generobre i det Indre, hvad Staten havde mistet i ydre Omfang«. (Smlgn. R. Haym: \textit{Wilhelm von Humboldt.} Berlin 1856, p. 262).} Det var under disse Forhold, at JOHANN
GOTTLIEB FICHTE i Berlin i Vinteren 1807—1808 holdt sine
»Taler til den tyske Nation«. Franskmændene holdt endnu Berlin
besat; Fichte's Ord overdøvedes ofte af de franske Trommer
paa Gaden, og ofte gik der det Rygte, at han var bleven
arresteret. Napoleon har dog ikke anset denne »Ideolog« for
farlig.

% --- p. 194 ---
\setcounter{footnote}{0}%
% CHECK p. 194: centred line: heading or verse
\opage{194}

\section*{II.}

Før \textit{vi} gaa ind paa Talernes Indhold, maa vi gøre lidt nærmere
Bekendtskab med Talerens Personlighed.

Fichte var en sachsisk Bondesøn, født i Aaret 1762. Hans
Forældre vare fattige, og Drengen hjalp dem i deres Profession,
Væveriet, eller passede Gæssene. Han udmærkede sig
tidligt ved sin klare Opfattelse og sin gode Hukommelse, og
det blev Anledningen til, at han opnaaede en videre Uddannelse
af sine Evner. En Adelsmand fra Omegnen kom en Søndag
for sent til Kirke, og da han gerne vilde vide, hvad Præsten
havde sagt, blev Gaasedrengen Johann Gottlieb hentet og
holdt hele Prædikenen for ham. Adelsmanden besluttede nu at
sørge for den opvakte Drengs Undervisning og hjalp ham frem
til at blive Student. Efter sin Velgørers Død havde Fichte nogle
trange Aar. Han led Nød og var ofte ved at fortvivle aldeles.
Kun den energiske Villie, som var Hovedtrækket i hans Karakter,
holdt ham oppe og bevarede ham for Selvmord. Det lykkedes
ham at opnaa en Huslærerplads i Zürich, og dette Schweizerophold
blev i flere Henseender betydningsfuldt for ham. Han
lærte frie, republikanske Forhold at kende, og han stiftede Bekendtskab
med Pædagogikens Reformator Pestalozzi. Han læste
i de følgende Aar Kant's Skrifter, der bleve afgørende for hans
hele aandelige Retning, og kom ogsaa under et Besøg i Kønigsberg
i personlig Berøring med den gamle Mester. I Aaret
1794 blev han Professor i Filosofi i Jena og udfoldede her en
omfattende og betydningsfuld Virksomhed. Hans Filosofi er
ikke let at fremstille i stor Korthed. Han fandt det selv overordentligt
vanskeligt at give en tilfredsstillende Fremstilling af
sin Filosofi; hele sit følgende Liv igennem arbejdede han paa
at faa et klart Udtryk for sin Tanke, uden at nogen Form helt
fyldestgjorde ham. Hans filosofiske Standpunkt er karakteriseret
ved, at han paa én Gang tager sit Udgangspunkt i Videnskabens
højeste Problemer og opfatter Tænkningen som Udslag af
Personlighedens inderste Væsen. Medens Henvisning til det
personlige Moment i Tænkningen for saa Mange betyder, at man
skal høre op at tænke, betød den for Fichte netop Nødvendigheden
af en saa meget des større Anstrengelse: det gjaldt om,
% --- p. 195 ---
\setcounter{footnote}{0}%
\opage{195}at Tænkningen kunde give et saa klart og konsekvent Udtryk,
som muligt, for Personlighedens Forhold til Erfaringen. Konsekvens
i Tanken og Karakter i Villien vare for ham uadskillelige.
Erfaringen havde for ham større Betydning i praktisk end
i teoretisk Henseende: Verden, som Erfaringen viser os den,
er blot det Materiale, vor Villie skal bearbejde, --- Grundlaget
for vor Pligt. --- Fra ortodoks Side blev han anklaget for Ateisme,
og ved et utvivlsomt Indgreb i den videnskabelige Lærefrihed
blev han afsat fra sit Professorat. Han levede nu i nogle
Aar i Berlin og blev derpaa preussisk Professor. Efter Jenaslaget
drog han til Königsberg, og da ogsaa denne By blev besat
af Fjenden, drog han over til København.

I København boede Fichte hos den berømte Jurist Anders
Sandøe Ørsted, som i sin Ungdom var en begejstret Tilhænger
af ham. Det har sin Interesse at se, at Fichte er kommen i
nær Berøring med udmærkede Repræsentanter for de nordiske
Landes Aandsliv. I Danmark vare begge Brødrene Ørsted hans
store Beundrere og forstode paa selvstændig Maade at bearbejde de
Paavirkninger, de modtoge af ham. Nordmanden Henrik Steffens
havde været hans Tilhører i Jena og omtalte ham med Begejstring
paa sine berømte Forelæsninger i København. Den svenske
Filosof Benjamin Høijer havde under sit Ophold i Dresden omgaaedes
Fichte og vist sig som en energisk Tænker under deres
hyppige Diskussioner. Et Sted i »Taler til den tyske Nation«
finder maaske sin Forklaring ved de Erfaringer, Fichte saaledes
havde gjort. Det maa tilføjes, at hans Indflydelse paa det danske
Aandsliv har været meget betydelig. Hans Personlighedslære
paavirkede dels gennem Ørstederne, dels gennem Steffens og
Sibbern den danske Tænkning paa en grundigere Maade end
senere den Hegelske Filosofi. Gennem Sibbern og Poul Møller
gaar der Tanketraade fra ham til Søren Kierkegaard, og
direkte har han grebet stærkt ind i Grundtvigs personlige Udvikling,
hvad jeg strax skal anføre et Vidnesbyrd om.

Da Freden i Tilsit var sluttet, vendte Fichte tilbage til Berlin
og holdt her i den følgende Vinter de Taler, vi nu skulle
beskæftige os med. Ved det nye Universitet i Berlin var han
meget virksom, men hans Virksomhed blev ikke  lang. Han
% --- p. 196 ---
\setcounter{footnote}{0}%
\opage{196}blev smittet af en Feber, hans Hustru havde paadraget sig ved
at pleje de Saarede, og døde i Begyndelsen af Aaret 1814.

Grundtvig skrev et Mindedigt over ham, i hvilket han, den
ortodokse Præst, udtaler, hvilken Vækkelse han skylder den for
Ateisme anklagede Filosof. To af de mest betegnende Vers af
dette Digt skulle her anføres:

\begin{verse}
Gottlieb Fichte! Ved din Kæmpegrav \\
Stander dybt bevæget jeg i Aanden: \\
% CHECK p. 196: centred line: heading or verse \\
Stirre maa jeg paa din Vandringsstav, \\
% CHECK p. 196: centred line: heading or verse \\
Som saa brat dig brødes udi Haanden. \\
Taarer falder og fra Øjelaag \\
Sagte ned paa sorte Kistelaag. \\
Røre maa sig over Grav min Tunge; \\
% CHECK p. 196: centred line: heading or verse \\
Men hvad kan, hvad tør den vel udsjunge?
\end{verse}

. . . . . . . . . . . . .

% CHECK p. 196: centred line: heading or verse

\begin{verse}
Fichte! Ja den Stund, jeg vandrer her, \\
% CHECK p. 196: centred line: heading or verse \\
Ej jeg kan, og ej jeg maa dig glemme, \\
Ej fornægte, at du varst mig kær, \\
At mig glædede din Kæmpestemme; \\
\textit{Den} mig vakte af min falske Ro, \\
Den mig vækkede til Tvivl og Tro, \\
Vækkede mig til at ihukomme \\
% CHECK p. 196: centred line: heading or verse \\
Sandheds Lys og Kærlighedens Blomme.
\end{verse}

\section*{III.}

Fichte spørger i sine »Taler« først og fremmest, om der
endnu gives et tysk Folk, og om dets fortsatte Bestaaen i dets
ejendommelige og selvstændige Væsen kan hævdes og fortjener
at hævdes. Den ydre Selvstændighed er tabt. Sejrherren bestemmer,
hvilken Del Tyskerne skulle tage i Historiens store
Begivenheder. Men den ydre Underkastelse bør ikke medføre
den indre. Aanden skal Erobreren ikke kunne kue. Kampen
med Vaabnene er forbi. Men der staar en anden Kamp tilbage:
for aandelig Selvstændighed, for Tankens og Karakterens
Udvikling:

»Lad os være paa vor Post overfor den Behagelighed, der
kan være ved at tjene; ti den røver endog vore Efterkommere
% --- p. 197 ---
\setcounter{footnote}{0}%
\opage{197}Haabet om fremtidig Befrielse. Dersom vor ydre Virken lægges
i hæmmende Lænker, saa lad os desto dristigere hæve vor Aand
til Tanken om Friheden, til Livet i denne Tanke, til kun at
ønske og begære denne ene Ting. Selvom Friheden for en Tid
forsvinder fra den synlige Verden, saa give vi den et Tilflugtssted
i det Inderste af vore Tanker, saalænge indtil den nye
Slægt voxer op omkring os, der faar Kraft til ogsaa i det Ydre
at fremstille denne Tanke\dots{} Lad blot ikke vor Aand kues
og underkastes og bringes i Fangenskab, ligesom vort Legeme!«

Der gives nemlig et Omraade, som Sejrherrens Skinsyge
ikke har lukket: det aandelige Livs og særligt Opdragelsens Omraade!
Der maa da de tilbageblevne Kræfter anvendes. At tabe
sig i Smerte og Sorg over det Tabte nytter ikke.

Lad fremfor Alt ikke Sejrherrens Overlegenhed føre de
Overvundne til ubetimelig Beundring! En saadan Beundring
vil let faa Udseendet af Frygt. Og den vil stride mod Folkets
gamle Maalestok for Storhed!

»Nej, brave tyske Landsmænd, saadan Uforstand forblive
langt fjernet fra vor Aand, --- saadan Besudling forblive fremmed
for vort Sprog, der er formet til Udtryk for Sandheden!
Lad Udlandet juble i Forbavselse over enhver ny Begivenhed,
--- lad det i hvert Aarti danne sig en ny Maalestok for Storhed,
skabe sig nye Guder og tale Gudsbespottelse i sin Iver
for at prise Mennesker! Lad os beholde vor gamle Maalestok
for Storhed: at stor er kun den, der med Begejstring fatter de
Ideer, som stedse kun bringe Held over Folkene! Og lad os
overlade Efterverdenen Dommen over de levende Mennesker!«

Haabet om en bedre Fremtid bygger Fichte paa det tyske
Folks Ejendommelighed. Det er et Urfolk, der bor paa sin egen
Grund fra umindelige Tider, og som har bevaret sit gamle Sprog
--- i Modsætning til de germanske Stammer, der bosatte sig
paa romersk Grund og dér uddannede de nylatinske Sprog.
Paa Sproget lægger Fichte en overordentligt stor Vægt i national
Henseende. »Hvilken umaadelig Indflydelse maa ikke Sprogets
Beskaffenhed have paa et Folks hele menneskelige Udvikling!
Sproget ledsager jo den Enkelte helt ind i hans Sinds
hemmeligste Dyb ved al hans Tænken og Villen. Det hæmmer
% --- p. 198 ---
\setcounter{footnote}{0}%
\opage{198}ham, eller det bevinger ham. Det forbinder hele denne Menneskemængde,
der taler det, til at udgøre en eneste Aand. I det
mødes det Sanselige og det Aandelige og smelte saaledes sammen,
at det ikke er muligt at sige, om Sproget selv er sanseligt
eller aandeligt.« --- Men fremfor Alt er det tyske Folk et
Folk med Sans for Inderlighed, Oprindelighed og Personlighed.
Hvad det tyske Folk har sat ind i Verden af aandeligt Arbejde,
vidner herom. \textit{Reformationen} var en Kamp for aandelig Frihed,
et Gennembrud af Personlighedens Trang. Netop fordi det
tyske Folk tog det gamle Kirkesystem med saa dyb Alvor, følte
det Trang til en Befrielse fra det; de Folkeslag, der affandt sig
med det gamle System paa en mere overfladisk Maade, kunde
ikke føle en saadan Trang til at befries fra det. Den af Leibniz
og Kant grundlagte tyske \textit{Filosofi} har --- i Modsætning til
andre Folks Filosofi --- særligt Blik for Livets indre Betingelser;
den opfatter ikke Sandheden som Noget, der kan modtages
blot ved ydre Paavirkning, men som Noget, der maa erhverves
gennem Udløsning af frie Kræfter. Ogsaa paa \textit{Opdragelsens}
Omraade --- især i Pestalozzi's Pædagogik --- ytrer sig den
samme Tendens til at gaa indefra udad, til at vække Kærligheden
og Trangen til Emnet og ikke paatvinge Kundskaben
udefra. Det er Pestalozzi's Pædagogik, Fichte støtter sig til i
det Udkast til en ny Nationalopdragelse, der udgør saa stor en
Del af Talernes Indhold. Den skal udvikle til Selvstændighed,
gøre Dannelsen til en personlig Ejendom, fra først til sidst sætte
den Unges Selvvirksomhed i Gang. ---

I Løbet af den Karakteristik, Fichte saaledes giver af sit
Folk for at indgyde det Mod og Lyst til at være sig selv, udtaler
han lejlighedsvis, at de Træk af ydre og indre Ejendommelighed,
han fremhæver, ogsaa gælde for de skandinaviske Folks
Vedkommende. Utvivlsomt er denne Bemærkning ikke fremkommen
uden Indflydelse af de Erfaringer, Fichte selv havde
gjort gennem Bekendtskabet med udmærkede Repræsentanter for
det nordiske Aandsliv. ---

Efter at have karakteriseret sit Folk indskærper Fichte den
betydningsfulde Forskel mellem Folk eller Fædreland og Stat.\footnote[1]{Den ottende Tale har til Overskrift: »Hvad et Folk er i Ordets højere \opage{199}Betydning, og hvad Fædrelandskærlighed er«, og kommer til det Resultat: »Folk og Fædreland ligger langt ud over Staten i dette Ords sædvanlige Betydning«.}

% --- p. 199 ---
\setcounter{footnote}{0}%
\opage{199}Hvad der hører til samme Folk eller Fædreland, har først og
fremmest samme Sprog:

»De, der tale samme Sprog, ere allerede forud for al menneskelig
Kunst ved selve Naturen forbundne ved en Mængde
usynlige Baand. De forstaa hverandre indbyrdes og formaa at
vinde stedse klarere gensidig Forstaaelse; de høre sammen og
udgøre et naturligt og uadskilleligt Hele. De kunne ikke optage
noget Folk af anden Herkomst og andet Sprog i sig og
blande sig med det, uden at volde sig selv Forvirring og mægtigt
forstyrre sin Udviklings sunde Fremgang. Af denne indre,
ved selve Menneskets aandelige Natur dragne Grænse bestemmes saa igen Boligernes ydre Begrænsning, og efter den naturligste
Opfattelse af Tingene ere ingenlunde de Mennesker, der
bo indenfor visse Bjerge og Floder, derfor ét Folk, men omvendt:
Menneskene bo sammen (skærmede af Floder og Bjerge,
naar Lykken har føjet det saa), fordi de allerede forud vare ét
Folk i Kraft af en langt højere Naturlov.«

Det er efter Fichte's religiøse Opfattelse tilsidst en af Livets
inderste Kilde udspringende Lov for Udvikling, der udgør det
enkelte Folks Væsen. Det Evige spalter sig i en Række af
specielle Livsformer, fordi det ikke kan udtømmes i nogen enkelt
Form. Denne særlige Udviklingslov eller specielle Livsform er
hvad man kalder et Folks Nationalkarakter. Den Enkelte, der
lever under denne Lov, vil ikke med klar Bevidsthed kunne
fatte Alt, hvad der ligger i den; men ved skjulte Traade hænger
den Enkelte gennem sin Nationalitet sammen med det oprindelige
Livs Kilde. Derfor strider det mod Aandslivets guddommelige
Orden, at et Universalmonarki skulde træde i Stedet
for de forskellige Nationaliteter. Ligesom Fichte tidligere tog
Ordet mod dem af sine Landsmænd, der glemte den nationale
Sorg i Beundring for Verdenserobreren, dette Overmenneske,
saaledes taler han her mod dem, der mente, at Nationaliteternes
Tid nu var forbi og fandt Trøst over deres Folks Undergang i
Tanken om at skulle tilhøre en Verdensstat. Der laa for Fichte
% --- p. 200 ---
\setcounter{footnote}{0}%
\opage{200}i denne Tanke en Raahed, som forgreb sig paa det Hellige, »et
Indgreb i den aandelige Verdens højeste Lov«, efter hvilken
Menneskehedens Væsen kun kan aabenbares i højst mangfoldige
Nuancer, dels enkelte Personligheder, dels Folkeslag.

Staten betyder i Modsætning til Folk og Fædreland den
sociale Ordning, der sikrer Retten og Freden og gør Handel og
Vandel mulig. Den er kun Middel og Betingelse eller ydre
Apparat for det ideale Indhold, der udfolder sig i Folket og
lever i Fædrelandskærligheden. ---

Det var Grundtankerne i de Ord, Fichte henvendte til sine
Landsmænd i Nødens Tid.

\section*{IV.}

Der fremtræder for den, der sammenligner Fichte's Taler
med de danske Nordslesvigeres Forhold, ikke blot slaaende Ligheder,
men ogsaa betydningsfulde Forskelle. Men i det Store
og Hele kunne de Danske i Slesvig styrke deres Mod og Haab
i Kampen ved den tyske Patriots Ideer. Han har jo udtrykkeligt
erklæret, at de ikke blot finde Anvendelse paa hans eget
Folk, men ogsaa paa Folkeslag som de nordiske.

Nordslesvigerne ere et Urfolk, baade hvad Jorden og hvad
Sproget angaar. De bo paa deres Fædres Jord, og de tale deres
Fædres Sprog. De have et Fædreland at bevare, selvom de ved
Erobring ere komne til at tilhøre en fremmed Stat. Fordi et Land
er erobret, blive dets Beboere ikke derfor --- det har Fichte
lært os --- en Del af det erobrende Folk. Der er endnu, efter
at Vaabnenes Kamp er forbi, en Kamp at føre om Sproget og
Nationaliteten, en Kamp, i hvilken nogle af de højeste og ædleste
Egenskaber, Mennesker kunne udvise, finde Anvendelse. Det
gælder Troskaben mod sig selv, Bevarelsen af den aandelige
Ejendommelighed.

Men Forholdene ere vanskeligere i Nordslesvig, end de vare
i Preussen paa Fichte's Tid. Han kunde henvise til, at der dog
var ét Omraade, paa hvilket Erobreren ikke havde tænkt paa
at lægge Haand, nemlig det aandelige Liv, særligt Opdragelsen.
Det var Napoleons brutale Foragt for »Ideologerne«, der lod
% --- p. 201 ---
\setcounter{footnote}{0}%
\opage{201}ham overse de aandelige Kræfters Betydning i Verden. Magthaverne
i Slesvig kende denne Betydning. Deres Skinsyge strækker
sig ogsaa ind paa dette Omraade. Man vil gennem Skolen
og ved at gøre det fremmede Sprog til eneste Undervisningssprog
tilintetgøre den erobrede Befolknings Nationalitet. Det
røber ganske vist en Forstaaelse for hvad aandelige Kræfter betyde,
som Napoleon ikke havde. Men det røber en langt større
Brutalitet at anvende Tvang paa det aandelige Omraade, end
helt at ignorere dette. Selv i Oldtidens paa Erobring grundede
Stater fik de Overvundne ofte Lov at beholde deres Sprog og
idethele deres aandelige Liv uhæmmet.

Til Erstatning nyde Nordslesvigerne godt af de Reformer,
der indførtes i Preussen i Nødens Tid efter Tilsiterfreden, saavel
som af de politiske Rettigheder, der skyldes den korte Frihedstid
efter 1848. De have borgerlige Rettigheder at hævde
som Medlemmer af den Stat, der har erobret dem, og de kunne
sende Repræsentanter til dens lovgivende Forsamlinger. Derigennem
have de Midler til at sætte visse Grænser for Tvangen
og til at bringe deres Tanker frem i det offentlige Livs Former.
Kampen bliver derved ikke blot en rent aandelig Kamp, men
tillige en Kamp om Retten.

Den Aand, der fører Magthaverne til at fare frem i Nord-Slesvig,
som de gøre, er unegteligt helt forskellig fra den Aand,
der gaar gennem Fichte's Taler. Det er Reaktionens, Bureaukratiets
og Militarismens Aand. Den meldte sig meget hurtigt
i Preussen, efter at den Begejstring, Fichte og en hel Række
af aandsbeslægtede Mænd havde vakt, var bleven brugt som et
politisk Middel i Frihedskrigen. Da Fichte's »Taler« ti Aar efter
hans Død skulde udkomme paany, mødte de Censurhindringer:
de vare nu blevne en farlig Bog. De to Sider i det tyske Væsen,
Idealiteten og Brutaliteten, træde her i Modsætning til hinanden.
Vi maa ikke forse os saaledes paa den sidste, at vi
glemme, hvor meget vi skylde den første. Vi maa tvertimod
glæde os over, at vi kunne appellere til det Ædleste og Bedste
i den tyske Folkeaand overfor den Fremfærd, der nu øves mod
vore fraskilte Landsmænd. Vi kunne indanke deres Sag for de
% --- p. 202 ---
\setcounter{footnote}{0}%
\opage{202}Ideers Domstol, til hvilke de Bedste i Tyskland toge deres Tilflugt,
da det var Nødens Tid for dem.

Man har undertiden ment, at der i Kulturen laa en Fare
for Nationaliteterne, idetmindste for de smaa Nationaliteter. Men
den virkelige Kultur sætter sin ædleste Blomst i den fine Sans
for Ejendommeligheder og Forskelle, i Forstaaelsen for det Nye
og Fremmede. Den vil derfor netop frede om de store og smaa
Nuancer indenfor det Menneskelige, forat Livet ikke skal blive
fattigere. Naar der er en Fare for Nationaliteterne, er det ikke, 
fordi der er for megen Kultur i Verden, men tvertimod, fordi 
der er for meget Barbari i Verden. Barbaren tramper det Fremmede
ned, fordi han ikke forstaar det, og den Kundskab og
Erfaring, han har om det Aandeliges Betydning, giver ham kun 
et Middel mere i Hænde til at øve Tvang. Med Barbariet er 
ingen Forhandling mulig; men det maa være muligt at komme 
til Forstaaelse med et Folk, der endnu i den Mand, der holdt
de begejstrede Taler i 1807—1808, ser en af sine bedste Sønner. 

% --- p. 203 ---
\setcounter{footnote}{0}%
\opage{203}\essay[Harnack om Kristendommens Væsen. (1901)]{HARNACK OM KRISTENDOMMENS VÆSEN\footnote[1]{En Anmeldelse i »Nordisk Tidsskrift for Filologi«. 1901.}.}
\streg

ADOLF HARNACK, Kristendommens Væsen. Autoriseret
Oversættelse ved Edv. Lehmann. København 1900, V. Pio's
Bogh. X + 219 S.

I det interessante Forslag til Reform af Undervisningen i
Græsk, som Wilamowitz-Moellendorff har indgivet til det preussiske
Kultusministerium, er bl. A. den Tanke fremkommen, at
Disciplene trinvis skulle indføres i den græske Literatur ligefra
Homer til Paulus. Hvis denne Tanke udføres, vil det blive Filologiens
Opgave at drøfte Urkristendommens Problem. Men
allerede nu er dette Problem et filologisk Problem, naar man
tager Begrebet Filologi i videre Betydning. Ligesom Studiet af
Platon og Værdsættelsen af hans Skrifter ikke blot er Filosofiens
Sag, men et strengt filologisk Grundlag ikke kan undværes, saaledes
forholder det sig analogt ogsaa med Studiet af de urkristelige
Skrifter, der maa ligge til Grund ved enhver Drøftelse af
»Kristendommens Væsen«: det hører ikke blot ind under Teologien,
men ogsaa under Filologien, naar det ikke skal svæve i
Luften. Det er en saadan Betragtning, der forklarer, hvorfor
den udmærkede Dogmehistorikers betydningsfulde Skrift anmeldes
i et filologisk Tidsskrift.

Det er forstaaeligt, at Harnack har følt Trang til at give
en positiv og sammenhængende Fremstilling af sin Opfattelse af
Kristendommens Væsen. Han har paa sin videnskabelige Løbe%
% --- p. 204 ---
\setcounter{footnote}{0}%
\opage{204}bane paa saa mange Punkter maattet indtage en kritisk og polemisk
Holdning overfor de herskende teologiske Meninger, at det
i Udenforstaaendes Øjne kunde faa Udseende af, at hans Standpunkt
var aldeles negativt. Ligesom den politiske Opposition ofte
bliver beskyldt for at være samfundsopløsende, saaledes beskyldes
den teologiske Opposition ofte for at være irreligiøs. Det
har sikkert været den kritiske Teolog en stor Tilfredsstillelse,
at han overfor en usædvanligt talrig Kreds af Berlineruniversitetets
Studerende har kunnet udvikle, hvad der for ham staar
som Kristendommens egentlige Væsen. Man mærker i hans
Bog den Varme og Begejstring, med hvilken han har tænkt og
talt over dette Emne. Det vil for Mange herhjemme være kærkomment,
at Bogen nu foreligger i dansk Oversættelse.

For at opdage, hvad Kristendommen egentligt er, maa man
ifølge Harnack se paa hele dens Historie. Urkristendommen er
kun denne Histories første Periode, og ingen enkelt Periode
kan fuldkomment og definitivt udtrykke Kristendommens Væsen.
Foruden det, Harnack kalder »Evangeliet«, Kærnen i Kristendommen,
er der i enhver Periode, ogsaa i den første, andre
medvirkende Elementer (»Koefficienter« kalder Harnack dem),
som skyldes Tidens og Kulturens Forhold til given Tid og Sted,
og som kunne indgaa en mere eller mindre inderlig Forbindelse
med den egentlige Kærne. Hvad Urkristendommen angaar,
bestaa disse Koefficienter, der maa udskilles, forat Evangeliet
kan vedblive at virke, især af Troen paa Historiens snare
Afslutning ved Messias' Genkomst og Troen paa Undere. De
centrale Elementer i Kristendommen fremtræde derimod klart
i Bjergprædikenen, især i Saligprisningerne og Fadervor. Hertil
henviser Harnack Enhver, der vil vide, hvad Kristendom egentligt
er. Dens ledende Tanker ere Ideerne om Guds Rige, som
skal komme, om Gud som Fader, om Menneskesjælens uendelige
Værdi, om en højere Retfærdighed og om Kærlighedsbuddet
som Lovens Fylde. Men det Vigtigste er, at der bag alle
disse Tanker staar en levende og mægtig Personlighed, dyb og
klar, vældig og mild som ingen anden. I hans Tale er den
inderligste Alvor forbunden med et aabent Blik for Livets Mangfoldighed.
I hans Lignelser spejler det Største sig i det Mindste.
% --- p. 205 ---
\setcounter{footnote}{0}%
\opage{205}Alle Modsætninger falde bort eller træde i lige Grad i det ene
Fornødnes Tjeneste. Her er ingen Sygelighed og ingen ekstatisk
Overspænding, og dog den stærkeste og inderligste Fordybelse
og Koncentration. Om hans personlige Udvikling vide vi Intet,
og vi vide idethele Intet om ham før hans tredivte Aar. Men
netop dette er, ifølge Harnack, et betegnende Træk, ti deraf
maa sluttes, at ingen Kriser og indre Storme ere gaaede forud
for hans offentlige Fremtræden. Alt hvad han siger, kommer
saa umiddelbart og saa selvfølgeligt, og staar i den Grad som
det, i hvilket han lever og aander, at det ikke kan være møjsommeligt
tilkæmpet.

Hvad dette sidste Punkt angaar, er den Slutning, Harnack
drager, dog ikke ganske sikker. Det gælder saa godt som om
alle Oldtidsskikkelser, at de staa for os i deres hele, fuldt udviklede
og harmoniske Form, uden at vi faa Noget at vide om,
hvorledes de ere blevne til det, de ere. Det kan i mange Tilfælde
komme af Mangel paa historiske Efterretninger; men det
hænger sikkert idethele sammen med, at man i Oldtiden ikke
havde vor psykologisk-historiske Interesse for at se Karaktererne
i deres Genesis. Oldtidens Psykologi var plastisk og intuitiv,
ikke genetisk og analyserende. Ligesaa lidt som man kan slutte,
at Sokrates uden indre Kampe har vundet den overlegne Selvbeherskelse,
med hvilken han fremtræder for os, ligesaa lidt
kan det være berettiget at drage en saadan Slutning med Hensyn
til Jesus af Nazareth.

Da Kristendommens Væsen ifølge Harnack kun forstaas ved
hele dens Historie, følger han den paa dens Vej fra Urkristendommen
gennem den apostoliske Tid, oldkirkelig Katolicisme,
romersk og græsk Katolicisme og Protestantisme. Han giver her
en Række ypperlige Karakteristiker. Af disse er maaske Skildringen
af den græske Kirke den mest vellykkede, og den vil
særligt interessere de Mange, der kende og elske den moderne
russiske Literatur. Man vil forstaa Leo Tolstoj's Stilling og Betydning
bedre gennem den. Dog burde der sikkert ogsaa være
henvist til de ejendommelige russiske Sekter, der frembyde nogle
af de mærkeligste religiøse Fænomener og belyse den Indflydelse,
Racekarakteren har paa det Religiøses Udformning.

% --- p. 206 ---
\setcounter{footnote}{0}%
\opage{206}Efter Reformationen er der ifølge Harnack ikke kommet
Noget frem, som kunde stille det religiøse Problem i en ny
Form. De Vanskeligheder, »Evangeliet« møder, ere og blive de
gamle. Denne Udtalelse maa undre enhver Læser, da det dog
er klart, at Harnacks eget Standpunkt, hele den Maalestok, han
anvender ved Sondringen mellem Kærne og Skal i Kristendommen,
netop skyldes Faktorer, der ere komne frem efter Reformationen.
Der er en stor Forskel mellem Luthers og Harnacks
Opfattelse af Kristendommens Væsen, og Harnack lægger ikke
Skjul derpaa, men giver en skarp Paavisning af Luthers Begrænsning.
Luther var for rask til at lave nye Landskirker;
han var for angst for »Sværmerne«, for ivrig til at bekæmpe
dem, for uvillig til at lære af dem. Og han stod ikke paa Højden
af den Erkendelse, der dog allerede var mulig i hans egen
Tidsalder. Særligt havde han ingen virkelig historisk Forstaaelse
af Urkristendommen. Utallige Problemer har han ikke kendt
og var derfor uden Evne til at skelne mellem Kærne og Skal,
Oprindeligt og Fremmed. --- Men naar denne Kritik af Reformatoren
staar fast, saa maa der dog paa det religiøse Omraade
være sket noget meget Væsentligt siden Reformationen! Forskellen
mellem Harnacks og Luthers Standpunkt maa skyldes
»Koefficienter«, der ere fremkomne efter Reformationen, og i
Kraft af sin egen Metode maatte Harnack være gaaet nærmere
ind paa »Evangeliets« Forhold til disse nye Kulturelementer,
end han har gjort. Og naar Harnack fra Urkristendommen udskiller
Mirakeltroen og den kiliastiske Messiastro som Koefficienter,
der ikke have med Evangeliets Kærne at gøre, saa foretager
han i Virkeligheden denne Udskillelse i Kraft af de moderne
Kulturelementer, der betinge hans Metode. Han staar
paa den moderne Naturvidenskabs, Historiekritiks og Psykologis
Grund --- i samme Øjeblik, som han erklærer dette Grundlag
for at være uden Betydning!

Her fremkommer nu et Problem af stor Betydning, som
Harnack i den foreliggende Bog ikke har fremdraget. Kan den
nytestamentlige Forestilling om Guds Rige helt skilles fra Forventningen
om dets snare, overnaturlige og ydre Realisering?
Og bestemmes ikke den urkristelige (nytestamentlige) Etik paa
% --- p. 207 ---
\setcounter{footnote}{0}%
\opage{207}de mest afgørende Punkter --- Forholdet til Familieliv, Statsliv,
Kunst og Videnskab --- netop ved denne Forventning? Hvis
dette er saa, maa den »kristelige« Livsanskuelse og Livsførelse,
der bliver mulig, efterat der paa den af Harnack angivne Maade
er skelnet mellem Kærne og Skal, blive af væsentligt anden
Beskaffenhed end den urkristelige.

Det skyldes sikkert den populære Karakter, Harnack har
givet sine Forelæsninger, at han ikke er gaaet ind paa denne
Side af Sagen. I sin store Dogmehistorie --- et af den moderne Videnskabs monumentale Værker --- udtaler han sig derimod om
det med al ønskelig Klarhed. Han anerkender der, at med de
apokalyptiske Forestillinger falder den egentlige nytestamentlige
Etik ogsaa bort. Derfor maa den kristelige Livsanskuelse suppleres,
og paa Spørgsmaalet om, hvorfra Suppleringen skal ske,
henviser han til --- Antiken og Goethe! Hvorledes denne Supplering
skal ske, saa at der kan komme et harmonisk Livshele
ud deraf, derom siger han Intet. Det har for mig været en
Skuffelse, at den nye Bog heller Intet bragte herom, ikke engang
berørte Problemet. Det forekommer mig at maatte være en
aldeles væsentlig Side af Sagen, naar det gælder at orientere
unge Videnskabsdyrkere overfor »Kristendommens Væsen«.

Hvad der efter Harnacks Overbevisning stadigt igen vil føre
tilbage til »Evangeliet«, til Kærnen i Kristendommen, udtaler
han ganske kort i sine Slutningsord: »Naar vi med fast Villie
hævde de Kræfter og Værdier, der paa vort indre Livs Højdepunkter
straale frem som vort højeste Gode, ja som vort egentlige
Selv, --- naar vi have Alvor og Mod nok til at lade dem
gælde som det Virkelige og forme vort Liv efter dem, --- og
naar vi saa se paa Gangen i Menneskehedens Historie, følge
dens opadstigende Udvikling og stræbende og tjenende opsøge
Aandernes Samfund i den, --- saa ville vi ikke synke hen i Lede
og Mismod, men vi ville vinde Vished om Gud, den Gud, som
Jesus Kristus kaldte sin Fader, og som ogsaa er vor Fader«.
--- Ja, hvad der her antydes, er Gangen i al religiøs Erfaring.
Religiøs Erfaring angaar netop det højeste Værdifulde, vi kende,
i dets Forhold til Virkeligheden. Men skulde den faste Overbevisning
om Værdiens Bestaaen ikke kunne formes og udtryk%
% --- p. 208 ---
\setcounter{footnote}{0}%
\opage{208}kes paa anden Maade og i andre Forestillinger end i dem, der blive tilbage efter den kritiske Sondring mellem Skal og Kærne
i Urkristendommen? Og ere ikke selve disse Forestillinger snarere
at kalde Billeder og Symboler end egentlige Tanker? --- Selvom disse Spørgsmaal besvares med Ja, er og bliver Kristendommen
den mægtigste aandelige Bevægelse, der har grebet ind i den evropæiske Menneskeheds Aandsliv, og dens Virkninger vil
ikke forsvinde. Men dens Forhold til Aandslivets andre Elementer er endnu mere indviklet, end selv den smukke og betydningsfulde Drøftelse, den her foreliggende Bog indeholder anerkender. Det er dog ikke de ringeste Bøger, man lukker
med et Spørgsmaal paa Læben, især naar de bidrage til, at Spørgsmaalet kan stilles klarere end før.

% --- p. 209 ---
\setcounter{footnote}{0}%
% CHECK p. 209: centred line: heading or verse
\opage{209}\essay[Tale for Georg Brandes. (1891)]{TALE FOR GEORG BRANDES.\footnote[1]{Ved Festen i Anledning af hans 25. Aars Jubilæum som Forfatter, i Oktober 1891.}}
\streg

Dette er ikke et Jubilæum som saa mange af dem, der
holdes i vore Dage, ved hvilke det eneste Mærkelige er, at
Aarene ere gaaede. Hvad der har samlet os her, er netop det,
at vor Æresgæst har formaaet at give de Aar, der nu ere gaaede hen siden hans første Optræden i Literaturen, et eget Præg gennem
sine Værker og sin Indflydelse paa det aandelige Liv. Han staar som den største Kraft i vor Literatur og som det ejendommeligste
Bidrag, vi i disse Aar have givet til Verdensliteraturen.
For denne Betragtning maa, forekommer det mig, paa
en Dag som denne alle andre Hensyn vige. Kritiken har havt
Ordet overfor ham, en mere eller mindre velvillig Kritik (især
en mindre velvillig); nu have vi kun at gøre med det, der vil
holde sig gennem Kritikens Skærsild og kan antages at blive
staaende som den egentlige Værdi af hans Arbejde.

Idet jeg søger at finde et enkelt Udtryk til at betegne, hvad
der staar for mig som Betydningen af Dr. Brandes' Virksomhed,
kommer jeg til at tænke paa Jægerhornet i »Drot og
Marsk«, som har den Evne, at det \textit{kalder Mænd}. Med et saadant
Horn er Dr. Brandes optraadt iblandt os.

Det \textit{kalder} \textit{Mænd}. Det betyder først, at det vækker dem
til Mænd, som endnu ikke ere det. Istedetfor Barnets snevre
% --- p. 210 ---
\setcounter{footnote}{0}%
\opage{210}Horizont aabnes der nu en videre, uendelig Horizont. De Skranker,
der hæmmede Blikket, falde bort; det opdages, hvor stor
Verden er, hvor utallige Muligheder der findes, hvor mange
Kræfter, der ere i fuld Virksomhed derude. Men Manden adskiller
sig fra Barnet ikke blot ved sit videre Blik, men ogsaa
ved sin større Selvstændighed og ved den større Klarhed over
sig selv. I Dr. Brandes' Værker finder jeg begge Bestræbelser,
baade den for at udvide Blikket og gøre os delagtige i den
fremmede Udvikling, og den for at føre os tilbage, lære os at
finde os selv. Han vil, at vi gennem Fordybelsen i det Fremmede
og Mangfoldige skulle berige os selv og tillige blive klarere
over vor egen ejendommelige Kraft og faa den vakt. Og hvis
nogen af de to Retninger har Overvægt over den anden, da er
det sikkert Bestræbelsen for at frigøre og udvikle den enkelte
Personlighed i dens Ejendommelighed. Det er denne Bestræbelse,
der ligger bag hans store Skildring af et halvt Aarhundredes
aandelige Historie, saaledes som den spejler sig i Literaturen.
Hans enestaaende Evne til at karakterisere, til at finde
det personlige Livs Kilder hos de Forfattere, han skildrer, har
en stadig Adresse til Læseren, paaminder om, at det at virke
ud af sin inderste Natur er det Højeste, man kan naa --- hvor
højt eller hvor ringe det end i sig selv maatte være. Der aabenbarer
sig i hans hele Virken en Higen efter at frigøre Individernes
Liv fra Alt hvad der hæmmer og forkvakler, og en
Tillid til det, der vil kunne udvikle sig, naar Frihed for Alvor
kommer til at raade paa det aandelige Omraade. I denne stadige
Higen og i denne urokkelige Tillid ser jeg den største
Adkomst til vor Tak. Mere end nogensinde gælder det sikkert
nu om, at den Livskilde, vi have i Friheden, ikke tilstoppes.
Udskejelserne og Karikaturerne, som ikke ville udeblive, komme
vi nok ud over.

Men ogsaa i en anden Betydning kalder Hornet Mænd.
Det har Bud til Mange, som selv vare Mænd samtidigt med
ham, og som arbejde paa andre Omraader, men følge hans Virken
med Sympati og paa deres Vis føle sig beslægtede med
den. En kraftig Rørelse paa et vigtigt Omraade af det aandelige
Liv forplanter sig til mange andre Omraader, ligesom Bølge%
% --- p. 211 ---
\setcounter{footnote}{0}%
\opage{211}slaget forplanter sig fra Søen ind i stille Bugter og Vige. Det
har da ogsaa gentagne Gange vist sig, hvorledes Mænd af meget
forskellig Retning og fra meget forskellige Omraader have kunnet
aflægge deres Vidnesbyrd om det Liv og den Vækkelse, Dr.
Brandes har formaaet at bringe ind i vort Aandsliv. Den Kreds,
som her er samlet, vidner om det Samme. Selv de, der kæmpe
under andre Faner, kunne dog anerkende den Betydning, det
har, at han er i Marken, og at hans Signal lyder.

Endeligt --- naar Hornet kalder Mænd, kan det ogsaa være
til Strid, det kalder, idet dets Toner lyde som en Udfordring.
Hvad der vækker til Liv, ægger ofte til Modsigelse, og hvor
selvstændige Mænd mødes, kan det ikke gaa af, uden at Tanke
kommer til at staa mod Tanke, Villie mod Villie. Et er, at Signalet
lyder til Fremrykning, et Andet, om Vejen, ad hvilken
det byder at gaa, er den rette, virkeligt fører fremad. Derom
kan der være Strid. Hyppigst og med størst Bitterhed er Dr.
Brandes bleven angreben af dem, der ingen Fremrykning ønske.
Jeg tror at have vist ham en stor Anerkendelse, idet jeg er gaaet
ud fra, at disse Angreb ikke havde taget mere paa ham, end
at han havde Kraft nok tilovers til at møde en Modstander, der
kom fra en anden Kant end den sædvanlige. --- Jeg skylder, i
Forbigaaende sagt, Dr. Brandes, at jeg er kommen ud af en
betænkelig Situation. En Kritik, jeg rettede mod nogle Udtalelser
af ham, blev, efter hvad jeg erfarede, tagen til Indtægt fra
en vis Side, der altid har havt et godt Øje til ham. Derved
var jeg ilde til Mode. Men da Dr. Brandes' Svar kom, blev
det fra samme Side taget til Indtægt mod mig, idet man var
glad ved at kunne slaa to Fluer (ganske vist ikke med ét Smæk).
Saaledes bortfaldt den betænkelige Situation. --- Idet jeg minder
om Fortidens Strid, er det ikke min Mening, at Striden blot
hører Fortiden til. Hvem véd, om der ikke kan blive Anledning
til ny Strid? Det er ikke af Stridslyst, jeg antyder denne Mulighed.
Men Baggesen har med Rette sagt: »I Aandens Verden
bør der være Strid!« Det er mageligere at slutte Trop og rulle
afsted som Del af en samlet Masse; men saaledes kommer man
ikke i Virkeligheden fremad: det har vor Æresgæst lært os!

Jeg for mit Vedkommende er kommen her for  at sige Dr.
% --- p. 212 ---
\setcounter{footnote}{0}%
\opage{212}Brandes Tak for, hvad han har virket for det aandelige Liv i
vort Fædreland. At vække levende Tilslutning og levende Modsigelse
og føre fremad paa begge Maader, det er Særkendet for
den store, den betydningsfulde Virken. Jeg vil paa denne Dag,
da Takken bringes ham fra de Unge i saa rigt og saa velfortjent
Maal, bringe ham en Hilsen fra Jevnaldrende, fra dem,
der havde de samme Lærere som han, og som Mænd have
fulgt hans Færd, snart med Tilslutning, snart med Modsigelse,
men aldrig med Ligegyldighed. Jeg tror at tale i Manges Navn,
naar jeg siger, at den Dag, Klangen af hans Horn ikke høres
mere, vil en af de betydeligste og ædleste Kræfter i det danske
Aandsliv i vort Arhundrede have ophørt at virke!

% --- p. 213 ---
\setcounter{footnote}{0}%
% CHECK p. 213: centred line: heading or verse
\opage{213}\essay[Tale for Henrik Ibsen. (1898)]{TALE TIL HENRIK IBSEN.}
\begin{center}{\small VED EN FEST I STUDENTERSAMFUNDET I ANLEDNING AF}\\
{\small HANS 70 AARS FØDSELSDAG.}\\
{\small DEN 2. APRIL 1898.}\end{center}
\streg

Kære Dr. Henrik Ibsen!

Det er Festernes og Triumfernes Tid for Dem. Men forud
ligger en lang Kamptid. De har kæmpet med Livets og Tidens
Problemer. De har kæmpet for at fyldestgøre Kunstens store
Krav og for at skabe nye Former for kunstnerisk Fremstilling.
Og De har kæmpet for at overvinde Læsernes og Tilskuernes
Modstand.

Lad mig først tale om denne sidste Kamp. At der blev
en Kamp, havde sin Grund i, at De har stillet større Fordringer
til Deres Læsere end nogen anden Forfatter i vor Tid. De
giver ikke fuldt færdige Løsninger eller Afgørelser. Der bliver
stedse Spørgsmaal og Gaader tilbage, og Læseren maa ofte lukke
Bogen med en Brod i sin Sjæl. Deres Digtning er gaadefuld,
som Livet er det. Enhver maa indarbejde i sin Personlighed,
hvad De giver. Fortolkningen eller Udfyldningen bliver hans
eget Værk og beror paa, hvilke Strenge der sættes i Bevægelse
hos ham. Naar Tæppet gaar ned for Deres Stykker, begynder
der et nyt Drama i Tilskuerens Sind; han nødes selv til at optræde,
skønt det ikke er paa nogen ydre Scene. De forlanger
et stort Arbejde af Læsere og Tilskuere. Derfor har man fundet
Dem streng og kold; man har ment, at De tvang Læser
% --- p. 214 ---
\setcounter{footnote}{0}%
\opage{214}eller Tilskuer ind i Kampen med Disharmonier og Problemer,
men holdt Dem selv udenfor. Og man ansaa Dem for livsfjendsk
og folkefjendsk, fordi De ikke gav harmoniske Løsninger og Afslutninger,
tydelige Svar, man kunde bringe hjem med fra Læsningen
eller Skuespillet. Derfor blev De et Modsigelsens Tegn.

Naar nu saa Mange fra saa mange Lejre samles om Dem
for at hylde Dem, saa er det ikke blot, fordi man har lært at
beundre og elske den store Kunst i Deres Værker. Det er ogsaa,
fordi man har faaet Syn for, at disse Deres Værker vidne om,
at De ikke har unddraget Dem Kampen med Livets Disharmonier
og med Tidens Problemer. Man har lært at betragte Dem
som Medkæmper i Tidens store Strid, bærende Deres Del af
Dagens Møje.

De har aldrig sværmet for Solidariteten. De har i den set
en Fare for Personligheden. Men Deres Værker vise, at De har
følt Dem solidarisk med dem, der lide og kæmpe. De har erfaret
Mørket og Tvivlen, Kuetheden og Bitterheden. I klare Billeder,
der aldrig ville glemmes, har De stillet os en Række Menneskeskikkelser
og Menneskeskæbner, der ere hentede fra Tidens Liv,
for Øje. Kunstneren skal ikke virke gennem Sætninger og Slutninger.
Han virker gennem Billeder og individuelle Skildringer,
og han maa overlade til Læser og Tilskuer at afgøre, om
det er Regler eller Undtagelser, Forbilleder eller Advarsler, der
gives gennem hans Fremstillinger. Saaledes giver jo Livet selv
os sine Fænomener og Begivenheder: vi maa selv finde ud,
om de ere Regel eller Undtagelse, om de skulle efterlignes eller
skyes.

Naar De saa sjeldent har talt Haabets Ord, --- De siger
etsteds:

% CHECK p. 214: centred line: heading or verse

\begin{verse}
Duer har jeg ingen af; \\
% CHECK p. 214: centred line: heading or verse \\
Duer er jo Haabets Fugle ---
\end{verse}

saa er det maaske, fordi De har ment, at Haab er et Lægemiddel,
der maa tildeles i smaa Doser, da Menneskene ere saa
tilbøjelige til at svælge i det. Eller det er maaske, fordi det
sande Haab maa vindes gennem tilkæmpet Forstaaelse. De følger
i Deres Værker Begivenheders og Karakterers Udvikling
Skridt for Skridt tilbage til det afgørende Punkt, hvor Grunden
% --- p. 215 ---
\setcounter{footnote}{0}%
\opage{215}til alt Følgende blev lagt. De sætter Fingeren paa det Sted,
hvor Sygdommen tog sin Begyndelse. Det er, som om De sagde
til os: den Strøm, der ikke kan stanses eller afledes, naar den
har voxet sig stærk, kunde maaske være bleven stanset eller
afledet ved sit første Udspring! Dette er Deres Maade at lære
os Haab paa: gennem Forstaaelse af Disharmoniernes Aarsag
lære vi, hvorledes Livet kan gøres harmonisk; gennem Forstaaelse
af Undergangen lære vi, hvorledes Livet kan bestaa i Kampen.
Da De rejste i Ægypten, saa De Tempelruiner og omstyrtede
Gudebilleder rage op af Ørkenens Sand ligesom Skeletter af
Karavaner, der vare omkomne ved Sandflugt. De spurgte Dem
selv: hvorfor gik alt dette til Grunde, saa at det end ikke, som
saa meget Andet, der maatte gaa til Grunde, lever i Menneskers
Erindring og Tanker? --- og Deres Svar var: denne Kultur
manglede Personlighed og derfor sand Livskraft. Denne Forstaaelse
bliver saa Grundlaget for Deres Haab:

\begin{verse}
Paa begravne Karavaner \\
% CHECK p. 215: centred line: heading or verse \\
Bygger jeg vor Fremtids Baner!
\end{verse}

--- Deres tredie Kamp, Kampen om den kunstneriske Form,
hænger sammen hermed. For at kunne fremstille Begivenhedernes
og Karakterernes Udvikling har De indført et nyt kunstnerisk Middel, den tilbagepegende, forklarende og opgørende Samtale.
De har indført Perspektivet i den dramatiske Kunst. Under
Replikernes stille, men faste Gang føres vi fra Øjeblikkets Situation
tilbage i Personernes Fortid. Karakterer og Begivenheder
udfoldes for os fra deres første Anlæg. Gennem Løngange føres vi til et Punkt, hvor store Udblik over Menneskelivet og dets
Kaar blive mulige. Hvilken Skat af psykologisk Iagttagelse og
praktisk Livsvisdom er der ikke nedlagt i disse Samtaler!

Den mægtige Billed- og Tankeverden, De har skabt, fængsler
vor Sans og vor Fantasi. Men den gives os ikke blot til
Beskuelse. Den bringer Røre i Sindene; den lader os forstaa
Skæbnens Gang i os og udenfor os; den vækker vor Selvvirksomhed
til Modstand, til Bortryddelsen af det, der kuer og
hæmmer. Derved forbereder den os til et nyt og kraftigt Tag
paa Livet og dets Forhold. ---

% --- p. 216 ---
\setcounter{footnote}{0}%
\opage{216}Lad os da nu takke \textit{Dem} for Deres trofaste Arbejde. De
har i de lange, ensomme Arbejdsaar kun spurgt om Et, om
hvad den sande Kunst krævede af Dem, og ved at være tro
mod Kunsten har De været tro mod Dem selv. Vi takker Dem
ikke blot for Deres Værker, men ogsaa fordi De ved Deres
urokkelige Fastholden af det store Maal, De havde sat Dem,
er bleven et Forbillede for Enhver, der arbejder for Sandhed
og Skønhed i en Tid, der saa ofte synes at være optagen af
helt andre Ting. At De har kunnet føre Kampen til Ende og
har fremtvunget Anerkendelsen af Deres Livs Gerning, er ikke
blot Noget, vi ønske Dem til Lykke med. Det er for os et
Grundlag for Haab, et Lysglimt til at forstaa Livet. Deres eget
Livs Drama fortoner sig henimod sin Afslutning lysere, end saa
mange af Deres digtede Dramaer gøre. Maatte dette Lysglimt
følge Dem, naar Festens Dage afløses af Arbejdets Dage, for
Dem de kæreste Dage, og maatte der endnu falde mange frodige
Arbejdsdage i Deres Lod!

% --- p. 217 ---
\setcounter{footnote}{0}%
% CHECK p. 217: centred line: heading or verse
\opage{217}\essay[Anmeldelse af Bjørnson's »Poul Lange og Tora Parsberg«. (1898)]{BJØRNSTJERNE BJØRNSON:\\ POUL LANGE OG TORA PARSBERG.}
\streg

»Undertiden stiger en Gud ned paa Jorden«, siger Poul
Lange i Bjørnsons nye Drama. Guden er her en Kvinde, der
hjælper Lange til at forstaa sig selv og til at kæmpe mod den
indre og ydre Modstand, han saa let bukker under for. Om
denne Kvinde siges det, og hun gentager det selv, at hun altid
forsøger sig paa det Umulige. Men et gammelt Ord siger, at
mod det Umulige kæmpe endog Guderne forgæves. Og det erfarer
Tora Parsberg ogsaa.

Skal Indholdet af det nye Drama angives i faa Ord, kan
det maaske angives paa denne Maade. Og herved er tillige udtrykt
det Middel, Digteren har brugt for at lade Poul Langes
Karakter træde frem for os, --- det eneste Middel, ved hvilket
en Karakter som hans kunde bringes til Aabenbarelse. Inderlig,
hensynsfuld, samvittighedsfuld, men tillige ængstelig overfor
Andres Dom, især overfor deres Latter og deres Haan, er han
et af de Mennesker, der kun lukke sig op efter en haard Kamp.
Han er en »undseelig Gjemmer«. I sin Begejstring for store
Ideer og Formaal er han tilbøjelig til at overse og ringeagte
de Midler, som Mennesker, der kalde sig selv praktiske, anse
nødvendige for at sætte Noget igennem. Har han saa selv af
Uerfarenhed eller i Nød ladet sig forlede til at benytte Midler,
der kunne lastes eller spottes, føler han hele sit Liv igennem stedse
% --- p. 218 ---
\setcounter{footnote}{0}%
\opage{218}paany Brodden derved. Ti paa Grund af hans Tilbøjelighed til
Selvreflektion fornys hans Fortid, især de mørke Partier af den,
stedse for ham i pinlig Erindring. Han er desuden vant til at
sammenligne sig med de største Aander og derfor til at bedømme sig selv strengt. Den lammende Selvkritik, der saa let
fører ham til Selvopgivelse, stammer ikke blot fra, at han er saa
lydhør for Andres Dom og Spot, men maaske især fra den
ideale Maalestok, han anvender paa sig selv. Han er tilsidst
kommen dertil, at han ikke betragter sig selv som normal, ---
at han frygter for at være »en Pøbel«. Denne kuede Selvfølelse
og den Angst, den røber, er en Hemmelighed, han ruger over.
Kun til enkelte Tider, naar han gaar op i Arbejde, eller naar
der ydes ham en hjertelig Haandsrækning udefra, kan han
glemme det indre Mørke og se lyst paa Fremtiden. Han har
en Trang til at handle efter sin egentlige, indre Natur, og han
handler i afgørende Øjeblikke ogsaa efter den; men der er i
ham en Understrøm af Selvreflektion og af matte Tvivl, der underhuler Grunden i hans ædle Natur, saa den Bygning, der kunde
have strakt sine Spir mod Himlen, undergraves og synke 
sammen.

Der maatte en Kvinde som Tora Parsberg til, forat en
saadan Natur kunde aabne sig. Hun har fulgt ham gennem
hans Liv, som Tilskuer paa nært Hold i hans Ungdom og Manddom; hun har været hans Fortroliges Fortrolige; og hun fik
fra hans Færd i mindre og større Kredse et Ideal af en Mand,
som hun bevarede i sit stærke og trofaste Hjertes Inderlighed,
indtil det Øjeblik kom, da hans Erklæring om, hvad hun nu
var bleven for ham, og hans Bøn om, at hun vilde forene sin
Skæbne med hans, gav hende Mod og Ret til at gribe ind i
hans Liv. I de store Samtaler med ham i første og tredie
Handling lærer hun ham at kende sig selv. Disse Samtaler forløbe
ikke saaledes, at Poul Lange selv føres til at sige, hvordan
han er; det er overvejende hende, der forkynder ham, hvad
der er hans egentlige Natur. Hun lader ham neppe tale ud,
saa sikker er hun i sin Forstaaelse af ham. Den dybeste Afgrund
i ham, den, det viste sig umuligt at slaa Bro over, har hun
dog ikke opdaget; Angstens, Selvkritikens og Svaghedens Dybde
% --- p. 219 ---
\setcounter{footnote}{0}%
\opage{219}i hans Natur anede hun ikke. Men ved sin inderlige Forstaaelse
og Medfølelse lærer hun ham, hvorledes netop det Svage og
Syge i hans Natur hænger sammen med det Fine, Ædle og
Store hos ham, hvorledes Forfølgelsen og Latteren, der rammer
ham, beviser hans Betydning, ikke --- som han frygter --- hans
Usselhed. Overfor Øjeblikspolitikernes fanatiske Mistænksomhed
og dumme Had viser hun ham, hvor fast et Tilhold han har
i sit Arbejde for Tanker og Opgaver, der betinge Slægtens
Fremtid og stedse paany ville fremstille sig for menneskelige
Samvittigheder.

Det er den ideale Poul Lange, saaledes som hans Billede
havde formet sig i hendes kraftige og dybe Sind, som hun fremholder
for den virkelige Paul Lange. Men det er nu en Gang
saa, at det alvorligste Problem i vort indre Liv er, om vi kunne
trøste os til at fastholde vort Jeg, som det er i vore bedste
Øjeblikke, eller om det Mørke, Svage og Lave, som vi ogsaa
kende i os, skulde være vort egentlige Jeg. Dette Problem
havde Poul Lange ikke kunnet løse selv; han faar nu i Tora
Parsberg den bedste Hjælper dertil, som tænkes kan. Maaske
er hendes Overlegenhed skildret med for stærke Træk. Hendes
Ord blive undertiden store, mægtige Toner, der fare hen
uden ret at individualiseres til Repliker efter den dramatiske
Kunsts strengeste Fordringer. Det drister jeg mig dog ikke til
at afgøre. Men Ingen vilde sikkert for større teknisk Fuldkommenheds
Skyld have undværet den vældige Energi, hvormed
de højeste Tanker, der kunne holde Livsmodet oppe, her strømme
frem. Det er Toner fra en Æolsharpe, der kunne erstatte Savnet
af adskilligt Filigranarbejde, man har bildt sig ind var Poesi.

Min Pen er ikke smidig nok til at skildre den Rigdom af
enkelte smaa Træk, med hvilke Digteren har formaaet at gøre
de to Hovedpersoner i Dramaet levende og anskuelige. Min
Bestræbelse er kun at fremdrage de Grundlinier, jeg har fundet
i den digteriske Karakteristik. Et Kunstværk er og bliver jo
dog alligevel tilsidst inkommensurabelt  for Analyse.

Et helt andet Billede af Poul  Lange end det, Tora Parsberg
drager frem, tegner sig af ham hos dem, der føle sig
% --- p. 220 ---
\setcounter{footnote}{0}%
\opage{220}modarbejdede eller skuffede ved hans politiske Optræden. Chefen
for det Ministerium, i hvilket han sidder, har længe vist sig
rænkefuld og troløs overfor ham. Det virkede saaledes paa ham, at han næsten blev sindssyg og var nærved at begaa Selvmord --- Nu har han indgivet sin Ansøgning om Afsked, og Spørgsmaalet
er, om han vil støtte sin gamle Chef overfor et Angreb i
Tinget. Han føler, at det stemmer med hans Natur at gøre det…
Ministerchefen er en af sit Folk højt fortjent Mand; han er desuden
en gammel Mand, som neppe har langt tilbage. Og fremfor
alt --- Poul Lange holder ikke af offentlige Afstraffelser; det stemmer
ikke med hans Natur at være Skarpretter. Han vil
helst positivt arbejde for det Nye og vender Ryggen til det Daarlige
og Usle. Han vil bevare sin Kraft til bedre Tider. Dog svarer han undvigende paa en Opfordring fra Hoffet om at forsvare
Ministerchefen og paa det Tilbud om Gesandtskabsposten
i London, som i Forbindelse dermed gøres ham. Og da en af
hans nærmeste Venner indstændigt opfordrer ham til at blive
borte paa det afgørende Møde i Tinget, lader han sig bevæge
til at love dette. Straks derefter finder den første store Samtale
med Tora Parsberg Sted, i hvilken hun aabner Lykken for
ham, vækker hans Selvtillid og styrker ham i Trangen til at
følge sin inderste Natur. Da han antyder, at man vil have ham
til at gøre Noget, der egentligt strider mod hans Natur, siger
hun: »Saa Andre vil sige dig, hvad som falder dig naturligt?«
Sejrsstolt forlader hun ham uden at ane, at hun har givet ham
en falsk Tryghed. Naar han ikke underretter Arne Kraft om,
at han efter nærmere Overvejelse tager sit Løfte om at blive
borte fra Mødet tilbage, uagtet der gaar fire Dage hen imellem
maa det vel forklares ved, at han i sin Lykkerus overser Nødvendigheden
deraf. Han, der ellers har for mange Tvivl, har her
for faa Betænkeligheder.

Efter at han i det afgørende Møde har forsvaret sin gamle
Chef som den, der vedvarende bedst formaar at samle Folket,
deltager han om Aftenen i en Fest, hvor mange af Stortingets
Medlemmer ere tilstede --- og her bliver han da selv Genstand
for en af de »offentlige Afstraffelser«, han ikke havde villet være
med til, naar det gjaldt Andre. Man vægrer sig ved at række
% --- p. 221 ---
\setcounter{footnote}{0}%
\opage{221}ham Haanden, kræver ham strengt til Regnskab og erklærer,
at han har mistet de Bedstes Tillid. Da ogsaa hans nærmeste
Ven deltager i denne Fordømmelse, synker han sammen.

For dem, der dømme ham, staar hans Handlemaade som
uforklarlig af andre Bevæggrunde end Stræben efter egen Fordel.
»Han véd nok, hvad han vil!« Sandheden, erklæres der,
maa i Politiken vente, til man har Brug for den. Der er i
hvert Øjeblik en enkelt Sag, som ligger for, og den, der hindrer
den, skal hensynsløst kastes til Side. Den mildeste Forklaring af hans Færd er i denne Kreds, at den udspringer af
sygelig Sentimentalitet.

Mesterligt er det i Dramaet gennem en Række af Enkeltheder,
Antydninger og Sammentræf paavist, hvorledes Langes
Karakter kunde tage sig ud for Øjeblikspolitikerne paa denne
Maade. Grunden til det Vrængebillede, der har dannet sig hos
dem, ligger dels i det Fine og Ædle i Langes Karakter, dels i
hans Svaghed, Vankelmod og Ubetænksomhed, dels i hans Modstanderes
udvortes og fanatiske Blik.

\stjerner

Efter det voldsomme Sammenstød er Paul Lange rystet i
sit Inderste. De gamle Tvivl dukke op; Angsten for at være
abnorm, for at være Pøbel rejser sig igen i hele sin Vælde,
som et ustanseligt Ekko af de fanatiske Ord, der have lydt til
ham. Mørket breder sig for ham paa alle Veje. Selv den lyse
Vej, Tora Parsberg havde aabnet ham, tør han ikke stole paa:
hvor kan han drage hende med ind i sin Skændsel, drage
hende ned i sit eget Mørke? Den anden store Samtale med
Tora Parsberg, der nu atter opsøger ham, hæver ham vel igen
op efter en stærk aandelig Kamp mellem de to Personligheder,
en Kamp mellem Haab og Fortvivlelse, mellem Liv og Død.
Hun forlader ham atter tryg, saa tryg, at hun end ikke forsøger
at faa ham til at udlevere Pistolerne, som hun véd, han
har i sit Gemme. Her viser sig atter Grænsen for hendes Forstaaelse
og for hendes Kraft; til det lurende Mismod og den
utrættelige Tvivl, der havde undergravet Grundlaget for hans
Personlighed, var hun ikke trængt ned, og hun havde vel heller
ikke kunnet overvinde disse dunkle Kræfter fuldstændigt. Neppe
% --- p. 222 ---
\setcounter{footnote}{0}%
\opage{222}har hun forladt ham, før disse Kræfter igen begynde at røre
sig. »Hvad maa hun dog i Grunden tænke? Ikke i Dag, men
den næste Dag? Naar Sejrsfølelsen er ovre? --- Ved Gud, jeg
taaler ikke den mindste Skygge af Ydmygelse mere!« Men i
samme Øjeblik kommer en ny Ydmygelse; et Telegram bringer
Efterretningen om, at hans Fjender have formaaet at tilføje
ham en ny offentlig Krænkelse. Da flyder Bægeret over, og
han retter det dræbende Skud  mod sit Bryst.

Dette Drama vil ikke forstaas af dem, der ikke have lært,
hvorledes det Svage og det Store, det Tunge og det Lette, det
Lave og det Ædle kan bo sammen i vor Natur saaledes, at et
Menneskes Handlemaade ikke kan forklares ud af en enkelt
Motivrække alene eller vurderes ud fra et eneste Synspunkt.
Det er en af Skyggesiderne ved det offentlige Liv, som det
former sig i vore Dage, at Forklaring og Bedømmelse af Mennesker
foregaar saa hurtigt og paa Grundlag af ydre, ved Øjeblikkets
Forhold bestemte Situationer. Digteren har vist os,
hvilke Tragedier deraf kan udspringe, naar det i Dybden ser
helt anderledes ud end paa Overfladen, uagtet Dybde og Overflade
dog høre nøje sammen.

En digterisk Fremstilling, der har naat sit Maal, har en
individuel Karakter, har skildret en enkelt bestemt Type; men
der kan ikke af den udledes almindelige Sandheder. Det vilde
være urigtigt at udlede af Paul Langes Skæbne, at Troskab
mod det Parti, man har sluttet sig til, ikke skulde være Pligt.
Dramaet peger kun hen paa, at det ikke kan være den højeste
Pligt; den Enkeltes Stræben efter Overensstemmelse med sin
inderste Natur kan føre til Brud med Partimoralen. Der gives
en Helligdom, som den Menneskekundskab og den Menneskevurdering,
der kan erhverves i det politiske Liv alene, ikke
aabner Adgang til. Stærke Naturer formaa at bevare denne
Helligdom og maaske endog tvinge Andre til at anerkende
den; svage og vaklende Naturer gaa under, fordi de lade Skrigene
udefra genlyde saa voldsomt i Helligdommens Indre, at
den indre Røst ikke kan høres. Men selvom Templets Mure vare
% --- p. 223 ---
\setcounter{footnote}{0}%
\opage{223}for svage, kunne de dog have indesluttet virkelige Gudebilleder ---
og det har Digteren her vist os.

Dersom Digteren har havt et Forbillede for sin Hovedfigur,
saa har han overfor ham ydet den Tjeneste, den døende Hamlet
bad sin Ven om:

\begin{verse}
Forklar mig og min Sag for dem, som ej forstaa, --- eee hele Sagen, det Mindste med det Største!
\end{verse}

% --- p. 224 ---
\setcounter{footnote}{0}%
% CHECK p. 224: centred line: heading or verse
\opage{224}\essay[Julius Langes Breve. (1902)]{JULIUS LANGES BREVE.}
\begin{center}{\small (FORORD TIL DEN TYSKE UDGAVE).}\\
(1902).\end{center}
\streg

Julius Lange har ikke blot Betydning som en udmærket
Forsker paa Kunsthistoriens Omraade, men ogsaa ved sin hele
lyse og varme Personlighed, der gennemtrængte al hans Færd
i Livet som i Videnskaben. Deraf kommer det, at hans kunsthistoriske
Skrifter ikke blot yde rig Belæring om Kunstens
Værker, deres Oprindelse og deres indbyrdes Forhold, men
ogsaa, dels gennem direkte Udtalelser, dels, og ikke mindst,
ved den Grundstemning, hvoraf Fremstillingen bæres, indeholde
en Skat af Livserfaring og sund Tænken over menneskelige Forhold.
Det vil derfor have sin store Interesse at lære ham at
kende fra en mere personlig Side, end hans videnskabelige Værker
gøre det muligt, --- at lære ham at kende, saaledes som
han var, naar han færdedes mellem sine Nærmeste og sine Venner. Og saaledes lære vi ham at kende i hans Breve.

Vi se Lange i disse Breve som det sunde, livskraftige og livsmodige
Menneske, betragtende Livet, ligesom han betragtede
Kunsten, med barnlig Umiddelbarhed og dog med Tænkerens
og det prøvede Menneskes dybe Forstaaelse. Der faldt mange
Skygger paa Langes Vej i de senere Aar, og hans Livsanskuelse
gav ham ingen dogmatisk Trøst. Med mandig Forstaaelse og
Energi fastholdt han det, der havde givet Livet Værdi for ham,
og arbejdede for det til det Sidste, midt under en smertefuld
og forstemmende Sygdom. Vi møde i disse Breve ikke mindst
% --- p. 225 ---
\setcounter{footnote}{0}%
\opage{225}talrige Exempler paa hans kostelige Humor, der ofte hjalp ham
selv og Andre over en mørk Time eller kastede Glans over en
festlig Sammenkomst.

Det at have kendt Lange personligt og at have hørt til
den Kreds, i hvilken han frit aabnede sin Aands og sit Lunes
Skatte, regnes af hans Venner for en af deres Livs bedste Ting. Gennem hans Breve kan nu Noget af dette Gode blive en større
Kreds til Del.

Dertil kommer, at Brevene naturligvis for en stor Del
handle om det, der laa Lange selv mest paa Hjertet, Kunsten
og dens Historie. Vi høre om hans Planer, hans Rejser og
hans Udkast. I improviserede Udtalelser kommer Meget af det
frem, som han senere indarbejdede og udførte videre i sit Livsværk, eller det, der skulde have været hans Livs store Værk,
men kun blev en Række af Brudstykker, --- »Billedkunstens
Fremstilling af Menneskeskikkelsen i Historiens forskellige Perioder«.


Lange har engang defineret Kunstværdi som den Værdi,
et Emne har for Kunstneren, og som det gennem hans Fremstilling
faar for Andre. Denne Definition (mod hvilken der
rent æstetisk maa gøres forskellige Indvendinger) er karakteristisk
for Lange som Kunsthistoriker. Gennem Studiet af Billedkunsten
i dens historiske Hovedperioder vilde han opdage,
hvorledes --- det vil sige, med hvilke Øjne og i hvilken Stemning
--- Menneskeskikkelsen er bleven opfattet. Og Opfattelsen
af Menneskeskikkelsen hviler igen paa Opfattelsen af Livet
og Livets Forhold. Saaledes fik Studiet af Kunsten for Lange
sin store menneskelige Betydning --- et væsentligt Bidrag til
Forstaaelse af Menneskelivet i dets Udvikling.

Men hin Definition passer desuden paa Langes egne Værker
og paa selve hans Personlighed. Paa hans Værker: ti gennem
dem faar Billedkunsten forøget Værdi for os formedelst
den Værdi, den har havt for Lange selv. Paa hans Personlighed:
ti gennem den styrkes og beriges vor Følelse ved Livet,
idet vi se, hvad Lange ved sin Energi, sin Humor, sin Varme
og sin klare Forstaaelse har faaet ud af Livet.

% --- p. 226 ---
\setcounter{footnote}{0}%
% CHECK p. 226: centred line: heading or verse
\opage{226}\essay[Goethe. (1899)]{GOETHE.}
\begin{center}(1899).\end{center}
\streg

Den 28. August 1899 er det halvandet hundrede Aar, siden
Goethe fødtes. Manges Tanker ville paa den Dag søge at klare
den Gæld, de staa i til denne store Genius. Hans Digtning har
vævet sig ind i vort Følelses- og Tankeliv, og hans Personlighed
staar for os som et mægtigt Forbillede for harmonisk, men
tillige dyb og alvorlig Livsførelse. Aldeles ublandede ere de Indtryk,
vi modtage af ham, naturligvis ikke. Der er mange af
hans Digte, som vi maa lade ligge, fordi de skyldes hans svage
Timer, eller fordi de bevæge sig i Former og Stemninger, hvis
Tid er forbi. Og der er i hans Livsførelse Mangt og Meget,
som maaske kan forstaas og undskyldes, men i hvert Tilfælde
ikke beundres. Og dog rager trods Tidens og Kritikens fordunklende
Magt den Dag idag hans Personlighed og hans Værk
op som en sollys Bjergtinde og betegner et af Højdepunkterne
for menneskeligt Aandsliv.

\section*{I.}

Hvor højt han rager op i Sammenligning med andre Mestre,
er det i en vis Forstand umuligt at afgøre. Store Aander kunne
ikke uden videre sammenlignes, fordi det i Forholdet mellem
dem ikke blot beror paa et Mere og Mindre, men ogsaa paa
Forskelle i Henseende til hvad og hvorledes: de høre til forskellige
Ordener, ligge ikke alle paa samme Linie, saa Afstan%
% --- p. 227 ---
\setcounter{footnote}{0}%
\opage{227}den mellem dem uden videre skulde kunne maales. Anlægge
vi den højeste Maalestok, ville vi dog vist i den første Række
indenfor Digtningens Verden kun træffe to store Navne: Homer
og Shakespeare. Deri vilde Goethe selv være enig med os. Da
man i den romantiske Periode forsøgte at stille Tieck op ved
hans Side som jevnbyrdig, fandt han det uretfærdigt, --- ligesaa
uretfærdigt, som hvis han selv vilde stille sig ved Siden af
Shakespeare. »Tieck«, sagde han, »er et Talent af stor Betydning;
men naar man vil stille ham lige med mig, farer man
vild. Jeg kan sige dette ligefrem, thi hvad angaar det mig?
Jeg har ikke skabt mig selv. Det var ligesom jeg vilde sammenligne
mig med Shakespeare, som jo heller ikke har skabt
sig selv, og som dog er et Væsen af højere Art, til hvem jeg
ser op og for hvem jeg bøjer mig i Ærefrygt.«

Her have vi strax et Træk, som er ejendommeligt for Goethe:
den sunde Selvfølelse, forenet med klar, paa dyb Forstaaelse
grundet Selvbegrænsning og med Forvisningen om, at der til
Grund for vor højeste Virken ligger Kræfter, som vi ikke selv
have givet os. I Begrænsningen laa for Goethe Mesterskabet, og
skønt den Energi, der virker indenfor Begrænsningen, er vor
egen, er Et med os selv, saa er den dog ikke vort eget Værk
--- saalidt som de Grænser, vi finde for vort Væsen, er vort
eget Værk. Goethes hele Opfattelse af Liv og Kunst bæres af
denne Tanke, som det er umuligt at udtrykke i et enkelt Ord.
Ordene blive idethele let for tunge, naar man skal tale om
Goethe. Han formaar at forene i sig Egenskaber, Bestræbelser
og Synsmaader, der hyppigst staa som uforenelige, og for hvis
Forening der ikke kan findes noget udtømmende Udtryk. Vi
maa stave os frem for at forstaa ham.

\section*{II.}

Hans Tro paa sig selv og hans Stræben efter at udvikle
sig og at drage Næring af Alt, hvad der møder ham paa hans
Vej, af Mennesker saavel som af Natur og Kunst, kan ofte se
ud som Egoisme, og dog ansaa han ikke sin Lyst og sin Higen
som berettiget afset fra Livets og Ideernes store Sammenhæng. 
% --- p. 228 ---
\setcounter{footnote}{0}%
\opage{228}Han betragtede ikke sig selv som Formaal og alt Andet og alle
Andre blot som Middel. Hvis man vil kalde ham Aandsaristokrat,
saa behøver dette Ord i hvert Tilfælde en nærmere Forklaring.
Med Stolthed kunde han paa sine gamle Dage sige, at
han i sit Kald som Forfatter aldrig havde spurgt, hvad den
store Masse vilde, og hvorledes han kunde gavne det Hele, men
kun stræbt at gøre sig selv indsigtsfuldere og bedre, forhøje sin
egen Personligheds Værdi og saa udtale, hvad han havde erkendt
som godt og sandt. Men det var ikke, fordi han negtede,
at et Værks Værdi tilsidst bestemmes ved dets varige Betydning.
Tvertimod: naar han ikke vilde lade »den store Masses«
Behag være afgørende Maalestok, var det netop, fordi den let
blændes af den øjeblikkelige Glans, medens det er de skabende
Aanders Kald i Ensomhed og ofte under Miskendelse at forme
og præge det, der skal have Værdi for lange Tider. Det er det,
Digteren i Forspillet til »Faust« udtaler overfor Teaterdirektøren
og Klownen:

% CHECK p. 228: centred line: heading or verse

\begin{verse}
Was glänzt, ist für den Augenblick geboren; \\
% CHECK p. 228: centred line: heading or verse \\
Das Aechte bleibt der Nachwelt unverloren.
\end{verse}

Hermed stemmer Goethes etiske Teori, i hvilken han har
Fichte til Forgænger og Carlyle til Efterfølger. »Det Gode«,
siger han i en Samtale med Eckermann, »er intet Produkt af
menneskelig Reflexion, men er medfødt og medskabt skøn Natur.
Mere eller mindre er det alle Mennesker medskabt, men
i høj Grad enkelte, ganske fortrinligt begavede Aander. Disse
have aabenbaret deres guddommelige Indre ved store Gerninger
eller Lærdomme, og det har saa ved Skønheden i sin Fremtræden
vundet Menneskenes Kærlighed og vældigt ledet til Ærefrygt
og Efterlignelse. Men ved Erfaring og Visdom kunde det
Sædeligt-skønnes og Godes Værdi komme til Bevidsthed, idet
det Slette i sine Følger godtgjorde sig som Noget, der forstyrrede
den Enkeltes og det Heles Lykke, hvorimod det Ædle og
Rette viste sig som Noget, der hidførte og fæstnede den særlige
og den almene Lykke.«

Der var for Goethe et stadigt Vexelforhold mellem de Udvalgte
og Menneskenes Mængde. Ligesom Værkers og Hand%
% --- p. 229 ---
\setcounter{footnote}{0}%
\opage{229}lingers Værdi bestemmes ved deres blivende Virkninger, der
kun erfares ved at vende Betragtningen fra den enkelte skabende
Aand udad til større Kredse, saaledes er paa den anden Side
ogsaa selve denne Aand Virkning af Fortiden. Goethe var opfyldt
af Taknemmelighed overfor betydelige Forgængere og Samtidige.
»Man taler stedse om Originalitet,« sagde han i en Samtale,
»men hvad vil det sige! Saasnart vi ere fødte, begynder
Verden at virke paa os, og det gaar fremdeles saaledes lige til
Enden. Og hvad kunne vi idethele kalde vort Eget, uden Energien,
Kraften, Viljen! Hvis jeg kunde sige, hvad jeg skylder
store Forgængere og Medlevende, saa vilde Intet blive tilbage.«
I et bekendt humoristisk Digt har han udviklet, hvorledes den
Enkeltes Natur er et Produkt af Forældres og Forfædres Natur,
saa at det til Slut kan blive et Spørgsmaal, hvad Originalt der
er ved hele Fyren.

\section*{III.}

Der er en indre Overensstemmelse mellem disse Tanker og
Goethes hele Naturopfattelse og Livsanskuelse.

Hvert enkelt Organ og hvert enkelt organisk Væsen, Plante
eller Dyr, er for hans aandfulde Betragtning i al sin Selvstændighed
og Ejendommelighed en Aabenbaring af Typer og Love,
der i mangfoldigt nuancerede Former raade i den hele Natur.
Selvstændighed og lovmæssig Sammenhæng ere for ham ikke
Modsætninger. Tvertimod: det enkelte Organs og det enkelte
Væsens Selvstændighed bliver netop mulig ved, at de staa som
Udtryk for og Virkninger af en omfattende Natursammenhæng.
Gennem stadige og zarte Overgange hænger den ene Form sammen
med den anden, det Største med det Mindste i Naturens
uendelige Rige. I sine naturvidenskabelige Afhandlinger og i
sine Digte om Planternes og Dyrenes Metamorfose har Goethe
udviklet disse Tanker.

 De genfindes i analog Form i hans religiøse Opfattelse. I
sit Digt »Prometheus« faar han Luft for den svulmende Selvfølelse,
der udsprang af Bevidstheden om den skabende Digterkraft,
der rørte sig i ham, og som gjorde ham uafhængig af
% --- p. 230 ---
\setcounter{footnote}{0}%
\opage{230}ydre Magter. Med titanisk Trods vender han sig her mod en- 
hver Guddom, der tænkes at gribe ind i Livet udefra. Men med
»Prometheus« maa sammenstilles et andet Digt, »Menneskehedens
Grænser«. Den titaniske Selvfølelse overfor alle ydre
Magter grunder sig paa Forvisningen om, at den Kraft, han besjæles af, stammer fra Tilværelsens store Sammenhæng, fra den 
dybeste Kilde til alt Liv og al Virken. Saaledes hænge Selv- 
stændighed og Begrænsning nøje sammen.

Den Tankens Mester, som især havde givet Goethe Impulser
til hans Livs- og Verdensanskuelse, var Spinoza. Med smukke
og varme Ord har Goethe i sin Selvbiografi udtalt, hvad han
skylder denne Filosof. Livet fordrer paa saa mange Punkter
Selvopgivelse, Resignation. Men de fleste Mennesker hjælpe
sig igennem ved at sætte en liden Nydelse i Stedet for en
anden: Beskæftigelser, Tilbøjeligheder, Kæpheste, Alt prøver man,
for tilsidst at udraabe, at Alt er forfængeligt. »Kun faa Mennesker
gives der, som forud ane en saadan uudholdelig Følelse,
og som resignere en Gang for alle for at undgaa alle partielle
Resignationer. Disse overbevise sig om det Evige, Nødvendige,
Lovmæssige, og søge at danne sig Begreber, der ere uomstyrtelige,
ja i Stedet for at ophæves ved Betragtningen af det Forgængelige,
netop bekræftes derved.« Den store Uegennyttighed
og den dybe Fred, der raader i Spinozas Tankegang, fordi han
anlægger sin Betragtning saaledes, at det Enkelte stedse ses i
sin Sammenhæng med det Hele, som begrænset og dog ejendommeligt,
vandt Goethes Hjerte for ham. Han gav sig ikke
ud for at være nogen videnskabelig Spinozafortolker, men tog
af hans Skrifter, hvad han kunde bruge. I Virkeligheden var
nu ogsaa den Opgave, Spinoza stillede for Erkendelsen: at se
det enkelte Væsen i al sin Ejendommelighed som Led i Tilværelsens
store lovmæssige Sammenhæng, af den Natur, at den
ikke kan løses rent videnskabeligt, men at dens Løsning beror
paa Evne til kunstnerisk Anskuen. Mellem det, Spinoza fordrede,
og det, Goethe i sine største Værker har ydet, er der
derfor et endnu større Slægtskab, end Goethe i sin Ungdom
var sig bevidst. Det var mere for Livets end for Digtningens
Kunst, han lærte af Spinoza. Han lærte en Resignationens Kunst,
% --- p. 231 ---
\setcounter{footnote}{0}%
\opage{231}der ikke gør kold og fattig, men betinger en stadig Ungdom
gennem Hengivelsen til den evigt unge Naturs store Sammenhæng.
Det var især hos Spinoza, Goethe fandt Ideen om en
Gud, der ikke »støder til Verden udefra«, men virker og rører
sig i det Inderste i Tilværelsen.

Goethe ansaa det for umuligt at finde et afsluttende Tankeudtryk
for den Maade, paa hvilken Elementer og Kræfter virke
sammen i hele Tilværelsen saavel som i det enkelte Væsen. Her
fandt han en af »Menneskehedens Grænser«. I Faust's Trosbekendelse,
der utvivlsomt er Goethes egen, hedder det, at ethvert
Navn her kun er »Klang og Røg« overfor den overstrømmende
Livsfylde, i hvilket Mennesket føler sig baaren af en
Magt, der er større end han selv. Den religiøse Følelse var for
Goethe en Trang til af Taknemlighed frivilligt at hengive sig til
et Højere, Renere, Ubekendt og finde et Udtryk for den evigt
Unævnte. (Marienbad-Elegien). Men Mennesket er efter hans
Mening ikke født til at løse Verdens Problemer, men til at søge,
hvor det uløselige Problem ligger, og saa holde sig indenfor det
Begribeliges Grænser.

\section*{IV.}

Lad os dog ikke glemme, at Goethe ikke var Filosof, skønt
saa mange af hans Digte høre til de tankevægtigste og dybsindigste,
der nogensinde ere skrevne.

Forholdet mellem hans Liv og hans Digtning stemmer med
hans Karakter idethele. Hvad der greb ham i Livet med Glæde
eller Sorg, Vemod eller Anger, fik Luft og fandt Form i hans
Digtning. Ogsaa her gik Selvhævdelsen og Hengivelsen i Et.
Hans Produceren var Hengivelse: ti Billederne myldrede ofte
frem hos ham i saadan overstrømmende Fylde, at han neppe fik
Tid til at skrive dem ned. Under hans Omstrejfen i den frie
Natur fremkaldte de skiftende Scener saavel som hans egne Oplevelser
og Tanker Digte, han nynnede for sig selv, men ikke
altid kunde erindre bagefter. Ikke blot Stoffet, ogsaa Formen
stillede Fordringer til ham, som han bøjede sig for. Men i og
med denne Selvhengivelse virkede Selvhævdelsen: ti den
% --- p. 232 ---
\setcounter{footnote}{0}%
\opage{232}
kunstneriske Frembringen var ham en Selvbefrielse, en Frigørelse
 for, hvad der overvældede, tyngede og knugede ham :

% CHECK p. 232: centred line: heading or verse

\begin{verse}
Und wenn der Mensch in seiner Qual verstummt, \\
Gab mir ein Gott zu sagen, was ich leide.
\end{verse}

I hans Naturdigte vibrere de zarteste Naturstemninger.
Som ingen Anden før og efter ham formaaede han at leve og
føle sig ind i Naturen. Han betegner (efter Ruysdael og Rousseau
) en ny Epoke i Naturfølelsens Historie. Man har, hvad
Udtryksformen angaar, særligt fremhævet, hvor mange nye sammensatte,
lykkeligt dannede Ord han har frembragt for at kunne
udtrykke Naturbilleder og Naturstemninger kort og malende,
Ord som »morgenschön«, »wellenathmend«, »Mondendämmerung«, »Nebelthal«, »Lebensfluten« o. s. v. Al Naturfølelse beror paa en uvilkaarlig Overførelse (eller som Tyskerne træffende
kalder det: Einfühlung) af menneskelig Stemning i Naturen.
Goethe formaaede at lade denne uvilkaarlige Overførelse have
sit frie Løb og dog klart blive sig Forskellen mellem Naturen
og Menneskefølelsen bevidst. I det prægtige Digt »Ilmenau«,
der afgiver saa mange Bidrag til hans Karakteristik, fremstilles
herlige Naturscener, medens dog Modsætningen mellem Digterens
Sind og den Natur, der stemmer ham, gør sig gældende.
Naturens Ro staar i Modsætning til hans egen Bekymring, og
han siger til Træerne i Skoven:

% CHECK p. 232: centred line: heading or verse

\begin{verse}
Ich sorge still, indesz ihr ruhig grünet.
\end{verse}

Han kan ikke glemme, hvor haardt Menneskene midt i denne
store Natur maa slide for at holde Livet oppe. Bondens, Kulsvierens
og Bjergmandens haarde Liv fremtræder for hans Blik,
skønt Medfølelsen taber Noget af sin Bitterhed ved Bevidstheden
om, hvad han som Hertugens Raadgiver har formaaet at gøre
for den Egn, i hvilken han færdes. Derhos udtaler han sin Bekymring
over den Maade, hans Værker virke paa; hvad han
har villet, har ikke virket, som han tænkte sig det:
\begin{verse}
Ich brachte reines Feuer vom Altar;\\
Was ich entzündet, ist nicht reine Flamme.
\end{verse}

% --- p. 233 ---
\setcounter{footnote}{0}%
\opage{233}Og dertil kommer Bekymringen for hans fyrstelige Ven,
hvis urolige Sind har saa vanskeligt ved at arbejde sig hen til
den Selvbegrænsning, som især er nødvendig for den, der skal
være Fører for Andre. Men Taagen, der har sænket sig over
Landskabet, svinder, Egnen ligger i Dagens fulde Klarhed, og
den mørke Stemning giver Plads for et lysende Haab.

Med Urette har man beskyldt Goethe for, at han paa egoistisk
Vis forvandlede alle Livsforhold til poetiske Emner eller
endog til Midler for poetisk Selvnydelse. Den Selvbefrielse, han
fandt i Digtningen, vidner netop imod denne Beskyldning. Der
gaar en stille Selvanklage gennem en hel Række af hans Digte,
--- Gøtz, Clavigo, Faust og flere. Han kunde ikke glemme, at
den store Opgave, hans Selvudvikling stillede ham, i hans Ungdom
havde ført ham til et Brud, der gjorde det af med en
ung Kvindes Livslykke. Friederikes Skæbne har Goethe følt
som sin Ulykke og som sin Skyld. Man hører dette gennem
Faust's Ord:

% CHECK p. 233: centred line: heading or verse

\begin{verse}
Bin ich der Flüchtling nicht? Der Unbehaus'te? \\
Der Unmensch ohne Zweck und Ruh, \\
% CHECK p. 233: centred line: heading or verse \\
Der wie ein Wassersturz von Fels zu Felsen brauste, \\
Begierig wüthend nach dem Abgrund zu! \\
% CHECK p. 233: centred line: heading or verse \\
Und seitwärts sie, mit kindlich dumpfen Sinnen, \\
Im Hüttchen auf dem kleinen Alpenfeld, \\
Und all ihr häusliches Beginnen \\
Umfangen in der kleinen Welt.
\end{verse}

\begin{verse}
Sie, ihren Frieden muszt' ich untergraben!\\
Du, Hölle, musztest dieses Opfer haben!
\end{verse}

Goethe har selv sagt, at hans Digte vare Lejlighedsdigte,
fremkaldte af Livet selv, først og fremmest hans eget Liv, og »forfattede
under umiddelbar Anskuen«. Og naar han tillige har
talt om sin alvorlige Maade at tage Livet paa (des Lebens
ernstes Führen), saa har han kunnet gøre det med den Bevidsthed,
at han har søgt at lade Livet komme til sin Ret ikke mindre
end Kunsten. Han fandt det træffende, naar det var sagt
om ham, at hans Stræben gik ud paa at give det Virkelige en
poetisk Form. Det var hans Livsopgave, og det var hans Alvor
% --- p. 234 ---
\setcounter{footnote}{0}%
\opage{234}med denne Livsopgave. Det er jo ogsaa en Opgave, som kommer
Livet selv til Gode, idet det Billede af Livet, Digterne
danne, igen virker tilbage paa det fortsatte Liv. Man har derfor
ikke Ret til at tale om en Modsætning mellem Mennesket
og Digteren. Det at være Digter og som Digter give et Billede
af Livet var Goethes Livsgerning \textit{som Menneske}\footnote[1]{Sammenlign, hvad Goethe siger om Schiller i Epilog zum Lied von der Glocke: Er wendete die Blüthe höchsten Strebens, / Das Leben selbst, an dieses Bild des Lebens.}, og bragte
den ham i Konflikter med andre menneskelige Forhold og Opgaver, Konflikter, der efterlode deres Braad i hans Indre, saa vilde
det i hvert Tilfælde ikke være udtømmende at forklare disse
Konflikter ved en Modstrid mellem Digteren og »Mennesket«.
Til Grund for dem laa en Strid mellem to \textit{menneskelige} Opgaver, en Strid, i hvilken det ikke er saa let at fælde Dom,
som det ofte er sket, men af hvilken Goethe iøvrigt efter eget
Sigende ikke gik ud uden Skrammer.

\section*{V.}

Studiet af Goethe er ikke mindst derved af saa stor Interesse
i psykologisk Henseende, at vi hos ham finde alle forskellige Grader i den Bevidsthed, ved hvilken Fantasien virker.
Han kunde forme med klart Overlæg, saavel som improvisatorisk,
og han kendte --- især i sin Ungdom --- den aldeles
uvilkaarlige, instinktmæssige Virken af Fantasien. Nogle af sine
Ungdomsværker har han digtet i en Slags søvngængeragtig Tilstand,
hvor han ikke vidste, hvad han gjorde, skønt han mærkede,
at det arbejdede i ham. Billederne fremstillede sig pludseligt
helt færdige for ham som »frie Guds Børn«, der sagde:
»Her er vi!« Det er, hvad han senere kaldte det Dæmoniske
i Geniet. Ligesaa lidt som Naturen, sammentigger Geniet Elementerne
til sine Helhedsværker. Der er fra først til sidst et
Vexelforhold mellem Helhed og Elementer, der trodser enhver
Analyse. Derfor er det ogsaa umuligt at udtrykke et Digterværks
Indhold i en enkelt Tanke, en enkelt Idé, --- ligesaa
% --- p. 235 ---
\setcounter{footnote}{0}%
\opage{235}umuligt, som det er at udtrykke Tilværelsens Væsen i en enkelt
Tanke. Et stort Digterværk er ligesaa vel som Naturen
inkommensurabelt for Tanken. »Jeg er«, siger Goethe, »af den
Mening, at jo mere inkommensurabel og for Tanken ufattelig
en poetisk Produktion er, des bedre. Der komme de og spørge,
hvilken \textit{Idé} jeg har søgt at legemliggøre i min Faust. Som
om jeg selv vidste det og kunde sige det!\dots{} At et Menneske
der fra svære Vildfarelser stedse stræber opad mod det Bedste,
kan forløses, det er vel en god Tanke, der er virksom og forklarer
Meget; men det er ingen \textit{Idé}, der skulde lægges til
Grund idethele og for hver enkelt Scene. Det vilde da ogsaa
være blevet en dejlig Ting, om jeg havde villet trække et saa
rigt, broget og mangfoldigt Liv, som jeg har fremstillet i Faust,
paa en eneste gennemgaaende \textit{Idés} magre Snor!«

Det var ikke Goethes Mening at ville miskende Tankens
Ret, hverken overfor Kunstens eller overfor Naturens Værker.
Men han hævdede Umuligheden af, at Tankens Arbejde kunde
bringes til Afslutning i en udtømmende Idé, og han frygtede
--- hvad den spekulative Filosofi, der udfoldede sig omkring
ham, ogsaa lod befrygte, --- at man skulde afslutte forhastet,
saa at abstrakte Tanker sattes i den umiddelbare levende Anskuelses
Sted.

Dog vilde man for Goethes hele Liv og Virken neppe kunne
finde noget bedre Udtryk end den Tanke, der efter hans egen
Indrømmelse udtrykker en væsentlig Side af Indholdet i hans
største Digt: Tanken om den stadige Stræben. Igennem hans
lange Liv gaar et uafbrudt Aandsarbejde, en stadig Higen efter
Fylde og Klarhed, og Hjertets Varme bevares helt op i Alderdommen.
Dermed skal Forskellen mellem den unge og den
gamle Goethe ikke negtes. Den rolige Kontemplation kom ganske
naturligt til at spille en større Rolle hos den gamle Mand end
hos den unge, der omtumledes i Trængselsperioden. Men allerede
i Ungdommen saa han midt i Bølgegangen op til de faste
Stjerner for at orientere sig og faa Hold paa sig selv. Og Stjernerne
blegnede ikke, da Bølgegangen lagde sig, saa at rolig Beskuen
blev mulig.

% --- p. 236 ---
\setcounter{footnote}{0}%
% CHECK p. 236: centred line: heading or verse
\opage{236}

\section*{VI.}

Selv den Største har sin Begrænsning, og skal han være
Noget for alle Tider, er der Adskilligt, man maa se bort fra eller
ændre. Hverken i den aandelige eller i den materielle Verden
ere Livsvilkaarene de samme til alle Tider. Opgaverne, der skulle
løses, og Hindringerne, der skulle overvindes, skifte med Tiderne. Det er da intet Under, at Meget forekommer os fremmed ---
selv hos den Største i den Generation, som indviede det Aarhundrede,
der nu er ved at løbe ud.

Man har undertiden villet betragte Goethe som en af Darwins
Forgængere. Men netop de for Darwins Lære mest karakteristiske
Punkter mangle hos Goethe: Ideen om Kampen for
Tilværelsen og Ideen om Arternes historiske Udvikling gennem
Afstamning. Goethes Naturopfattelse er præget af en begejstret
Optimisme. Han ser Naturlivet som en harmonisk Udfoldelse af
evige Typer. Vel er der stor Forskel mellem Vilkaarene for
Livet idethele; men der er for Goethe ingen bitter Modstrid tilstede
mellem, hvad det Heles Bestaaen kræver, og hvad det
enkelte Væsen stræber efter. Naturen er en Moder, der skjuler
sig bag sine Børn, men trods sin Usynlighed ikke lader sin
Moderkærlighed uden Vidnesbyrd. Naar den hæmmer Udviklingstrangen
paa et enkelt Punkt, er det kun for blidt at lede
Udviklingen over til højere Former.

\begin{verse}
Doch hier hält die Natur, mit mächtigen Händen, die Bildung\\
An und lenket sie sanft in die Vollkommenheit.
\end{verse}

For Darwin viser Naturen os ingen harmonisk Leg af Former
og Kræfter. Vejen til nye Former gaar gennem Hæmning,
Lidelse og Undergang. Der falder nu stærkere Skygger over
Naturlivet: det er kun paa Overfladen, at Naturen viser os Lys
og Fred. Naar Goethe i Ilmenaudigtet følte Naturens Ro i Modsætning
til sin egen Sorg, saa lærer den nye Naturopfattelse, at
ogsaa Naturens Liv, ikke mindre end Menneskets, kan hæmmes
og kues, saa at det bliver vanskeligere end forhen at »føle sig
ind« i Naturen. Alt medens Træerne staa tilsyneladende roligt
og grønnes, kæmpe de indbyrdes om Lys og Luft og Plads til
at brede sig.

% --- p. 237 ---
\setcounter{footnote}{0}%
\opage{237}Ogsaa paa selve det menneskelige Omraade ere vi blevne
nødte til at lægge større Vægt paa Strid og Modsætninger, end
Goethe og hans Tid gjorde. Der var for Goethe en dyb Analogi
mellem Natur, Kunst og Samfundsliv. Hans Ideer om disse
»tre store Verdensegne«, som han kalder dem, udviklede sig
samtidigt og i Overensstemmelse med hverandre. Men ogsaa
paa det sociale Omraade falde Skyggerne nu stærkere, og ere
Tvivlene blevne mere skærende. Skal Muligheden af en harmonisk
Løsning fastholdes, behøves der en kraftig Vilje og et
større Indskud af Haab og Frejdighed, end Goethes rolige Beskuen
medførte eller ansaa for nødvendigt. Goethes Harmonilære
mangler ikke Energi. Men der stilledes den ikke den Opgave
at spænde over saa store Modsætninger, som stilles et Forsøg
i lignende Retning i vore Dage. Hans Livsanskuelse bliver
i vore Øjne naiv.

Og dog --- maaske har Naiviteten trods Alt Ret. Den sande
Storhed fremtræder med Naivitetens Præg, og overfor Livet og
dets Problemer gælder det maaske tilsidst, at vi maa vedblive
at være som Børn. Selvom en Faust, hvis den digtedes i vore
Dage, maatte indeholde andre og tildels dyberegaaende Problemer
end det goetheske Værk, saa vilde Grundtankerne dog sikkert
vedblive at være de samme. Goethe staar endnu som vor
Lærer og vil vedblive at staa saaledes i lange Tider. De »Guds
Børn«, der fremtraadte for hans Tanke og hans Fantasi, ere
endnu ikke blevne Oldinge. Atter og atter vende vi tilbage til
hans rige Billed- og Tankeverden, til hans Naturstemninger og
hans Livserfaringer. Naar vi fejre hans Mindefest, er det ikke,
fordi hans Aand er bleven et blot Minde for os; den følger os
i vor Stræben efter at føle og forstaa Livet, og Alt, hvad der
efter hans Tid er fremkommet paa Kunstens og Tankens Omraade,
har ikke ladet os vokse fra denne Fører.
\end{document}
